\documentclass[11pt,a4paper]{article}
\pdfmapfile{+lm.map}
\pdfmapfile{+cm.map}
\pdfmapfile{+cmextra.map}
\pdfmapfile{+symbols.map}
\pdfmapfile{+euler.map}
\usepackage[T1]{fontenc}
\usepackage[utf8]{inputenc}
\usepackage{lmodern}
\usepackage[margin=27mm]{geometry}
\usepackage{amsmath,amssymb,amsthm,mathtools}
\setcounter{MaxMatrixCols}{20}
\usepackage{booktabs,longtable,array,enumitem}
\usepackage{xcolor}
\usepackage{microtype}
\usepackage{xurl}
\usepackage[colorlinks=true,linkcolor=blue!45!black,citecolor=blue!45!black,urlcolor=blue!45!black]{hyperref}
\newtheorem{theorem}{Theorem}[section]
\newtheorem{proposition}[theorem]{Proposition}
\newtheorem{lemma}[theorem]{Lemma}
\newtheorem{corollary}[theorem]{Corollary}
\newtheorem{sourcefact}[theorem]{Cited source statement}
\theoremstyle{definition}
\newtheorem{definition}[theorem]{Definition}
\newtheorem{remark}[theorem]{Remark}
\setlength{\emergencystretch}{2em}
\setlength{\parskip}{3pt}
\title{K17: Step 17, the $K$-groups of $\mathbb Z$,\\
Soul\'e characters and the pro-$\ell$ kernels of the tripod}
\author{Research run K17---checked deductions and explicit open gates}
\date{21 November 2026\\Timed checkpoint; research in progress}
\begin{document}
\maketitle
\begin{abstract}
The exact quantitative Selmer theorem behind step 17 is recovered from
the pinned source. Its rational-field specialization, together with
etale $K$-theory and arithmetic duality, bounds the orders of
$K_{4k}(\mathbb Z)$; a reconstruction at fixed levels yields exponent
seven. We separate the exact Vandiver equivalence, full-domain Soul\'e
indices, and the stronger integral generation claims. Complete
deductions settle the pro-$2$ kernel-field equality with unrestricted
real places and every odd prime satisfying Vandiver, hence every
prime below $2^{31}$. The proof combines a general pole-separation
criterion with the integral modular-symbol theorems of Fukaya--Kato
and Wake--Wang-Erickson; no simple-root hypothesis is needed. More generally, the one-variable Greenberg
condition that the plus Iwasawa module is finite suffices, through a
bounded rational-coordinate argument on the entire arithmetic group. At every
irregular prime satisfying this condition the integral Lie abelianization has nonzero
finite primary torsion in infinitely many even weights, so integral odd
generation fails. The examples at $37$ and $691$ locate explicit defects.
At $2$ a central-extension
parity induction constructs adapted Newton generators in every odd
weight and proves integral free generation in all weights. A separate
integral Steinberg-homology source theorem removes the rank-eleven term
from a finite $K_{12}$ criterion; a boundary-degree indexing discrepancy
in the older cell data is explicitly isolated. The other remaining
cases retain explicit gates. Cote 149, scans 79--85, is read through
exact limiting Kummer sequences. A marked generic Galois diagram
determines the geometric tangent subset. Counterexamples rule out
automorphism-invariant recognition from the isolated normalizer and
from the full relative diagram with the generic marking omitted.

This checkpoint is not the final package. Every completed new proof
has a whole-proof reading recorded in the check ledger. Literature
theorems, exact arithmetic computations and source inspections have
distinct check types. No new Lean build, formal certificate or external
peer review is claimed. Unresolved assertions appear separately.
\end{abstract}
\tableofcontents
\clearpage
\part{Checked deductions and rechecked source statements}

\section{Recovering the arithmetic statement behind step 17}
\label{sec:step17}

\paragraph{Source contract and provenance.}
Alp\"oge's \emph{Remark}, p.~4, step~17 and footnote~6, names
$H^1(\mathcal O_E[1/\ell],\mathbb Q_\ell/\mathbb Z_\ell(-e))$
and a weakening of Vandiver, but supplies neither quantifiers nor a theorem
number for that weakening. The attached map identifies candidate modules only.
We recovered those modules from the public mirror
\url{https://github.com/CocoSgt/riemann-hypothesis-lean}, at commit
\texttt{7cbf7a0c32fca1badf05931879a37c3f8a2b03b4}.
The retrieved \path{src/Main.lean} is byte-identical to the supplied file;
its SHA-256 is recorded in the evidence manifest. Git blob identifiers were
checked for the manifest-based downloads. These are source inspections, not
a new Lean build or a claim to have authenticated the entire mirror. Under
the prompt's stipulated theorem contract, we use the recovered unconditional
Alp\"oge theorem below. The arithmetic translation and consequences are the
new deductions checked here.

In this section a source path beginning with \texttt{AV/} abbreviates
\path{src/Development/AutomorphicArithmetic/AutomorphicVarieties/}.
The source facts have precise declaration locators; no theorem is inferred
merely from a file name or from a proposition's definition.

\subsection{The local conditions and the fields}

Put $D_\ell(e)=\mathbb Q_\ell/\mathbb Z_\ell(-e)$ and, for a number field $E$,
\[
 \mathcal S_\ell(E,e)=
 \ker\left(H^1_{\mathrm{et}}(\operatorname{Spec}\mathcal O_E[1/\ell],D_\ell(e))
       \longrightarrow\prod_{v\mid\ell}H^1(E_v,D_\ell(e))\right).
 \tag{S1}
\]
Thus the selected classes are unramified outside $\ell$ and locally trivial
at every place above $\ell$. They are a subgroup of the displayed $H^1$,
not that entire group. The source definitions are
\path{src/Formalism/GaloisCohomology/Statement.lean}, declarations
\texttt{LData}, \texttt{qzTw}, \texttt{Sha}, and \texttt{shaLen}.

For a finite extension coefficient ring $O\subset\overline{\mathbb Q}_\ell$
which is finite over $\mathbb Z_\ell$ and is a DVR, write
\[
 D_O=O\otimes_{\mathbb Z_\ell}(\mathbb Q_\ell/\mathbb Z_\ell),
 \qquad e(O)=\operatorname{length}_O(O/\ell O).
 \tag{S2}
\]
For a Dirichlet character $\psi$ modulo $N$ with values in $O^\times$, define
$\mathcal V_{\ell,N}(e,\psi)$ to be the subgroup of
$H^1(G_{\mathbb Q},D_O(\psi^{-1}\chi_\ell^{-e}))$
whose restriction is zero on every finite inertia group and on every
decomposition group at a prime dividing $\ell N$. Here $\chi_\ell$ is the
cyclotomic character. Set
\[
 v_{\ell,N}(e,\psi)
 =\operatorname{length}_O\mathcal V_{\ell,N}(e,\psi)/e(O).
 \tag{S3}
\]
In (S3) the quotient of lengths is used only when the module has finite
length. Literally, the source's \texttt{Vn} applies
\texttt{ENat.toNat} to its extended-natural lengths (which sends
infinity to zero); every estimate used below simultaneously proves
finiteness, so (S3) then gives exactly its value. These are
\texttt{Coef}, \texttt{Vlen}, and \texttt{Vn} in
\path{AV/HeightAndVanishing/DirichletPrimeToLPartMulLPowerPart.lean}, together
with \texttt{QO}, \texttt{eO}, \texttt{shaFam}, and \texttt{ShaU} in
\path{AV/GaloisLibrary/SelmerTransport/ConjStableShaFam.lean}.
The scalar character is defined in
\path{AV/CharacterLevelBounds/ShaFamMono.lean}, \texttt{rho}.

\begin{proposition}[K17-FIELD-DESC: the field interpretation]
Let $e>0$, let $\ell$ be odd, and suppose $E/\mathbb Q$ is a cyclic
extension of degree prime to $\ell$, cut out by $\psi$. Suppose $O$ contains
the values of $\psi$. Restriction gives an isomorphism
\[
 \mathcal V_{\ell,N}(e,\psi)
 \simeq
 \left(O\otimes_{\mathbb Z_\ell}\mathcal S_\ell(E,e)\right)^{\psi},
 \tag{S4}
\]
where the right side is the $\psi$-eigenspace for the natural action of
$\operatorname{Gal}(E/\mathbb Q)$ on cohomology, including its action on
the Tate coefficients. The same description applies to character components
inside a larger abelian extension of degree prime to $\ell$.
\end{proposition}
\begin{proof}
Write $\Delta=\operatorname{Gal}(E/\mathbb Q)$. All positive-degree
cohomology of $\Delta$ with $\ell$-primary coefficients vanishes: it is
annihilated by $|\Delta|$, which acts invertibly on those coefficients.
Inflation--restriction consequently identifies the global $H^1$ on the
left with the indicated eigenspace of $H^1(G_E,D_O(-e))$. The same argument
applies to decomposition and inertia groups, since their quotient groups
have order dividing $|\Delta|$. Hence restriction detects the vanishing of
each local restriction.

An unramified local class at a place $w\nmid\ell$ of $E$ lies in
\[
 H^1_{\mathrm{unr}}(E_w,D_O(-e))
 =D_O/(\operatorname{N}w^{-e}-1)D_O=0.
\]
The last equality holds because $\operatorname{N}w^{-e}-1$ is a nonzero
element of $\mathbb Q_\ell$ and multiplication by any nonzero scalar from
$\mathbb Z_\ell$ is surjective on the divisible module $D_O$.
Thus all extra decomposition conditions away from $\ell$, including
those dividing $N$, are automatic after the inertia conditions. The local
conditions at $\ell$ are exactly (S1). Restriction therefore gives (S4).
Finite flatness of $O$ over $\mathbb Z_\ell$ commutes with the cochain
description and the eigenspace idempotent, so no scalar-extension defect
remains. This proves the assertion.
\end{proof}

For CM fields the source's ``plus'' notation is twist-dependent. If $c_E$
is complex conjugation, then
\[
 \mathcal S_\ell(E,e)^+
 =\{x\in\mathcal S_\ell(E,e):c_E x=(-1)^e x\}.
 \tag{S5}
\]
It is not always the $+1$ eigenspace. See
\path{AV/LengthEqAddEigen.lean}, \texttt{ShaPlus}, \texttt{ShaMinus},
\texttt{shaLen\_eq\_add}, and
\path{AV/PartsSign.lean}, \texttt{mk\_mem\_ShaPlus\_iff}.
The direct sum of these two eigenspaces uses that $2$ is invertible;
the source treats $\ell=2$ separately.

The fields used in the larger twisting argument are not arbitrary number
fields. The source's \texttt{Host} fixes a Dirichlet character $\chi$ of
modulus $m$, an integer $a\ge2$, and primes $q\equiv3\pmod8$,
$d\equiv7\pmod{16}$ with $(qd,m)=1$. They determine
\[
 E=F_\chi\,\mathbb Q(\mu_{2^{a+2}})^+
      \mathbb Q(\sqrt{-q})\,\mathbb Q(\sqrt{-d}),
 \qquad N=\operatorname{lcm}(2^{a+2},m)qd,
 \tag{S6}
\]
where $F_\chi$ is the fixed field of $\ker\chi$. This is the fixed-field
description encoded by \texttt{IsCompField} and
\texttt{isCompField\_ofChar} in \path{AV/StmtFields.lean}, and by
\texttt{Host.H} in
\path{AV/HeightAndVanishing/DirichletPrimeToLPartMulLPowerPart.lean}.
Indeed the four congruence conditions in that definition fix respectively
the real $2$-power cyclotomic field, $F_\chi$, and the two quadratic fields;
the intersection of the corresponding subgroups fixes their compositum.
The selected character components have $\psi(-1)=(-1)^e$.

For completeness, the exact finite-parameter organization is in
\path{AV/HeightAndVanishing/LayerLe.lean}, where \texttt{PointHV}
is the proposition defined to assert the following statement:
for each $\chi$ and function $j:\mathbb N\to\mathbb N$, there are
$C,r_1,r_2\in\mathbb N$, with $r_1,r_2$ prime, such that, for every host
with character $\chi$, both auxiliary primes smaller than $q,d$, every
$t\ge2$ satisfying \texttt{Host.Big}, and every $B\ge0$ satisfying
the specified \texttt{LayerLe} bounds at the finite set \texttt{Bchi},
the weighted sum of the lengths in (S5), over all odd primes and
$1\le e<t$, is bounded by $(Nt)^C+CB$.
Here \texttt{Host.Big} means $2t+1<q,d$ and each of
$m,\operatorname{ord}\chi,p_K(\chi),h_K(\chi),p_1(\chi)$ is less than
$q,d$. The finite exceptional set is
\[
 B_\chi=\{3\}\cup
 \{\text{odd prime divisors of }
   (\operatorname{ord}\chi)m p_K(\chi)h_K(\chi)p_1(\chi)\}.
\]
The three auxiliary quantities are the least primes/class number specified
in \path{AV/StmtFields.lean}, \texttt{pK}, \texttt{hStar}, \texttt{pOne}.
We record this definition to identify the fields and classes in the
note. The closed source proof composition in the following host-family
subsection supplies its theorem inputs; the stated layer bounds remain
explicit premises. The deductions about $\mathbb Z$ below use the unconditional,
quantified theorem next, and do not assume an unverified numerical value
for any of these constants or layer bounds.

\subsection{The unconditional character theorem}

The exact admissibility condition \texttt{GD} is in
\path{AV/HeightAndVanishing/Counts.lean}. In ordinary notation it says:
$e\ge1$, $N$ is even, $\ell$ is odd, the order of $\psi$ is prime to
$\ell$, the factor $\psi_\ell$ supported at the $\ell$-part of the
conductor satisfies $\psi_\ell^2=1$, and
$\psi(-1)=(-1)^e$. There is an auxiliary imaginary quadratic field $K$
such that
\[
 \ell\nmid |D_K|h_K,\quad (|D_K|,N)=1,\quad
 K\not\subset\mathbb Q(\mu_\ell)F_{\psi/\psi_\ell},\quad
 (\ell-1)\nmid2e. \tag{S7}
\]
The third condition is exactly the source's subgroup noncontainment
\texttt{eta}; it has not been silently deleted. The source definition
of \texttt{chiP} agrees with $\psi/\psi_\ell$ under its prime-to-$\ell$
order hypothesis. In particular these conditions force $\ell\ge5$.

\begin{sourcefact}[Alp\"oge: K17-S17, the recovered quantitative theorem]
For every imaginary quadratic $K$ there exist $c_5,A\in\mathbb N$ with
the following property. Fix $e,N$, a subgroup
$H\subset(\mathbb Z/N\mathbb Z)^\times$, and distinct primes $p,p'$
not dividing $30N|D_K|$. For each prime $\ell$ in a finite set $L$ choose
a permitted coefficient ring $O_\ell$. If at least one character with
kernel $H$ satisfies (S7) and the preceding admissibility conditions at
each $\ell\in L$, then all the associated lengths are finite and
\[
 \sum_{\ell\in L}\log\ell
   \sum_{\ker\psi=H}v_{\ell,N}(e,\psi)
 \le 2c_5\varphi([(\mathbb Z/N\mathbb Z)^\times:H])
          (N(e+1))^A(\log p+\log p'). \tag{S8}
\]
The character sum is over the characters valued in $O_\ell$; existence
of one such faithful character provides its whole orbit.

Exact locators: \path{AV/UnitaryThreeGroups/KatoUse.lean},
\texttt{CorGK}; and \path{AV/SmallClassNumbers/CorGKOf.lean},
\texttt{UpperTheory.V.SmallClassNumbers.corGK}, an unconditional theorem.
\end{sourcefact}

\paragraph{Proof reconstruction and the cited inputs.}
The arithmetic input \texttt{PropGK} supplies, for an auxiliary prime $p$,
a nonzero integer $n_p$ such that
\[
 \log|n_p|\le c_5\varphi([G:H])(N(e+1))^A\log p,
 \qquad
 \sum_{\ker\psi=H}v_{\ell,N}(e,\psi)\le2v_\ell(n_p)
 \quad(\ell\ne p). \tag{S9}
\]
Its source proof is \texttt{propGK} in
\path{AV/UnitaryThreeGroups/OrbEquiv.lean}. It combines
\texttt{OneChar} (the Selmer-length estimate by rank-three/rank-two
automorphic congruences) with \texttt{CongruenceNumberRankThree}
(a nonzero norm of Frobenius-trace differences). The latter provides the
integer and its archimedean size, while the former bounds its local
valuation. The faithful-character orbit has size $\varphi([G:H])$.
Their unconditional inputs are assembled in
\texttt{SmallClassNumbers.corGK\_of}, \texttt{litN}, and \texttt{corGK}
of the file cited above: local and global Galois cohomology, the Kato
input, adjoint vanishing, automorphy, and congruence-number estimates.
These are cited Alp\"oge theorems under the source contract; they are
not independent new automorphy proofs in this paper.

Here is the complete elementary passage from (S9) to (S8).
Use (S9) twice, with $p$ and $p'$. For each $\ell$ at least one applies,
so the nonnegative length sum is at most
$2v_\ell(n_p)+2v_\ell(n_{p'})$. For a nonzero integer $n$,
\[
 \sum_{\ell\in L}v_\ell(n)\log\ell\le\log|n|
\]
because the product of the indicated prime powers divides $|n|$.
Summing and using the two archimedean estimates proves (S8), including
the auxiliary-prime cases. This is precisely
\path{AV/UnitaryThreeGroups/CorGSumLog.lean}, \texttt{corGK}.
The complete source proofs of this passage and of the orbit/norm
translation were read in this run.

\subsection{The rational-field specialization}

\begin{theorem}[K17-S17-Q]
There exist absolute $C,A\in\mathbb N$ such that for every positive even
integer $e$ and every finite set of primes $L$ satisfying
$\ell\ge5$ and $(\ell-1)\nmid2e$,
\[
 \mathcal S_\ell(\mathbb Q,e)\text{ is finite},\qquad
 \sum_{\ell\in L}\log\ell\,
   \operatorname{length}_{\mathbb Z_\ell}\mathcal S_\ell(\mathbb Q,e)
 \le C(e+1)^A. \tag{S10}
\]
In particular the group is zero whenever $\log\ell>C(e+1)^A$
and the displayed prime conditions hold.
\end{theorem}
\begin{proof}
In (S8) take $K=\mathbb Q(\sqrt{-3})$, $N=2$, $H=G$, the trivial
character, and auxiliary primes $p=7,p'=11$. The discriminant is $-3$
and the class number is one. The character conditions are immediate;
evenness of $e$ gives the parity condition. For $\ell\ge5$ the field
$K$ is not contained in $\mathbb Q(\mu_\ell)$, since the latter is
unramified at $3$ whereas $K$ is ramified there. This checks (S7).
All these checks are also made by \texttt{gd} in
\path{AV/SmallClassNumbers/ClassNumberThree.lean}.

Choose the bottom coefficient ring. It is isomorphic to
$\mathbb Z_\ell$ by \texttt{exists\_coef} in
\path{AV/SmallClassNumbers/LenEqZeroIff.lean}. Then $D_O=D_\ell(0)$
and $e(O)=1$, since $O/\ell O\simeq\mathbb F_\ell$ is a simple
$O$-module. The extra local condition at $2$ is automatic by the
unramified local calculation in the proof of K17-FIELD-DESC.
Thus the length in (S8) is the length in (S10).

The right side of (S8) is
$2c_5 2^A(e+1)^A(\log7+\log11)$, bounded by
$36c_5 2^A(e+1)^A$. Take $C=36c_5 2^A$.
This is the proof \texttt{bound\_of} in
\path{ClassNumberThree.lean}; its unconditional input and conclusion
are \texttt{corGK} and \texttt{bound} in
\path{SmallClassNumbers/CorGKOf.lean}.
For the last assertion use $L=\{\ell\}$. A nonzero finite
$\mathbb Z_\ell$-module has length at least one, contradicting
the strict inequality on $\log\ell$. The source independently states
this cocycle formulation as \texttt{Large} and proves
\texttt{large\_of} in \path{LenEqZeroIff.lean}.
\end{proof}

\paragraph{What is and is not numerical here.}
The source has also a proposition \texttt{CorG} with a displayed
$e^7N^{18}|D_K|^5$ shape. A definition of that proposition is not a
proof of it. We have used the established \texttt{CorGK}, with its
existential exponent, throughout. Neither ``100ish'' in footnote~6
nor an unspecified constant is a certified degree range. The next
section gives every vanishing and cyclicity conclusion that actually
follows from (S10), and identifies the remaining numerical gate.

\subsection{A closed theorem for the Case~I host family}
\label{sec:host-case-one}

\paragraph{K17-HOST-CASE-I: a theorem, beyond the field definitions.}
There is a closed source proof of the host-family assertion
\texttt{CountI}.  In full namespace notation its proof term is the
composition
\[
 \texttt{UpperTheory.V.U3.countI}\;
 \texttt{UpperTheory.V.SmallClassNumbers.corGK}.
\]
The first theorem has type
\(\texttt{U3.CorGK}\longrightarrow\texttt{HV.CountI}\);
the second is the closed theorem already used for the rational-field
specialization.  This composition uses the established, field-dependent
polynomial theorem \texttt{CorGK}.  It does not require the stronger
uniform discriminant and level exponents appearing in the separately
defined proposition \texttt{HV.CorG}.

Here is the precise mathematical reading.  Fix a character \(\chi\) in
the source's \texttt{PChar} domain, of modulus \(m\), and put
\(o=\operatorname{ord}(\chi)\).
Let \(p_K\) be the least prime at least seven congruent to three modulo
four and coprime to \(m\); set
\[
 K_\chi=\mathbb Q(\sqrt{-p_K}),\qquad
 h_\chi=|\operatorname{Cl}(K_\chi)|.
\]
Let \(p_1\) be the least prime congruent to one modulo the odd part of
\(o\), and coprime to \(2mp_K\).  The finite set used by the source is
\[
 \mathcal B_\chi=\{3\}\ \cup\
 \{\ell:\ell\text{ is an odd prime dividing }omp_Kh_\chi p_1\}.
\]
There exist natural numbers \(C\) and distinct primes \(r_1,r_2\),
depending only on \(\chi\), with the following property.

Choose a host \(h=(\chi,a,q,d)\), where \(a\ge2\), \(q,d\) are primes,
\(q\equiv3\pmod8\), \(d\equiv7\pmod{16}\), and
\(\gcd(qd,m)=1\).  These are exactly the source's
\texttt{PrimePair} conditions.  Its level and associated field are
\[
 N=\operatorname{lcm}(2^{a+2},m)qd,
 \qquad
 E_h=F_\chi\,\mathbb Q(\mu_{2^{a+2}})^+
          \mathbb Q(\sqrt{-q})\,\mathbb Q(\sqrt{-d}).
\]
Require
\[
 r_i<q,\quad r_i<d\quad(i=1,2),
\]
and, for \(t\in\mathbb N\), the explicit host-size conditions
\[
 2t+1<q,d,
 \qquad y<q,d\quad
 (y\in\{m,o,p_K,h_\chi,p_1\}).
\]
These are precisely \(h.\texttt{Big}(t)\).  Let \(e\in\mathbb N\) with \(0<e<t\), and let
\(L\) be any finite set of odd primes such that
\[
 \ell\notin\mathcal B_\chi,
 \qquad (\ell-1)\nmid 2e
 \qquad(\ell\in L).
\]
This pair of conditions is exactly \(h.\texttt{CaseI}(\ell,e)\).

For each \(\ell\), take any coefficient ring \(\mathcal O_\ell\) of the
source, and let \(X_h(\mathcal O_\ell,e)\) be the characters modulo
\(N\) that factor through \(\operatorname{Gal}(E_h/\mathbb Q)\) and
satisfy \(\psi(-1)=(-1)^e\).  Use \(V_{\ell,N}(e,\psi)\) for the
cohomology group with the everywhere-inertia and selected decomposition
conditions defined in the source reconstruction above.  Then all these
groups, for \(\ell\in L\) and \(\psi\in X_h(\mathcal O_\ell,e)\),
have finite length, and
\begin{equation}\label{eq:host-case-one-bound}
 \sum_{\ell\in L}\log\ell\!
 \sum_{\psi\in X_h(\mathcal O_\ell,e)}
 \frac{\operatorname{length}_{\mathcal O_\ell}
                   V_{\ell,N}(e,\psi)}
      {\operatorname{length}_{\mathcal O_\ell}
                   (\mathcal O_\ell/\ell\mathcal O_\ell)}
 \le (Nt)^C.
\end{equation}
The constants precede the host, the twist, the coefficient choices and
the finite prime set.  Both even and odd \(e\) are permitted here;
the character parity changes with \(e\).  No assertion about the whole
ambient \(H^1\) is substituted for these local kernels.

\paragraph{Proof and source check.}
Apply \texttt{CorGK} to the fixed imaginary quadratic field
\(K_\chi\), obtaining constants \(c_5,A\).  Choose
\(r_1\ge30mp_K+1\) and \(r_2\ge r_1+1\), both prime.
The prime divisors of the host level are among those dividing \(2mqd\).
The displayed inequalities therefore ensure
\(r_i\nmid30N|\operatorname{disc}(K_\chi)|\).
The source theorem \texttt{gdOfHost}, with its complete proof checked,
shows that every character in \(X_h(\mathcal O_\ell,e)\) satisfies the
\texttt{GD} hypotheses for this fixed \(K_\chi\).  Its field exclusion
is proved by ramification at \(p_K\): this prime ramifies in
\(K_\chi\), while the cyclotomic and character fields occurring in the
opposing compositum are unramified there.  The class-number and character
order exclusions are exactly supplied by \(\ell\notin\mathcal B_\chi\).

Group the characters by the kernel \(H\) of their character on
\((\mathbb Z/N\mathbb Z)^\times\).  Each quotient by such a kernel is
cyclic.  Applying \texttt{CorGK} to each resulting character orbit gives
finiteness and a bound with orbit factor \(\varphi([G:H])\), where
\(G=(\mathbb Z/N\mathbb Z)^\times\).  The source's finite character
count proves
\[
 \sum_H\varphi([G:H])\le |G|=\varphi(N)\le N.
\]
Indeed the characters with a fixed kernel are the faithful characters
of that cyclic quotient, numbering \(\varphi([G:H])\), and different
kernels give disjoint sets of characters.  Over a characteristic-zero
algebraically closed coefficient field there are exactly \(|G|\)
characters of the finite abelian group \(G\).  This also proves the
inequality without assuming the individual rings \(\mathcal O_\ell\)
already contain every root of unity.

Consequently the left side of \eqref{eq:host-case-one-bound} is at most
\[
 2c_5(\log r_1+\log r_2)\,N\,[N(e+1)]^A.
\]
Choose \(k\in\mathbb N\) with
\(2c_5(\log r_1+\log r_2)\le16^k\).
Since \(N\ge16\), \(e+1\le t\), and \(t\ge1\), this is at most
\((Nt)^{k+1+A}\).  Set \(C=k+1+A\).
These are the choices and inequalities in the checked source proof.

The exact implementation locators, relative to
\path{src/Development/AutomorphicArithmetic/AutomorphicVarieties/}, are
\path{UnitaryThreeGroups/CountI.lean}, theorem
\texttt{UpperTheory.V.U3.countI};
\path{SmallClassNumbers/CorGKOf.lean}, theorem
\texttt{UpperTheory.V.SmallClassNumbers.corGK};
\path{HeightAndVanishing/GdOfHost.lean}, theorem \texttt{gdOfHost};
\path{HeightAndVanishing/Counts.lean}, definition \texttt{CountI};
and \path{HeightAndVanishing/DirichletPrimeToLPartMulLPowerPart.lean},
definitions \texttt{Host}, \texttt{Big}, \texttt{X}, and \texttt{CaseI}.
The auxiliary primes and bad set are defined in \path{StmtFields.lean}.
The complete \texttt{U3.countI} proof, including its quantifier order,
auxiliary-prime exclusions, finite orbit partition and final absorption
inequality, was independently read.  No fresh Lean compilation is
claimed.

\paragraph{K17-HOST-POINT: the full point assertion also closes.}
Further retrieval supplies the remaining inputs to
\texttt{pointHV\_of\_counts}.  Consequently \texttt{PointHV} has a
closed proof by composition of the recovered source theorems.  Its
arithmetic layer-bound premise remains part of the theorem; closing
the source proof does not set that bound to zero.

For precision, fix \(\chi\) and any function
\(j:\mathbb N\longrightarrow\mathbb N\).  There are
\(C\in\mathbb N\) and primes \(r_1,r_2\), depending only on
\(\chi,j\), such that the following holds.  Take any host \(h\) as
above, with \(r_i<q,d\), any \(t\ge2\) satisfying
\(h.\texttt{Big}(t)\), and any real number \(B\ge0\).
For each \(\ell\in\mathcal B_\chi\), assume the following layer
condition: for every cyclotomic \(\ell\)-power extension
\[
 F/\mathbb Q(\mu_{\ell m}),\qquad
 [F:\mathbb Q(\mu_{\ell m})]=\ell^{j(\ell)},
\]
all the groups \(\mathcal S_\ell(F,e)\), \(1\le e<t\), have finite
length, and
\begin{equation}\label{eq:host-layer-premise}
 \sum_{e=1}^{t-1}\operatorname{length}_{\mathbb Z_\ell}
                 \mathcal S_\ell(F,e)\le B.
\end{equation}
Here \emph{cyclotomic \(\ell\)-power extension} means
\(F=\mathbb Q(\mu_{\ell m})(\mu_{\ell^b})\) for some \(b\), as
encoded by \texttt{IsLayer}; the statement is invariant under the
number-field presentations used in the source.

Write
\[
 \mathcal S_\ell(E_h,e)^{+_e}
 =\{x\in\mathcal S_\ell(E_h,e):c_{E_h}x=(-1)^e x\},
\]
where the complex-conjugation action includes the Tate coefficients.
Then these eigenspaces have finite length for every odd prime \(\ell\)
and every \(1\le e<t\).  For every finite set \(L\) of odd primes,
\begin{equation}\label{eq:host-point-bound}
 \sum_{\ell\in L}\log\ell\sum_{e=1}^{t-1}
 \operatorname{length}_{\mathbb Z_\ell}
       \mathcal S_\ell(E_h,e)^{+_e}
 \le (Nt)^C+CB.
\end{equation}
There is no Case~I restriction on \(L\) in this assertion.  This is
the precise field-level, finite-parameter statement defined by
\texttt{PointHV}, now with its source inputs supplied.

\paragraph{Closed source inputs and proof.}
Set \(h_N=\texttt{SmallClassNumbers.litN}\), and obtain
\(h_H\) by applying
\path{published_hypotheses_from_patching_alt} to \(h_N\) and the
closed source theorem \texttt{descWeak}.  The following checked
assembly supplies every argument of \texttt{pointHV\_of\_counts}.
The directory prefix in the next list is
\path{ModularityArgument/ThresholdArguments/DenseThresholdCases/}.
\begin{enumerate}
 \item \path{HVFinalU.lean}, \texttt{pt\_s}, gives Poitou--Tate
 duality from \(h_N\); the local compact finiteness and duality
 arguments are the fields \texttt{hN.coh.lfK} and
 \texttt{hN.coh.ltK}.
 \path{KFinal.lean}, \texttt{lemmaK\_litN\_alt}, gives
 \(\exists C_K,\ \texttt{LemmaK}(C_K)\).
 \item \texttt{CountI} is supplied by the closed composition proved
 above.  \texttt{CountII} and \texttt{CountIII} are supplied by
 \path{AdjointLiftingOfCohomology.lean},
 \path{countII_and_countIII_of_printsH_of_sp4In}, using
 \path{Sp4In0.lean}, \texttt{sp4In0\_alt}, at \(h_H\).
 \item \path{DnEndsTop.lean}, \texttt{ends0\_alt}, supplies the
 reduced unitary lifting statement at \(h_H\).  Applying
 \path{LinkedUnlE.lean},
 \path{count_v_of_printed_bound_k_reduced_unitary_lift}, with
 \(h_H\) and the preceding \texttt{LemmaK} witness gives
 \texttt{CountV}.
 \item Outside that directory, \path{LevelCohomology/PropZ.lean}
 proves \texttt{X.propZ} without additional arguments.
 \texttt{countZ X.propZ} gives \texttt{CountZ}.
\end{enumerate}
The construction of \(h_H\) is in
\path{PublishedHypothesesFromPatching.lean}; the exact same closed
\(h_N,h_H\) are used by \texttt{SmallClassNumbers.corGK}.
The theorem \texttt{pointHV\_of\_counts} is in
\path{HeightAndVanishing/JointWinOfCounts.lean}.  Applying it to the
inputs above proves \texttt{PointHV}.  This is a forward proof
composition; no converse implication from the later \texttt{VWin}
theorem is being used.

For clarity, the final arithmetic combination can also be checked
directly.  The source's \texttt{sumOut} bounds the contribution from
odd primes outside \(\mathcal B_\chi\) by \((Nt)^{C_0}\).
For each \(\ell\in\mathcal B_\chi\), its \texttt{sumIn}, using
\eqref{eq:host-layer-premise}, bounds that prime's contribution by
\((Nt)^{C_\ell}+C_\ell B\).  All these bounds include the requisite
finiteness assertions.  Set
\[
 D=\sum_{\ell\in\mathcal B_\chi}C_\ell,
 \qquad b=|\mathcal B_\chi|,
 \qquad C=C_0+D+b+bD.
\]
For \(x=Nt\ge2\), the total is at most
\[
 x^{C_0}+b x^D+bDB
 \le (b+1)x^{C_0+D}+CB
 \le x^{C_0+D+b}+CB
 \le x^C+CB,
\]
because \(b+1\le2^b\le x^b\).  This proves
\eqref{eq:host-point-bound}.  The source proves these steps in
\path{HeightAndVanishing/FinalLe.lean}, \texttt{sumIn};
\path{HeightAndVanishing/SumOut.lean}, \texttt{sumOut}; and
\path{HeightAndVanishing/PointArith.lean}, \texttt{point}.
The layer condition is \texttt{LayerLe} in
\path{HeightAndVanishing/LayerLe.lean};
\texttt{Lchi} in \path{PrimeSetBoundDataOfPrimitiveChar.lean}
is the full cyclotomic field of level \(\ell m\), and
\texttt{IsLayer} is in \path{StmtFields.lean}.

Every assembly declaration just cited, and the complete proofs of
\texttt{sumIn}, \texttt{sumOut}, \texttt{countIV}, and \texttt{point},
were read for this source reconstruction.  The imported arithmetic
and automorphic theorems are used under the same stipulated source
contract as \texttt{CorGK}; this is not a fresh build of the full
Lean project.  The explicit proof composition is recorded separately
as an uncompiled review fragment.

% K17 supporting section, 2026-11-21 UTC.
% Authoring/check model: GPT-6 Astra Pro; effort: not exposed.
% Bibliography entries are supplied separately in k_groups_bibitems.tex.
% Uses the parent's theorem/sourcefact/proof environments and amsmath/amssymb.

\section{The precise consequences for \texorpdfstring{$K$}{K}-groups}
\label{sec:k-groups}

\subsection{The known comparison with Vandiver}

Let \(\ell\) be an odd prime, put
\[
 S_\ell=\operatorname{Spec}\mathbb Z[1/\ell],\qquad
 A_\ell=\operatorname{Cl}(\mathbb Z[\mu_\ell])\otimes\mathbb Z_\ell,
\]
and let \(A_\ell^{[a]}\) denote the eigenspace for the \(a\)-th power
of the Teichm\"uller character, with \(a\) taken modulo \(\ell-1\).
For a finite abelian group \(M\), write \(M\{\ell\}\) for its
\(\ell\)-primary subgroup.

\begin{sourcefact}[Known comparison: BG-KUR-15, BG-WEI-82, BG-WEI-109]
For every \(k\ge1\),
\[
 K_{4k}(\mathbb Z)\{\ell\}
 \simeq H^2(S_\ell,\mathbb Z_\ell(2k+1)),                 \tag{K1}
\]
and
\[
 K_{4k}(\mathbb Z)\{\ell\}=0
 \quad\Longleftrightarrow\quad A_\ell^{[-2k]}=0.         \tag{K2}
\]
The two groups in (K2) are cyclic simultaneously.
Consequently Vandiver's conjecture at \(\ell\) is equivalent to either
of the following conditions:
\[
\begin{split}
 &K_{4k}(\mathbb Z)\{\ell\}=0
       \quad(1\le k\le(\ell-3)/2);\\
 &K_{4k}(\mathbb Z)\{\ell\}=0
       \quad\text{for every }k\ge1.
\end{split}
\tag{K3}
\]
\end{sourcefact}

Here (K1) uses the modern integral comparison edge map, followed by
localization; it is not an unspecified factorial-normalized Chern
character. Exact locators are Weibel \cite[VI.8.2, VI.8.3,
VI.10.9]{Weibel2013} and Kurihara
\cite[Lemma~1.2, Corollary~1.5, Proposition~3.7, Remark~3.11]{Kurihara1992}.
Historically, Kurihara's Proposition~3.7 gave one direction for
\(K\)-groups; his Remark~3.11 made the converse conditional on the
then-conjectural comparison. The modern unconditional equivalence
should therefore not be attributed to an unconditional 1992
comparison theorem.

We record the short character check. Over
\(\mathbb Z[\mu_\ell,1/\ell]\) the omitted prime is principal and the
Brauer group is zero. Kummer theory identifies \(H^2(\mu_\ell)\)
with the class group modulo \(\ell\). Tensoring with
\(\mu_\ell^{\otimes2k}\), and taking invariants under the group of
order \(\ell-1\), selects the exponent \(-2k\). Thus the dimensions
modulo \(\ell\) agree. Finiteness and Nakayama give both the zero
and cyclicity assertions. This argument does not identify the full
finite groups or their orders.

The trivial eigenspace is zero: extension after norm to
\(\mathbb Q\) acts on it by \(\ell-1\), while
\(\operatorname{Cl}(\mathbb Z)=0\). The other even characters are
exactly the exponents \(-2k\) for
\(1\le k\le(\ell-3)/2\). This proves (K3), including the empty range
at \(\ell=3\). For a fixed \(k\), (K2) concerns only one component.
The periodic assertion is about \emph{vanishing}, with period
\((\ell-1)/2\) in \(k\); it is not periodicity of the full finite
groups. In particular \(k=0\), where \(K_0(\mathbb Z)=\mathbb Z\),
is excluded. As a useful coverage consequence, vanishing of
\(K_{4k}\) for \(1\le k\le R\) would imply Vandiver for
every odd \(\ell\le2R+3\) (VK-COVER).

\paragraph{Background cases that do not use K17-S17.}
All regular odd primes have zero primary part in every positive
degree \(4k\). The 2-primary part is also zero in every such degree:
specialize Rognes--Weibel \cite[Theorem~0.7, proved as
Theorem~7.11]{RognesWeibel2000} to \(\mathbb Z[1/2]\), where
\(s=1,t=u=0\), and use localization and finiteness.
Furthermore,
\[
 K_4(\mathbb Z)=0,\qquad K_8(\mathbb Z)=0.              \tag{K4}
\]
The first is Rognes \cite[Theorem~0.1, p.~268]{Rognes2000};
the second is Dutour Sikiri\'c--Elbaz-Vincent--Kupers--Martinet
\cite[Theorem~6.4]{DEKM2026}, published in the January 2026 issue.
An older table listing \(K_8\) as unknown is obsolete.
Hart--Harvey--Ong \cite[\S1, item~(2)]{HartHarveyOng2017}
verify Vandiver for every odd prime \(\ell<2^{31}\).
These facts exclude all such prime divisors of every
\(|K_{4k}(\mathbb Z)|\). Also (K2) and (K4) give vanishing whenever
\(k\equiv0,1,2\pmod{(\ell-1)/2}\); the zero residue uses the
trivial class-group component.

\subsection{The negative Selmer group is exactly the required dual}

Write
\[
 D_\ell(e)=\mathbb Q_\ell/\mathbb Z_\ell(-e),\qquad
 \mathcal S_\ell(e)=
 \ker\!\left(H^1(S_\ell,D_\ell(e))
              \longrightarrow H^1(\mathbb Q_\ell,D_\ell(e))\right).
\]
All cohomology is continuous. The notation agrees with the
rational-field negative Selmer group in K17-S17-Q.

\begin{theorem}[VK-PT and VK-LOCAL]
For odd \(\ell\) and \(k\ge1\), there is a canonical isomorphism
\[
 \bigl(K_{4k}(\mathbb Z)\{\ell\}\bigr)^\vee
       \simeq\mathcal S_\ell(2k).                     \tag{K5}
\]
With the bottom coefficient ring \(O=\mathbb Z_\ell\), the
source's group
\(\operatorname{ShaU}(\chi_\ell^{-2k};2\ell)\)
is this same group.
\end{theorem}

\begin{proof}
Put \(r=2k+1\), \(T=\mathbb Z_\ell(r)\), and
\(D=T^\vee(1)=\mathbb Q_\ell/\mathbb Z_\ell(1-r)\).
At finite level set
\(A_n=T/\ell^nT=\mathbb Z/\ell^n(r)\), with reduction
transition maps, and let
\(D_n=A_n^\vee(1)=\mathbb Z/\ell^n(1-r)\), with the
dual inclusion maps into \(D_{n+1}\).
The finite Poitou--Tate sequence and local duality
\cite[I, Theorem~4.10; II, Theorem~2.9]{MilneADT}
give the segment
\[
\begin{split}
 H^1(\mathbb Q_\ell,A_n)&\longrightarrow H^1(S_\ell,D_n)^\vee
 \longrightarrow H^2(S_\ell,A_n)\\
 &\longrightarrow H^2(\mathbb Q_\ell,A_n)
 \longrightarrow H^0(S_\ell,D_n)^\vee\longrightarrow0.
\end{split}
\tag{K6}
\]
The modified real-place terms vanish since \(\ell\) is odd.
Local duality identifies \(H^2(\mathbb Q_\ell,A_n)\) with
\(H^0(\mathbb Q_\ell,D_n)^\vee\), and the last map is dual to
restriction of invariants. The global and \(\ell\)-local
cyclotomic images are both \(\mathbb Z_\ell^\times\). Their
invariants in \(D_n\) are therefore equal, including when they are
nonzero. Hence the last map in (K6) is an isomorphism.
Exactness and duality therefore give compatible exact sequences
\[
 0\longrightarrow H^2(S_\ell,A_n)^\vee
 \longrightarrow H^1(S_\ell,D_n)
 \longrightarrow H^1(\mathbb Q_\ell,D_n).
\]
This checks the direction of the dual injection before passing
to a limit.

The finite-coefficient arithmetic cohomology groups are finite,
so their inverse systems satisfy the Mittag--Leffler condition.
In particular Jannsen's exact sequence
\cite[(3.1)]{Jannsen1988}
\[
 0\longrightarrow\varprojlim\nolimits^1 H^1(S_\ell,A_n)
 \longrightarrow H^2_{\rm cont}(S_\ell,T)
 \longrightarrow\varprojlim H^2(S_\ell,A_n)
 \longrightarrow0
\]
has zero left term. Pontryagin duality identifies the dual of
this inverse limit of finite groups with
\(\varinjlim H^2(S_\ell,A_n)^\vee\).
Filtered direct limits are exact, and continuous cohomology
with discrete torsion coefficients commutes with this filtered
union: both globally and locally,
\(\varinjlim H^1(D_n)=H^1(D)\).
Taking the direct limit of the finite exact sequences proves
\[
 H^2_{\rm cont}(S_\ell,T)^\vee
 \simeq\ker\!\left(H^1(S_\ell,D)
                   \longrightarrow H^1(\mathbb Q_\ell,D)\right).
\]
Now (K1) gives (K5). Thus neither a local \(H^2\) correction
nor an unexamined inverse-limit term remains.

It remains to compare the explicit source local conditions.
For \(q\ne\ell\), \(D_\ell(e)\) is unramified and
\[
 H^1_{\rm unr}(\mathbb Q_q,D_\ell(e))
 =D_\ell(e)/(q^{-e}-1)D_\ell(e)=0 \qquad(e>0).         \tag{K7}
\]
Indeed \(q^{-e}-1\) is a nonzero \(\ell\)-adic integer and
multiplication by it is onto the divisible group
\(\mathbb Q_\ell/\mathbb Z_\ell\).
The geometric Frobenius convention replaces it by \(q^e-1\)
and gives the same result. Thus local triviality at 2 is
automatic for a class unramified away from \(\ell\).
Local triviality at \(\ell\) implies the inertia condition there.
Conversely, the source's inertia conditions give an
unramified-away-\(\ell\) class, and its decomposition condition
at \(\ell\) makes that class lie in the displayed kernel.
The odd-primary real conditions add nothing.

Finally the recovered definitions give
\(Q_O=O\otimes_{\mathbb Z_\ell}
(\mathbb Q_\ell/\mathbb Z_\ell)\) and
\(e(O)=\operatorname{length}_O(O/\ell O)\).
The bottom ring supplied by \texttt{exists\_coef} is isomorphic
to \(\mathbb Z_\ell\), so \(Q_O=D_\ell(0)\) and \(e(O)=1\).
The scalar twist is exactly \(\chi_\ell^{-e}\).
This proves the asserted normalization and group identification.
\end{proof}

\subsection{Removing every excluded prime}

\begin{theorem}[VK-EXCEPTIONS]
If \(e>0\) is even and \(\ell\) is odd with
\((\ell-1)\mid2e\), then
\[
 K_{2e}(\mathbb Z)\{\ell\}=0.                         \tag{K8}
\]
\end{theorem}

\begin{proof}
The character \(\omega^{-e}\) has order at most two.
The trivial component is already zero. Otherwise
\(-e\equiv(\ell-1)/2\pmod{\ell-1}\).
Evenness of \(e\) forces \(\ell\equiv1\pmod4\).
Put \(L=\mathbb Q(\mu_\ell)\),
\(F=\mathbb Q(\sqrt\ell)\), and \(H=\operatorname{Gal}(L/F)\).
For \(x\in A_\ell^{[-e]}\), the norm and ideal-extension maps
satisfy
\[
 jN(x)=\sum_{h\in H}h(x)=\frac{\ell-1}{2}x.           \tag{K9}
\]

We need only the elementary estimate \(h_F<\ell\).
Put \(M=\lfloor\sqrt\ell/2\rfloor\).
Minkowski's bound \cite[Theorem~4.3, p.~70]{MilneANT}
provides an integral ideal of norm at most \(M\) in each class.
If \(a_F(n)\) is the number of ideals of norm \(n\), then
\(a_F(n)\le d(n)\): at a split prime the coefficient of
the \(a\)-th power is \(a+1\), at an inert prime it is zero or
one, and at a ramified prime it is one. Multiplicativity gives
the claim. Hence
\[
 h_F\le\sum_{n=1}^M a_F(n)
 \le\sum_{n=1}^M d(n)
 =\#\{(a,b)\in\mathbb Z_{>0}^2:ab\le M\}
 \le M^2\le\frac{\ell}{4}<\ell.                      \tag{K10}
\]
Here \(\ell\ge5\) in the nontrivial quadratic case.
Thus \(\operatorname{Cl}(F)\{\ell\}=0\), so \(N(x)=0\).
The scalar \((\ell-1)/2\) is a unit on an \(\ell\)-primary
group; (K9) gives \(x=0\). Apply (K2) with \(k=e/2\).
No integral injectivity of an ideal-extension map is assumed.
\end{proof}

\subsection{The global order and cyclicity estimates}

\begin{theorem}[VK-POLY and VK-CYCLIC-RANGE]
The recovered theorem K17-S17-Q implies absolute
\(C,A\in\mathbb N\) such that, for every \(k\ge1\),
\[
 \log|K_{4k}(\mathbb Z)|\le C(2k+1)^A.               \tag{K11}
\]
Put \(B_k=C(2k+1)^A\). For every prime \(\ell\),
\[
 v_\ell\bigl(|K_{4k}(\mathbb Z)|\bigr)
       \le\left\lfloor\frac{B_k}{\log\ell}\right\rfloor. \tag{K12}
\]
In particular \(\log\ell>B_k\) implies vanishing of the primary
part, while \(2\log\ell>B_k\) makes it either zero or
\(\mathbb Z/\ell\).
\end{theorem}

\begin{proof}
The positive even \(K\)-group is finite. Its 2-primary part is
zero by the cited theorem; the 3-primary part is zero by the
trivial-character case of (K2). All primes excluded by
\((\ell-1)\mid4k\) vanish by (K8).
Apply K17-S17-Q to the finite set of remaining prime divisors
of its order. By (K5), the normalized source lengths are
exactly
\(v_\ell(|K_{4k}(\mathbb Z)|)\). The finite weighted sum is
therefore the logarithm of the order, proving (K11).
Singleton estimates give (K12). A finite
\(\mathbb Z_\ell\)-module of length at most one is zero or
\(\mathbb Z/\ell\), proving the last assertion.
\end{proof}

For odd \(\ell\), Kurihara's cyclicity comparison also makes
\(A_\ell^{[-2k]}\) cyclic in the last range. It does not
bound that class-group component's order by \(\ell\);
the comparison controls its number of generators, not its
whole order.

\subsection{Extracting exponent seven from the source proof}

\begin{theorem}[VK-EXP7]
Keeping the source's fixed levels and using the exact
highest-weight dimensions, the bound can be sharpened to
\[
 \exists C_7\in\mathbb N\quad
 \log|K_{4k}(\mathbb Z)|\le C_7(2k)^7
 \quad\text{for every }k\ge1.                       \tag{K13}
\]
The constant has not been numerically evaluated.
\end{theorem}

\begin{proof}
We identify the representation-theoretic input and then give
the full count-to-norm argument. All automorphy, cohomology,
and representation-classification inputs remain the cited
source theorems used in K17-S17.

In \path{WeylC.lean}, indices start at zero and
\[
 \nu_j=j+1+H-m_{n-1-j},\qquad N_{\rm det}=n+H.
\]
The helper \path{ExistsScalar.lean:existsScalar} calls
\path{Scalar.lean:scalar}, which constructs the polynomial
module generated by a nonzero highest-weight vector of
weight \(\nu\). This module is precisely one copy of
the irreducible \(L_\nu\). To check this assertion, decompose
the finite polynomial representation into irreducibles.
Every nonzero component of the highest-weight vector has
highest weight \(\nu\); in the isotypic component
\(L_\nu\otimes W\), that vector is a highest-weight vector
in \(L_\nu\) tensored with one vector of \(W\). Its orbit
span is one copy of \(L_\nu\).
The implementation takes the span over all matrices.
It equals the invertible-matrix orbit span: any linear
functional vanishing on the latter gives a polynomial that
vanishes on the Zariski-dense invertible locus, hence on all
matrices. Thus this exact dimension applies to the actual
constructed module, not just an alternative informal choice.

The helper \path{ExistsTens.lean:exists_tens} gives the
required irreducible coefficient representation \(\Xi\)
after determinant twist \(-N_{\rm det}\), with
\[
 \dim\Xi\le(\dim L_\nu)^{[F:\mathbb Q]}
\]
for the totally complex fields used here. The infinitesimal
character is the same one used by the source's cohomology
theorem. The standard highest-weight classification and Weyl
dimension formula are cited background
\cite[\S\S27.2--27.3 and \S28.1]{EtingofGL2024}:
\[
 \dim L_\nu=\prod_{0\le i<j<n}
       \frac{\nu_i-\nu_j+j-i}{j-i}.                  \tag{K14}
\]

The rank-three trace-count source uses
\(m=(e-1,-1,-e-1)\), \(H=e+1\).
It gives
\[
 \nu=(2e+3,e+4,5),\qquad
 \lambda=\nu-(e+4)(1,1,1)=(e-1,0,1-e).
\]
Formula (K14) is \(e^3\), and an imaginary quadratic
field has two complex embeddings, so \(\dim\Xi_3\le e^6\).
The rank-two branch uses \(m=(e,-e)\), \(H=e\), giving
\(\nu=(2e+1,2)\) and \(\lambda=(e-1,-e)\).
Its one-embedding dimension is \(2e\); hence
\(\dim\Xi_2\le4e^2\).

Fix \(K=\mathbb Q(\sqrt{-3})\), \(N=2\).
The proof of \texttt{cardK} counts pairwise nonisomorphic
cuspidal representations by total cohomology at
\(\operatorname{KLev}(K,3,6)\).
The theorem \path{CohBoundFree.lean:cohBoundFree} therefore
gives a fixed positive integer \(B_0\) such that this total
cohomology is at most \(B_0\dim\Xi_3\).
It follows, by the unchanged
\path{CountSys.lean:exists_eigData} construction, that
the rank-three trace data have at most \(B_0e^6\) elements.

For the rational rank-two branch the level is
\(N|\operatorname{disc}K|=6\).
The source's quadratic-base-change argument uses six fixed
imaginary quadratic fields with pairwise coprime discriminants.
For concreteness one may take discriminants
\(-3,-7,-11,-19,-23,-31\).
At any prime \(p\), five of them avoid \(p\).
For each field choose a fixed positive integer \(B_i\)
bounding total cohomology at
\(\operatorname{KLev}(K_i,2,18)\) by \(B_i\dim\Xi_2\).
These seven cohomology constants are independent of
\(e,p,\ell\).

The proof of
\path{CountQuadraticOfHyp.lean:count_quadratic_of_hypotheses}
then applies unchanged. Its helper
\path{RootSum.lean:rootSum} introduces at most \(2^n\)
choices of roots of Frobenius eigenvalues and a Boolean
factor for two possible lifts. For \(n=2\) the multiplier
is \(8\). Thus, if \(I_p\) is the five-element subset
avoiding \(p\), the rational trace-data cardinality is at most
\[
 8\sum_{i\in I_p}4B_i e^2
 \le32\left(\sum_{i=1}^6B_i\right)e^2.
\]
This indexes whole eigensystems, retaining the multiplicities
needed by the norm product. Replacing them by distinct
single-prime trace values would not justify this argument.
Both trace-conversion helpers preserve indexing cardinality.
With
\[
 D=B_0+32\sum_{i=1}^6B_i,
\]
their disjoint union therefore has at most \(De^6\) elements.

Let \(p\ge7\), \(p\nmid6\), choose \(w\mid p\) in \(K\), and put
\(q=Nw=p^f\), \(1\le f\le2\). For the trivial character
orbit, the source constructs the integer
\[
 n_p(e)=\prod_{s}
       q^e\bigl(t_s-(q^{-e}+1+q^e)\bigr)
       \in\mathbb Z\setminus\{0\}.                  \tag{K15}
\]
The data are Galois stable, each scaled factor is integral,
and the purity bound \(|t_s|\le3\) makes every factor nonzero.
The exact source theorem is
\path{FrobeniusTraceNormNonzeroLogBound.lean:}
\texttt{frobenius\_trace\_norm\_nonzero\_log\_bound}.
It gives
\[
\begin{split}
 \log|n_p(e)|
 &\le De^6(2e\log q+1)\\
 &\le5De^7\log p,
\end{split}
\tag{K16}
\]
since \(q\le p^2\), \(e\ge1\), and \(\log p\ge1\).
The source's unchanged \texttt{OneChar} and orbit argument
then bounds the normalized negative-Selmer length by
\(2v_\ell(n_p(e))\), for \(\ell\ne p\) in the source's
admissible range. Its other arithmetic hypotheses are
exactly those checked by \path{ClassNumberThree.lean:gd}.

Take \(e=2k\) and \(p=7\). All excluded primary parts are
zero by (K8) and the known 2-primary theorem.
The one additional primary part \(\ell=7\) is zero because
7 is regular. Thus every prime satisfies
\[
 v_\ell(|K_{4k}(\mathbb Z)|)\le2v_\ell(n_7(2k)).
\]
Consequently
\[
 |K_{4k}(\mathbb Z)|\mid|n_7(2k)|^2,\qquad
 \log|K_{4k}(\mathbb Z)|
 \le10D\log7\,(2k)^7<20D(2k)^7.
\]
Taking \(C_7=20D\) proves (K13).
\end{proof}

This is a mathematical sharpening of the recovered proof,
not a declaration that the separate proposition
\path{HeightAndVanishing/Counts.lean:CorG} has been proved.
The exact source proof witnesses, without these improvements,
give exponent \(133\): the rank-three exponent is \(36\),
the six rank-two exponents are \(16\) each, and the norm
estimate adds one. Replacing sums of exponents by maxima
already gives \(37\). The exponent \(7\) additionally uses
the fixed levels and the exact coefficient dimensions.
No claim of a new compiled Lean theorem is made.

\paragraph{An exact small-degree computation target: VK-TRACE-GCD.}
Use (K15) also at \(p=11\). Since the 11-primary part
already vanishes,
\[
 |K_{4k}(\mathbb Z)|
 \mid \gcd\bigl(|n_7(2k)|,|n_{11}(2k)|\bigr)^2.       \tag{K17}
\]
Indeed each valuation is bounded by twice each of the two
integer valuations. Here \(q_7=7\), whereas \(q_{11}=121\).
A complete computation of these trace data and their gcd
could prove \(K_{12}=0\) at \(e=6\), or \(K_{16}=0\)
at \(e=8\), if the gcd has only already-excluded prime factors.
Neither trace computation has been completed in this run.
Formula (K17) identifies the finite calculation that would
be needed, including its normalization and multiplicities.

\subsection{An exact polynomial certificate for a new vanishing degree}

The remaining finite computation can be expressed without choosing
algebraic eigenvalues individually. Retain the entire indexed trace
multiset from the proof of VK-EXP7, including repeated values.

\begin{proposition}[VK-POLY-CERT]
Let $e=2k>0$ and let $p\ge7$ be a regular auxiliary prime. Put
$q=\operatorname{Nm}(w)$ for the chosen place of
$\mathbb Q(\sqrt{-3})$ above $p$, and define
\[
 F_{p,e}(X)=\prod_{s\in S_{p,e}}(X-q^e t_s),\qquad
 E_{p,e}=q^{2e}+q^e+1.
\]
Then $F_{p,e}$ is monic in $\mathbb Z[X]$, its value at $E_{p,e}$
is nonzero, and
\[
 |K_{2e}(\mathbb Z)|\mid |F_{p,e}(E_{p,e})|^2.
 \tag{K18a}
\]
More generally, let $P_{p,e}\in\mathbb Z[X]$ be monic and nonzero,
with $F_{p,e}\mid P_{p,e}$ in $\mathbb Q[X]$. Remove the maximal
power of $X-E_{p,e}$ from $P_{p,e}$, obtaining $P^*_{p,e}$.
Then
\[
 P^*_{p,e}(E_{p,e})\ne0,\qquad
 |K_{2e}(\mathbb Z)|\mid |P^*_{p,e}(E_{p,e})|^2.
 \tag{K18b}
\]
For two such auxiliary primes the order divides the square of the
gcd of these two nonzero evaluations. If every prime factor of this
gcd is below $2^{31}$, the already verified Vandiver range implies
$K_{2e}(\mathbb Z)=0$.
\end{proposition}

\begin{proof}
The source properties \texttt{TrData.int} and
\texttt{TrData.conj} in \path{AutomorphicVarieties/AutomorphicForms/Classifies.lean} say that every
$q^e t_s$ is an algebraic integer and that every Galois automorphism
permutes the indexed multiset. The elementary symmetric
coefficients are therefore rational algebraic integers. This gives
the asserted integral monic polynomial. Its evaluation satisfies
\[
 n_p(e)=(-1)^{|S_{p,e}|}F_{p,e}(E_{p,e}).
\]
The bound $|t_s|\le3$ and $q\ge7$ imply
$E_{p,e}>3q^e$, so every evaluated factor is nonzero.
The all-prime divisibility in VK-EXP7, using regularity to remove
the auxiliary prime itself, gives (K18a).

Monic polynomial division gives $P_{p,e}=F_{p,e}Q$ with
$Q\in\mathbb Z[X]$. Each exact division by the monic linear
polynomial $X-E_{p,e}$ also stays in $\mathbb Z[X]$.
Since $F_{p,e}(E_{p,e})\ne0$, $F_{p,e}$ has no such linear
factor. Unique factorization in $\mathbb Q[X]$ shows that it still
divides $P^*_{p,e}$; monic division makes the quotient integral.
This proves (K18b). Taking valuations at each prime proves the
gcd assertion. The final assertion uses the already cited vanishing
of every primary part at primes below $2^{31}$.
\end{proof}

Thus a proved, multiplicity-preserving inclusion of the source
trace system in a larger integral Hecke model is sufficient.
Removing the exact scalar root above can avoid a separate
Eisenstein projection. Taking the squarefree part is not permitted:
distinct representations may have the same trace at one prime,
and their multiplicities occur in the source integer.

Both $7$ and $13$ split in $\mathbb Q(\sqrt{-3})$ and are regular.
The following exact evaluation points were checked by integer
arithmetic; no corresponding Hecke polynomials were computed.
\[
\begin{array}{c|r|r|r}
K\text{-group}&e&p&E_{p,e}\\\hline
K_{12}&6&7&13\,841\,404\,851\\
K_{12}&6&13&23\,298\,089\,949\,291\\
K_{16}&8&7&33\,232\,936\,334\,403\\
K_{16}&8&13&665\,416\,609\,998\,910\,563
\end{array}
\]
The required highest weights are $(5,0,-5)$ and $(7,0,-7)$.
The open input is a complete integral Hecke model at the source's
specified levels, together with the coverage and normalization
proof. A dimension count or a spectrum computed from a
conjectural trace formula does not supply that input. The
literature and available-software checks, and this precise finite
certificate reduction, are recorded in
\path{evidence/work/k_numeric_followup.md}.

\subsection{What the Host estimates leave open}
\label{sec:host-uniformity-gate}

\paragraph{K17-HOST-GREENBERG-GATE.}
The established Step~17, Host and window estimates do not by themselves
remove the hypothesis that the one-variable plus Iwasawa module is
finite.  The following argument checks the full quantifiers and gives
explicit \emph{compatibility countermodels}: abstract modules and
artificial trace polynomials satisfying the numerical properties under
discussion.  They are not arithmetic counterexamples, actual
$K$-groups, or actual automorphic trace systems.  The conclusion
concerns the proposed inference from these estimates; it does not
exclude a different use of the stipulated arithmetic source contract.

\paragraph{The constants and the tower parameter.}
The source's \texttt{CountI} first fixes $\chi$, then chooses
$C,r_1,r_2$, and bounds the indicated sum of normalized lengths by
$(Nt)^C$ over its admissible hosts and Case~I primes.
The full \texttt{PointHV} has the stronger but still ordered form
\[
 \begin{gathered}
 \forall\chi\ \forall j:\mathbb N\longrightarrow\mathbb N\
 \exists C,r_1,r_2\ \forall h,t,B:\\
 \bigwedge_{\ell\in\mathcal B_\chi}
     \operatorname{LayerLe}(\chi,\ell,j(\ell),t,B)\\
 \Longrightarrow\quad
 \sum_{\ell}\log\ell\sum_{1\le e<t}
       \operatorname{length}\mathcal S_\ell(E_h,e)^{+_e}
       \le (Nt)^C+CB.
 \end{gathered}
 \tag{HU.1}
\]
All sums over primes in (HU.1) mean the source's arbitrary finite sets
of odd primes.  The admissibility conditions and $B\ge0$ are as in the preceding
Host theorem.  For each fixed $\ell$, $j(\ell)$ is one fixed
cyclotomic layer depth.  Replacing it by larger depths permits the
constants to change.  Moreover $B$ is an input estimate for the
finite-layer groups, and the output retains it; it is not an
independently obtained uniform bound on the tower.

More explicitly, \texttt{LayerLe} ranges over presentations of
$\mathbb Q(\mu_{\ell m})$ and cyclotomic $\ell$-power extensions of
degree $\ell^{j(\ell)}$.  The definition \texttt{IsLayer} requires
both that cyclotomic condition and that degree.  Its conclusion
\texttt{VVat} bounds the sum of strict Selmer lengths for $1\le e<t$.
The exact source locators are
\path{HeightAndVanishing/Counts.lean:CountI},
\path{HeightAndVanishing/LayerLe.lean:LayerLe,PointHV},
\path{StmtFields.lean:IsLayer}, and
\path{LengthEqAddEigen.lean:VVat}, relative to
\path{src/Development/AutomorphicArithmetic/AutomorphicVarieties/}
in the pinned source \cite{K17PinnedLeanMirror}.
The final combination is
\path{HeightAndVanishing/PointArith.lean:point}.

There is also a field-theoretic restriction on this uniformity.
For a fixed character of modulus $m$ and order $o$, the host is
\[
 E_h=F_\chi\,\mathbb Q(\mu_{2^{a+2}})^+
           \mathbb Q(\sqrt{-q})\,\mathbb Q(\sqrt{-d}),
 \qquad N=\operatorname{lcm}(2^{a+2},m)qd.
 \tag{HU.2}
\]
Every field adjoined to $F_\chi$ in (HU.2) has degree a power of
two.  Hence $[E_h:F_\chi]$ is a power of two.  Since the fields
are abelian over $\mathbb Q$, for each odd prime $p$ the restriction
map identifies the $p$-primary parts of
$\operatorname{Gal}(E_h/\mathbb Q)$ and
$\operatorname{Gal}(F_\chi/\mathbb Q)$.  Their order is the fixed
number $p^{v_p(o)}$, independent of $a,q,d$.
Consequently this host family does not contain an unbounded
$p$-cyclotomic tower.  Such a specialization changes $\chi$ and
therefore permits the constants to change.  When $p$ divides the
cyclotomic modulus, it also lies in $\mathcal B_\chi$ and its layer
estimate occurs explicitly in (HU.1).

Restriction and norm through the extension in (HU.2) multiply a
$p$-primary group by a $p$-adic unit.  Whenever the local conditions
are preserved, this transports an upper bound to a direct summand;
it does not make the bound decrease.  Indeed the host condition
$q,d>2t+1$ forces $N>4t^2$.
Even a bound polynomial in conductor with universal exponents
would allow positive Iwasawa $\lambda$: at conductor $p^{n+1}$
it is exponential in $n$, whereas the model below grows only
linearly in $n$.

\paragraph{The downstream window theorem.}
The checked endpoint has the definition
\[
 \texttt{VWin}:\quad
       \forall\chi,\ \texttt{VWinAt}(\chi,401).
\]
Its source is \path{CharacterBoxes/ExistsPageN.lean}.
Here the directory prefix is
\path{src/Development/AutomorphicArithmetic/}.
For each fixed $\chi$, each $0<\varepsilon<1/52$, and each function
$a(x)$ with
\[
 x^{2+8\varepsilon}\le 2^{a(x)}
      \le 2x^{2+8\varepsilon},
\]
it chooses $J\ge1$ and $X_0$.  For every $X\ge X_0$ and every
pair $(q,d)$ satisfying \texttt{WinPair} at
$(\varepsilon,X^{1/J})$, it then chooses $Y$ with
$X^{1/J}\le Y$ and $Y^2\le X$.  For every $Y\le x\le Y^2$
and every field satisfying \texttt{IsChiField}, its conclusion is
\[
 \sum_{\ell}\log\ell
     \sum_{1\le e<\lfloor x^{401}\rfloor}
           \operatorname{length}\mathcal S_\ell(E,e)
      \le \sqrt{\exp(x^\varepsilon)}
      =\exp(x^\varepsilon/2).
 \tag{HU.3}
\]
This includes both eigenspaces and the prime two.  The proof is
\path{CMFieldExponentialBounds/AdmissibleOfPageN.lean},
declarations \texttt{vwinAt\_of\_joint} and
\texttt{vwin\_of\_joint}; its complete argument and the
complete endpoint definitions were read and their retrieved source
blobs verified at the pinned commit \cite{K17PinnedLeanMirror}.
The intermediate definitions \texttt{ThmBab,ThmHV,JointWin} are in
\path{PrimeSetBoundDataOfPrimitiveChar.lean}.
The full proof \path{HeightAndVanishing/ThmHV.lean:thmHV}
chooses the fixed depth
$j(\ell)=\lfloor\log_\ell R_\ell\rfloor+3$ before applying
\texttt{PointHV}.

\begin{sloppypar}
The endpoint has a closed source composition under the same stipulated
contract as the preceding Host result.  Apply
\path{UpperTheory.V.v_window_of_literal_automorphy_n}
to \path{UpperTheory.V.SmallClassNumbers.litN}.
The latter is in \path{SmallClassNumbers/CorGKOf.lean}; the former
is the final theorem of \path{WindowBoundOfLiteratureLevel.lean}.
Its directory is
\path{ModularityArgument/ThresholdArguments/}\allowbreak%
\path{DenseThresholdCases/}.
For the two premises in \texttt{vwin\_of\_joint}, the intervening
\path{HVCore.lean:of_four_alt}, in that same directory,
uses two explicit providers.

First, \texttt{Win.thmBab\_of\_lemmaK} gives \texttt{ThmBab}
using Poitou--Tate, the local inputs, Ferrero--Washington and
the established \texttt{LemmaK} witness.
Second, \texttt{Odd.thmBcFree\_of} gives \texttt{ThmBcFree}
when applied to \texttt{ML.iv} with the Poitou--Tate and unit inputs.
The complete body of \texttt{of\_four\_alt} was read.
Then \path{HeightAndVanishing/JointWinOfCounts.lean:vwin_of_counts}
uses \texttt{jointWin\_of\_counts} to produce
$\texttt{JointWin}(401)$ and \texttt{thmBc\_iff} to convert
\texttt{ThmBcFree} into $\texttt{ThmBc}(401)$.
Thus the conditional premises of \texttt{vwin\_of\_joint} are
supplied by explicit source declarations; they are not omitted
assumptions.  The compatibility argument below would also apply
if (HU.3) were merely imposed as a stronger numerical test.
\end{sloppypar}

For every fixed $\varepsilon>0$, every fixed polynomial in $x$
is eventually smaller than $\exp(x^\varepsilon/2)$: taking
logarithms reduces this to $\log x=o(x^\varepsilon)$.
Thus (HU.3) permits such polynomial contributions even without
selecting favorable windows.  Letting $\varepsilon$ tend to zero
at a fixed $x$ would reverse the quantifiers: both $J$ and $X_0$,
as well as the prescribed hosts, may change.

\paragraph{An explicit nonfinite module.}
Let $p$ be odd, let $\Gamma\cong\mathbb Z_p$ have generator $\gamma$,
put $u=1+p$, and set
\[
 \Lambda=\mathbb Z_p[[\Gamma]],\qquad
 M=\Lambda/(\gamma-u).
 \tag{HU.4}
\]
With $\gamma=1+T$, this is
$\mathbb Z_p[[T]]/(T-p)\cong\mathbb Z_p$.
It is infinite, $\Lambda$-torsion and $\mathbb Z_p$-torsion free,
with $\mu=0$ and $\lambda=1$.
Give the Tate character of $\Gamma$ the value $u$ at $\gamma$.
For $e\ge1$,
\[
 M(e)_\Gamma\cong
       \mathbb Z_p/(u^{e+1}-1),\qquad
 \operatorname{length}_{\mathbb Z_p}M(e)_\Gamma
       =1+v_p(e+1).
 \tag{HU.5}
\]
For completeness, if $b=1+p^s c$ with $s\ge1$ and $c$ a unit,
the leading term of $b^p-1$ is $p^{s+1}c$ and every later
binomial term has larger valuation.  Here $p$ is odd.
Induction gives $v_p(u^{p^n}-1)=n+1$, and raising such a number
to an exponent prime to $p$ does not change that valuation.
Writing $k=p^nk_0$ proves
$v_p(u^k-1)=1+v_p(k)$ and hence (HU.5).

These finite lengths satisfy
\[
 \begin{split}
 \log p\ \operatorname{length}M(e)_\Gamma
     &\le \log p+\log(e+1),\\
 \sum_{e=1}^{t-1}\operatorname{length}M(e)_\Gamma
     &=(t-1)+v_p(t!)
       \le (t-1)+\frac{t}{p-1}<\frac{3t}{2}.
 \end{split}
 \tag{HU.6}
\]
The factorial estimate follows by summing
$\lfloor t/p^j\rfloor$ over $j\ge1$.
Taking $t=\lfloor x^{401}\rfloor$ therefore gives a contribution
$O(x^{401})$, fully compatible with (HU.3).
Selecting a single $\Delta$-character progression only decreases
the sum.  Allowing polynomially many auxiliary character components
still gives a fixed polynomial in $x$.
Nevertheless
\[
 M/(\gamma^{p^n}-1)M\cong\mathbb Z_p/p^{n+1}\mathbb Z_p,
 \tag{HU.7}
\]
whose lengths tend to infinity.  Reduction gives surjective
transition maps on these quotients, and their inverse limit is $M$.
Thus finite specialization lengths, their polynomial growth and
$\mu=0$ do not force the module to be finite.

\paragraph{Compatibility with a nonzero degree-twelve group.}
This example can respect the relevant parity and exceptional-character
exclusions.  Take the prime
\[
 p=2147483659=2^{31}+11.
\]
Let $\Delta\cong(\mathbb Z/p\mathbb Z)^\times$, write $\omega$
for its Teichm\"uller character, and give $M$ the even
$\Delta$-character $\theta=\omega^{-6}$.
Because $|\Delta|=p-1$ is a $p$-adic unit, the groups
$S_e=(M(e)_\Gamma)^\Delta$ satisfy
\[
 S_e=
 \begin{cases}
  M(e)_\Gamma,&e\equiv6\pmod{p-1},\\
  0,&\text{otherwise}.
 \end{cases}
 \tag{HU.8}
\]
In particular $S_6\cong\mathbb Z/p\mathbb Z$.
The character $\theta$ is neither trivial nor quadratic.
The nonzero positive twists start at $e=6$ and obey
\[
 \log|S_e|\le\log p+\log(e+1)<22+e\le e^7.
\]
Consequently the abstract finite groups $G_k=S_{2k}$ have
\[
 \log|G_k|\le(2k)^7\quad(k\ge1),\qquad G_3\ne0.
 \tag{HU.9}
\]
They satisfy the quantitative bound with constant one and have
prime support above $2^{31}$.  They are a compatibility model
for the numerical statements, not an identification with
$K_{4k}(\mathbb Z)$.

Primality was verified by exact trial division up to $46340$.
An additional Lucas certificate uses
$p-1=2\cdot3\cdot149\cdot2402107$ and witness $2$.
Each factor is prime by trial division,
$2^{p-1}=1\pmod p$, and the residues
$2^{(p-1)/q}\pmod p$, in the displayed factor order, are
\[
 2147483658,\quad2021977949,\quad1547536445,\quad509190093.
\]
Each residue minus one is coprime to $p$.
The exact executable and output are
\path{work/host_strength_model.py} and
\path{work/host_strength_model_results.json}.

\paragraph{Artificial trace polynomials with a nontrivial gcd.}
The two-prime divisibility certificate would become decisive if its
actual evaluated gcd were computed to be one.  Its general numerical
properties do not force this.  To see this directly, for $q=7,13$
put
\[
 E_q=q^{12}+q^6+1,\qquad
 b_q=(-E_q^2\bmod p)\in\{0,\ldots,p-1\},\qquad
 F_q(X)=X^2+b_q.
 \tag{HU.10}
\]
These polynomials are monic integral, their root multisets are
Galois stable, and each root $\alpha$ has
$|\alpha|^2=b_q<p<9q^{12}$.  Hence their normalized roots satisfy
$|\alpha/q^6|<3$.  Their evaluations are nonzero and divisible by
$p$.  The exact data are
\[
 \begin{array}{c|r|r}
 q&E_q&b_q\\ \hline
 7&13841404851&1328819029\\
 13&23298089949291&1446269942
 \end{array}
\]
and
\[
 \begin{split}
 F_7(E_7)&=191584488250615151230,\\
 F_{13}(E_{13})&=542800995285254312397672623,\\
 \gcd(F_7(E_7),F_{13}(E_{13}))
       &=15032385613=7p.
 \end{split}
 \tag{HU.11}
\]
All values and inequalities were checked with exact integers.
The degree-two index sets satisfy the counting upper bounds.
Thus monicity, integrality, Galois stability, the trace
absolute-value bound, nonvanishing, bounded cardinality and
two-prime square divisibility are compatible with a common prime
above the verified range.  These are artificial polynomials:
no claim identifies their roots with automorphic representations
or with a single global compatible system.  Such arithmetic
information is additional to the estimates just tested.

Finally, the available natural-constant bound for $K_{12}$ has
right side $C6^7\ge279936$.  To deduce vanishing from the
verified prime cutoff alone, one would need
$\log|K_{12}|<\log(2^{31})$, or a sufficiently strong
prime-by-prime replacement.  The former bound does not give the
latter.  Thus the remaining gates are actual Iwasawa-module
finiteness, a computed complete trace certificate, or the
separately stated finite-rank homology vanishings.  No new
arithmetic vanishing follows from the Host or window estimates
alone.

\subsection{A smaller integral homology gate for \texorpdfstring{$K_{12}$}{K12}}

There is an independent route to the first undecided complete vanishing
degree through integral Steinberg homology. The following localized
form of a known resolution removes the rank-eleven term. We first
record the earlier two-rank implication with its announced source
input explicit, and then give a criterion that retains the questioned
rank-eight endpoint as a hypothesis.

\begin{sourcefact}[K17-ST2-LOCAL]\label{prop:steinberg-two-local}
For every $n\ge3$, put $R_n=\mathbb Z[1/(n+1)!]$ and let
$\mathrm{St}_n=\widetilde H_{n-2}(\mathcal T_n(\mathbb Q),\mathbb Z)$.
Then
\[
 H_2(\mathrm{SL}_n(\mathbb Z);\mathrm{St}_n\otimes R_n)=0,
 \qquad
 H_2(\mathrm{GL}_n(\mathbb Z);\mathrm{St}_n\otimes R_n)=0.
\]
\end{sourcefact}

This is already included in Miller--Patzt--Putman,
Remark~1.18, p.~7 \cite{MillerPatztPutman2025}: the cited
codimension-two proof applies whenever every prime at most $n$
is invertible. Their Lemma~25.3, Claim~2, pp.~56--57, explains
that localization argument for the degree-one model; the
uniform degree-two assertion is stated in Remark~1.18, not in
Claim~2 itself. Our ring $R_n$ satisfies the stated condition.
Moreover their Corollary~E, p.~7, already proves degree-two
vanishing with arbitrary coefficients for $n\ge8$, so the
rank-$11$ case used below holds even over $\mathbb Z$.

Before locating this newer source we independently checked the
weaker $R_n$ form directly from Br\"uck--Miller--Patzt--Sroka--Wilson,
Proposition~8.5 and Lemmas~8.7 and~8.9, pp.~57--59
\cite{BrueckEtAl2024}. That complete recheck remains in the
evidence, with its actual audit events. It is not claimed as a
new result of this run and is not reproved here. The conditional
$K_{12}$ reduction that uses it is given in full next.

\begin{proposition}[K17-K12-TWO-RANKS]\label{prop:k12-two-ranks}
Put $R=\mathbb Z[1/13!]$. Assume in addition the rank-eight endpoint
$H_5(\mathrm{GL}_8;\mathrm{St}_8\otimes R)=0$
announced in DEKM's Theorem~4.5 and Table~5. If
\[
 H_4(\mathrm{GL}_9(\mathbb Z);\mathrm{St}_9\otimes R)=0,
 \qquad
 H_3(\mathrm{GL}_{10}(\mathbb Z);\mathrm{St}_{10}\otimes R)=0,
                                                               \tag{K23}
\]
then $K_{12}(\mathbb Z)=0$.
Equivalently, it suffices to establish the corresponding localized
vanishings of $H_{12}(\mathrm{Vor}_{\mathrm{GL}_9(\mathbb Z)})$
and $H_{12}(\mathrm{Vor}_{\mathrm{GL}_{10}(\mathbb Z)})$.
Neither of the two displayed hypotheses is established here.
The additional announced endpoint is not independently verified by
the posted enumeration: the dimension issue and the fully explicit
three-rank criterion appear in Proposition~\ref{prop:k12-three-ranks}.
\end{proposition}

\begin{proof}
Use Quillen's rank spectral sequence
\[
 E^1_{a,b}=H_b(\mathrm{GL}_a(\mathbb Z);\mathrm{St}_a\otimes R)
 \ \Longrightarrow\ H_{a+b}(BQ;R).
\]
The exact construction and its Voronoi comparison are recorded in
Dutour Sikiri\'c--Elbaz-Vincent--Kupers--Martinet,
Proposition~1.5 and \S6.1; their Table~5 is on author-preprint p.~17
\cite{DEKM2026}. In total degree $13$, that table, with its rank-eight endpoint retained
as the additional hypothesis, gives zero in ranks $2\le a\le8$ except for
\[
 E^1_{7,6}=R.
\]
Ranks $0$ and $1$ also give zero in this degree: the former is the
point $BQ_0$, and $\mathrm{GL}_1(\mathbb Z)=C_2$ has no positive
homology over $R$. The terms of ranks $9$ and $10$ vanish by (K23). The rank-$11$
term is zero by Source statement~\ref{prop:steinberg-two-local}.
The rank-$12$ term is $H_1(\mathrm{GL}_{12};\mathrm{St}_{12}\otimes R)$,
which vanishes by the localized Church--Putman result; the rank-$13$
term is the zero Steinberg coinvariant group of Lee--Szczarba.
Both precise forms are recorded in Kupers, Lemmas~2.1--2.2 and the
complete localization argument on pp.~2--3
\cite{Kupers2019}. Higher ranks have negative second degree and do
not occur. Thus the only nonzero $E^1$ term on this total-degree
diagonal is the single free rank-one module $R$.

The rational homology of $BQ$ in degree $13$ has rank one.
Indeed, Borel's rational $K$-group calculation and the rational
infinite-loop-space decomposition give an exterior generator in
degree $1$ and polynomial generators in degrees $6,10,14,\ldots$.
The only monomial of degree $13$ is the degree-one generator times
the square of the degree-six generator. These are the same rational
inputs used in the proof of Proposition~6.1 of
\cite{DEKM2026}; the additional degree requires just the
displayed monomial count.

It follows that no differential can have nonzero image in the
surviving rank-one term. More explicitly, inductively that term is
a submodule of $R$ until an incoming nonzero differential occurs.
Any nonzero submodule of a rank-one module over the domain $R$ has
rank one; such an incoming image would kill the only rational
contribution in total degree $13$, contradicting the computed
rational rank. Hence all incoming images are zero. Outgoing
differentials replace the term by a submodule and preserve its
torsionfreeness. The limiting diagonal has only this one term, so
$H_{13}(BQ;R)$ is torsion free.

The Hurewicz map
\[
 K_{12}(\mathbb Z)\otimes R
   =\pi_{13}(BQ)\otimes R\longrightarrow H_{13}(BQ;R)
\]
is injective. For an exact source justification, apply Arlettaz's
stable Hurewicz theorem to the $0$-connected spectrum
$\Sigma\mathbf K(\mathbb Z)$: Theorem~6.6 and Definition~6.5 of
\cite[pp.~40--42]{Arlettaz2000} bound the kernel in degree $13$
by a product of exponents of the first $12$ positive stable stems.
Only primes at most $(12+3)/2$, hence only $2,3,5,7$, divide that
product. The ordinary Hurewicz map of
$BQ=\Omega^\infty\Sigma\mathbf K(\mathbb Z)$ followed by homology
suspension is the stable Hurewicz map, so its kernel is a subgroup
of that same kernel. All these primes are units in $R$.

The finite group $K_{12}(\mathbb Z)\otimes R$ now injects into a
torsion-free module, so it is zero. The primary parts at all primes
inverted in $R$ vanish as well: the $2$-primary part is zero, and
$3,5,7,11,13$ are regular primes. Therefore $K_{12}(\mathbb Z)=0$.
Finally, Proposition~1.5 of \cite{DEKM2026} identifies
$H_q(\mathrm{GL}_n;\mathrm{St}_n)$ with
$H_{q+n-1}(\mathrm{Vor}_{\mathrm{GL}_n})$ after inverting the
indicated finite-stabilizer orders. At $(n,q)=(9,4)$ and $(10,3)$
the Voronoi degree is $12$, proving the equivalent formulation.
\end{proof}

The earlier Voronoi remark listed ranks $9,10,11$ as missing
\cite[Remark~6.6, author-preprint p.~19]{DEKM2026}.
Source statement~\ref{prop:steinberg-two-local} removes rank $11$.
The remaining rank-$9$ and rank-$10$ computations give this alternative
to the complete trace-polynomial certificate only after the announced
rank-eight endpoint is supplied. The next proposition retains it
explicitly; the still undecided group is not asserted to vanish.

% Source-index correction: independent three-rank sufficient condition.
% The earlier two-rank statement may be retained only with its announced DEKM input explicit.
% Complete self-check; independent review requested.
\subsubsection{Retaining the rank-eight endpoint}
\label{sec:k12-index-safe}

The dimension convention in the reported Voronoi computation requires
an explicit qualification.  If the minimal rank-one forms span a space
of dimension $r$, the associated projective cell has dimension $r-1$.
Consequently
\[
 H_q(\mathrm{GL}_N;\mathrm{St}_N\otimes R)
 \quad\hbox{uses projective degree }q+N-1
 \quad\hbox{and cone rank }q+N.
 \tag{KI.1}
\]
DEKM's Theorem~4.5 announces vanishing at the endpoints used in its
Table~5, but the displayed enumeration and posted data reach cone rank
$12$ at $N=8$ and cone rank $11$ at $N=9,10$
\cite[\S\S1.1--1.2, Definition~2.1, Tables~3--4,
Theorem~4.5]{DEKM2026}.
The direct orientation argument therefore reaches Steinberg degrees
$4,2,1$, respectively.  In particular it does not independently verify
$H_5(\mathrm{GL}_8;\mathrm{St}_8)=0$.
No claim that this announced vanishing is false is made here.

\begin{proposition}[K17-K12-THREE-RANKS]
\label{prop:k12-three-ranks}
Put $R=\mathbb Z[1/13!]$.  If
\[
\begin{split}
 H_5(\mathrm{GL}_8(\mathbb Z);\mathrm{St}_8\otimes R)&=0,\\
 H_4(\mathrm{GL}_9(\mathbb Z);\mathrm{St}_9\otimes R)&=0,\\
 H_3(\mathrm{GL}_{10}(\mathbb Z);\mathrm{St}_{10}\otimes R)&=0,
\end{split}                                                     \tag{KI.2}
\]
then $K_{12}(\mathbb Z)=0$.
Equivalently, it suffices to prove the three localized projective
Voronoi homology vanishings in degree $12$ at ranks $8,9,10$.
The rank-eight hypothesis is retained explicitly because its announced
source endpoint has the dimension-verification gap just described.
\end{proposition}

\begin{proof}
Use Quillen's rank spectral sequence
\[
 E^1_{a,b}
 =H_b(\mathrm{GL}_a(\mathbb Z);\mathrm{St}_a\otimes R)
 \Longrightarrow H_{a+b}(BQ;R).
\]
On total degree $13$, the independent low-rank computation gives
$E^1_{7,6}=R$ and zero in ranks $2,\ldots,6$:
the projective degree is $12$, and the original
Elbaz-Vincent--Gangl--Soul\'e Theorem~4.3 gives precisely these values
at ranks $5,6,7$; the projective complexes at ranks at most $4$
have dimension less than $12$
\cite[Theorem~4.3, p.~597; \S3.4]{ElbazGanglSoule2013}.
Ranks $0,1$ give no positive-degree contribution.
The three terms of ranks $8,9,10$ vanish by (KI.2).
The rank-$11$ degree-two term vanishes by
Miller--Patzt--Putman, Corollary~E, even integrally
\cite[Corollary~E, p.~7]{MillerPatztPutman2025}.
The rank-$12$ degree-one term and rank-$13$ degree-zero term
vanish by the localized Church--Putman and Lee--Szczarba inputs
recorded in \cite[Lemmas~2.1--2.2]{Kupers2019}.
There are no higher-rank terms on this diagonal.

Thus the only term on total degree $13$ is a single copy of $R$.
Rationally $H_{13}(BQ)$ has rank one: the only degree-$13$ monomial
in its rational infinite-loop-space homology is the degree-one
exterior generator times the square of the degree-six polynomial
generator.  No incoming differential can have nonzero image in the
rank-one term, since such an image has rank one and would kill the
only rational contribution.  Outgoing differentials only replace it
by a submodule of $R$.  Hence $H_{13}(BQ;R)$ is torsion free.

Arlettaz's stable Hurewicz bound, applied to
$\Sigma\mathbf K(\mathbb Z)$ in degree $13$, makes
$K_{12}(\mathbb Z)\otimes R\to H_{13}(BQ;R)$ injective:
the relevant first twelve stable stems have prime divisors at most
$7$, all inverted in $R$
\cite[Definition~6.5 and Theorem~6.6, pp.~40--42]{Arlettaz2000}.
The ordinary Hurewicz map followed by homology suspension is the
stable map, so its kernel is contained in the same killed kernel.
Since $K_{12}(\mathbb Z)$ is finite, its localization is zero.
Its primary parts at the inverted primes also vanish: the $2$-part
is known to vanish, and $3,5,7,11,13$ are regular.
This proves $K_{12}(\mathbb Z)=0$.  Formula (KI.1) proves the
equivalent Voronoi formulation.
\end{proof}


\subsubsection{A finite double-building sufficient condition}

The decomposition complex of Miller--Patzt--Putman gives a further
conditional form of the two-rank test. This uses homology at two
specific ranks instead of a connectivity conjecture at all ranks.

\begin{proposition}[K17-K12-DOUBLE]\label{prop:k12-double-building}
Let $R=\mathbb Z[1/13!]$, and let $\mathcal T^2(\mathbb Z^N)$
be the double Tits building of compatible flags of proper direct
summands. If
\[
 \widetilde H_{10}(\mathcal T^2(\mathbb Z^9);R)=
 \widetilde H_{11}(\mathcal T^2(\mathbb Z^9);R)=
 \widetilde H_{11}(\mathcal T^2(\mathbb Z^{10});R)=0,
                                                               \tag{DB}
\]
then
\[
 H_4(\mathrm{GL}_9;\mathrm{St}_9\otimes R)=
 H_3(\mathrm{GL}_{10};\mathrm{St}_{10}\otimes R)=0.
\]
Together with the explicit additional hypothesis
$H_5(\mathrm{GL}_8;\mathrm{St}_8\otimes R)=0$, these imply
$K_{12}(\mathbb Z)=0$ by Proposition~\ref{prop:k12-three-ranks}.
The three hypotheses in \textup{(DB)} are not proved here.
\end{proposition}

\begin{proof}
For $i\ge-1$, use the augmented decomposition complex
\[
 S_i(\mathbb Z^N)=
 \bigoplus_{V_1\oplus\cdots\oplus V_{i+2}=\mathbb Z^N}
       \mathrm{St}(V_1)\otimes\cdots\otimes\mathrm{St}(V_{i+2}),
\]
where each direct summand has positive rank. Thus $S_{-1}=\mathrm{St}_N$.
Its differential multiplies adjacent Steinberg apartment classes with
alternating signs. By \cite[Theorem~23.3, p.~49]{MillerPatztPutman2025},
\[
 H_i(S_\bullet(\mathbb Z^N))
   \cong\widetilde H_{i+N-1}(\mathcal T^2(\mathbb Z^N)).
\]
Flat localization preserves this identity.
Theorem~D$'$ of the same source, p.~64, supplies exactness at
$S_{-1},S_0,S_1$ for $N\ge4$. Hence \textup{(DB)} gives exactness
of
\[
 S_d\otimes R\longrightarrow S_{d-1}\otimes R\longrightarrow\cdots
 \longrightarrow S_0\otimes R\longrightarrow\mathrm{St}_N\otimes R
 \longrightarrow0
\]
with $(N,d)=(9,4)$ or $(10,3)$.
The first rank-$9$ hypothesis alone also gives the same exactness
for $(N,d)=(9,3)$; this will recover a needed rank-$9$ boundary
without relying on the announced DEKM endpoint.

Add the kernel of the leftmost differential as a term in degree
$d+1$. This is an exact resolution by $R$-free modules; projectivity
over the group ring is not needed. Its hyperhomology spectral sequence
converges to $H_*(\mathrm{GL}_N(\mathbb Z);\mathrm{St}_N\otimes R)$.
For $0\le p\le d$, its first page is
\[
 E^1_{p,q}=
 \bigoplus_{\substack{r_1+\cdots+r_{p+2}=N\\r_j\ge1}}
 H_q\left(\prod_j\mathrm{GL}_{r_j}(\mathbb Z);
                   \bigotimes_j(\mathrm{St}_{r_j}\otimes R)\right).
                                                               \tag{DB1}
\]
Indeed, $\mathrm{GL}_N(\mathbb Z)$ is transitive on ordered integral
decompositions with fixed ranks, their stabilizers are the displayed
products, and Shapiro's lemma applies. This is the single-rank
construction in the complete proof of
\cite[Lemma~23.5, pp.~50--51]{MillerPatztPutman2025}.
The extra kernel term is in column $d+1$ and cannot occur in total
degree $d$.

The low-rank homology data in
\cite[Theorems~4.3 and~8.1]{ElbazGanglSoule2013}, together with
the documented cone ranks through $12$ at rank $8$ in
\cite[Proposition~1.5, Tables~3--5]{DEKM2026}, localized in $R$, say
that for $1\le r\le8$ and $0\le q\le4$ the only nonzero groups
are
\[
 H_0(\mathrm{GL}_1;\mathrm{St}_1\otimes R)=R,
 \qquad
 H_3(\mathrm{GL}_3;\mathrm{St}_3\otimes R)=
 H_3(\mathrm{GL}_4;\mathrm{St}_4\otimes R)=R.
\]
Here the rank-one statement also follows directly from
$\mathrm{GL}_1(\mathbb Z)=C_2$. These ranges require only cone
rank $12$ at rank $8$, hence are within the displayed enumeration.
The additional rank-$9$ degree-three endpoint would require cone
rank $12$ there, beyond the posted rank-$11$ enumeration; it is
derived next from \textup{(DB)} instead. Every module in the
displayed low-rank ranges is free over $R$, so no Tor terms occur
in the required iterated K\"unneth formulas.

First take $(N,d)=(9,3)$. The hypothesis
$\widetilde H_{10}(\mathcal T^2(\mathbb Z^9);R)=0$ gives
exactness at $S_2$. Thus the partial resolution starts at $S_3$,
and its additional kernel lies in degree $4$. Every factor in
\textup{(DB1)} has rank at most $8$. On $p+q=3$, a nonzero
K\"unneth term could have only rank-one factors, giving
$q=0,p=3$ and total rank $5$, or one rank-three or rank-four
factor in degree three, giving $q=3,p=0$ and total rank at most
$4+1=5$. Neither can give rank $9$. Therefore
\[
 H_3(\mathrm{GL}_9;\mathrm{St}_9\otimes R)=0. \tag{DB2}
\]
The lower rank-$9$ groups vanish by the documented cone-rank
range through $11$ (and independently the degree-one and
degree-two groups vanish by Corollary~E of
\cite{MillerPatztPutman2025}).

For $(N,d)=(9,4)$, a nonzero K\"unneth contribution to
\textup{(DB1)} on $p+q=4$ could have only rank-one factors,
or one rank-three or rank-four factor in degree three and rank-one
factors otherwise. In the first case $q=0,p=4$ and the total rank
is six. In the second $q=3,p=1$ and the total rank is at most
$4+1+1=6$. Neither has rank nine.

For $(N,d)=(10,3)$ the rank-nine factors contribute nothing through
degree three, now using \textup{(DB2)}. The same list of possibilities applies. Rank-one
factors alone give $q=0,p=3$ and total rank five. A single
rank-three or rank-four factor gives $q=3,p=0$ and total rank at
most $4+1=5$. Neither has rank ten.

Thus every first-page term in the relevant total degree is zero in
both spectral sequences. They give the required degree-four and
degree-three Steinberg vanishings. Under the additional explicit
rank-eight hypothesis, Proposition~\ref{prop:k12-three-ranks}
gives the final conditional conclusion. The announced rank-eight
endpoint is not independently inferred from the posted enumeration.
\end{proof}

% K17: low-perfection-excess simpliciality and partial target-cell pruning.
% Author /root/vandiver_k; GPT-6 Astra Pro; effort not exposed.
% 2026-11-21 UTC. This section proves no full K12 vanishing.

\subsection{Simpliciality and partial orientation pruning of the remaining cells}
\label{subsec:k12-orientation}

For a well-rounded lattice $L$ of rank $N$, write $p(v)=vv^t$, let
$r$ be the dimension of the real span of its minimal-vector quadratic
forms, and choose one representative of each minimal pair.
The corresponding projective cell has dimension $r-1$.
The maximal index $\iota(L)$ is the maximum of
$[L:\sum_{j=1}^N\mathbb Z v_j]$ over independent minimal vectors.
The following deduction uses the known theorem that every perfect cone
of dimension at most nine is basic and therefore simplicial
\cite[Theorem~1.1(i), pp.~644--646]{DutourHulekSchurmann2015}.
Here the cited dimension is the dimension of the cone's linear span.

\begin{proposition}[K17-K12-CIRCUIT-SIMPLE]
\label{prop:k12-low-excess-simplicial}
If $r-N\leq4$, the minimal-vector quadratic forms are linearly
independent. In particular, every projective degree-twelve cell in
rank nine or ten has exactly thirteen minimal-vector pairs.
\end{proposition}

\begin{proof}
Put $d=r-N\geq0$. If the forms are dependent, choose an
inclusion-minimal dependent family $p(v_1),\ldots,p(v_k)$.
It has rank $k-1$, and its unique relation up to scaling has every
coefficient nonzero. Set $W=\operatorname{span}_{\mathbb R}
\{v_1,\ldots,v_k\}$ and $m=\dim W$. Separate the positive and negative
coefficients to obtain an equality of positive semidefinite sums.
The image of either sum is the span of its supporting vectors, so
both supporting sets span $W$. Each contains at least $m$ vectors;
therefore $k\geq2m$.

Let $\sigma_W$ be generated by all minimal-vector rays lying in $W$,
and let $r_W$ be its dimension. This is a face: if $P$ is orthogonal
projection onto $W^\perp$, the functional $q\mapsto\operatorname{tr}(Pq)$
is nonnegative on every generator and vanishes precisely on those
in $W$. The space $W$ is rational. Equivalently, $\sigma_W$ is the
perfect cone of $L\cap W$ with its restricted form; its minimum
equals that of $L$. The chosen circuit gives $r_W\geq k-1$.

Choose $N-m$ minimal vectors whose images form a basis of
$(L\otimes\mathbb R)/W$. Their projected quadratic forms are linearly
independent, since in suitable quotient coordinates they are the
diagonal elementary matrices. Projection annihilates the span of
$\sigma_W$, so
\[
 r\geq r_W+N-m\geq k-1+N-m\geq N+m-1.
\]
Thus $m\leq d+1$, and also
\[
 r_W\leq r-N+m=d+m\leq2d+1\leq9.
\]
The cited theorem makes $\sigma_W$ simplicial. Distinct minimal pairs
give distinct extremal rays: a positive sum of rank-one positive
semidefinite matrices has rank one only if every summand has the
same image. The circuit is therefore a dependence among distinct
extremal rays of $\sigma_W$, a contradiction.
\end{proof}

\begin{proposition}[K17-K12-ORIENT-LOWINDEX]
\label{prop:k12-low-index-orientation}
A projective degree-twelve cell has an orientation-reversing integral
stabilizer element in either of the following cases:
\begin{enumerate}
\item its lattice has rank nine and maximal index one;
\item its lattice has rank ten and maximal index at most two.
\end{enumerate}
Consequently these cells contribute no generators to the localized
Voronoi chain groups over $R=\mathbb Z[1/13!]$.
\end{proposition}

\begin{proof}
By Proposition~\ref{prop:k12-low-excess-simplicial}, there are thirteen
minimal pairs, and their quadratic forms are independent. First suppose
the maximal index is one. In a basis of minimal vectors write their
column matrix as $V=[I_N\ A]$. Every maximal minor has value $0,1$,
or $-1$, so $V$ is totally unimodular. The full integral relation
lattice is the row lattice of
\[
 B=[-A^t\ I_c],\qquad c=13-N,
\]
which is also totally unimodular. Heller's bound gives at most
$c(c+1)/2$ nonzero columns distinct up to sign
\cite[Theorem~4.2, pp.~1358--1360]{Heller1957}.
For $c=4$ or $3$ this bound is ten or six, respectively.
Thirteen columns therefore force either two zero columns or two
nonzero columns equal up to sign. After a sign change of a minimal
vector if necessary, their transposition fixes every row of $B$.
It descends to an integral involution of $L$, exchanges two minimal
rays, and fixes the others. Its determinant on their independent
quadratic span is $-1$.

It remains to treat rank ten with maximal index two. Let $L_0$ be
the lattice generated by all minima. If $L_0=L$, the gcd of the
maximal minors of $V$ is one. Since their nonzero values have
absolute value one or two, one minor has absolute value one.
Again use its columns as a basis and write $V=[I_{10}\ A]$.
The rank-three relation matrix $B=[-A^t\ I_3]$ has maximal minors
of absolute value at most two, by complementary minors. The exact
bimodular bound is
\[
 c(2,3)=\frac{3^2+3}{2}+3=9
\]
\cite[Theorem~1, p.~4]{LeePaatStallknechtXu2022}.
This counts nonzero columns distinct up to sign and permits
nonprimitive columns. Thirteen columns again force a signed-equal
pair or two zeros, producing the preceding integral ray transposition.

Suppose instead that $L_0\ne L$. Its index divides every independent
ten-vector index, hence $[L:L_0]=2$, and each independent ten minima
is a basis of $L_0$. Relative to one such basis, $V=[I_{10}\ A]$
and $B=[-A^t\ I_3]$ are totally unimodular. There are at most six
nonzero signed-column classes of $B$ and at most one zero class.
Orient the representatives so columns in each nonzero class are equal.
Distinct nonzero classes have distinct nonzero reductions in
$\mathbb F_2^3$: equal supports without a common overall sign would
give a $2\times2$ minor of absolute value two.

Write $L=L_0+\mathbb Z(w/2)$, where $w$ gives a nonzero class in
$L_0/2L_0$, and choose its coordinate lift $a\in\mathbb F_2^{13}$.
Integral surjectivity of $V$ gives
\[
 \ker(\overline V:\mathbb F_2^{13}\longrightarrow L_0/2L_0)
 =\operatorname{row}(\overline B).
\]
A permutation $P$ within the signed-column classes fixes the integral
relation lattice and defines an automorphism of $L_0$. It extends to
$L$ exactly when
\[
 Pa-a\in\operatorname{row}(\overline B).
\]
If a class has at least three columns, two corresponding entries of
$a$ agree, and their transposition fixes $a$. The same applies to a
two-column class with equal entries. Either yields the required
orientation reversal.

In the remaining case every class has size at most two and the
entries of $a$ in each doubled class are opposite. Thirteen columns
then force exactly seven classes, six nonzero and one zero, with six
doubles and one singleton. Their six nonzero labels omit a unique
nonzero $b\in\mathbb F_2^3$. Choose a linear functional $\lambda$
with $\lambda(b)=1$ and, if the singleton has a nonzero label $a_0$,
also $\lambda(a_0)=0$. This is possible because $a_0\ne b$.
A nonzero functional is one on four vectors of $\mathbb F_2^3$;
removing $b$ leaves exactly three present labels of value one,
all doubled. Let $P$ swap these three pairs. Since their glue bits
are opposite,
\[
 Pa-a=\lambda\overline B.
\]
Thus $P$ preserves the glue and extends to an integral automorphism
of $L$. It acts on the quadratic-ray basis as three transpositions,
so its determinant is $-1$.

In all cases the sum of the rays is a fixed positive radial vector;
passing to the projective cell preserves the orientation sign.
An orientation-reversing stabilizer makes the cell generator equal
to its negative in orientation coinvariants. Since two is invertible
in $R$, that generator is zero.
\end{proof}

These are partial removals of the two target chain groups. The
remaining cells have maximal index at least two in rank nine and
at least three in rank ten. Neither the proposition nor simpliciality
alone proves that the remaining homology groups vanish.

% K17: an explicit chain-level survivor. No homology or K12 nonvanishing.
\subsection{A concrete obstruction to eliminating every target cell by orientation}
\label{subsec:k12-survivor}

The preceding orientation argument has a genuine limit at maximal
index three. We give explicit data and an exact certificate for that
limit. The result concerns a chain generator; it gives no nonzero
homology class.

\begin{proposition}[K17-K12-ORIENT-SURVIVOR]
\label{prop:k12-orientation-survivor}
There is a rank-ten well-rounded lattice whose perfect cone has
thirteen independent minimal-vector rays and maximal minimum-basis
index three, and whose entire integral cell stabilizer preserves
orientation. Its projective degree-twelve cell therefore contributes
an orientation generator over $\mathbb Z[1/13!]$.
\end{proposition}

\begin{proof}
First take $L=\mathbb Z^9$ with normalized quadratic form $Q=G/36$,
where the positive integral Gram matrix is
\begingroup
\setlength{\arraycolsep}{3pt}
\[
G=\begin{pmatrix}
36&7&7&2&7&0&-9&0&-9\\
7&36&7&0&0&-11&0&-11&4\\
7&7&36&4&-7&0&0&-7&-4\\
2&0&4&36&0&0&-12&0&-6\\
7&0&-7&0&36&0&0&-11&-7\\
0&-11&0&0&0&36&4&-7&-15\\
-9&0&0&-12&0&4&36&-7&-15\\
0&-11&-7&0&-11&-7&-7&36&-11\\
-9&4&-4&-6&-7&-15&-15&-11&36
\end{pmatrix}.
\]
\endgroup
Its complete minimum set, of normalized norm one, consists of the
pairs of $e_1,\ldots,e_9$ and
\begin{align*}
a&=(-1,-1,0,-1,-1,-2,-2,-2,-3),\\
b&=(-1,0,-1,-1,-1,-1,-2,-2,-3),\\
c&=(0,-1,-1,-1,-2,-2,-2,-3,-3).
\end{align*}
The exact arithmetic verification of positivity and completeness is
given below. Put $V=[I_9\ a\ b\ c]$. Exact elimination gives rank
twelve for its twelve quadratic rays.

The integral relation lattice of $V$ has the basis
$B_0=[-A^t\ I_3]$, where $A=[a\ b\ c]$. A unimodular row change,
column reordering, and signs identify its signed columns with
\[
\begin{split}
 &e_1,\ 2e_1,\ 3e_1;\\
 &te_1+e_2 \quad(t=0,-1,-2);\\
 &te_1+e_3 \quad(t=0,-1);\\
 &te_1+e_2+e_3 \quad(t=0,-1,-2,-3).
\end{split}
\tag{*}
\]
The certificate checks this identification integrally. Its
$3\times3$ minors have maximal absolute value three and include a
unit minor. Complementary minors show that the maximal minimum-basis
index of $V$ is three.

An integral cell stabilizer permutes the primitive minimal vectors
up to sign. It preserves their integral relation lattice and therefore
induces a member of $\operatorname{GL}_3(\mathbb Z)$ preserving the
signed column set $(\ast)$. Conversely a signed-column symmetry
preserves the relation lattice and descends to the quotient lattice.
All columns are distinct up to sign, so each induced matrix determines
the signed permutation uniquely.

The pair $\{\pm3e_1\}$ is the unique pair with content three.
After multiplying by $-I$, any symmetry $U$ therefore fixes $e_1$.
Modulo $\mathbb Ze_1$, the three nonzero signed ray classes have
different multiplicities three, two, and four. Each is fixed.
The quotient action must fix the rays of $e_2,e_3,e_2+e_3$, hence is
$\epsilon I$ with $\epsilon\in\{1,-1\}$. Thus
\[
 U=\begin{pmatrix}1&x&y\\0&\epsilon&0\\0&0&\epsilon\end{pmatrix}.
\]
For $\epsilon=1$, preservation of the finite intervals forces $x=y=0$.
For $\epsilon=-1$, the intervals of sizes three and two force
$x=2,y=1$; their sum is three, which also preserves the last interval.
The full group is consequently
\[
 \{I,-I,J,-J\},\qquad
 J=\begin{pmatrix}1&2&1\\0&-1&0\\0&0&-1\end{pmatrix}.
\]
The antipodal matrix acts trivially on pairs. The matrix $J$ fixes the
three axis pairs and reverses the intervals of sizes three, two, and
four. This is respectively one, one, and two transpositions, four
in total. Every stabilizer therefore acts by an even permutation of
the independent quadratic rays. It preserves their orientation and
that of their projectivization.

Finally use $L'=L\perp\mathbb Z$ with normalized form $Q\perp[1]$,
or integral Gram matrix $G\perp[36]$. Its minimum pairs are the
twelve preceding pairs and the last coordinate vector; a vector
with both components nonzero has normalized norm at least two.
The thirteen quadratic rays are independent. The relation matrix
has the twelve columns $(\ast)$ and a unique zero column.
Every signed-column symmetry fixes that zero column and has even
permutation on the others, by the preceding classification. A sign
on the final vector fixes its quadratic ray. Hence every integral
stabilizer preserves orientation. A full minimum-vector basis must
include the last coordinate vector, so its maximal index is still
three. The cone dimension is thirteen and its projective degree is
twelve, as required.
\end{proof}

\paragraph{Exact certificate for the minimum set.}
The standalone standard-library program
\path{k12_survivor_certificate.py} uses only hard-coded integer
data, rational arithmetic, and integer square roots. It computes and
checks $Q=L_0D L_0^t$, with $L_0$ unit lower triangular and diagonal
\[
\begin{gathered}
1,\quad \frac{1247}{1296},\quad \frac{725}{774},\quad
\frac{12863}{13050},\quad \frac{337943}{373027},\quad
\frac{14603131}{16221264},\\
\frac{1271478295}{1577138148},\quad
\frac{263690369113}{411958967580},\quad
\frac{120536209775}{3164284429356}.
\end{gathered}
\]
These positive values certify positive definiteness. To enumerate
all integer vectors of norm at most one it uses
\[
 Q[z]=\sum_i D_i\left(z_i+\sum_{j>i}(L_0)_{ji}z_j\right)^2
\]
and descends from the last coordinate. If the remaining budget is
$R\geq0$ and the current rational shift is $a/b$, $b>0$, put
\[
 M=\left\lfloor\sqrt{\left\lfloor(R/D_i)b^2\right\rfloor}\right\rfloor.
\]
The complete interval of possible next integer coordinates is
\[
 \left\lceil\frac{-M-a}{b}\right\rceil
 \leq z_i\leq
 \left\lfloor\frac{M-a}{b}\right\rfloor.
\]
This follows by multiplying the next quadratic inequality by $b^2$;
its left side is an integer square. Thus every possible branch is
visited and none is lost to rounding. The executed exact search
visits 326 nodes and returns precisely the twenty-four signed
vectors stated above. It also directly verifies their norms, the
quadratic rank, the relation-matrix identification, and all maximal
relation minors. The complete JSON output accompanies the program.

The program independently checks the stabilizer list by considering
every ordered image of a fixed unimodular relation-column basis among
the twenty-four signed columns. It obtains exactly four matrices,
all with even pair permutations, agreeing with the interval argument.
The original floating-point search was used only to discover $G$;
it is not an input to this standalone exact verification.

\paragraph{Provenance and remaining problem.}
The relation configuration was found among the extremizers in
\cite[\S5.2 and Tables~1,3, p.~85]{AverkovSchymura2024}. The explicit
existence proof above does not depend on completeness of that
classification. This generator is not removed by orientation
coinvariants. Whether it is a cycle, a boundary, or part of a nonzero
homology class is a separate chain-complex question. In particular,
the example establishes neither nonvanishing of the target homology
nor nonvanishing of $K_{12}(\mathbb Z)$.

\begin{proposition}[K17-K12-SURVIVOR-CYCLE]
\label{prop:k12-survivor-cycle}
Let $\tau$ be the rank-nine projective degree-eleven cell and
$\sigma=\tau\perp\mathbb Z$ the rank-ten projective degree-twelve
cell constructed in Proposition~\ref{prop:k12-orientation-survivor}.
Their orientation generators are cycles over $R=\mathbb Z[1/13!]$.
This statement does not assert that either cycle is a nonboundary.
\end{proposition}

\begin{proof}
Number the minimum-pair representatives as
$v_1=e_1,\ldots,v_9=e_9,v_{10}=a,v_{11}=b,v_{12}=c$.
The cone on these twelve independent rays is simplicial. Its facets
are therefore obtained by deleting one ray. Every deletion still
spans the rank-nine vector space, as is checked directly from these
vectors (equivalently, every column of the rank-three relation
matrix is nonzero).

For each removed index $i$, the table gives two remaining indices
$j,k$, a sign $\epsilon$, and an integral functional $\varphi$.
\[
\begin{array}{c|cc|r|l}
i&j&k&\epsilon&\varphi(z)\\\hline
1&4&12&1&-z_1+z_4\\
2&3&5&1&z_2+z_3-z_5\\
3&2&5&1&z_2+z_3-z_5\\
4&1&5&-1&z_1-2z_4+z_5\\
5&1&7&1&z_1+z_5-z_7\\
6&2&4&-1&z_2+z_4-z_6\\
7&1&5&-1&z_1+z_5-z_7\\
8&1&9&1&z_1+z_8-z_9\\
9&1&8&-1&z_1+z_8-z_9\\
10&1&11&1&z_1\\
11&1&10&1&z_1\\
12&1&4&1&z_1-z_4
\end{array}
\]
Direct integer dot products verify
\[
\varphi(v_j)=1,\qquad \varphi(v_k)=-\epsilon,\qquad
\varphi(v_h)=0\quad(h\notin\{i,j,k\}).
\]
Put $u=\epsilon v_k-v_j$ and $g=I+u\varphi$.
Since $\varphi(u)=-2$, one has $g^2=I$, and its entries are integers.
It is therefore an integral lattice automorphism. It sends
$v_j\mapsto\epsilon v_k$ and $v_k\mapsto\epsilon v_j$ and fixes every
other representative in the facet. Thus it exchanges two independent
quadratic rays and fixes the others, reversing the facet orientation.
Every facet term is individually zero in orientation coinvariants
over $R$, because two is invertible. This proves $\partial\tau=0$.

For $\sigma$, deleting one of the first twelve rays gives the
corresponding facet of $\tau$ joined with the final coordinate ray.
The block-diagonal extension $g\perp1$ is an orientation-reversing
stabilizer of that facet. Deleting the last coordinate ray leaves
rank nine inside the rank-ten ambient lattice and hence belongs to
the boundary subcomplex of the relative Voronoi complex. All terms
of $\partial\sigma$ are therefore zero. No comparison of orbit
multiplicities is needed, since every full-rank term vanishes
separately.

The standalone exact verifier
\path{k12_survivor_boundary_certificate.py} checks all twelve
integral identities, the involutions, the stated actions, and the
full ranks using only the standard library and rational arithmetic.
It was executed successfully; its complete matrices are saved in
the accompanying JSON certificate.
\end{proof}

The next issue for these explicit cycles is their image under the
incoming Voronoi boundary maps. Their existence gives no evidence
that those images are nonzero in homology.

% Checked exact coface deduction; whole independent audits requested.
\subsection{The two explicit surviving orientation cycles are boundaries}
\label{subsec:k12-survivor-coface}

The orientation survivors above do not obstruct the vanishing of their
homology groups: each has a concrete incoming boundary with coefficient
one.  This is a calculation for the two displayed cells, not a
calculation of the entire Voronoi complexes.

\begin{proposition}[K17-K12-SURVIVOR-BOUNDARY]
\label{prop:k12-survivor-boundary}
Let $\tau$ and $\sigma$ be the cells in
Proposition~\ref{prop:k12-survivor-cycle}, and put
$R=\mathbb Z[1/13!]$. There are oriented Voronoi cells $\eta$ at
rank nine in projective degree twelve and $\eta^+$ at rank ten
in projective degree thirteen for which
\[
             \partial\eta=\tau,\qquad
             \partial\eta^+=\sigma
             \quad\hbox{over }R.                         \tag{CB.1}
\]
The two explicit cycles therefore represent zero in localized
Steinberg homology. No assertion about the remaining classes in
either homology group is made.
\end{proposition}

\begin{proof}
Retain the numbering $v_1,\ldots,v_{12}$ of the minimum-pair
representatives of $\tau$, and set
\[
 w=e_8+e_9,\qquad Q_\eta=H/720,
\]
where
\begingroup
\setlength{\arraycolsep}{2.3pt}
\[
H=\begin{pmatrix}
720&138&140&40&140&0&-180&0&-180\\
138&720&98&0&0&-220&0&-199&80\\
140&98&720&67&-140&0&0&-140&-53\\
40&0&67&720&0&0&-240&0&-93\\
140&0&-140&0&720&-48&-43&-220&-61\\
0&-220&0&0&-48&720&34&-56&-300\\
-180&0&0&-240&-43&34&720&-66&-264\\
0&-199&-140&0&-220&-56&-66&720&-360\\
-180&80&-53&-93&-61&-300&-264&-360&720
\end{pmatrix}.
\]
\endgroup
The standalone exact certificate described below verifies that
$Q_\eta$ is positive definite and that its entire minimum set is
\[
              \{\pm v_1,\ldots,\pm v_{12},\pm w\},
              \qquad \min(Q_\eta)=1.                     \tag{CB.2}
\]
The full minimum set determines a Voronoi cell by the
minimal-class/cell correspondence
\cite[\S1.2.1, p.~4; \S2.1, p.~6]{DEKM2026}.
The thirteen quadratic rays are linearly independent. Hence their
cone is a simplicial Voronoi cone of dimension thirteen and its
projectivization $\eta$ is a degree-twelve cell. Its facets are
exactly the thirteen single-ray deletions. Deleting $w$ gives
$\tau$. Deleting any $v_i$ still gives full vector rank nine.

For each $i$, use the explicit odd integral reflection $g_i$ and
functional $\varphi_i$ in the proof of
Proposition~\ref{prop:k12-survivor-cycle}. Every printed functional
vanishes on $w$. In particular the two expressions involving its
nonzero coordinates are $z_1+z_8-z_9$, which vanish at
$(z_1,z_8,z_9)=(0,1,1)$. Consequently
\[
 g_i(w)=w .
\]
The automorphism $g_i$ therefore stabilizes the facet of $\eta$
obtained by deleting $v_i$, exchanges two of its original
quadratic rays and fixes all its other rays. These rays are
independent, so it reverses the facet orientation. Every one of
these twelve facets is zero in orientation coinvariants over $R$.
The remaining facet $\tau$ has no orientation-reversing stabilizer,
by Proposition~\ref{prop:k12-orientation-survivor}.

We must also check that $\eta$ itself is an allowed orientation
generator. Its stabilizer permutes its facets and preserves the
property of having an orientation-reversing integral stabilizer,
since stabilizers of equivalent cells are conjugate. Thus it
fixes the unique facet without that property, namely $\tau$.
It consequently fixes the complementary ray $\widehat w$ and
restricts on the other twelve rays to a symmetry of $\tau$.
That symmetry has even ray permutation. The permutation on all
thirteen independent rays is therefore even as well. Every
integral stabilizer of $\eta$ preserves its orientation.

For completeness the boundary coefficient is an actual unit,
not an unspecified stabilizer quotient. In the convention of
\cite[\S1.2.1, equations~(1)--(2), p.~4]{DEKM2026}, the coefficient
of an orbit representative is the sum of the oriented incidence
signs over the actual facets in that orbit. There is exactly one
actual facet of $\eta$ which can be equivalent to $\tau$: every
other facet has an odd stabilizer, whereas $\tau$ does not.
The incidence of a facet of a simplex is $\pm1$. With the
orientation on the cone of $\eta$ given by the ordered rays
$(\widehat v_1,\ldots,\widehat v_{12},\widehat w)$ and that on
$\tau$ given by the first twelve rays, the cited convention
makes this coefficient $+1$. Equivalently, the equivariant
cellular boundary has transfer index
$[\Gamma_\eta:\Gamma_\eta\cap\Gamma_\tau]=1$, because the facet
is fixed by $\Gamma_\eta$; the subsequent inclusion into
$\Gamma_\tau$ contributes no multiplier on degree-zero
coinvariants. There is no additional target-stabilizer factor.
This proves the first equality in (CB.1).

For the second, take the orthogonal extension
$Q_\eta\perp[1]$ on $\mathbb Z^{10}$. Its minimum pairs are
the thirteen in (CB.2), embedded in the first nine coordinates,
and $\pm e_{10}$: a vector with both components nonzero has
norm at least two. Its fourteen quadratic rays are independent,
and the projective degree is thirteen. Call this cell $\eta^+$.
Deleting $\widehat w$ gives $\sigma$. Deleting any original
$\widehat v_i$ gives a full-rank facet whose odd stabilizer is
$g_i\perp1$. Deleting $\widehat e_{10}$ leaves vector rank nine
inside the rank-ten ambient lattice, so that facet belongs to
the relative boundary subcomplex and contributes zero.

The unique full-rank facet of $\eta^+$ with a nonzero orientation
generator is therefore $\sigma$. Every stabilizer of $\eta^+$
preserves it and fixes its complementary ray, so the preceding
argument also proves that $\eta^+$ is orientation preserving.
The same incidence calculation gives coefficient $\pm1$ on
$\sigma$ and no other nonzero term. Choose the orientation of
$\eta^+$ to make that coefficient $+1$. This proves (CB.1).
\end{proof}

\paragraph{Complete exact verification and its limits.}
The self-contained program
\path{k12_survivor_coface_certificate.py} contains the displayed
integer matrix, the thirteen vectors and all twelve reflection
rows as literal integer data. It uses only rational arithmetic,
integer square roots and the Python standard library.
It factors $Q_\eta=L D L^t$ exactly, checks every diagonal entry
of $D$ is positive and checks the matrix identity itself.
For enumeration it uses
\[
 Q_\eta[z]=\sum_i D_i
       \left(z_i+\sum_{j>i}L_{ji}z_j\right)^2.
\]
Descending in $i$, if the remaining budget is $B$ and the shift
is $a/b$ with $b>0$, the next coordinate is restricted by
\[
 (b z_i+a)^2\le (B/D_i)b^2.
\]
An integer square root therefore gives the exact complete
interval of possible integers $z_i$; every such integer is
visited. The executed search visits $388$ nodes and finds
exactly the twenty-six signed vectors in (CB.2), with no
other nonzero integer vector of norm at most one.
The program checks their norms directly, the rank thirteen of
their quadratic rays, and the full vector and quadratic ranks
of every other facet. It then checks the twelve integral
involutions, their transpositions of facet rays, and the
identities $\varphi_i(w)=0$. The complete rational pivots,
minimum list and witness matrices are retained in the
accompanying JSON output.

Floating-point linear programming was used only to discover a
candidate; it is not an input to this final verifier. The boundary
argument uses the formal Voronoi chain convention, not the
questioned endpoint enumeration in Theorem~4.5 of the same
source. These calculations kill the two specified cycles.
They neither establish the whole rank-nine or rank-ten target
homology vanishing nor supply the missing rank-eight
$H_5(\mathrm{GL}_8;\mathrm{St}_8\otimes R)$ input to the
$K_{12}$ criterion.

% Complete restricted classification deduction; independent audit requested.
\subsection{A complete rank-ten reduction through maximal index three}
\label{subsec:k12-trimodular-classification}

The explicit coface calculation combines with a published finite
classification to remove a whole class of degree-twelve generators.
The assumption that the minimum vectors generate the ambient lattice
is essential to the integral argument.

\begin{lemma}[K17-K12-FIXED-COFACE]
\label{lem:k12-fixed-coface}
Let $\rho$ be an orientation-preserving, full-rank simplicial
Voronoi cell, with minimum-pair representatives $v_1,\ldots,v_s$.
For each $i$, suppose $g_i\in\mathrm{GL}(L)$ preserves the other
$s-1$ quadratic rays and reverses the orientation of their span.
Suppose a primitive $w\in L$ satisfies $g_iw=\pm w$ for every $i$.
If a positive form has exactly
\[
       \{\pm v_1,\ldots,\pm v_s,\pm w\}
\]
as its minimum set, and these $s+1$ quadratic rays are independent,
then $\rho$ is a boundary with coefficient one over any coefficient
ring in which two is invertible.
\end{lemma}

\begin{proof}
The displayed minimum set defines a simplicial coface $\xi$.
Deleting $w$ gives $\rho$. Every other full-rank facet has the
odd stabilizer $g_i$, which fixes the additional quadratic ray;
a facet of deficient vector rank is zero in the relative complex.
Thus $\rho$ is its unique full-rank facet with a surviving
orientation generator. A stabilizer of $\xi$ must preserve that
facet, fix the remaining ray and act evenly on its rays.
Therefore $\xi$ is orientation preserving. In the incidence
sum over actual facets there is exactly one term equivalent to
$\rho$, with sign $\pm1$; choose the orientation of $\xi$ to
make the sign positive. This proves the assertion, with the
same no-transfer-factor justification as in
Proposition~\ref{prop:k12-survivor-boundary}.
\end{proof}

The admissibility of the enlarged minimum set is an explicit
hypothesis of this lemma. The fixed-ray equations alone do not
prove the existence of that coface. They reduce the search to
the finite union of rational subspaces
$\bigcap_i\ker(g_i-\epsilon_i I)$, with $\epsilon_i\in\{1,-1\}$,
after which a complete minimum certificate is still required.

\begin{proposition}[K17-K12-TRIMODULAR-GENERATED]
\label{prop:k12-trimodular-generated}
Let $\rho$ be a rank-ten projective degree-twelve Voronoi cell.
Assume that its minimum vectors generate its ambient lattice $L$
and that every basis chosen from its minimum vectors has lattice
index at most three. Then its orientation generator is either
zero or, up to sign and integral equivalence, the explicit
generator $\sigma$ of
Proposition~\ref{prop:k12-orientation-survivor}.
In either case it lies in the incoming boundary image over
$R=\mathbb Z[1/13!]$.
\end{proposition}

\begin{proof}
By Proposition~\ref{prop:k12-low-excess-simplicial}, the cell has
exactly thirteen minimum pairs and their quadratic rays are
independent. Choose representatives and let
\[
 V:\mathbb Z^{13}\longrightarrow L\cong\mathbb Z^{10}
\]
be their column matrix. The generation assumption makes $V$
surjective. Choose a $\mathbb Z$-basis of its rank-three kernel
and put the basis vectors in the rows of a matrix
$B\in M_{3,13}(\mathbb Z)$.

We record the normalization needed here. For every set $I$ of
ten column indices, complementary Pl\"ucker coordinates give
\[
       \det V_I=\pm\det B_{I^c}.                         \tag{TG.1}
\]
Indeed, over $\mathbb Q$ these two systems of complementary
minors are proportional, since the row space of $B$ is the
kernel of $V$. Both Pl\"ucker vectors are primitive integral
vectors. The maximal minors of $V$ have gcd one because $V$
is surjective. The row lattice of $B$ is saturated, since its
quotient is the free lattice $L$, so its maximal minors also
have gcd one. The proportionality scalar is consequently a
unit $\pm1$; the remaining signs are the complementary-index
signs. This proves (TG.1). In particular every three-by-three
minor of $B$ has absolute value at most three. No assumption
that some minimum basis initially has index one is needed.

If two nonzero columns of $B$ agree up to sign, the corresponding
signed transposition of $\mathbb Z^{13}$ fixes every row of $B$.
The same is true for the ordinary transposition of two zero
columns. It descends to an integral automorphism of
$\mathbb Z^{13}/\operatorname{row}(B)=L$, exchanges the two
minimum rays and fixes the others. Since the quadratic rays
are independent, this reverses orientation. We may therefore
assume that all nonzero columns are distinct up to sign and
there is at most one zero column.

For a full-rank integer three-row matrix whose maximal minors
have absolute value at most three, the maximal number of
distinct columns, with both signs and zero permitted, is
$25$. More precisely, the published values are
\[
 h(1,3)=13,\qquad h(2,3)=19,\qquad h(3,3)=25,
\]
and there are exactly three unimodular equivalence classes of
centrally symmetric lattice polytopes attaining the last value
\cite[\S\S5.1--5.2, Tables~1 and~3, p.~85]{AverkovSchymura2024}.
This is a published finite-classification input; that complete
classification was not rerun here.

If $B$ had no zero column, adjoining the negatives of its
thirteen columns and zero would give twenty-seven distinct
columns without increasing the maximal minor bound, a
contradiction. Hence $B$ has exactly one zero column and twelve
nonzero columns distinct up to sign. Their convex hull with
their negatives is a full-dimensional centrally symmetric
lattice polytope $P$. Determinants are multilinear, so their
maximum on $P$ is attained at vertices and remains at most
three. Every lattice point of $P$ can therefore be adjoined as
a column without increasing that bound. Since the twelve
signed pairs and zero already give twenty-five points, the
bound forces
\[
       P\cap\mathbb Z^3=\{0,\pm b_1,\ldots,\pm b_{12}\},
       \qquad \Delta(P)=3.                              \tag{TG.2}
\]
Thus $P$ is one of the three published extremal polytopes.
Any affine unimodular equivalence between these centrally
symmetric polytopes has zero translation part, since their
unique centers are zero; it is therefore a linear integral
basis change of the relation lattice.
Here nonprimitive columns are allowed, as required by the
definition of $h$.

The zero column is a rational coloop of $V$. Since the minima
generate $L$, it splits integrally: if $v$ is that column of
$V$ and $W$ is the span of the other vectors, then
\[
 L=(L\cap W)\oplus\mathbb Zv ,
\]
and the other twelve minima generate $L\cap W$, of rank nine.
The cell is therefore the integral join of its rank-nine core
with one coordinate ray. The core is determined, up to integral
equivalence and choices of signs and order, by its relation
matrix: it is the quotient of $\mathbb Z^{12}$ by that
matrix's row lattice. This reduces the claim to the three
extremal records.

Here are the complete checks for those records. We use the
order in the authors' published data file
\path{dim_3_delta_3_extremal.txt}.

\emph{First record.}
Its signed columns can be represented by
\[
B_1=\begin{pmatrix}
0&0&1&1&1&2&2&2&3&3&3&3\\
0&1&0&0&1&0&0&1&0&0&1&1\\
1&0&-1&0&0&-1&0&0&-1&0&-1&0
\end{pmatrix}.
\]
Using a unit three-column minor to solve its relations gives
the core representatives $e_1,\ldots,e_9,a_1,b_1,c_1$, where
\[
\begin{split}
a_1&=(-1,-1,-1,-2,-2,-2,-3,-2,-3),\\
b_1&=(0,-1,0,0,-1,0,0,-1,-1),\\
c_1&=(-1,-1,-2,-2,-2,-3,-3,-3,-3).
\end{split}
\]
Let $u_1,\ldots,u_6$ be the rows of
\[
\begin{pmatrix}
0&0&0&1&1&1&1&1&1\\
0&0&1&0&1&1&1&1&1\\
0&0&1&1&0&1&1&1&1\\
0&1&1&1&1&1&1&1&1\\
1&0&1&1&1&1&1&1&1\\
1&1&1&1&1&2&2&2&2
\end{pmatrix}.
\]
Direct integer multiplication gives the exact matrix identity
\[
 \sum_{j=1}^6u_ju_j^t
     =\sum_{j=1}^5e_je_j^t+c_1c_1^t.                   \tag{TG.3}
\]
None of the $u_j$ is zero or one of the twelve proposed pairs.
If precisely those twelve pairs were minima of norm one,
every $Q[u_j]$ would be strictly greater than one, while
(TG.3) makes their sum six. This is impossible. The first
record cannot be a complete core minimum set.

\emph{Second record.}
Its signed column set is the one denoted $(*)$ in
Proposition~\ref{prop:k12-orientation-survivor}, up to the
explicit integral basis change and ordering checked there.
Solving its unit-minor relations gives exactly
$e_1,\ldots,e_9,a,b,c$ from that proposition. The core is
therefore $\tau$ and its one-coloop join is $\sigma$.
Proposition~\ref{prop:k12-survivor-boundary} supplies its
incoming boundary with coefficient one.

\emph{Third record.}
Use
\[
B_3=\begin{pmatrix}
0&0&1&1&1&1&2&2&2&2&3&3\\
0&1&-1&-1&0&0&-1&-1&0&0&-1&0\\
1&0&-1&0&-1&0&-1&0&-1&0&-1&0
\end{pmatrix}.
\]
The unimodular matrix
\[
 U=\begin{pmatrix}-1&0&-3\\0&-1&1\\0&0&1\end{pmatrix}
\]
preserves its signed column set and induces the pair permutation
$(1\,11)(3\,9)(5\,7)$. It therefore preserves the relation
lattice and descends to an integral core automorphism with
an odd permutation of its twelve quadratic rays. Extending
by the identity on the coloop gives an odd stabilizer of
the rank-ten cell. Its orientation generator is zero.

These three cases exhaust (TG.2). Thus every surviving generator
in the stated class is $\sigma$, up to sign and integral
equivalence, and has the asserted boundary.
\end{proof}

\begin{corollary}
Every localized class in
$H_3(\mathrm{GL}_{10}(\mathbb Z);\mathrm{St}_{10}\otimes R)$
has a projective degree-twelve cycle representative using only
cells of maximal minimum-basis index at least four or cells
whose minimum vectors generate an index-three proper sublattice
and whose maximal minimum-basis index is three.
\end{corollary}

\begin{proof}
Subtract the explicit incoming boundaries for the
minimum-generated cells covered by the proposition. Cells of
maximal index at most two already vanish by
Proposition~\ref{prop:k12-low-index-orientation}. For a remaining
cell of maximal index three, the index of the sublattice generated
by all minima divides every minimum-basis index, and therefore
divides three. If it is proper, that index is three. No assertion
that the remaining cells vanish is used.
\end{proof}

\paragraph{Certificate and retained classification input.}
The standalone exact verifier
\path{k12_trimodular_classification_certificate.py}
checks the lattice points, maximal minors and a unit minor in
each of the three archived extremal records. It checks (TG.3),
the complete relation reconstruction in the second case, and
the displayed odd integral symmetry in the third. Its input
matrices and vectors are literal integer data, and its output
is retained as JSON. It does not recompute the published
extremal coverage: the values of $h$ and completeness of the
three-record list remain the explicitly cited primary input.
The reduction also leaves the indicated ambient-glue and
larger-index cases open; it is not a computation of the entire
rank-ten homology group or of $K_{12}(\mathbb Z)$.

% Exact restricted classification candidate; independent whole-proof audit requested.
\subsection{The twenty-six remaining ternary-glue cycles}
\label{subsec:k12-ternary-glue}

The proper-sublattice case in Proposition~\ref{prop:k12-trimodular-generated}
also admits a finite classification. The following statement identifies
actual cells and their boundaries. It does not assert that their homology
classes are nonzero.

\begin{theorem}[K17-K12-TERNARY-CYCLES]
\label{thm:k12-ternary-cycles}
Let $\rho$ be a rank-ten projective degree-twelve Voronoi cell whose
largest minimum-basis index is three and whose minimum vectors generate
a proper sublattice $L_0$ of its ambient lattice $L$.
Then either its orientation coinvariant is zero, or it is integrally
equivalent to exactly one of the twenty-six templates in
Table~\ref{tab:k12-ternary-templates} below. Each listed template is
realized by a positive form with exactly thirteen minimum pairs, has
orientation-preserving stabilizer, and is a cycle over
$R=\mathbb Z[1/13!]$.

Consequently, after subtracting the previously certified boundary
$\sigma$, every degree-twelve cycle supported on cells of maximal
minimum-basis index at most three is represented by an $R$-linear
combination of these twenty-six cycles. No linear independence or
nonvanishing of their homology classes is asserted.
\end{theorem}

\paragraph{Template notation.}
Write the following lists as columns:
\[
\begin{split}
\mathcal B_1&=((0,0,1),(0,1,0),(1,0,-1),(1,0,0),
                  (1,1,-1),(1,1,0)),\\
\mathcal B_2&=((0,0,1),(0,1,0),(0,1,1),(1,0,0),(1,0,1)),\\
\mathcal B_3&=((0,0,1),(0,1,0),(0,1,1),(1,0,0)),\\
\mathcal B_4&=((0,0,1),(0,1,0),(1,-1,-1),(1,0,0)),\\
\mathcal B_5&=((0,0,1),(0,1,0),(1,0,0)).
\end{split}                                                    \tag{TG.4}
\]
A mask $m\in\{1,\ldots,7\}$ denotes the subset
$C(m)=\{c\in\mathbb F_3:\text{bit }c\text{ of }m\text{ is one}\}$.
Thus $1,3,5,6,7$ denote $\{0\},\{0,1\},\{0,2\},\{1,2\}$
and $\mathbb F_3$, respectively.
For a mask list on $\mathcal B_i$, repeat its $j$th column once for
each $c\in C(m_j)$ and give that position color $c$.
The last table column specifies additional zero Gale columns:
$-$ means none, $0$ means one of color zero, $1$ means one of color
one, and $01$ means two of colors zero and one.
The resulting $3$-by-$13$ matrix is $B$, and the color vector is $c$.
Define
\[
 L_0=\mathbb Z^{13}/\operatorname{row}(B),\qquad
 L=L_0+\mathbb Z\frac{\bar c}{3}\ \subset L_0\otimes\mathbb Q.
                                                               \tag{TG.5}
\]
The thirteen images of the coordinate vectors give the template rays.
Every $B$ used here has a unit maximal minor, so its row lattice is
saturated and $L_0$ is free of rank ten.

\begin{table}[ht]
\centering
\begin{tabular}{ccl|ccl}
Template&Mask list&Zero colors&Template&Mask list&Zero colors\\\hline
$1.1$&$113777$&$-$&$1.11$&$131757$&$1$\\
$1.2$&$131377$&$1$&$1.12$&$131773$&$1$\\
$1.3$&$131637$&$01$&$1.13$&$131775$&$1$\\
$1.4$&$131657$&$01$&$1.14$&$131777$&$-$\\
$1.5$&$131673$&$01$&$1.15$&$133337$&$1$\\
$1.6$&$131675$&$01$&$1.16$&$133537$&$1$\\
$1.7$&$131677$&$0$&$1.17$&$133557$&$1$\\
$1.8$&$131677$&$1$&$1.18$&$133567$&$1$\\
$1.9$&$131737$&$1$&$1.19$&$133577$&$-$\\
$1.10$&$131756$&$01$&$1.20$&$173537$&$-$\\\hline
$2.1$&$31737$&$01$&$2.4$&$71733$&$01$\\
$2.2$&$33337$&$01$&$2.5$&$71737$&$0$\\
$2.3$&$71337$&$01$&$2.6$&$71737$&$1$
\end{tabular}
\caption{The surviving proper-index-three templates. The first index
specifies $\mathcal B_1$ or $\mathcal B_2$ in (TG.4).}
\label{tab:k12-ternary-templates}
\end{table}

\begin{proof}
The index $[L:L_0]$ divides every minimum-basis index and divides an
index equal to three. Since it is proper, it is three. Every independent
ten minima are consequently a basis of $L_0$. Choose one of these
bases and write the minimum matrix as $V=[I_{10}\ A]$.
Every square minor of $A$ is, up to sign, a maximal minor of $V$;
thus $V$ and its integral relation matrix $[-A^t\ I_3]$ are totally
unimodular. There are thirteen pairs by
Proposition~\ref{prop:k12-low-excess-simplicial}.

The convex hull of the nonzero signed relation columns is an
origin-symmetric lattice polytope of determinant parameter one.
The complete published classification has five such classes
\cite[\S5.2, Table~2, p.~85]{AverkovSchymura2024}.
Their exact vertex records are the signed lists (TG.4).
The integer-point checks in the accompanying verifier show that each
record contains no additional nonzero lattice points besides its
vertices. Therefore every nonzero signed relation-column class is
exactly one of the columns in one of (TG.4), and every column in that
list occurs. This use of the five-class coverage is a cited primary
input; the complete polytope classification is not claimed rerun.

Write $L=L_0+\mathbb Z(\bar c/3)$. Reduction of the exact presentation
of $L_0$ gives
\[
 \operatorname{row}(B\bmod3)\longrightarrow\mathbb F_3^{13}
      \longrightarrow L_0/3L_0\longrightarrow0.                 \tag{TG.6}
\]
Changing a representative $c$ amounts to adding a row of $B$ modulo
three; replacing its nonzero glue generator also allows multiplication
by $-1$.
Within each nonzero signed Gale class change signs so that its
columns are equal. A transposition of two positions with the same
color fixes $c$, preserves the integral relation lattice and descends
to an automorphism of $L$ exchanging two independent quadratic rays.
It reverses orientation. Thus, for a surviving cell, each such class
has at most three positions, with distinct colors.

On zero Gale columns the coordinate signs may be changed independently.
Among three such positions two colors have the same type, zero or
nonzero. A suitably signed transposition then preserves the glue and
reverses orientation. Thus there are at most two zero columns; with
two their colors must have different types. After sign changes the
possibilities are exactly those recorded above the table.

\emph{The short-support obstruction.}
For any representative of the glue choose its thirteen coefficients
$a_j$ in $\{-1,0,1\}$. Put
$x=\sum a_jv_j/3$ and $t=\#\{j:a_j\ne0\}$.
Since $x\notin L_0$, it is nonzero, and every $x-a_jv_j$ with
$a_j\ne0$ lies outside $L_0$. Normalize the minimum to one.
Completeness of the minimum set implies $Q[x-a_jv_j]>1$.
Direct expansion gives the identity
\[
 \sum_{a_j\ne0}\bigl(Q[x-a_jv_j]-1\bigr)
       =(t-6)Q[x].                                             \tag{TG.7}
\]
The left side is strictly positive. Therefore $t\ge7$.
All twenty-seven row shifts in (TG.6) are tested in the verifier.
A color state with a representative of support at most six cannot
be a complete minimum set. This is an elementary quadratic identity,
not an unverified numerical positivity test.

\emph{All possible stabilizer tests and orbit identifications.}
For a fixed list $\mathcal B=(b_1,\ldots,b_s)$, enumerate every integral
matrix $U$ preserving its signed set. Choose a unit three-column minor;
the images of those three columns must be an ordered signed triple
from the same set with determinant $\pm1$. Testing all such triples
is finite and exhaustive. Write
\[
 Ub_j=\epsilon_j b_{\pi(j)},\qquad\epsilon_j\in\{1,-1\}.
\]
For every $\lambda\in\mathbb F_3^3$ test
\[
 \epsilon_j C(m_j)
     = C(m_{\pi(j)})+\lambda\mathbin{\cdot}b_{\pi(j)}
       \quad\text{for all }j.                                 \tag{TG.8}
\]
If this holds, match each color uniquely by
$c\mapsto\epsilon_jc-\lambda\cdot b_{\pi(j)}$.
The resulting signed coordinate permutation $P$ satisfies
$BP=UB$ and $Pc-c\in\operatorname{row}(B\bmod3)$.
Its transpose preserves the relation lattice, and hence so does $P$.
It descends to $L_0$ and extends to $L$ by (TG.6).
The sign of its ray permutation is exactly its orientation sign.

Conversely, every integral stabilizer of the template permutes its
primitive minimum vectors up to signs. It induces an integral
automorphism of the relation lattice, hence one of the matrices $U$
just enumerated. Its action on the glue line has scalar $1$ or $-1$;
in the latter case compose with global negation, which fixes every
quadratic ray. Signs on zero columns can then be normalized without
changing the orientation permutation. Thus (TG.8) tests every
possible orientation sign. It also gives all equivalences of color
states, together with their row shifts. Distinct published polytope
classes cannot become equivalent in this process.

Here is the full finite enumeration, with masks ranging over
$1,\ldots,7$ and total multiplicity, including zeros, equal to thirteen:
\[
\begin{array}{c|r|r|r|r|r}
\mathcal B_i&\text{states}&\text{odd symmetry}&\text{support}\le6
       &\text{remaining states}&\text{orbits}\\\hline
1&63450&21978&24300&17172&20\\
2&2310&1410&0&900&6\\
3&14&14&0&0&0\\
4&14&14&0&0&0\\
5&0&0&0&0&0\\\hline
\text{total}&65788&23416&24300&18072&26
\end{array}                                                     \tag{TG.9}
\]
The exact verifier exhausts the stated finite loops, verifies every
matrix and color condition with integers, and computes the complete
orbits under (TG.8) and row shifts. Their representatives are precisely
Table~\ref{tab:k12-ternary-templates}. No matrix optimization is used
in this finite classification.

\emph{Exact realization of every retained template.}
The companion data file supplies an integer Gram matrix $H$ and a
positive denominator $d$ for each template, in an explicit basis of
$L$. Its change-of-basis numerator $D$ satisfies
$\det(D/3)=1/3$ and $Dv_j^{\rm ambient}=3v_j^{L_0}$; the glue class
is checked directly modulo three. Thus these are forms on the
specified ambient lattices, not forms only on $L_0$.
For each $H/d$ the verifier constructs and checks its rational
$LDL^t$ decomposition with strictly positive pivots. Descending
enumeration using that identity visits every vector of norm at most
one, with integer square-root bounds and no rounding. Each search
returns exactly the twenty-six signed prescribed vectors. Their
quadratic span has rank thirteen. The searches use between $325$
and $467$ nodes. Complementary minors relative to $L_0$ are all
$0,\pm1$, so the largest basis index in $L$ is exactly three.
These checks prove realization and completeness for all twenty-six
templates. The exhaustive test (TG.8) has no odd stabilizer for any
of them, as asserted.

\emph{Their boundary is zero.}
Each cell is simplicial, so its facets delete a single minimum ray.
A deficient-vector-rank facet is zero in the relative complex.
For a full-rank facet, its twelve minimum vectors still contain a
basis of $L_0$, since every independent ten are such a basis.
It therefore has a rank-two totally unimodular Gale matrix and
the same index-three glue. Heller's bound gives at most three
nonzero signed Gale classes
\cite[\S1, the formula $h(1,m)=m^2+m+1$]{AverkovSchymura2024}.
If it had no odd same-color transposition, each class could contain
at most three positions and the zero class at most two. This accounts
for at most $3\cdot3+2=11$ positions, fewer than twelve.
Thus every full-rank facet has an odd integral stabilizer.
Its orientation coinvariant is zero since two is a unit in $R$.
The cellular incidence sum therefore vanishes.

Finally, minimum-generated cells of maximal index at most three
are already zero or the certified boundary $\sigma$ by
Proposition~\ref{prop:k12-trimodular-generated}; maximal index at most
two vanishes by Proposition~\ref{prop:k12-low-index-orientation}.
These observations prove the final cycle reduction.
\end{proof}

\paragraph{Replayable certificate and precise boundary gate.}
The standard-library verifier
\path{k12_ternary_glue_certificate.py} uses the frozen integer
input \path{k12_ternary_glue_form_data.json} and writes its
complete check report as JSON. It checks the five archived records,
the entire finite color enumeration, all twenty-six orbit and
stabilizer conclusions, and every rational minimum certificate.
The only retained classification input is the published completeness
of the five determinant-one polytopes.
To remove this entire contribution to the $K_{12}$ chain gate it is
now sufficient to put these twenty-six explicitly specified cycles
in the incoming boundary image. Such fillings are not supplied by
the classification theorem. Cells of maximal basis index at least
four remain outside its scope.

% Exact twenty-six-coface deduction; independent whole-proof audit requested.
\subsection{Coefficient-one fillings of all twenty-six cycles}
\label{subsec:k12-ternary-fillings}

The finite gate in Theorem~\ref{thm:k12-ternary-cycles} has explicit
fillings. The proof uses actual cofaces, including a check that the
formerly deficient-rank facets stay in the relative boundary.

\begin{proposition}[K17-K12-INDEX-THREE-BOUNDARIES]
\label{prop:k12-index-three-boundaries}
For each of the twenty-six cells $\rho_i$ in
Table~\ref{tab:k12-ternary-templates}, there is a projective
degree-thirteen rank-ten Voronoi cell $\eta_i$ with
\[
                       \partial\eta_i=\rho_i
\]
over $R=\mathbb Z[1/13!]$, after choosing its orientation.
Consequently every localized class in
$H_3(\mathrm{GL}_{10}(\mathbb Z);\mathrm{St}_{10}\otimes R)$
has a degree-twelve cycle representative supported only on cells
whose largest minimum-basis index is at least four.
\end{proposition}

\begin{proof}
Fix one of the twenty-six templates and use its explicit ambient
lattice basis from the preceding certificate. Write its thirteen
minimum representatives as $v_1,\ldots,v_{13}$.
The frozen data file \path{k12_ternary_coface_data.json}
specifies a primitive vector $w$, an integral symmetric matrix $H$
and a positive denominator $d$. The exact verifier checks a positive
rational $LDL^t$ decomposition of $H/d$ and enumerates all integer
vectors of norm at most one. The result is exactly
\[
             \{\pm v_1,\ldots,\pm v_{13},\pm w\}.              \tag{TG.10}
\]
It also verifies that the fourteen quadratic rays are independent.
Thus (TG.10) gives a full-rank simplicial coface $\eta$ of $\rho$.
The twenty-six exhaustive searches require between $327$ and $477$
nodes; no finite-box cutoff or floating approximation is used in
these checks. The correspondence between complete minimum sets and
Voronoi cells is the same as in
\cite[\S1.2.1, p.~4; \S2.1, p.~6]{DEKM2026}.

For every old-ray deletion except an explicitly marked relative
facet, the data provide distinct indices $a,b$ different from the
deleted index $j$, a sign $\epsilon\in\{1,-1\}$ and an integral
linear functional $\varphi$ in the ambient lattice coordinates.
Put
\[
       g_j=I+(\epsilon v_b-v_a)\varphi.                         \tag{TG.11}
\]
The verifier checks every entry of the following identities:
\[
\begin{gathered}
g_j^2=I,\qquad g_jv_a=\epsilon v_b,\qquad
g_jv_b=\epsilon v_a,\\
g_jv_k=v_k\quad(k\notin\{j,a,b\}),\qquad g_jw=w.
\end{gathered}                                                  \tag{TG.12}
\]
In particular $g_j$ is an integral automorphism of the actual ambient
lattice. On the corresponding facet of $\eta$ it exchanges exactly
two independent quadratic rays and fixes every other ray. This
facet has an odd stabilizer and its orientation coinvariant vanishes.

The marked old-ray deletions are precisely the zero-column positions
in the original Gale matrix. For each, the verifier separately checks
that the other twelve old vectors together with $w$ have vector rank
nine. These facets remain in the relative boundary of the rank-ten
complex. This check is necessary: merely adding a ray fixed by all
the other reflections would not prevent a former coloop deletion
from acquiring full rank.

The remaining facet, obtained by deleting $w$, is the original
orientation-preserving cell $\rho$. It is therefore the unique
full-rank facet with nonzero orientation coinvariant. Any integral
stabilizer of $\eta$ must preserve this facet: vector rank and the
existence of an odd stabilizer are preserved by equivalence. It fixes
the remaining ray and has even permutation sign on $\rho$.
Consequently $\eta$ itself is orientation preserving.

The boundary coefficient is the sum of signed incidences over actual
facets in the target orbit
\cite[\S1.2.1, equations~(1)--(2), p.~4]{DEKM2026}.
Only the single facet $\rho$ can be in its orbit, and its simplicial
incidence is $\pm1$. Equivalently, the stabilizer of $\eta$ fixes
that facet, so the source transfer index is one; inclusion on
degree-zero orientation coinvariants does not multiply by the
target-stabilizer index. Choose the orientation to obtain
$\partial\eta=\rho$.

The verifier performs these checks for every one of the twenty-six
templates, without a missing or duplicated index. Minimum-generated
cells of maximal index at most three were already zero or the
boundary $\sigma$; cells of maximal index at most two were zero.
Subtract the displayed fillings with the coefficients of the
corresponding generators from any degree-twelve cycle. The resulting
cycle has no cell of maximal index at most three in its support.
This proves the homology representative statement.
\end{proof}

\paragraph{Exact replay and scope.}
The standard-library program
\path{k12_ternary_coface_certificate.py} reads the frozen
coface data and the preceding template data, verifies all rational
minimum and integral matrix identities, and writes a report for all
twenty-six fillings. The prior finite classification supplies the
orientation-preserving property of each original cell; the proof
above supplies the unique-facet and incidence-one conclusions.
The result removes the whole maximal-index-at-most-three contribution
to the rank-ten chain gate. It does not assert vanishing of the
remaining cells of larger index, of the entire Steinberg homology
group, or of $K_{12}(\mathbb Z)$.

\subsection{Inflation and a restricted rank-eight endpoint reduction}
\label{subsec:k12-coloop-rank8}

The orthogonal-sum operation used for the explicit coface is the
inflation construction of \cite[\S6.1]{ElbazGanglSoule2013}. We record
its chain-level consequence and its limits before giving a further
deduction for the rank-eight endpoint.

\begin{lemma}[Inflation, with the incidence convention fixed]
\label{lem:k12-inflation}
Let $C_d(N;R)$ be the oriented cellular coinvariants of the relative
projective Voronoi complex. Inflation defines a degree-one map
\[
 J_N:C_d(N;R)\longrightarrow C_{d+1}(N+1;R),\qquad
 J_N(\tau)=\tau\perp\mathbb Z.
\]
With the cone-span orientation convention of
\cite[\S1.2.1, equations (1)--(2)]{DEKM2026}, and the new quadratic
ray appended to an oriented basis, one has
$\partial J_N=-J_N\partial$. Thus an explicit equation
$\partial b=\tau$ gives $\partial(-J_Nb)=J_N\tau$, with coefficient
one. If two is invertible in $R$, then $J_{N+1}J_N=0$.
\end{lemma}

\begin{proof}
Normalize a form witnessing the minimum set of $\tau$ to minimum one.
Its orthogonal sum with $[1]$ has exactly the old minimum pairs and the
new coordinate pair: a vector with both components nonzero has norm
at least two. The resulting projective cell is the cone on $\tau$.
For a facet $\rho$ of $\tau$, the basis obtained by appending an old
missing ray to the oriented basis of $J_N\rho$ differs by one
transposition from the corresponding basis for $J_N\tau$.
Its incidence is therefore the negative of the old incidence.
The additional base facet has vector rank $N$ in the rank-$(N+1)$
ambient lattice and is zero in the relative complex. Old relative
boundary faces remain relative boundary faces after inflation.
This proves the differential identity, also for nonsimplicial cells.

The construction precedes passage to coinvariants. An old integral
equivalence $g$ extends to $g\perp1$, and an old orientation reverser
extends to one for the inflated cell. Thus the map descends, with no
stabilizer-index multiplier; actual equivalent facets are summed with
their ordinary incidence signs on both sides. Finally the integral
permutation of the two added coordinate axes fixes the old quadratic
span and transposes two independent new rays. It reverses orientation,
proving the double-inflation cancellation after inverting two.
\end{proof}

A rational coloop is a minimum ray whose removal lowers the vector
rank. It is \emph{integrally split} if, for its primitive representative
$v$ and the span $W$ of the other minima,
$L=(L\cap W)\oplus\mathbb Zv$. Such a cell is an inflation, even when
the original witness form has nonzero cross terms. If the minima
generate $L$, every rational coloop splits integrally. In general the
glue in an overlattice can prevent this.

Lemma~\ref{lem:k12-inflation} is not a contracting homotopy for the
subcomplex spanned by inflated cells. A cycle in the image of $J_N$
need not lift to a cycle, since its lifted boundary can lie in
$\ker J_N$. Consequently lower-rank homology vanishing alone does not
eliminate all cycles supported on such cells. The known small-rank
splitting in \cite[Theorem~6.1]{ElbazGanglSoule2013} is a separate
statement; \cite[\S6.3]{ElbazGanglSoule2013} also describes the failure
of a splitting at the next rank.

\begin{proposition}[K17-K12-EXCESS5]
\label{prop:k12-excess-five}
Let $L$ be a well-rounded lattice of rank $N\geq6$ whose perfect cone
$\sigma$ is nonsimplicial and has perfection rank $N+5$.
There is a rational subspace $W$ of dimension six such that the
minimum rays in $W$ span a nonsimplicial cone of dimension eleven.
Outside $W$ there are exactly $N-6$ minimum pairs, whose images form
a basis modulo $W$; each is a rational coloop.

If $N\geq8$ and the minima generate $L$, the cone has an integral
orientation-reversing stabilizer. In particular, a nonsimplicial
rank-eight cell of projective degree twelve with minimum-generated
ambient lattice is zero in oriented coinvariants over
$R=\mathbb Z[1/13!]$.
\end{proposition}

\begin{proof}
Choose a circuit of $k$ minimum quadratic rays and let their vector
span be $W$, of dimension $m$. Its quadratic rank is $k-1$.
The positive and negative parts of a circuit relation are equal
positive semidefinite sums. Comparing kernels shows that both
supports span $W$, so $k\geq2m$.

The cone $\sigma_W$ of all minima in $W$ is a perfect cone for
$L\cap W$ and is a face of $\sigma$: a positive semidefinite
functional with kernel $W$ exposes it. It is nonsimplicial, since
distinct rank-one positive semidefinite generators are extremal and
the selected circuit lies among them. Write $r_W=\dim\sigma_W$.
Choosing $N-m$ minima independent modulo $W$ and projecting their
quadratic rays gives
\[
 N+5\geq r_W+N-m,
 \qquad r_W\geq k-1\geq2m-1.
\]
Thus $m\leq6$ and $r_W\leq m+5\leq11$.
By \cite[Theorem~1.1(ii)]{DutourHulekSchurmann2015}, the only
nonsimplicial perfect cone of dimension at most ten is the $D_4$
cone. That exception would have $m=4$ and $r_W=10$, contradicting
$r_W\leq m+5$. Hence $r_W=11$ and $m=6$. The chosen circuit has
exactly twelve rays.

Choose minima $v_1,\ldots,v_t$ independent modulo $W$, where
$t=N-6$. Since $\dim\sigma=11+t$, their quadratic rays form a
basis modulo $\operatorname{span}(\sigma_W)$. Choose a vector-space
complement $U$ and use their images as its basis. For another
minimum vector $v\notin W$, express $vv^t$ in this quadratic basis,
modulo a form supported on $W$. If the $U$-coordinates of $v$ are
$a_1,\ldots,a_t$, the $U\times U$ block is diagonal on the right.
Therefore $a_i a_j=0$ for $i\ne j$, so exactly one coordinate,
say $a_j=a$, is nonzero. The coefficient of $v_jv_j^t$ is $a^2$.
The $W\times U$ block then shows that the $W$-component of $v$
is $a$ times that of $v_j$. Thus $v=av_j$, and equal minimum norms
give $a=\pm1$. There are precisely the stated $t$ outside pairs.

If the minima generate $L$, let $M$ be generated by those in $W$.
The complete minimum list gives
\[
 L=M\oplus\bigoplus_{j=1}^t\mathbb Zv_j,
 \qquad M=L\cap W.
\]
For $N\geq8$, interchange $v_1,v_2$ and fix $M$ and all remaining
$v_j$. This is an integral lattice automorphism. It fixes the
eleven-dimensional quadratic span of $\sigma_W$ and transposes
two independent outside rays, hence reverses the cell orientation.
The coinvariant therefore vanishes after inverting two.
\end{proof}

This proposition alone leaves simplicial rank-eight cells and nonsimplicial
cells with nontrivial ambient glue. The following theorem removes the
latter as well. Neither result asserts the rank-eight
$H_5$ vanishing, enumerate the complete degree-twelve chain group,
or compute its incoming boundary matrix. The source's low-dimensional
perfect-cone classification is invoked as cited background; its
enumeration was not rerun in this deduction.

\subsection{Eliminating every nonsimplicial rank-eight endpoint cell}
\label{subsec:k12-r8-nonsimplicial}

The preceding excess-five reduction can be completed in rank eight,
including the possible ambient glue. The additional inputs are the
classification of short perfection relations and the minimum-basis
index bounds. These are invoked as cited theorems, with the later
corrections and completion of the index bound included.

\begin{theorem}[K17-K12-R8-NONSIMPLICIAL]
\label{thm:k12-r8-nonsimplicial}
Let $L$ be a well-rounded lattice of rank eight. Suppose the cone
$\sigma$ generated by all its minimum quadratic rays has dimension
thirteen and is nonsimplicial. Then $\sigma$ has an integral
orientation-reversing stabilizer.

Consequently every nonsimplicial rank-eight projective degree-twelve
Voronoi cell is zero in oriented coinvariants over any ring in which
two is invertible. There is no assumption that the minima generate
the ambient lattice over $\mathbb Z$.
\end{theorem}

\begin{proof}
Normalize the minimum to one. We first reconstruct the equality case
of the circuit argument to specify the core to which the source
classification will apply. Choose a circuit of $k$ distinct minimum
quadratic rays and let $W$ be the span of their underlying vectors,
of dimension $m$. Its quadratic rank is $k-1$. Splitting its relation
into positive and negative parts gives equal positive semidefinite
sums. Equality of their kernels shows that each sign support spans
$W$, so $k\geq2m$.

Let $\sigma_W$ be the cone of all minimum rays in $W$, of dimension
$r_W$. It is the actual perfect cone for $L\cap W$, and it is a face
of $\sigma$, exposed by a positive semidefinite functional with
kernel $W$. It is nonsimplicial: distinct rank-one positive
semidefinite rays are extremal in a cone generated by such rays,
and the chosen circuit lies among them. Choose $8-m$ minima
independent modulo $W$. Projecting their quadratic rays gives
\[
 13\geq r_W+8-m,
 \qquad r_W\geq k-1\geq2m-1.
\]
Thus $m\leq6$ and $r_W\leq m+5\leq11$.
By \cite[Theorem~1.1(ii)]{DutourHulekSchurmann2015}, the only
nonsimplicial perfect cone of dimension at most ten is the $D_4$
cone. Its vector rank four and quadratic rank ten contradict
$r_W\leq m+5$. Hence
\[
             m=6,\qquad r_W=11,\qquad k=12.             \tag{R8.1}
\]
Each sign support therefore consists of six independent minima.

Choose minima $v,w$ independent modulo $W$. The classes of
$vv^t,ww^t$ form a basis of the quotient of the thirteen-dimensional
quadratic span by the eleven-dimensional core span. Split the vector
space as $W\oplus U$, using the images of $v,w$ as a basis of $U$.
For another outside minimum $x$, with $U$-coordinates $(a,b)$,
comparison of its $U\times U$ quadratic block with a combination
of $vv^t,ww^t$ forces $ab=0$. If $a\ne0$, the $W\times U$ block
then gives $x=av$; the other case gives $x=bw$. Equality of the
minimum norms makes the nonzero scalar $\pm1$. Thus $v,w$ are the
only two outside minimum pairs and are rational coloops.

\medskip\noindent
\emph{The integral six-dimensional core.}
Let $M_0$ be generated by the twelve circuit vectors. It has minimum
one, the two sign supports are independent six-element systems of
minima, and their union generates $M_0$. The circuit is
perf-irreducible in the sense of
\cite[Definition~2.4 and \S3]{BergeMartinet2009}: a splitting into
two perfection relations would contradict its minimality.
The hypotheses of \cite[Corollary~6.9]{BergeMartinet2009} therefore
hold. In rank six that result leaves $D_6$ or the Watson
configuration of \cite[Theorem~5.1]{BergeMartinet2009}.

The $D_6$ alternative is impossible. Its thirty minimum pairs would
also be minimum pairs of $L$, and their quadratic rays span
$\operatorname{Sym}^2(W)$ of dimension twenty-one, contradicting
(R8.1). Indeed, the roots $e_i+e_j,e_i-e_j$ span the off-diagonal
quadratic coordinates by subtraction and the diagonal coordinates
by their sums. We may therefore choose independent minima
$f_1,\ldots,f_6$ such that
\[
 e=\frac{f_1+\cdots+f_6}{3},\qquad
 M_0=\sum_{i=1}^6\mathbb Zf_i+\mathbb Ze,
 \qquad
 \text{circuit pairs }\{f_i,e-f_i:1\leq i\leq6\}.        \tag{R8.2}
\]
Set $F_0=\sum_i\mathbb Zf_i$ and $M=L\cap W$. The frame has
$[M_0:F_0]=3$. The six-dimensional minimum-basis index bound in
\cite[Theorem~5.2]{MartinetIndex2001} gives
\[
                   3[M:M_0]=[M:F_0]\leq4.
\]
Hence $M=M_0$. Saturation of the circuit lattice has been proved,
rather than assumed in the application of the short-relation theorem.

\medskip\noindent
\emph{The full minimum set of the core.}
Use the $f_i$ as rational coordinates and let $S$ be the span of
$f_if_i^t$ and $(e-f_i)(e-f_i)^t$. The circuit gives $\dim S=11$,
so this is the entire span of $\sigma_W$. If a nonzero vector $x$
has $xx^t\in S$, its off-diagonal coordinates satisfy
\[
 x_jx_k=c-d_j-d_k\quad(j\ne k),
 \qquad c=\frac13\sum_{i=1}^6d_i.                       \tag{R8.3}
\]
In fact, if the coefficient of $(e-f_i)^2$ is $b_i$, take
$d_i=b_i/3$; the other six rays contribute only to the diagonal.
Alternating four entries of (R8.3), for distinct $j,k,l,m$, gives
\[
                       (x_j-x_m)(x_k-x_l)=0.
\]
If $x_a\ne x_b$, every coordinate outside $\{a,b\}$ has a common
value. If both $x_a,x_b$ differed from this value, applying the last
identity with two outside indices would contradict it. Thus all but
at most one coordinate of $x$ coincide. Write $x=a\mathbf1+bf_i$.
Equation (R8.3) forces $d_j=d$ for $j\ne i$ and $d_i=d-ab$, hence
$c=2d-ab/3$ and
\[
                         a^2=-ab/3.
\]
If $a=0$, $x$ lies on the line of $f_i$. Otherwise $b=-3a$,
so it lies on the line of $e-f_i$. The all-equal case gives only
zero. Consequently the twelve lines in (R8.2) are the only nonzero
rank-one lines in $S$. Equal minimum norms show that their signed
representatives are the complete core minimum set.

\medskip\noindent
\emph{Core symmetries and orientation.}
Every permutation $\pi$ of the $f_i$ fixes $e$ and preserves $M$
and its twelve minimum pairs. On the twelve-dimensional free ray
space it acts by two copies of $\pi$, with determinant
$\operatorname{sgn}(\pi)^2=1$. It fixes the unique relation
\[
             \sum_i f_if_i^t-\sum_i(e-f_i)(e-f_i)^t=0.   \tag{R8.4}
\]
Therefore its determinant on the eleven-dimensional quotient $S$
is $+1$.

Define $F(f_i)=e-f_i$. Then $F(e)=e$ and $F^2=1$, so $F$ is an
integral automorphism of $M$. It exchanges the two ray systems by
six transpositions, giving determinant $+1$ on the free ray space,
and negates (R8.4). Its determinant on $S$ is consequently $-1$.
Every $\pi F$ also reverses orientation. Extending such a map by
the identity on $v,w$ keeps this sign. Swapping $v,w$ while fixing
$M$ has sign $-1$ whenever it preserves the ambient lattice.
These are cone automorphisms; no orthogonality of the original
witness form is needed.

\medskip\noindent
\emph{Bounding the ambient glue.}
The lattice generated by all minima is
\[
                  L_{\min}=M\oplus\mathbb Zv\oplus\mathbb Zw.
\]
Write $g=[L:L_{\min}]$. The independent minimum frame
$f_1,\ldots,f_6,v,w$ has index $3g$ in $L$.
The completed rank-eight bound in
\cite[appendix, pp.~6--8]{MartinetIndexErrata2014} gives index at
most nine outside $E_8$, and
\cite[Theorem~10.1]{MartinetIndex2001} shows that index nine itself
occurs only for $E_8$. The latter lattice is impossible here: its
complete minimum quadratic span has dimension thirty-six, whereas
$\dim\sigma=13$. It follows that
\[
                           3g\leq8,
                  \qquad g\in\{1,2\}.                  \tag{R8.5}
\]
If $g=1$, extend $F$ by fixing $v,w$ to obtain the desired odd
integral stabilizer.

\medskip\noindent
\emph{Every admissible binary glue.}
For $g=2$, write
\[
 L=L_{\min}+\mathbb Z\frac{u+av+bw}{2},
       \qquad u\in M,\quad a,b\in\{0,1\}.
\]
The outside bits cannot both vanish, since $L\cap W=M$.
If $a=b=1$, interchanging $v,w$ and fixing $M$ preserves the
glue class and reverses orientation.

Suppose exactly one outside bit is one, say $a=1,b=0$.
Since $[M:F_0]=3$ is odd, the six $f_i$ form a basis of $M/2M$.
Represent $u$ by a binary word of weight $k$. Modulo $2M$ one has
\[
 e=\sum_i f_i,
 \qquad F(u)=u+k\sum_i f_i.
\]
For even $k$, $F$ fixes the glue class. For $k=3$, it complements
the six bits; choose a permutation $\pi$ sending the complementary
three-element support back to the original support. Then $\pi F$
fixes the class. In both cases the extension by the identity on the
outside directions is an integral odd stabilizer of $L$.

The remaining weights $k=1,5$ are impossible. A weight-one class
is represented by a minimum $f_i$, and a weight-five class by a
minimum $e-f_i$. Hence $u\equiv m\pmod{2M}$ for some core minimum
$m$. Both $(m+v)/2$ and $(m-v)/2$ belong to $L$ and are nonzero,
while
\[
 Q\!\left[\frac{m+v}{2}\right]
 +Q\!\left[\frac{m-v}{2}\right]
       =\frac{Q[m]+Q[v]}2=1.
\]
This contradicts their having norms at least one. The other
one-bit case is identical. All possible ambient glues have now
been treated, proving the theorem.
\end{proof}

\begin{corollary}
\label{cor:k12-r8-simplicial-gate}
In the localized rank-eight projective degree-twelve Voronoi
chain group, only simplicial cells can contribute. Equivalently,
each surviving cell in that degree has exactly thirteen minimum
quadratic rays.
\end{corollary}

\paragraph{Certificate and retained inputs.}
The complete standard-library verifier
\path{k12_rank8_watson_certificate.py} checks the integral
Watson model, its frame index three and quadratic rank eleven,
the quotient orientation sign, and all $192$ binary glue classes
with nonzero outside bits. It supplies odd maps for $168$ classes
and minimum representatives proving the other $24$ impossible.
Its JSON output is retained. The no-extra-rank-one argument above
is symbolic and does not depend on finite sampling.

The short-relation and minimum-index classifications remain
explicit cited inputs. The known EVGS rank-six data provide a
separate corroboration: among $329$ codimension-ten records there
is exactly one with twelve pairs, the Watson model. This database
check is not required by the proof and does not rerun its coverage.
The remaining simplicial rank-eight cells and their boundary
matrices have not been eliminated. In particular this theorem
does not assert the rank-eight $H_5$ vanishing or
$K_{12}(\mathbb Z)=0$.

\begin{corollary}[K17-K12-R8-INDEX5]
\label{cor:k12-r8-index-five}
Every potentially contributing rank-eight projective degree-twelve
Voronoi cell has exactly thirteen minimum pairs, and its maximal
minimum-basis index is at most five. More precisely, for each minimum
frame $F=\sum_{i=1}^8\mathbb Zf_i$, the quotient $L/F$, up to signs,
permutations of the frame and changes of quotient generators, is one
of the following types. Superscripts on entries denote repetition.
\[
\begin{array}{c|l}
 L/F & \text{coefficients of generators on the frame}\\ \hline
 0 & \text{none}\\
 C_2 &(1^m,0^{8-m})/2,\quad 5\leq m\leq8\\
 C_3 &(1^m,0^{8-m})/3,\quad 7\leq m\leq8\\
 C_4 &(1^5,2^2,0)/4,\ (1^5,2^3)/4,\
       (1^6,2^2)/4,\ (1^7,2)/4\\
 C_2\times C_2 &(1,1,1,1,1,0,0,0)/2,\
                 (0,0,0,1,1,1,1,1)/2\\
 C_5 &(1^5,2^3)/5
\end{array}
\]
This is a list of necessary quotient types, including possible
ambient glue, rather than a classification of the complete cells.
\end{corollary}

\begin{proof}
Theorem~\ref{thm:k12-r8-nonsimplicial} removes the nonsimplicial
cells. A remaining cell has thirteen independent quadratic rays,
hence exactly thirteen minimum pairs.

Use the complete minimum-frame quotient classification in
\cite[\S11 and Table~11.1, author PDF pp.~45--48]{MartinetIndex2001}.
Its setup specifies that each type has a unique minimal minimum
configuration and that the column $s$ records its number of pairs.
A type supported on $n\leq8$ frame coordinates forces at least
\[
                           s+(8-n)                     \tag{R8.6}
\]
minimum pairs in the ambient rank-eight lattice: the support
lattice already contains $s$ pairs and the remaining frame vectors
supply $8-n$ further pairs.

The completed index bound, from
\cite[appendix]{MartinetIndexErrata2014} and
\cite[Theorem~10.1]{MartinetIndex2001}, excludes orders above eight
outside $E_8$; that exception has $120$ pairs. The complete quotient
list contains no order seven. The six index-six types first occur
in rank eight and have pair counts
\[
                     31,\ 27,\ 120,\ 28,\ 14,\ 36.
\]
For index eight, the primitive rank-seven type is $E_7$ with
$63$ pairs, and the six primitive rank-eight types have counts
\[
                     48,\ 75,\ 120,\ 56,\ 48,\ 32.
\]
All exceed the allowable count in (R8.6). Thus every minimum-basis
index is at most five.

For the remaining orders, apply (R8.6) to the same complete table.
Order two with support four forces $12+4=16$ pairs, leaving
supports five through eight. Order three with support six forces
$12+2=14$, leaving supports seven and eight. At order four, every
row except the four cyclic types and one elementary type displayed
above forces at least sixteen pairs. The first displayed cyclic
type has support seven, and the other four types have support
eight. At order five, the distributions $(4,4)$ and $(6,2)$ force
sixteen pairs, leaving $(5,3)$. This gives precisely the stated
necessary list.

For the signed-coordinate assertion in the order-five row, one can
use the direct calculation in \cite[\S9, case $d=5$, p.~31;
Proposition~9.1 and proof, p.~33]{MartinetIndex2001}.
After signs and a change of the single quotient generator, it lists
exactly the distributions $(6,2),(5,3),(4,4)$ and proves that the
first and last force sixteen pairs. Thus this assertion does not
rely on interpreting the broader monomial-code equivalence in
Definition~7.4 as a signed change of frame.
\end{proof}

No assertion is made that each remaining quotient type extends to
a cell with exactly thirteen rays, or that the remaining cells or
their homology classes vanish. The published quotient-type
classification is an explicit retained input.

\subsection{Excluding proper index-five glue at rank eight}
\label{subsec:k12-r8-glue-five}

\begin{proposition}[K17-K12-R8-GLUE5]\label{prop:k12-r8-glue-five}
Let \(L\) be a rank-eight lattice whose complete minimum cone is
simplicial of dimension thirteen.  If \(L_0\) is the sublattice
generated by its minimum vectors, then \([L:L_0]\ne5\).
\end{proposition}

\begin{proof}
There are exactly thirteen minimum pairs.  By
Corollary~\ref{cor:k12-r8-index-five}, every basis selected from
the minimum vectors has index at most five in \(L\).
For an index-five frame, signs, permutation of the frame and a
change of quotient generator put the numerator in the form
\[
                         (1,1,1,1,1,2,2,2)/5.              \tag{G5.1}
\]
For clarity, this exact signed-coordinate statement follows from
the direct calculation in
\cite[\S9, case \(d=5\), p.~31; Proposition~9.1 and full proof,
p.~33]{MartinetIndex2001}.
The only possible coefficient distributions are
\((6,2),(5,3),(4,4)\); the first and last force sixteen
minimum pairs.  The same count is recorded in
\cite[complement to \S7, p.~4]{MartinetIndexErrata2014}.

Suppose \([L:L_0]=5\).  Every independent eight-element minimum
frame has index in \(L\) a positive multiple of five and at most
five.  It is therefore a \(\mathbb Z\)-basis of \(L_0\).
Fix one such frame, and let \(V\) be the \(8\)-by-\(13\) matrix
of representatives of the minimum pairs in its coordinates.
The matrix contains \(I_8\), and every maximal minor is zero or
\(\pm1\).  Any smaller nonzero square minor can be padded to
a maximal minor using complementary identity columns.
Thus \(V\) is totally unimodular.
Its rational column matroid is the same as its column matroid
over \(\mathbb F_2\), since every nonzero minor is \(\pm1\).
It has rank eight and nullity five.  It has no loops or
parallel elements, since the minimum pairs are nonzero and
distinct up to sign.

Fix one nonzero generator of \(L/L_0\), and represent it by
\(u/5\), where \(u\in L_0/5L_0\).
In every minimum basis \(B\) of \(L_0\), all eight coordinates
of \(u\) are nonzero and their product is a nonsquare in
\(\mathbb F_5\).  Indeed (G5.1) has product \(2^3\), a
nonsquare.  Signs are squares modulo five, permutations
preserve the product, and a change of quotient generator
multiplies the product by \(c^8\), which is a square.
The property therefore holds for the one fixed generator
in every basis.

Let \(v\) be a column outside a chosen basis, and let
\(S=\{j:v_j\ne0\}\).  Each \(v_j\) lies in \(\{0,1,-1\}\),
by the maximal-minor calculation.  For each \(i\in S\),
replacing the \(i\)th basis vector by \(v\) gives another
minimum basis.  If \(u\) has old coordinates \(a_j\), its
new coordinates are
\[
 a'_i=a_i/v_i,\qquad
 a'_j=a_j-v_j a_i/v_i\quad(j\ne i).
                                                               \tag{G5.2}
\]
For \(j\in S\), put \(r_j=a_j/v_j\in\mathbb F_5^\times\).
Nonvanishing of the new coordinates for every pivot implies
that these residues are pairwise distinct.  Hence \(|S|\le4\).
Let \(\chi_5\) denote the quadratic character of
\(\mathbb F_5^\times\).
Comparing the nonsquare products in the two bases and
canceling the unchanged coordinates gives
\[
 \chi_5\left(\prod_{j\in S}r_j\right)
 =
 \chi_5\left(r_i\prod_{\substack{j\in S\\j\ne i}}(r_j-r_i)\right)
 \qquad(i\in S).                                             \tag{G5.3}
\]
All sign factors disappear because \(-1\) is a square
modulo five.

No subset of size two or three satisfies (G5.3).
Multiplying all residues by \(c\ne0\) multiplies both
products by \(c^{|S|}\), preserving equality or inequality
of their characters.  The two multiplicative orbits of
two-element subsets are represented by \(\{1,4\}\) and
\(\{1,2\}\).  Pivoting at \(1\), the old and new products
are respectively \(4,3\) and \(2,1\), of opposite quadratic
characters.  The four three-element subsets form one
multiplicative orbit, represented by \(\{1,2,3\}\).
At pivot \(1\), its products are \(1,2\), again of opposite
characters.

Support one would make \(v\) a duplicate minimum pair;
support zero is impossible.  Hence every nonbasis column,
relative to every minimum basis, has support exactly four.
Every fundamental circuit consequently has five elements.
Any circuit is fundamental for a basis obtained by extending
that circuit with one element removed.  Therefore every
circuit of this binary matroid has size five.

Two distinct circuits cannot intersect.  Their symmetric
difference would be a nonzero binary cycle of positive even
size at most eight.  Every binary cycle is a disjoint union
of circuits: choose a minimal nonempty cycle in its support,
subtract it over \(\mathbb F_2\), and repeat.
If every circuit has size five, a nonzero cycle of even size
must contain at least two disjoint circuits and have at
least ten elements.  This is a contradiction.

All circuits are therefore pairwise disjoint.  Their
characteristic vectors are linearly independent by
disjointness, and span the binary kernel by the same
cycle decomposition.  Nullity five would require five
disjoint five-element circuits, hence at least twenty-five
columns.  There are only thirteen.  The contradiction
proves the proposition.
\end{proof}

The conclusion excludes proper ambient index five for the
remaining simplicial rank-eight cells.  It leaves
minimum-generated cells of maximal minimum-basis index five
and proper glue of the other allowed orders.  It does not
assert vanishing of their chains or homology.

\subsection{The full-domain Soul\'e index}

\begin{theorem}[VK-SOULE and VK-SOULE-NORM]
Let \(\ell\) be odd and \(m\ge3\) odd, and put
\[
 N_\ell=\operatorname{Gal}
   (\overline{\mathbb Q}/\mathbb Q(\mu_{\ell^\infty})).
\]
For the standard Soul\'e character \(\kappa_m\) on this
full cyclotomic kernel,
\[
 [\,\mathbb Z_\ell(m):\kappa_m(N_\ell)\,]
   =|H^2(S_\ell,\mathbb Z_\ell(m))|
   =|K_{2m-2}(\mathbb Z)\{\ell\}|.                   \tag{K18}
\]
\end{theorem}

\begin{proof}
Fix compatible roots \(\zeta_n\) of order \(\ell^n\).
The standard character is given by Kummer action on roots of
\[
 E_{m,n}=
 \prod_{a\in(\mathbb Z/\ell^n)^\times}
       (1-\zeta_n^a)^{a^{m-1}}.                    \tag{K19}
\]
Kurihara \cite[p.~233, (4)]{Kurihara1992} defines the
cyclotomic class \(c(1)\in H^1(S_\ell,\mathbb Z_\ell(m))\)
at finite level by corestriction of
\((1-\zeta_n)\otimes\zeta_n^{\otimes(m-1)}\).
Restriction--corestriction to \(N_\ell\) gives the sum
of its conjugates. Factoring out
\(\zeta_n^{\otimes(m-1)}\) produces precisely (K19).
Thus \(\operatorname{res}_{N_\ell}c(1)=\kappa_m\);
there is no factorial in this normalization.

The group \(H^1(S_\ell,\mathbb Z_\ell(m))\) is free of
rank one and \(c(1)\ne0\), as recalled on the same page.
Let \(\alpha\) be a primitive generator. Its reduction in
\(H^1(S_\ell,\mathbb F_\ell(m))\) is nonzero by the
coefficient exact sequence. Write
\(\Gamma=\mathbb Z_\ell^\times=\Delta\times\Gamma_1\),
where \(|\Delta|=\ell-1\) and \(\Gamma_1=1+\ell\mathbb Z_\ell\).
Since \(m\) is odd, its residual character on \(\Delta\)
is nontrivial. Inflation--restriction gives
\[
 H^1(\Gamma,\mathbb F_\ell(m))
 =H^1(\Gamma_1,\mathbb F_\ell(m))^\Delta=0.           \tag{K20}
\]
Here \(\Gamma\) is abelian, so \(\Delta\) acts trivially
on \(\Gamma_1\) and by the nontrivial character on the
coefficient module.

If \(\alpha|_{N_\ell}\) had image contained in
\(\ell\mathbb Z_\ell(m)\), its reduction would restrict
to zero. Inflation--restriction for the cyclotomic quotient
would put that nonzero reduction in the zero group (K20),
a contradiction. Apply this first in the maximal extension
unramified outside \(\ell\); the full group \(N_\ell\)
surjects onto its cyclotomic kernel. The image of
\(\alpha|_{N_\ell}\) is a compact, hence closed subgroup
of \(\mathbb Z_\ell(m)\). Since it contains a unit,
it is the whole module.
Writing \(c(1)=a\alpha\), \(a\ne0\), now gives
\[
 [\mathbb Z_\ell(m):\kappa_m(N_\ell)]
 =[H^1(S_\ell,\mathbb Z_\ell(m)):\mathbb Z_\ell c(1)].
\]
Kurihara's Proposition~5.1 gives an inequality; his
Remark~5.2 explicitly states equality, citing Bloch--Kato,
\S6, (6.8)--(6.9). We use that equality as a known
theorem. It identifies the last index with \(|H^2|\).
The modern comparison (K1) completes (K18).
\end{proof}

This reconciles the index formula with Ishii's precise
full-domain statement \cite[Theorem~1.1, p.~2]{IshiiKummer2024}:
every odd-degree character is nonzero; the character is
surjective when \(m\equiv1\pmod{\ell-1}\); and all odd-degree
characters are surjective exactly when Vandiver holds at
\(\ell\). For a fixed degree the condition is the one
component in (K2).

Changing the chosen radical does not change the character
on \(N_\ell\), which fixes the roots of unity.
Replacing \(1-\zeta_n^a\) by \(\zeta_n^a-1\) changes a
Kummer unit by a sign, whose odd-primary Kummer class is
zero. Changing the compatible root generator by
\(b\in\mathbb Z_\ell^\times\) changes the scalar coordinate
by the unit \(b^{-m}\), so the index is unchanged.
Half-products or real norms may introduce powers of 2,
also units here. Ichimura--Kaneko's universal-series
coordinate \cite[\S1.1, pp.~313--314]{IchimuraKaneko1989}
contains \((1-\ell^{m-1})^{-1}\kappa_m/m!\):
the first factor is a unit, whereas \(m!\) need not be.
It must not be inserted into the definition (K19).
The original Ichimura--Sakaguchi paper was bibliographically
identified but not directly read in this run; no original
theorem number from it is asserted.

The proof uses oddness of \(\ell\), especially (K20), and
does not extend the index formula to 2. Nor does it
identify the image of \(\kappa_m\) on a deeper Galois
filtration subgroup with its image on all of \(N_\ell\).
Consequently it proves no integral free-generation theorem.

\subsection{A reproducible certificate at 37}

The criterion of Buhler--Harvey
\cite[\S4, author preprint pp.~4--5]{BuhlerHarvey2011}
with the printed exponent corrected by Hart--Harvey--Ong
\cite[\S7, author preprint p.~13]{HartHarveyOng2017}
says to test, for every irregular index \(a\) at \(\ell\),
\[
 V_{\ell,a}=
 \prod_{c=1}^{(\ell-1)/2}
   (\zeta^{c/2}-\zeta^{-c/2})^{c^{\ell-1-a}}\pmod q,\qquad
 V_{\ell,a}^{(q-1)/\ell}\not\equiv1\pmod q,           \tag{K21}
\]
where \(q\equiv1\pmod\ell\) is prime, \(\zeta\) has order
\(\ell\) modulo \(q\), and the half exponents are interpreted
modulo \(\ell\).
For \(\ell=37\), the residues of
\(B_2,B_4,\ldots,B_{34}\) modulo 37 are
\[
31,16,15,16,4,17,32,22,31,15,15,17,12,29,2,0,2.
\]
Thus the only irregular index is \(a=32\).
Taking \(q=149\), \(\zeta=107\), and its square root
\(s=16\) in the order-37 subgroup, we have
\(\zeta^{19c}=s^c\), since \(2\cdot19\equiv1\pmod{37}\).
Thus the computation with factors \(s^c-s^{-c}\) implements
the corrected half-exponent formula exactly and gives
\[
 16^{37}\equiv1,\quad16\not\equiv1,\quad
 V_{37,32}=48,\quad48^4=142\not\equiv1\pmod{149}.
\]
Trial division certifies the primes and hence the root's
exact order. This proves Vandiver at 37 by the cited
criterion, and consequently
\[
 K_{4k}(\mathbb Z)\{37\}=0\quad(k\ge1).
\]
The accompanying exact-arithmetic script compares two
Bernoulli algorithms and two modular-product algorithms,
and checks the equality of every square-root factor with
the corrected half-exponent factor.
The code and the application of the criterion were read
in full after running. This is a reproducible numerical
certificate with a cited arithmetic criterion, not a new
formal proof certificate (VK-37).

\subsection{The remaining numerical question}

The known primary-part exclusions give the sufficient gate
\[
 |K_{4k}(\mathbb Z)|<2^{31}
 \quad\Longrightarrow\quad K_{4k}(\mathbb Z)=0.       \tag{K22}
\]
A nontrivial finite group has a prime divisor of its order,
and all smaller primes have been excluded. Thus either a
small enough explicit bound or the exact trace gcd in
(K17) could close a fixed-degree case. Merely knowing that
an absolute constant in (K11) or (K13) exists does not.
Indeed the source's particular coarse choice
\(C=36c_5\,2^A\) satisfies \(C\ge36>\log(2^{31})\)
whenever \(c_5\ge1\); that bound cannot meet the gate even
with its smallest positive integer constant.

For a precise logical illustration, choose a prime
\(q>2^{31}\), put \(h=(q-1)/2\), and define abstract groups
\[
 M_k=
 \begin{cases}
 \mathbb Z/q,& k\equiv3\pmod h,\\
 0,&\text{otherwise}.
 \end{cases}
\]
They have no 2-primary or below-\(2^{31}\) primary parts,
\(M_1=M_2=0\), the prescribed fixed-prime vanishing
periodicity, and
\(\log|M_k|\le\lceil\log q\rceil(2k+1)^7\), but
\(M_3\ne0\). These are only a model of the specifically
listed size and numerical patterns. They are not actual
\(K\)-groups or cyclotomic class data and give no
counterexample to Vandiver (VK-NUM-GATE).

The current primary-source review found no proof of
\(K_{12}(\mathbb Z)=0\) or \(K_{16}(\mathbb Z)=0\).
This reports the scope of the search, not an exhaustive
negative theorem about the literature.
Kupers \cite[Remark~2.6, p.~4]{Kupers2019} gives an
integral Church--Farb--Putman route to \(K_{12}\).
The subsequent result of Br\"uck--Miller--Patzt--Sroka--Wilson
\cite[Theorem~B]{BrueckEtAl2024} proves rational
codimension-two vanishing, and does not state the integral
and higher-codimension input required by that route.
The later Miller--Patzt--Putman theorem independently removes rank $11$.
Proposition~\ref{prop:k12-three-ranks} retains ranks $8,9,10$, explicitly
isolating the older computation's unverified rank-eight endpoint.
No new specified degree vanishing is claimed from the
existential constants extracted here.

\section{Tripod kernels and integral Soul\'e generators}
\label{sec:tripod-integral}

\subsection{Correct field formulation and the odd regular theorem}

For a prime $p$, write
\[
 X=\mathbb P^1_{\mathbb Q}\setminus\{0,1,\infty\},\qquad
 \rho_p:G_{\mathbb Q}\longrightarrow
 \operatorname{Out}\bigl(\pi_1^{\mathrm{et}}(X_{\overline{\mathbb Q}})^{(p)}\bigr),
\]
\[
 F=\mathbb Q(\mu_p),\quad K=\mathbb Q(\mu_{p^\infty}),\quad
 \Omega_p^*=\overline{\mathbb Q}^{\ker\rho_p},
\]
and let $\Omega_p$ be the maximal pro-$p$ extension of $F$ unramified outside $p$. Thus the desired equality is
\[
 \Omega_p^*=\Omega_p,
 \quad\text{equivalently}\quad
 \ker\rho_p=\operatorname{Gal}(\overline{\mathbb Q}/\Omega_p).
\]
The kernel is a subgroup of $G_{\mathbb Q}$; it is not itself the maximal ramification Galois group. At $p=2$ the ramification convention at infinity is specified in Theorem~\ref{thm:pro2-kernel}.

Put $G_p=\operatorname{Gal}(\Omega_p^*/K)$, with the Ihara weight filtration $F^mG_p$, and
\[
 \mathfrak g_p=\bigoplus_{m\geq1}F^mG_p/F^{m+1}G_p.
\]
For odd $m\geq3$, choose $\sigma_m\in F^mG_p$ whose Soul\'e value generates the actual image ideal $\kappa_m(F^mG_p)$. Let
\[
 \psi:\operatorname{Lie}_{\mathbb Z_p}(s_3,s_5,\ldots)
       \longrightarrow\mathfrak g_p,
 \qquad s_m\longmapsto\sigma_m,
\]
and let $\Psi$ be the corresponding homomorphism from the free pro-$p$ group on these generators.

\textbf{Sharifi's exact theorem.} For an \emph{odd regular} prime $p$, $\Psi$ is onto. If rational Deligne--Ihara holds, then $\psi$ and $\Psi$ are isomorphisms and $\Omega_p^*=\Omega_p$ \cite[Theorem~1.1, author's preprint p.~2]{Sharifi2002}. Brown supplies the rational conjecture \cite{Brown2012,Brown2019}. Therefore both the kernel-field equality and integral odd free generation hold for odd regular primes. This resolves the correctly formulated known row LM6~A.44 without step~17. Ishii's recent primary account restates exactly this scope \cite[Theorem~1.1 and Remark~1.1.1, p.~2]{IshiiKernel2026}.

\subsection{A separate deduction for the prime \texorpdfstring{$2$}{2}}
\label{sec:pro2-proof}

The following result is a deduction developed and checked in this run. We did not locate it as a stated theorem in the literature searched, and we make no novelty claim. Its proof uses established arithmetic and rational Deligne--Ihara results together with the complete argument below. It does not use the proposed step--17 estimate. The initiating researcher, two other agents, and the parent researcher each read the whole argument; the check ledger distinguishes those reviews from the cited literature inputs and from formal verification.

\begin{theorem}[K17-P2-MOMENT]\label{thm:pro2-kernel}
Let $X=\mathbb P^1_{\mathbb Q}\setminus\{0,1,\infty\}$ and $\Pi=\pi_1^{\mathrm{et}}(X_{\overline{\mathbb Q}})^{(2)}$. Let $\Omega_2$ be the maximal pro-$2$ extension of $\mathbb Q$ unramified at every finite prime other than $2$, with complexification at the real place allowed. If $\Omega_2^*$ is the fixed field of the kernel of the outer representation
\[
 \rho_2:G_{\mathbb Q}\longrightarrow\operatorname{Out}(\Pi),
\]
then $\Omega_2^*=\Omega_2$.
\end{theorem}

The convention at infinity is part of the statement. The theorem concerns faithfulness of a pro-$2$ action; it does not assert that its integral weight Lie algebra is freely generated by odd Soul\'e elements.

\paragraph{Factorization through the maximal ramification group.}
The image of $\rho_2$ is pro-$2$. Here is a direct justification that does not invoke an odd-prime formulation. The geometric group $\Pi$ is free pro-$2$ of rank two, and its Frattini quotient is $\mathbb F_2(1)^2$. The Galois action on this quotient is trivial. For any finite $2$-group $Q$, the kernel of
\[
 \operatorname{Aut}(Q)\longrightarrow
 \operatorname{Aut}(Q/\Phi(Q))
\]
is a $2$-group. To see this directly, fix a basis of $Q/\Phi(Q)$ of size $d$. Its lifts form a set of cardinality $|\Phi(Q)|^d$, and every such tuple generates $Q$ by the defining property of the Frattini subgroup. The displayed kernel acts freely on this set, since an automorphism fixing a generating tuple is the identity. Its order therefore divides a power of $2$. Use the characteristic finite quotients obtained from the descending Frattini series of $\Pi$. They are cofinal among its finite quotients, and an automorphism trivial on $\Pi/\Phi(\Pi)$ induces an automorphism trivial on the Frattini quotient of each of them. Consequently the corresponding kernel in $\operatorname{Aut}(\Pi)$ is pro-$2$. Inner automorphisms lie in that kernel, so the kernel in $\operatorname{Out}(\Pi)$ is also pro-$2$.

At every odd prime, the marked curve $(\mathbb P^1,\{0,1,\infty\})$ has good reduction with disjoint marked sections. Prime-to-residue-characteristic specialization therefore makes inertia act trivially on its outer pro-$2$ fundamental group. Equivalently, one may use the all-primes factorization recalled in Hain--Matsumoto, \S8.4, p.~27 \cite{HainMatsumoto2001}. Combining unramifiedness outside $2$ with the pro-$2$ image gives an induced continuous representation of
\[
 P=\operatorname{Gal}(\Omega_2/\mathbb Q).
\]
It suffices to prove that representation faithful.

\paragraph{The arithmetic source.}
Gras--Jaulent's Theorem/Definition~0(iv) and Theorem~2 give
\[
 P\simeq C_2*\mathbb Z_2
\]
for the unrestricted-at-infinity group \cite[pp.~1--3]{GrasJaulent2008}. The specialization of their $2$-rationality hypothesis to $\mathbb Q$ is elementary. A rational number whose valuations at all odd primes are even is, modulo rational squares, $\epsilon2^e$ with $\epsilon\in\{1,-1\}$ and $e\in\{0,1\}$. Being a square in $\mathbb Q_2$ forces $e=0$ by valuation and $\epsilon=1$ because $-1$ is not a $2$-adic square. The global and local $2$-primary root-of-unity groups both have $2$-rank one. Thus their criterion applies.

Kronecker--Weber identifies the full abelianization with the cyclotomic character
\[
 \chi:P\longrightarrow\mathbb Z_2^\times.
\]
Write $P=\langle c\rangle*\langle\gamma\rangle$, where $c^2=1$ and $\langle\gamma\rangle\simeq\mathbb Z_2$. Then $\chi(c)=-1$. Replacing $\gamma$ by $c\gamma$ if necessary is an automorphism of the free product and ensures that
\[
 u=\chi(\gamma)\in1+4\mathbb Z_2.
\]
Since $\chi$ is the abelianization, $u$ is a topological generator of $1+4\mathbb Z_2$.

The subgroup $P_i=\ker(\chi\bmod4)$, corresponding to $\mathbb Q(i)$, is free pro-$2$ on $\gamma$ and $c\gamma c$. One can verify this without an infinite-rank subgroup theorem: let the involution exchange the two generators $a,b$ of a free pro-$2$ group. Its semidirect product with $C_2$ has the universal property of $C_2*\mathbb Z_2$, with $a=\gamma$ and $b=c\gamma c$. Its index-two kernel is the asserted free group. Set
\[
 h=(c\gamma c)\gamma^{-1},\qquad H=\ker\chi.
\]
Then $(\gamma,h)$ is another free pro-$2$ basis of $P_i$, and $H$ is the kernel of the map $P_i\to\mathbb Z_2$ sending $\gamma$ to $1$ and $h$ to $0$. For $x\in\mathbb Z_2$, write
\[
 h_x=\gamma^xh\gamma^{-x}.
\]

\paragraph{Finite quotients and measure-valued Magnus coefficients.}
Let $N_n$ be the closed normal subgroup of $P_i$ generated by $\gamma^{2^n}$. Then
\[
 P_i/N_n\simeq C_{2^n}*\mathbb Z_2.
\]
The kernel $H_n$ of its map to $C_{2^n}$ is free pro-$2$ on the $2^n$ elements $h_j$, $j\bmod2^n$. Indeed, the free pro-$2$ group on those symbols, with the cyclic permutation action, has semidirect product with $C_{2^n}$ possessing precisely the universal property of $C_{2^n}*\mathbb Z_2$.

The subgroups $N_n$ decrease to $1$. In any finite continuous $2$-group quotient of $P_i$, the image of $\gamma$ has finite $2$-power order, so $N_n$ lies in its kernel for all sufficiently large $n$. Compactness also makes the map to the inverse limit surjective: compatible quotient cosets form a nested family of nonempty closed subsets of $P_i$. Hence
\[
 P_i\simeq\varprojlim_nP_i/N_n,
 \qquad
 H\simeq\varprojlim_nH_n.                                  \tag{P2.1}
\]
The transition maps fold the finite bases by reduction of $j$ modulo $2^n$. The elements $h_x$ topologically generate $H$. To check this last point explicitly, the kernels of $H\to H_n$ decrease to $1$, and compactness implies that every open normal subgroup contains one of them. Every finite quotient of $H$ therefore factors through an $H_n$, whose indicated generators are images of the $h_x$.

The finite-rank integral Magnus maps are injective. A primary locator is \cite[\S5.1, p.~21, (5.1)]{MinacPasiniQuadrelliTan2020}, whose mod-$2$ statement implies the integral injection. There is also a short direct verification. The augmentation ideal of $\mathbb F_2[Q]$ for a finite $2$-group $Q$ is nilpotent: choose a central element $z$ of order two, use $(z-1)^2=0$, and induct on $Q/\langle z\rangle$. Thus any map from a finite-rank free pro-$2$ group to $Q$ is recovered by substituting $X_j\mapsto\overline g_j-1$ into the mod-$2$ Magnus series. An element with trivial expansion maps to $1$ in every finite quotient and is therefore $1$.

The degree-$r$ coefficients of the compatible integral expansions in (P2.1) belong to
\[
 M_r=\varprojlim_n\mathbb Z_2[(\mathbb Z/2^n)^r]
     =\mathbb Z_2[[\mathbb Z_2^r]],
\]
the bounded $\mathbb Z_2$-valued measures on $\mathbb Z_2^r$. With multiplication given by external product and concatenation, they define an injective homomorphism
\[
 \theta:H\longrightarrow1+\prod_{r\geq1}M_r,
 \qquad \theta(h_x)=1+\delta_x.                             \tag{P2.2}
\]
Injectivity follows from the finite-rank maps and (P2.1). Consequently every nontrivial element has a first finite nonzero Magnus degree. Completion in the $(2,I)$ topology creates no element with all homogeneous coefficients zero: the completed coefficient algebra is the displayed product.

\paragraph{The actual rational image.}
Let $B$ be the proalgebraic closure over $\mathbb Q_2$ of the rational outer image induced by the tripod representation, and let $U$ be its prounipotent radical. Hain--Matsumoto, \S8.3 and Proposition~8.3, Theorem~8.4, pp.~25--26, gives an extension
\[
 1\longrightarrow U\longrightarrow B\longrightarrow\mathbb G_m\longrightarrow1
\]
whose reductive quotient is induced by $\chi$, and a continuous map
\[
 \rho_U:H\longrightarrow U(\mathbb Q_2)                    \tag{P2.3}
\]
with Zariski-dense image \cite{HainMatsumoto2001}. The rational action is induced functorially from the pro-$2$ action; a trivial pro-$2$ outer automorphism gives a trivial rational outer automorphism. Their Appendix~A.6, A.9 and A.10 records the relevant comparison. We have used the actual arithmetic image, not an unrelated motivic group.

Brown's rational Deligne--Ihara theorem and Hain--Matsumoto's graded comparison give, with positive arithmetic weights,
\[
 \operatorname{gr}\operatorname{Lie}U
 \simeq\operatorname{Lie}_{\mathbb Q_2}(s_3,s_5,s_7,\ldots),
 \qquad\operatorname{wt}(s_m)=m.                           \tag{P2.4}
\]
The exact outer-action statement for every prime is Brown \cite[\S2.1, Conjecture~1, p.~2; \S2.3, Theorem~2.1 and following paragraph, p.~4]{Brown2019}; the original theorem is \cite[Theorem~1.1 and its consequence, p.~950]{Brown2012}.

We spell out the completion step. In a separated complete positively filtered Lie algebra whose associated graded is free with only finitely many generators below each weight, choose lifts of the free graded generators. The resulting map from the weight-completed free Lie algebra induces an isomorphism on associated gradeds. Induction through successive finite weight quotients proves it is an isomorphism on each quotient; inverse limits give an isomorphism of the completed Lie algebras. A Levi splitting of the negatively weighted extension, as in \cite[Lemma~3.7, p.~9]{HainMatsumoto2001}, permits weight-homogeneous choices. Thus the completed universal enveloping algebra of $\operatorname{Lie}U$ is the weight-completed free associative algebra on the $s_m$. Its finite weight quotients are finite-dimensional.

Let $\ell_m:U(\mathbb Q_2)\to\mathbb Q_2(m)$ be the $m$th abelian coordinate, and put $a_m=\ell_m(\rho_U(h))$. Conjugation by $\gamma$ acts on that coordinate by $u^m$, because the unipotent part of a lift acts trivially on abelianization. Hence
\[
 \ell_m(\rho_U(h_x))=a_mu^{mx}.                             \tag{P2.5}
\]
Every $a_m$ is nonzero. Otherwise that coordinate would vanish on all the topological generators $h_x$, hence on $H$, contradicting the density in (P2.3) and the nonzero $m$th abelian quotient in (P2.4). No integral primitive normalization is asserted here.

\paragraph{Substitution at a finite weight cutoff.}
Suppose $g\neq1$ lies in the kernel of (P2.3). Let $r$ be its first nonzero Magnus degree, and let
\[
 0\neq\mu=\mu_r(g)\in\mathbb Z_2[[\mathbb Z_2^r]].
\]
In the completed free associative algebra put $Y_x=\rho_U(h_x)-1$. Its constant term is zero and its part of generator length one is
\[
 \sum_{m\geq3,\ m\text{ odd}}a_mu^{mx}s_m.
\]
Here generator length is distinct from arithmetic weight and from geometric depth.

Substitution $\delta_x\mapsto Y_x$ is well-defined and continuous coefficient by coefficient. At target weight at most $N$, only Magnus degrees at most $\lfloor N/3\rfloor$ contribute. The remaining coefficient functions on compact powers of $\mathbb Z_2$ are continuous and bounded in $\mathbb Q_2$. Integration against a bounded $\mathbb Z_2$-measure is defined by uniform approximation by locally constant functions, after scaling into $\mathbb Z_2$ if necessary. The same approximation proves continuity, multiplicativity and Fubini's identity. This construction needs no lattice uniform across all weights. The substitution and $\rho_U$ agree on every $h_x$ and are continuous homomorphisms on each finite weight quotient, so topological generation makes them agree on $H$. No universal property of a free pro-$2$ group is applied to a noncompact target.

The coefficient of the ordered word $s_{m_1}\cdots s_{m_r}$ in the image of $g-1$ is therefore
\[
 \left(\prod_{j=1}^ra_{m_j}\right)
 \int_{\mathbb Z_2^r}\prod_{j=1}^ru^{m_jx_j}\,d\mu(x).
\]
Lower Magnus degrees are zero, and higher ones have generator length greater than $r$. Since $\rho_U(g)=1$ and every $a_m\neq0$, we obtain
\[
 \int_{\mathbb Z_2^r}\prod_{j=1}^ru^{m_jx_j}\,d\mu(x)=0
 \quad(m_j\geq3\text{ odd}).                              \tag{P2.6}
\]
The variables remain ordered; this is not a symmetrized moment.

\paragraph{Analytic separation.}
The Amice transform
\[
 A_\mu(T_1,\ldots,T_r)
   =\int_{\mathbb Z_2^r}\prod_{j=1}^r(1+T_j)^{x_j}\,d\mu(x)
\]
lies in $\mathbb Z_2[[T_1,\ldots,T_r]]$ and converges on the open unit polydisc. It is injective as a transform of bounded measures: its coefficients are the moments against products of binomial polynomials, the multivariate Mahler basis. Equivalently it is the canonical isomorphism of the multivariable Iwasawa algebra with the power-series algebra.

For odd $m\geq3$ put $t_m=u^m-1$. These points lie in $4\mathbb Z_2$, and the distinct points $t_{1+2^j}$ converge to $u-1$. Distinctness follows from the infinite order of $u$. Equation (P2.6) says that $A_\mu$ vanishes on the full Cartesian grid of these points. The one-variable nonarchimedean identity theorem applies on the closed disc of radius $1/2$, contained in the radius-one convergence disc. Fixing grid values in all but one variable and applying that theorem successively in each variable shows $A_\mu=0$. Therefore $\mu=0$, a contradiction.

It follows that $\rho_U$ is injective on $H$. The kernel of the pro-$2$ tripod action is contained in $H$ because geometric abelianization records $\chi$. Its elements are also trivial under the induced rational outer action, so they are $1$. This proves Theorem~\ref{thm:pro2-kernel}.

\paragraph{What this proof does and does not decide.}
The proof uses the explicit arithmetic presentation $C_2*\mathbb Z_2$, not merely abstract freeness of $H$. It settles the field equality with the stated convention at infinity. It does not prove integral odd generation of the weight Lie algebra, and it does not transfer to irregular primes without a substitute for its linear-moment separation. The latter obstruction is made explicit below.


\subsection{Four meanings of an integral generator}
\label{sec:four-integral-meanings}

The following claims differ and must be kept separate:
\begin{enumerate}
\item[(A)] A nonzero integral odd class $\sigma_m\in\mathfrak g_{p,m}$ exists.
\item[(B)] The full-domain character
\[
 \kappa_m:G_{\mathbb Q(\mu_{p^\infty})}^{\mathrm{ab}}
       \longrightarrow\mathbb Z_p(m)
\]
is onto.
\item[(C)] An element of the weight subgroup $F^mG_p$ has $\kappa_m$-value $1$ in a fixed integral normalization.
\item[(D)] The map $\psi$ above is an isomorphism of graded $\mathbb Z_p$-Lie algebras.
\end{enumerate}

For odd $p$, (A) is already supplied by Soul\'e nonvanishing and Sharifi's integral construction \cite[Proposition~2.3, pp.~5--6]{Sharifi2002}. It is primitive relative to the \emph{actual image ideal}, which need not be $\mathbb Z_p$. At $p=2$, existence also follows from Brown and the actual lattice: multiply a nonzero rational odd generator by a sufficiently large power of $2$ to put it in $\mathfrak g_{2,m}$, then choose a representative in $F^mG_2$. This establishes existence, independently of the kernel theorem and of step~17.

For odd $p$, all of (B), for odd $m\geq3$, hold if and only if Vandiver holds at $p$. More precisely, each character is nonzero, and it is onto if $m\equiv1\pmod{p-1}$; simultaneous surjectivity is equivalent to Vandiver \cite[Theorem~1.1(1)--(3), p.~2]{IshiiKummer2024}. The theorem is attributed there to the classical Soul\'e and Ichimura--Sakaguchi work. We verified this precise primary-paper locator; we did not verify an original Ichimura--Sakaguchi theorem number and do not invent one. On all of $G_{\mathbb Q}$ the class is a cocycle, while on the cyclotomic kernel it is an ordinary continuous homomorphism.

There is a decisive restriction issue in (C). Sharifi's integral stable-derivation comparison gives, on $F^mG_p$,
\[
 \kappa_m=(p^{m-1}-1)(m-1)!\lambda_m^{(p)},
 \qquad \lambda_m^{(p)}\text{ integral}
 \quad\cite[(2.10), p.~7]{Sharifi2002}.
\]
Thus
\[
 \kappa_m(F^mG_p)\subseteq(m-1)!\mathbb Z_p(m).              \tag{I.1}
\]
At $p=3,m=5$, this image is contained in $3\mathbb Z_3$. Yet $3$ is regular, so (B) and (D) hold. Consequently full-domain surjectivity, weight-subgroup primitivity, and integral generation cannot be conflated. Formula (I.1) does not dispute the cohomological index formula for the full-domain cyclotomic class; it concerns a further restriction.

For (D), odd regular primes are settled positively by Sharifi--Brown. The prime $37$ gives a negative answer despite Vandiver and the corresponding vanishing of all the $K_{4k}$-groups at that prime. The proof follows.

\subsection{A counterexample to integral odd generation at \texorpdfstring{$37$}{37}}
\label{sec:integral37}

Let $\mathfrak h_p=\operatorname{im}\psi$. McCallum--Sharifi's notation in \S9 agrees with this choice: their ``Deligne conjecture in degree $i$'' means injectivity of $\operatorname{gr}^i\psi$ \cite[p.~24]{McCallumSharifi2003}. Brown gives precisely this injectivity, because the rational map is injective and every graded piece of the integral free source is torsion-free.

\begin{proposition}[T-INTEGRAL37]\label{prop:37-counterexample}
At $p=37$, Vandiver holds, all $K_{4k}(\mathbb Z)\otimes\mathbb Z_{37}$ vanish for $k\geq1$, and all full-domain odd Soul\'e characters are onto. Nevertheless
\[
 (\mathfrak g_{37}^{\mathrm{ab}})_{32}\neq0,
 \qquad
 \mathfrak h_{37,32}\subsetneq\mathfrak g_{37,32}.
\]
In particular $\mathfrak g_{37}$ cannot be generated over $\mathbb Z_{37}$ by homogeneous elements all of odd weight, and is not the free graded integral Lie algebra on one generator in every odd weight at least three.
\end{proposition}

\begin{proof}
The only irregular index for $37$ is $32$. Vandiver at $37$ can be verified by the small exact modular certificate accompanying this report, using the criterion of Buhler--Harvey \cite[\S4, pp.~4--5]{BuhlerHarvey2011}: take $q=149$ and $\zeta=107$, whose square root in the order-$37$ subgroup is $16$. The half-exponent normalization is corrected in Hart--Harvey--Ong, \S7, p.~13; the detailed certificate above uses exactly that correction. The required product is $V=48\pmod{149}$ and $V^{(149-1)/37}=142\not\equiv1\pmod{149}$. The computation uses exact integer arithmetic and was independently cross-checked; it is not a formal proof certificate. The modern K-theoretic comparison recorded elsewhere in this report gives the asserted $K_{4k}$ vanishing. The full-domain character assertion follows from \cite[Theorem~1.1(3)]{IshiiKummer2024}.

McCallum--Sharifi's Theorem~7.5 proves that their pairing at index $32$ for $37$ is nonzero. Their Corollary~8.3, p.~23, makes a coefficient explicit: $e_{3,32}$ is nonzero, represented by $-11$ up to a nonzero normalization scalar. Their Theorem~9.7, p.~28, applies when Vandiver holds, $m$ is the least irregular even index, rational Deligne holds through degree $m$, and some $e_{i,m}\neq0$ has $i<m/2$. All hypotheses hold with $(p,m,i)=(37,32,3)$, by Brown and $3<16$.

For clarity, the proof of that cited theorem identifies the mechanism. In degrees below $32$, the integral map is bijective. Thus the weight-$32$ commutator part of $\mathfrak g$ equals that of $\mathfrak h$. The nonzero cup coefficient, through their Proposition~9.6, prevents bijectivity at $32$; Brown's injectivity makes the defect a failure of surjectivity. Their theorem consequently gives both displayed conclusions.

If a graded Lie algebra is generated by homogeneous odd-weight elements, its abelianization is generated by their odd-weight images and has no even-weight component. The nonzero degree-$32$ abelianization contradicts this. This proves the proposition.
\end{proof}

The authors' Magma computation in Theorem~7.5 is used as a published theorem; we did not rerun it. The deduction above, including every hypothesis of Theorem~9.7, was read in full. It does not say that $\mathfrak h_{37}$ is nonfree: rational injectivity makes it a free integral subalgebra. It says that its degree-$32$ lattice is proper. The graded pieces of $\mathfrak g_{37}$ remain torsion-free; their Lie abelianization can have $37$-torsion.

Most importantly, Proposition~\ref{prop:37-counterexample} does not refute $\Omega_{37}^*=\Omega_{37}$. An arithmetic source and its faithful image can have these same extra integral directions. This is a counterexample to the proposed integral-generation implication, not a counterexample to the kernel-field equality.

Sharifi also identifies the first possible defect under Vandiver: if $\Psi$ is not onto and $m$ is its first defective weight, then $m$ must be even and $p\mid B_m$ \cite[Theorem~4.1, p.~10]{Sharifi2002}. The odd-weight alternative would violate Vandiver. This explains why removing the plus class-group obstruction does not eliminate all integral-generation obstructions.

\subsection{Irregular primes: the sharper remaining gate}
\label{sec:irregular-gate}

Write $H_p=\operatorname{Gal}(\Omega_p/K)$ and $M_p=H_p^{\mathrm{ab}}$, with the conjugation action of $\Gamma=\operatorname{Gal}(K/F)\simeq\mathbb Z_p$ and $\Lambda=\mathbb Z_p[[\Gamma]]$. A useful correction is that $H_p$ is already free pro-$p$ even at irregular primes. Iwasawa's original Theorem~2 states that the maximal $S$-ramified pro-$p$ group over a cyclotomic $\mathbb Z_p$-field is free exactly when its $\mu$-invariant is zero \cite[pp.~273--274]{Iwasawa1981}. Ferrero--Washington applies to $K/\mathbb Q$, which is abelian. This does not make the finite-base group $\operatorname{Gal}(\Omega_p/F)$ free.

Let $A_\infty$ be the direct limit, under extension of ideals, of the $p$-parts of the class groups in the cyclotomic tower. Kummer duality gives
\[
 M_p^+\simeq(A_\infty^-)^\vee(1)
 \quad\cite[Corollary~3.3.14 and Proposition~3.4.7]{SharifiIwasawaNotes}.
                                                               \tag{I.2}
\]
The finite-level minus class groups inject into $A_\infty^-$, by the same source's Proposition~3.4.1. At an irregular prime that minus class group is nonzero, so $M_p^+\neq0$. With $\mu=0$, this is a nonzero finite-rank free $\mathbb Z_p$-module but a torsion $\Lambda$-module. The distinction between $p$-torsion and $\Lambda$-torsion is essential.

Every odd linear Soul\'e coordinate kills $M_p^+$: complex conjugation acts by $+1$ on the source and $-1$ on the target $\mathbb Q_p(m)$, so equivariance gives twice the value equal to zero. Thus the degree-one separation used in the pro-$2$ proof fails on a concrete nonzero module. A lift of a nonzero class in $M_p^+$ can have a nonzero first source Magnus coefficient while all target generator-length-one coordinates vanish. Its image may still be nontrivial at higher length. Abstract freeness of $H_p$, weak Leopoldt, or $\mu=0$ alone therefore cannot repeat that proof.

A sharper continuation is to control the full nonlinear lifts of the $\Gamma$-relations defining this torsion module. Their first nonlinear terms are cup products; the degree-$32$ result at $37$ detects one such direction through a bracket. Kernel equality requires detection of every remaining direction and its iterated commutators, or an equivalent separation theorem for the full noncommutative relation ideal. A class-group order or characteristic ideal alone does not specify those nonlinear terms. The following subsections supply an all-weight separation proof at $37$, its general pole criterion, and the modular-symbol input that makes Vandiver alone sufficient. In particular the field equality holds for every prime below $2^{31}$; the remaining all-prime gate is explicit. Neither Proposition~\ref{prop:37-counterexample} nor (I.2) gives a negative answer.

% Checked deduction: author, root, vandiver and archive whole-proof readings completed.
% Candidate and first independent audit: work/irregular37_pole_candidate.md.
\subsection{A nonlinear separation argument at \texorpdfstring{$37$}{37}}
\label{sec:kernel37-poles}

The failure of integral odd generation at $37$ does not decide the kernel-field question. The following argument detects its extra even arithmetic direction through a pole in a quadratic moment function. It uses the published cup-product computation, but supplies a separate group-theoretic and analytic proof of faithfulness. The proof is a deduction developed in this run, not a theorem located in the literature or a consequence of step~17.

\begin{theorem}[K17-37-KERNEL]\label{thm:kernel37-poles}
For the maximal pro-$37$ extension $\Omega_{37}$ of $\mathbb Q(\mu_{37})$ unramified outside $37$, the tripod kernel field satisfies
\[
 \Omega_{37}^*=\Omega_{37}.
\]
This is compatible with $(\mathfrak g_{37}^{\mathrm{ab}})_{32}\ne0$ and failure of integral odd generation.
\end{theorem}

We first prove the separation statement needed for the arithmetic application.

\paragraph{K17-POLE-ONE: the compact orbit basis and separation criterion.}
Let $p$ be odd, let $\Gamma=\langle\gamma\rangle\simeq\mathbb Z_p$, and put $\Lambda=\mathbb Z_p[[\Gamma]]$. Suppose $H$ is free pro-$p$ with a continuous $\Gamma$-action, and that its abelianization has the compatible decomposition
\[
 H^{\mathrm{ab}}=
 \left(\bigoplus_{a\in A}\Lambda e_a\right)\oplus\mathbb Z_p e_y,
 \qquad \gamma e_y=\alpha e_y,\qquad \alpha\in1+p\mathbb Z_p.       \tag{N.1}
\]
Here $A$ consists of all odd residue classes modulo $p-1$. A commuting finite group $\Delta$ of order $p-1$ acts on $e_a$ by character $a$ and on $e_y$ by an even character. Choose lifts $x_a,y\in H$, and set
\[
 x_{a,t}=\gamma^t x_a\gamma^{-t},\qquad t\in\mathbb Z_p.
\]
These are a free compact generating space for $H$, in the following precise sense.

Let $F_X$ be the inverse limit of the finite-rank free pro-$p$ groups on
\[
 X_n=(A\times\mathbb Z/p^n)\amalg\{y\},
\]
with maps folding the residue indices. This limit has the universal property for continuous maps from $X=(A\times\mathbb Z_p)\amalg\{y\}$ into a pro-$p$ group: after composition with any finite quotient of the target, such a map factors through an $X_n$. Restricting that property to abelian targets gives
\[
 F_X^{\mathrm{ab}}=\mathbb Z_p[[X]]=
 \Lambda^{|A|}\oplus\mathbb Z_p.
\]
Thus the natural map $F_X\to H$ is an isomorphism on abelianizations. It is onto: a proper closed subgroup of a pro-$p$ group lies in a proper open subgroup and then in a maximal open subgroup, so cannot surject on the Frattini quotient. Since $H$ is free pro-$p$, projectivity supplies a section. Its closed image surjects on $F_X/\Phi(F_X)$ and must therefore be all of $F_X$. The map is an isomorphism. This argument uses neither infinite-rank Hopficity nor an assumed interchange of abelianization and inverse limits.

The compatible finite-rank Magnus embeddings consequently give an injective measure-valued Magnus expansion of $H$. Assign degree $1$ to the orbit letters and degree $2$ to $y$. Every nontrivial element has a first finite nonzero degree. In any degree, each word pattern has a bounded $\mathbb Z_p$-measure as coefficient in the orbit parameters of its letters of degree $1$.

\paragraph{A quadratic meromorphic moment.}
Suppose there is a continuous dense map $\rho:H\to U(\mathbb Q_p)$ into the actual rational tripod image, whose Lie algebra is the completed free Lie algebra on odd generators $s_m$, $m\ge3$. Suppose $\gamma$ is represented by a Levi element with cyclotomic value $u$, a topological generator of $1+p\mathbb Z_p$. Thus it scales $s_m$ by $u^m$, and $\Delta$ acts by character $m$.

Let $\ell_m$ be the $m$th abelian coordinate. Equivariance gives
\[
 \ell_m(\rho(x_{a,t}))=
 \begin{cases}a_m u^{mt},&m\equiv a\pmod{p-1},\\0,&\text{otherwise}.
 \end{cases}
\]
Each $a_m$ is nonzero by density and the compact generating space: $y$ has no odd linear coordinate by parity. Rescale each free letter by $a_m$, weight by weight over $\mathbb Q_p$. The linear part of $\rho(x_{a,t})-1$ now has coefficient $u^{mt}$ on its matching letters. No uniformly integral change of lattice is asserted.

Set
\[
 E=\gamma y\gamma^{-1}y^{-\alpha}\in\overline{[H,H]}.
\]
Its degree-two coefficients involving two orbit letters give bounded measures $\mu_{a,b}$ and transforms
\[
 C_{a,b}(T_1,T_2)=
 \int_{\mathbb Z_p^2}(1+T_1)^t(1+T_2)^v\,d\mu_{a,b}(t,v)
 \in\mathbb Z_p[[T_1,T_2]].
\]
Assume that for one ordered pair $(a_0,b_0)$,
\[
 C_{a_0,b_0}(0,0)\in\mathbb Z_p^\times.                         \tag{N.2}
\]
We will verify this hypothesis arithmetically at $37$.

Let $B_{m,n}$ be the coefficient of $s_ms_n$ in $\rho(y)-1$. Since the latter has no length-one part, comparison in length two yields
\[
 (u^{m+n}-\alpha)B_{m,n}
   =C_{a,b}(u^m-1,u^n-1),\qquad a\equiv m,\ b\equiv n\pmod{p-1}.
                                                                    \tag{N.3}
\]
Indeed the two sides of $E=\gamma y\gamma^{-1}y^{-\alpha}$ give the factor $u^{m+n}-\alpha$, while only the two-orbit-letter Magnus terms can contribute to target length two. Terms of $\overline{[H,H]}$ involving $y$ have assigned degree at least three.

Write
\[
 D(T_1,T_2)=(1+T_1)(1+T_2)-\alpha.
\]
Equation (N.3) represents the length-two kernel of $y$ by $C_{a,b}/D$. Divisions will only be performed in a fraction field after clearing denominators. The numerator in (N.2) is a unit. The divisor is genuine in the open unit polydisc, since $T_2=\alpha/(1+T_1)-1$ and $\alpha-1\in p\mathbb Z_p$.

\paragraph{Ordered tilings and pole separation.}
Suppose $g\ne1$ belongs to $\ker\rho$, and let $d$ be its first nonzero assigned Magnus degree. A source word of degree $d$ specifies a subset
\[
 J\subseteq\{1,\ldots,d-1\},\qquad
 \text{no two elements of }J\text{ adjacent}.
\]
Each $j\in J$ denotes one $y$ filling target positions $j,j+1$; the other positions are singleton orbit letters. For each pattern and its singleton types, let $A_J$ be the Amice transform of its coefficient measure. It depends only on the unpaired variables.

Fix the residue classes of an ordered target word of length $d$. Its vanishing coefficient gives, at $T_i=u^{m_i}-1$,
\[
 \sum_J A_J(T_i:i\text{ unpaired})
       \prod_{j\in J}
       \frac{C_{a_j,a_{j+1}}(T_j,T_{j+1})}{D_j}=0,
 \qquad D_j=(1+T_j)(1+T_{j+1})-\alpha.                         \tag{N.4}
\]
There are finitely many tilings. Multiply by all $D_j$, $1\le j<d$. The resulting integral analytic identity holds even at grid points where a denominator vanishes, by (N.3). Smaller source degrees vanish by choice of $d$, and larger ones have target length greater than $d$; thus no omitted source terms contribute.

The finite-weight substitution is valid exactly as in the pro-$2$ proof. At every arithmetic weight cutoff only finitely many source degrees contribute. The coefficient functions on compact orbit spaces are bounded and continuous; integration, uniform locally constant approximation, and Fubini give a continuous multiplicative map. It agrees with $\rho$ on the generators and hence on $H$.

If $\alpha=u^w$ for a positive integer $w$, restrict all test weights to exceed $w$; otherwise no exclusion is needed. This removes every possible zero denominator on the test grid. For each odd residue class $a$, the remaining tails $m=a+(p-1)p^n$ still give distinct points converging to $u^a-1$ in $p\mathbb Z_p$. Applying the analytic identity theorem successively in all variables makes the denominator-cleared version of (N.4) an identity of power series. Hence (N.4) holds in the fraction field of
\[
 R=\mathbb Z_p[[T_1,\ldots,T_d]][1/p].
\]
This ring has a uniform bound on coefficient denominators. Replacing it by $\mathbb Q_p[[T_1,\ldots,T_d]]$ would destroy the argument by allowing inverses of $D_j$ with unbounded denominators.

Choose a nonzero coefficient pattern with the maximum number of $y$ letters, and denote its tiling by $J$. Assign target types $(a_0,b_0)$ at each of its pairs and the specified source type at each singleton. This fixes an identity (N.4) with $A_J\ne0$. Consider
\[
 I_J=(D_j:j\in J)\subset R.
\]
This is prime: its disjoint equations substitute $T_{j+1}=\alpha/(1+T_j)-1$, leaving a power-series ring in the unpaired variables and one variable from each pair, with $p$ inverted. The substitutions are legitimate in the $(p,T)$ topology because their constant term is in $p\mathbb Z_p$.

For $k\notin J$, $D_k\notin I_J$. Its neighboring positions belong to different blocks; the remaining variables of these blocks are independent, and their product is not identically $\alpha$. Thus $D_k$ is a unit in $R_{I_J}$. Multiply (N.4) by $\prod_{j\in J}D_j$ and reduce modulo $I_J$ in this localization. A tiling omitting any pair of $J$ vanishes. A tiling properly containing $J$ has more $y$ letters and its coefficient is zero by maximality. Only $J$ remains:
\[
 A_J\prod_{j\in J}C_{a_0,b_0}(T_j,T_{j+1})=0
 \quad\text{in }\operatorname{Frac}(R/I_J).
\]
Every $C$ factor is a unit. The unpaired-variable ring embeds in $R/I_J$, so $A_J$ remains nonzero, a contradiction. If $J$ is empty, maximality leaves just its singleton coefficient, giving the same contradiction directly. This proves injectivity under (N.1)--(N.2).

\paragraph{The arithmetic module at $37$.}
Now set $p=37$, $F=\mathbb Q(\mu_{37})$, $K=\mathbb Q(\mu_{37^\infty})$, $P=\operatorname{Gal}(\Omega_{37}/F)$, and $H=\operatorname{Gal}(\Omega_{37}/K)$. Choose commuting lifts of $\Delta=\operatorname{Gal}(F/\mathbb Q)$ and $\Gamma=\operatorname{Gal}(K/F)$ using Sharifi's procyclic construction \cite[\S2, p.~3]{Sharifi2002}. If the chosen generator has cyclotomic value $u_0$, replace it by its $v$th power, where $v=\log(1+37)/\log(u_0)\in\mathbb Z_{37}^{\times}$. The resulting $\gamma$ still commutes with $\Delta$ and has $\chi(\gamma)=1+37$.

Iwasawa's theorem and Ferrero--Washington give freeness of $H$ \cite[Theorem~2, pp.~273--274]{Iwasawa1981}. Put $M=H^{\mathrm{ab}}$. Weak Leopoldt gives $\operatorname{rank}_{\Lambda}M=18$. Kummer duality gives that $M^+$ is $\Lambda$-torsion \cite[Theorems~3.3.3--3.3.4, Proposition~3.4.7]{SharifiIwasawaNotes}. Also $M$ is $\mathbb Z_{37}$-torsion-free because $H$ is free pro-$37$.

The minimal generators in McCallum--Sharifi, \S8, pp.~22--23, consist of $\gamma$, eighteen odd $x_m$ for $3\le m\le37$, and a single even generator $y$ of character $32$ \cite{McCallumSharifi2003}. Vandiver at $37$ and its unique irregular index $32$ were checked above. The exact sequence
\[
 0\longrightarrow M_{\Gamma}\longrightarrow P^{\mathrm{ab}}
   \longrightarrow\Gamma\longrightarrow0
\]
follows from the split semidirect product $P=H\rtimes\Gamma$. Consequently $M/(37,T)M$ has exactly eighteen odd directions and one even direction. Nakayama gives eighteen generators for $M^-$. Its rank is eighteen, since $M^+$ has rank zero. A surjection $\Lambda^{18}\to M^-$ has rank-zero kernel, which must vanish as a submodule of a free module over the domain $\Lambda$. Thus
\[
 M^-\simeq\bigoplus_{a\text{ odd}\bmod36}\Lambda e_a.           \tag{N.5}
\]

Similarly $M^+$ is cyclic, generated by the class of $y$, and is nonzero because its quotient modulo $(37,T)$ has dimension one. The explicit relation in McCallum--Sharifi, Corollary~8.3, p.~23, begins
\[
 37y-3[\gamma,y]-11[x_3,x_{29}]-[x_7,x_{25}]+\cdots=0
                                                                  \tag{N.6}
\]
in descending $p$-central degree two. Linearizing the full relator in the cyclotomic kernel and projecting onto the plus part gives an annihilator
\[
 f(T)y=0,\qquad f(T)\equiv37-3T\pmod{(37,T)^2}.                \tag{N.7}
\]
Here is the filtration justification for the coefficient congruence. Before imposing arithmetic relations, linearize the finite-rank free presentation group into $A_{\mathrm{free}}\rtimes\Gamma$, where $A_{\mathrm{free}}$ is the free $\Lambda$-module on the non-cyclotomic generators. Its filtration
\[
 (p,T)^{n-1}A_{\mathrm{free}}\rtimes p^{n-1}\Gamma
\]
respects commutators and $p$th powers; use
$(1+T)^{p^{n-1}}-1\in(p,T)^n$. Thus descending $p$-central degree three contributes only $(p,T)^2$ to each coefficient. Read the $y$ coefficient there, then pass to $M^+$. The odd commutators vanish and odd linear terms disappear in this projection.

Weierstrass preparation makes $f$ a unit times $T-\beta$, where $\beta\in37\mathbb Z_{37}$ and $v_{37}(\beta)=1$. Hence $M^+$ is a quotient of $\Lambda/(f)\simeq\mathbb Z_{37}$. It is nonzero and $\mathbb Z_{37}$-torsion-free, so
\[
 M^+\simeq\mathbb Z_{37},\qquad\gamma e_y=(1+\beta)e_y.        \tag{N.8}
\]
Equations (N.5)--(N.8) prove (N.1), with $\alpha=1+\beta$.

The lifts can preserve the preceding Frattini classes: project their abelian classes with the integral $\Delta$-idempotents and lift back through $H\to M$. The projections fix their prescribed residual characters. Each odd projected class generates its free rank-one $\Lambda$-eigenspace, and the even class generates $M^+$. A change of finite free presentation generators with the same Frattini classes is induced by an automorphism: the endomorphism is surjective by the Frattini criterion and injective by Hopficity of a finitely generated profinite group. Such a change fixes descending $p$-central degree two: both $p$th powers and commutators retain their classes. Unit rescalings within the one-dimensional eigenspaces only rescale the displayed nonzero coefficients by units. Thus these changes preserve the information used from (N.6).

\paragraph{Recovering the unit quadratic numerator.}
For the compact basis supplied by (N.1), write $\varphi$ for the $\Gamma$-automorphism of $H$ and $E=\varphi(y)y^{-\alpha}$. The group $P$ has the presentation on $\gamma,x_a,y$ with one defining relation
\[
 R=\gamma y\gamma^{-1}\varphi(y)^{-1}.
\]
To verify the presentation, assign images to these generators in any pro-$p$ target satisfying $R$. Send the compact orbit generator $x_{a,t}$ to $\gamma^t x_a\gamma^{-t}$. The compact free universal property gives a map from $H$; its $\Gamma$-compatibility is automatic on the orbits and is exactly $R$ on $y$. Hence it extends to $H\rtimes\Gamma$.

In descending $p$-central degree two the relation has class
\[
 [\gamma,y]+(1-\alpha)y-E_2,
\]
where $E_2$ is the quadratic Magnus term after folding the orbit parameters to their augmentations. Its ordered $x_a x_b$ coefficient is, up to the commutator sign convention, $C_{a,b}(0,0)$. Because there is one defining relation, any relation's degree-two initial term is a scalar multiple of this one: conjugation fixes that class and $p$th powers of the relator vanish in that degree. No assertion about freeness of the full relation module is required.

In (N.6) the coefficients of $[\gamma,y]$ and $[x_3,x_{29}]$ are both units modulo $37$. It follows that $C_{3,29}(0,0)$ is a unit. This is (N.2), with $(a_0,b_0)=(3,29)$.

\paragraph{The actual rational image and conclusion.}
Hain--Matsumoto's actual-image comparison and Brown's rational Deligne--Ihara theorem supply the dense map $H\to U(\mathbb Q_{37})$ and its completed free odd Lie algebra \cite[Proposition~8.3, Theorem~8.4]{HainMatsumoto2001}\cite[\S2.1--2.3]{Brown2019}. Arrange that the image of $\gamma$ lies in a Levi torus by successively solving in positive weights, since $u^w-1\ne0$ in $\mathbb Q_{37}$ for every $w>0$. This converges in the weight completion. The commuting finite group $\Delta$ then also lies in that torus: its possible unipotent part would centralize $\gamma$, whose unipotent centralizer is trivial by the same weight argument.

All hypotheses of the separation argument are now verified, so the actual rational action is injective on $H$. The kernel of the pro-$37$ tripod action lies in $H$ because abelianization records the cyclotomic character, and it is trivial under the induced rational action. Thus it is trivial on $P$, proving Theorem~\ref{thm:kernel37-poles}.

This proof deals with the single even arithmetic direction at $37$. It does not supply an all-irregular-primes theorem. With several even directions or higher $\Lambda$-torsion rank, the corresponding matrix and higher-order pole separation must be proved separately.

% CHECKED DEDUCTION: full self-reading, two independent whole-proof audits, and parent review passed.
% Existing checked p=2 and p=37 arguments remain separate.
\subsection{Several even directions and the range \texorpdfstring{$p<1000$}{p<1000}}
\label{sec:multieven-poles}

The preceding pole argument extends to several even arithmetic directions and to characteristic polynomials with repeated roots. The arithmetic hypothesis is most cleanly stated directly in Galois cohomology.

\begin{theorem}[K17-POLE-GENERAL]\label{thm:pole-general}
Let $p$ be an odd prime satisfying Vandiver, let $F=\mathbb Q(\mu_p)$, and let $P$ be the Galois group of its maximal pro-$p$ extension unramified outside $p$. Assume that the cup products of the odd-character summand of $H^1(P,\mathbb F_p)$ span $H^2(P,\mathbb F_p)$. Then
\[
 \Omega_p^*=\Omega_p.
\]
No assumption that the relevant Iwasawa polynomials have degree one, split over $\mathbb Q_p$, or have distinct roots is required.
\end{theorem}

Here odd means the $-1$ eigenspace for complex conjugation, equivalently the sum of the odd Teichm\"uller characters. Cup/relator duality says that the hypothesis supplies a nonzero odd--odd commutator coefficient in each nonzero even relation eigenspace. This is a statement about the given arithmetic group and its full cup product; it does not require a chosen basis of cyclotomic units.

\paragraph{The arithmetic module and its companion bases.}
Set $K=\mathbb Q(\mu_{p^\infty})$, $H=\operatorname{Gal}(\Omega_p/K)$, $\Gamma=\operatorname{Gal}(K/F)=\langle\gamma\rangle$, $\Lambda=\mathbb Z_p[[\Gamma]]$, and $M=H^{\mathrm{ab}}$. Choose $\gamma$ commuting with a lift of $\Delta=\operatorname{Gal}(F/\mathbb Q)$ and with $\chi(\gamma)=u=1+p$, as above.

Iwasawa's theorem and Ferrero--Washington make $H$ free pro-$p$, so $M$ is $\mathbb Z_p$-torsion-free. Weak Leopoldt gives $\operatorname{rank}_{\Lambda}M=(p-1)/2$. Kummer duality makes $M^+$ a torsion $\Lambda$-module; it is finitely generated over $\mathbb Z_p$ because it is pseudo-isomorphic to $(X_\infty^-)^\iota(1)$, whose $\mu$-invariant is zero. Finite kernels and cokernels preserve this finite generation. Thus $M^+$ is finite free over $\mathbb Z_p$ \cite[Theorem~2]{Iwasawa1981}\cite[Theorems~3.3.3--3.3.4, Proposition~3.4.7, Corollary~3.4.8]{SharifiIwasawaNotes}.

Under Vandiver the minimal generators of $P$ in McCallum--Sharifi, \S8, pp.~22--23, are $\gamma$, one $x_a$ for each odd residue class $a$, and one $y_e$ for each occurring even character $e$ \cite{McCallumSharifi2003}. From
\[
 0\longrightarrow M_\Gamma\longrightarrow P^{\mathrm{ab}}
  \longrightarrow\Gamma\longrightarrow0
\]
and Nakayama, $M^-$ has $(p-1)/2$ generators over $\Lambda$ and the same rank. The kernel of the resulting surjection from a free module has rank zero and is zero. Hence
\[
 M^-\simeq\bigoplus_{a\text{ odd}\bmod p-1}\Lambda e_a.       \tag{G.1}
\]
Each nonzero even component $T_e=M^{(e)}$ is cyclic over $\Lambda$ and finite free over $\mathbb Z_p$. Put $\lambda_e=\operatorname{rank}_{\mathbb Z_p}T_e$ and choose a cyclic generator $v_e$. If $q_e(z)$ is the characteristic polynomial of $\gamma$ on $T_e$, then
\[
 q_e(z)\equiv(z-1)^{\lambda_e}\pmod p.                         \tag{G.2}
\]
Indeed the continuous action of the pro-$p$ group $\Gamma$ on $T_e/pT_e$ has finite $p$-group image and is unipotent.

The vectors $v_e,\gamma v_e,\ldots,\gamma^{\lambda_e-1}v_e$ form a $\mathbb Z_p$-basis. Cayley--Hamilton makes their span stable under $\gamma$ and $\gamma^{-1}$, since $q_e$ has unit constant coefficient. The span is compact and closed, so contains the entire $\Lambda$-span of $v_e$. It therefore surjects onto $T_e$, and a surjection between free $\mathbb Z_p$-modules of the same finite rank is an isomorphism.

Choose actual lifts $x_a,y_e\in H$ and set
\[
 x_{a,t}=\gamma^t x_a\gamma^{-t},\quad t\in\mathbb Z_p,
 \qquad y_{e,j}=\gamma^j y_e\gamma^{-j},\quad0\le j<\lambda_e.
\]
These constitute a free compact basis of $H$. To recall the complete argument, take the inverse limit of the free pro-$p$ groups on the finite quotients of this compact generating space. Its abelianization is precisely (G.1) together with the indicated finite bases of the $T_e$. The natural map to $H$ is an isomorphism on abelianizations, hence is onto by the Frattini criterion for closed subgroups. Projectivity of $H$ supplies a section; its closed image is Frattini-surjective in the source and hence equals the source. The map is an isomorphism.

We may make the lifts exact nonlinear $\Delta$-eigenlifts by Sharifi's averaging lemma \cite[Lemma~2.1, pp.~3--4]{Sharifi2002}:
\[
 \delta x_a\delta^{-1}=x_a^{\omega(\delta)^a},\qquad
 \delta y_e\delta^{-1}=y_e^{\omega(\delta)^e}.
\]
This preserves the prescribed abelian eigenspace classes. Commutation of $\gamma$ and $\Delta$ gives the same eigenproperty for all their indicated conjugates.

\paragraph{The companion error and its cup coefficient.}
Write
\[
 q_e(z)=z^{\lambda_e}-\sum_{j=0}^{\lambda_e-1}c_{e,j}z^j
\]
and, with the product taken in increasing $j$, define
\[
 E_e=(\gamma^{\lambda_e}y_e\gamma^{-\lambda_e})
          \left(\prod_jy_{e,j}^{c_{e,j}}\right)^{-1}.
\]
The characteristic-polynomial relation in $M$ puts $E_e$ in the closed commutator subgroup of $H$. Its quadratic Magnus coefficients in two odd orbit letters give bounded measures and integral Amice transforms $C_{e;a,b}(T_1,T_2)$. We claim that for each $e$, there is an ordered odd pair $(a_e,b_e)$ such that
\[
 C_{e;a_e,b_e}(0,0)\in\mathbb Z_p^\times.                       \tag{G.3}
\]
The following relation-space argument supplies this comparison even when there are several even characters.

Let $F_{\mathrm{free}}$ be free pro-$p$ on $\gamma,x_a,y_e$, retaining only $j=0$ as an even generator. Substitute the actual conjugate words for all $y_{e,j}$ and $x_{a,t}$ in the compact free expansion of $E_e$. The group $P=H\rtimes\Gamma$ has one defining relator per even character:
\[
 R_e=\gamma^{\lambda_e}y_e\gamma^{-\lambda_e}
       \left(\prod_jy_{e,j}^{c_{e,j}}\right)^{-1}E_e^{-1}.       \tag{G.4}
\]
This is verified by the universal property. An assignment of $\gamma,x_a,y_e$ into a pro-$p$ target satisfying (G.4) gives images of all compact generators of $H$. Compatibility with $\gamma$ is automatic on the odd orbits and along each finite even chain except at its last element; the last compatibility is exactly (G.4). It therefore extends uniquely to the semidirect product.

Let $R$ be the closed relation subgroup. The exact eigenlifts let $\Delta$ act on $F_{\mathrm{free}}$ by fixing $\gamma$ and raising each other generator to its prescribed Teichm\"uller unit power. These are finite-order automorphisms; the quotient to $P$ is equivariant. Abelianize the kernel of $F_{\mathrm{free}}\to\Gamma$. Its free relative abelianization is
\[
 A_{\mathrm{free}}=
   \bigoplus_a\Lambda X_a\oplus\bigoplus_e\Lambda Y_e.
\]
The image of $R_e$ is exactly $q_e(\gamma)Y_e$. The word $E_e$ is a closed commutator word in the compact free basis, so has zero relative abelianization before imposing any arithmetic relations. Normal conjugation acts on this module through $\Gamma$. Hence the whole closed image is
\[
 R_{\mathrm{lin}}=\bigoplus_eq_e(\gamma)\Lambda Y_e,             \tag{G.5}
\]
a free $\Lambda$-module, with its indicated character $e$ on each summand.

Linearization gives a surjection
\[
 R/\overline{R^p[R,F_{\mathrm{free}}]}
       \longrightarrow R_{\mathrm{lin}}/(p,T)R_{\mathrm{lin}}.  \tag{G.6}
\]
The source is generated by the number of defining relators and the target has exactly that dimension. Thus (G.6) is an equivariant isomorphism, and the class of $R_e$ has character $e$. The cup-normalized character-$e$ relation is therefore a nonzero scalar multiple of $R_e$ in this quotient. All relators lie in the Frattini subgroup; passing to descending $p$-central degree two annihilates $R^p$ and $[R,F_{\mathrm{free}}]$.

The companion product in (G.4) contains only $y_e$ and its $\gamma$-conjugates and contributes no odd--odd bracket in that degree. Its odd--odd coefficient is consequently minus the augmentation coefficient of $E_e$, namely $-C_{e;a,b}(0,0)$ up to the commutator convention. The assumed odd cup surjectivity and the usual minimal-presentation cup/relator duality prove (G.3) \cite[\S8, pp.~21--23]{McCallumSharifi2003}. No freeness or asphericity claim about the full relation module is being used.

\paragraph{Quadratic target coefficients and partial fractions.}
Use the actual continuous dense rational representation $H\to U(\mathbb Q_p)$ and the completed free odd Lie algebra supplied by Hain--Matsumoto and Brown \cite{HainMatsumoto2001,Brown2019}. Conjugate rationally so that $\gamma$ is in a Levi torus; the successive positive-weight equations are solvable because $u^w-1\ne0$. The commuting $\Delta$ then also acts by its weight characters. Normalize each odd free letter $s_m$ so that the linear coefficient of $\rho(x_{a,t})-1$ on its matching residue class is $u^{mt}$. The necessary scalars are nonzero by density. Every $y_{e,j}$ has target generator length at least two by parity.

For an ordered pair $(m,n)$ put $z=u^{m+n}$, and let $B_{e,j;m,n}$ be the coefficient of $s_ms_n$ in $\rho(y_{e,j})-1$. The actual conjugate definition gives
\[
 B_{e,j;m,n}=z^jB_{e,0;m,n},
 \qquad
 q_e(z)B_{e,0;m,n}=C_{e;a,b}(u^m-1,u^n-1),                    \tag{G.7}
\]
where $a\equiv m$, $b\equiv n\pmod{p-1}$. The second equality is the length-two component of $E_e$. These coefficients have character $e$, so vanish when $a+b\not\equiv e$. For those residue classes (G.7) makes $C_{e;a,b}$ vanish on the entire Cartesian weight grid; the successive analytic identity theorem forces $C_{e;a,b}=0$ as a power series. Thus wrong-character numerators vanish identically, and the unit pair in (G.3) has $a_e+b_e\equiv e$.

Before dividing, exclude all exceptional weight sums. There are finitely many positive integers $w$ such that $u^w$ is a root of some $q_e$, since $u$ has infinite order. Restrict every test weight $m$ to exceed all such $w$. Then all pair denominators are nonzero. In each fixed odd residue class, the tails $m=a+(p-1)p^n$ still accumulate $p$-adically at $a$, and suffice for every analytic identity below.

Let $L/\mathbb Q_p$ be a finite extension splitting all $q_e$, with valuation ring $\mathcal O_L$ and uniformizer $\pi$. Every root $\alpha$ satisfies $\alpha-1\in\pi\mathcal O_L$ by (G.2). If its multiplicity is $n_{e,\alpha}$, ordinary partial fractions give the two bases of the same space of strictly proper rational functions:
\[
 \left\{\frac{z^j}{q_e(z)}:0\le j<\lambda_e\right\},
 \qquad
 \left\{\frac1{(z-\alpha)^k}:1\le k\le n_{e,\alpha}\right\}_{\alpha}.
\]
Apply the resulting constant invertible $L$-linear change to the degree-two Magnus letters $y_{e,j}$. This is a linear change in the associative coefficient algebra, not a change of actual pro-$p$ generators. At each fixed degree, only finitely many scalar coefficients occur, so measures retain a finite common denominator bound. The first nonzero assigned degree is preserved.

In the new letters $Y_{e,\alpha,k}$, the leading target kernels are
\[
 \frac{C_{e;a,b}(T_1,T_2)}
      {\bigl((1+T_1)(1+T_2)-\alpha\bigr)^k}.                   \tag{G.8}
\]
All even letters have assigned degree two and the orbit letters degree one. The measure-valued Magnus map remains injective after scalar extension and this invertible homogeneous change.

\paragraph{Colored pair tilings and Laurent coefficient separation.}
Suppose $g\ne1$ is in the kernel and let $d$ be its first nonzero assigned degree. A coefficient pattern is an ordered tiling of $d$ target positions by singleton odd letters and paired even letters. Each pair is colored by $(e,\alpha,k)$. For fixed target residue classes, its vanishing target coefficient gives
\[
 \sum_{\mathcal P} A_{\mathcal P}(T_i:i\text{ singleton})
 \prod_{(j,e,\alpha,k)\in\mathcal P}
 \frac{C_{e;a_j,a_{j+1}}(T_j,T_{j+1})}{D_{j,\alpha}^k}=0,
 \quad D_{j,\alpha}=(1+T_j)(1+T_{j+1})-\alpha.                \tag{G.9}
\]
All $A_{\mathcal P}$ are bounded $L$-valued measure transforms. Only the indicated first source degree contributes: a larger degree has larger target length, and every smaller degree is zero. Substitution at each finite arithmetic weight cutoff is continuous and multiplicative by integration of bounded continuous coefficient functions and uniform locally constant approximation, as in the preceding complete proofs.

There are finitely many patterns and denominators. Clearing denominators in (G.9) gives a bounded-coefficient analytic identity on the full Cartesian tails of the allowed weight grids. Successive one-variable identity theorems give an identity of power series over $L$. Its bounded-denominator ring is
\[
 R_L=\mathcal O_L[[T_1,\ldots,T_d]][1/\pi],
\]
not the ring permitting arbitrary unbounded coefficient denominators.

Choose a nonzero pattern with the maximum number of even letters. Fix its pair positions $J$ and its $(e,\alpha)$ colors. Among nonzero patterns with those positions and colors choose a pole-order vector $(k_j)$ maximal in the componentwise partial order. This is possible because there are finitely many patterns. Assign the unit target pair $(a_e,b_e)$ from (G.3) at each chosen pair and its source orbit type at each singleton.

For $j\in J$ introduce $S_j=D_{j,\alpha_j}$. The pairs are disjoint and the substitutions
\[
 T_{j+1}=\frac{\alpha_j+S_j}{1+T_j}-1
\]
give an isomorphism of formal $\mathcal O_L$ power-series rings: their constant terms are topologically nilpotent. Retain the unpaired variables and one left variable from each pair as base variables. Let $K_0$ be the fraction field of their bounded-denominator power-series ring. Equation (G.9) embeds into Laurent series over $K_0$ in the $S_j$, with finite lower bounds on all exponents.

Every other denominator is analytic in these variables. At the same pair position a different root gives $S_j+(\alpha_j-\alpha')$ with nonzero constant in $K_0$. At a different position its two entries lie in different blocks, whose remaining base variables are independent; their product is not identically its root. Its constant term is again nonzero in $K_0$. These denominators are therefore invertible in $K_0[[S_j:j\in J]]$. A different even character at a selected position contributes zero for the assigned target types, since their sum fixes $e$. Each selected $C$ remains a nonzero unit at $S=0$.

Extract the coefficient of $\prod_{j\in J}S_j^{-k_j}$. A pattern missing a selected pair or using a different root there has no negative exponent in that variable. A pattern properly containing $J$ has more even letters and zero coefficient by maximality. For patterns with the same positions and root colors, an order vector $k'$ can contribute only if $k'_j\ge k_j$ for every $j$, since all remaining series have nonnegative exponents. Componentwise maximality excludes every strict such vector. Thus only the selected pattern contributes, and its coefficient is
\[
 A_{\mathcal P}\prod_{j\in J}
 C_{e_j;a_{e_j},b_{e_j}}\big|_{S=0}\ne0.
\]
The unpaired-variable coefficient ring embeds in $K_0$, and all factors are units. This is a contradiction. When there are no even pairs, the same calculation reduces to odd-moment separation. Hence the actual rational action is injective on $H$. Cyclotomic abelianization detects any kernel outside $H$, and the rational action is induced by the pro-$p$ outer action. Theorem~\ref{thm:pole-general} follows.

\paragraph{K17-CUP-PARITY: from $p$-units to odd Galois generators.}
Let $A_F$ be the $p$-class group and
\[
 E=\mathcal O_F[1/p]^\times/(\mathcal O_F[1/p]^\times)^p.
\]
For a $p$-unit $v$, $v/c(v)$ is a global unit with all complex conjugates of absolute value one, hence a root of unity. An elementary proof is that the monic polynomials of its positive powers have bounded degree and bounded integer coefficients, so only finitely many roots occur and two powers coincide. Thus $E^-$ is exactly the one-dimensional space generated by $\mu_p$; it is nonzero because $F$ does not contain $\mu_{p^2}$.

The Kummer sequence and Vandiver give
\[
 0\longrightarrow E\longrightarrow H^1(P,\mu_p)
  \longrightarrow A_F[p]\longrightarrow0,
 \qquad E^+=H^1(P,\mu_p)^+.
\]
There is no Brauer contribution, since $F$ is totally imaginary with just one omitted finite prime, and
\[
 H^2(P,\mu_p^{\otimes2})=A_F\otimes\mu_p
\]
is a plus space. Mixed parity cups therefore vanish. The alternating cup on the one-dimensional $E^-$ vanishes because $p$ is odd. Surjectivity of the cup on all $p$-units consequently implies surjectivity from $E^+\otimes E^+$. Untwisting gives precisely the odd $H^1(P,\mathbb F_p)$ cup hypothesis of Theorem~\ref{thm:pole-general}. Under cup/relator duality, $x_a$ has character $\omega^a$, its Kummer dual has character $\omega^{1-a}$, and a relation of character $\omega^e$ corresponds to twisted target character $\omega^{2-e}$. Thus the indices satisfy $a+b\equiv e\pmod{p-1}$, as used above.

\begin{corollary}[K17-KERNEL-LT1000]\label{cor:kernel-lt1000}
For every odd prime $p<1000$, one has $\Omega_p^*=\Omega_p$.
\end{corollary}
\begin{proof}
Buhler--Harvey verify Vandiver for every prime below $163{,}577{,}356$, in particular for this entire range \cite[abstract and \S1, p.~1]{BuhlerHarvey2011}. Sharifi's Theorem~1.2 and Corollary~5.9 prove that the $p$-unit cup pairing is surjective for every odd $p<1000$ \cite[pp.~2 and~27]{Sharifi2007}. The parity lemma supplies the exact odd-generator hypothesis, so Theorem~\ref{thm:pole-general} applies.
\end{proof}

The cited computational theorems are used as published results: Sharifi's Theorem~5.7 is an explicit Magma modular-symbol computation, and the Vandiver range is the published Buhler--Harvey verification. These computations were not rerun here. The deduction imposes no separate numerical $\lambda$-invariant or simple-root condition. For an arbitrary remaining prime, the explicit sufficient conditions are Vandiver and the stated odd cup surjectivity. This argument does not settle their validity for all primes, and it does not prove integral odd generation of the image Lie algebra.

% CHECKED extension: complete self-reading and two independent refinement audits passed.
\subsection{A weaker remaining arithmetic condition}
\label{sec:pole-nonvanishing}

The proof of Theorem~\ref{thm:pole-general} uses the unit cup coefficient only to ensure that a selected numerator stays nonzero on a pole divisor. This gives a weaker sufficient condition, expressed in the actual quadratic companion errors.

Retain the construction of $H$, its compact basis, $q_e(z)$, and the integral numerator functions $C_{e;a,b}(T_1,T_2)$ from the preceding proof. That construction, up to the assertion (G.3), uses Vandiver but no cup-surjectivity hypothesis. Let $L/\mathbb Q_p$ split the finitely many polynomials $q_e$, let $\mathcal O_L$ be its valuation ring, and let $\pi$ be a uniformizer. Every root $\alpha$ lies in $1+\pi\mathcal O_L$.

\begin{theorem}[K17-POLE-NV]\label{thm:pole-nv}
Let $p$ be an odd prime satisfying Vandiver. Suppose that for every occurring even character $e$ and every distinct root $\alpha$ of $q_e$, there is an ordered pair of odd characters $(a,b)$, with $a+b\equiv e\pmod{p-1}$, such that
\[
 C_{e;a,b}\!\left(T,\frac{\alpha}{1+T}-1\right)\ne0
 \quad\text{in }\mathcal O_L[[T]][1/\pi].                    \tag{NV}
\]
Then $\Omega_p^*=\Omega_p$. The pair $(a,b)$ may depend on $\alpha$, and no assumption on the multiplicities of the roots is needed.
\end{theorem}

Equivalently, for each $(e,\alpha)$ at least one matching numerator is not divisible by
\[
 D_\alpha=(1+T_1)(1+T_2)-\alpha
\]
in $\mathcal O_L[[T_1,T_2]][1/\pi]$. The equivalence follows from the formal change of variables $S=D_\alpha$ and $T_2=(\alpha+S)/(1+T_1)-1$; its constant displacement is topologically nilpotent. Thus the condition concerns restriction to the divisor, rather than just nonvanishing somewhere in the two-variable domain.

\begin{proof}
All the arithmetic module, compact free basis, actual rational realization, companion identity, and partial-fraction steps of the preceding proof remain valid without (G.3). They give degree-two formal letters with leading target kernels
\[
 \frac{C_{e;a,b}(T_1,T_2)}{D_\alpha^{\,k}},
 \qquad 1\le k\le n_{e,\alpha}.
\]
The finite scalar extension and homogeneous change of letters preserve the bounded-denominator coefficient measures and the first nonzero assigned Magnus degree. Restrict the test weights to the same sufficiently large Cartesian tails to avoid all exceptional denominator zeros.

Suppose a nontrivial kernel element has first assigned degree $d$. Among its nonzero colored tilings choose one with the maximum number of pairs. Fix its pair positions and $(e,\alpha)$ colors, and then choose a componentwise maximal pole-order vector $(k_j)$. At each chosen pair use the odd target types supplied by (NV) for that particular $(e_j,\alpha_j)$; at each singleton retain its specified source orbit type.

The vanishing target coefficients and the successive analytic identity theorem give exactly (G.9) in the bounded-denominator fraction ring. At the disjoint selected pairs put $S_j=D_{j,\alpha_j}$ and retain the unpaired variables and one left variable from each pair as base variables. Let $K_0$ be the fraction field of the bounded-denominator power-series ring in these base variables. The same denominator classification embeds (G.9) in
\[
 K_0[[S_j:j\in J]][S_j^{-1}:j\in J].
\]
Other-root denominators at a selected pair have nonzero constant in $K_0$; denominators at different positions connect different independent blocks and also have nonzero constant. Wrong-character numerators vanish identically, as already proved using the Cartesian weight grids.

Extract the coefficient of $\prod_{j\in J}S_j^{-k_j}$. A pattern missing a selected pair or using a different root there cannot supply the corresponding negative exponent. A pattern with additional pairs has zero coefficient by maximality of the pair count. Among patterns with the same positions and colors, only a componentwise greater or equal pole-order vector could contribute, and the chosen maximality excludes every strict such vector. Therefore the coefficient is precisely
\[
 A_{\mathcal P}\prod_{j\in J}
 C_{e_j;a_j,b_j}\!\left(T_j,\frac{\alpha_j}{1+T_j}-1\right).
\]
The selected coefficient $A_{\mathcal P}$ is nonzero. Each other factor is nonzero by (NV), the variables belonging to distinct pairs are independent, and the base power-series ring is a domain. Their product is consequently nonzero in $K_0$, contradicting the extracted zero coefficient. When there are no pairs, this is the already checked odd-moment separation argument.

Thus the actual rational action is injective on $H$. The cyclotomic character detects any kernel outside $H$, and the rational action is induced by the pro-$p$ outer action. This proves the asserted equality.
\end{proof}

\paragraph{Comparison with the cup condition.}
A numerator with unit constant coefficient is a unit in $\mathbb Z_p[[T_1,T_2]]$ and remains a unit under every indicated divisor substitution. Hence the cup condition in Theorem~\ref{thm:pole-general} implies (NV). The converse analytic implication fails: the formal numerator $C=p$ has no nonzero reduction modulo $p$, but restricts to a nonzero constant on every divisor. This is an example of analytic data, not an arithmetic prime. No additional prime satisfying (NV) but failing the cup hypothesis has been established here.

Repeated roots explain why the divisor qualification matters. A globally nonzero numerator divisible by $D_\alpha$ cancels a highest pole order; it need not suffice for the present separation argument. Condition (NV) rules out such cancellation on at least one odd pair for each $(e,\alpha)$.

\paragraph{Ordinary basis choices.}
For a fixed $\gamma$, replacing a cyclic lift $y_e$ by another lift of the same abelian class changes its length-two target coefficients by a bounded analytic commutator transform $A$. Consequently
\[
 B'_e=B_e+A,\qquad C'_e=C_e+q_e(z)A,
 \qquad z=(1+T_1)(1+T_2),
\]
so their restrictions at $z=\alpha$ agree. Replacing the cyclic generator by $b_e(\gamma)v_e$, where $b_e$ is a unit modulo $q_e$, instead gives
\[
 B'_e=b_e(z)B_e+A,\qquad
 C'_e=b_e(z)C_e+q_e(z)A.
\]
The restricted numerators differ by $b_e(\alpha)\ne0$. Changing an odd $\Lambda$ basis multiplies its normalized linear moments by a power-series unit; this divides the corresponding normalized numerator by unit factors in its two variables and again preserves (NV). Scalar extension preserves nonzero power series. These facts make the condition independent of these routine choices; the cyclotomic generator may simply remain fixed by $\chi(\gamma)=1+p$.

\paragraph{Arithmetic scope of the refinement.}
The criterion records exactly where the separation proof uses its
quadratic input. The next subsection supplies the stronger cup
condition at every Vandiver prime, so (NV) introduces no independent
arithmetic hypothesis in that range. A failure of (NV) would not by
itself prove a nontrivial kernel: a source direction might still be
detected at a larger target generator length. An argument beyond
Vandiver would also have to accommodate the different module structure.

% Verified chain from Vandiver to the cup hypothesis.
% Exact FK locators refer to the author preprint dated 15 August 2012.
\subsection{Vandiver alone supplies the cup hypothesis}
\label{sec:vandiver-kernel}

The extra cup hypothesis in Theorem~\ref{thm:pole-general} follows from
Vandiver by combining the integral results of Fukaya--Kato and
Wake--Wang-Erickson. The resulting statement is stronger than the numerical
range obtained from Sharifi's earlier cup computations.

\begin{theorem}[K17-VANDIVER-KERNEL]\label{thm:vandiver-kernel}
Let $p$ be an odd prime. If
\[
 p\nmid h\bigl(\mathbb Q(\mu_p)^+\bigr),
\]
then the odd-character cup products span $H^2(P,\mathbb F_p)$ and
\[
 \Omega_p^*=\Omega_p.
\]
The hypothesis is equivalently
\[
 K_{4k}(\mathbb Z)\{p\}=0
 \quad\text{for every }k\ge1,
\]
or equivalently the same vanishing for $1\le k\le(p-3)/2$.
These are equivalent formulations of the sufficient hypothesis;
no converse from the kernel-field equality is asserted.
\end{theorem}

\begin{proof}
For $p\ge5$, put $F_r=\mathbb Q(\mu_{p^r})$,
$F_\infty=\bigcup_rF_r$, and
\[
 X=\varprojlim_r\operatorname{Cl}(F_r)\{p\},\qquad
 \mathfrak X=\operatorname{Gal}(M/F_\infty),
\]
where $M$ is the maximal abelian pro-$p$ extension of $F_\infty$
unramified outside $p$. Thus $X$ is a quotient of $\mathfrak X$.
Let $\omega$ be the Teichm\"uller character and
$\Lambda=\mathbb Z_p[[\operatorname{Gal}(F_\infty/F_1)]]$.

First, Vandiver at $F_1$ implies $X^+=0$. Here is a proof of the
persistence statement needed for this implication. Suppose $L/K$ is a
cyclic extension of degree $p$ ramified at precisely one finite prime,
which is totally ramified, and suppose $\operatorname{Cl}(K)\{p\}=0$.
Let $E/L$ be the Hilbert $p$-class field. It is Galois over $K$, and
$G=\operatorname{Gal}(E/K)$ is a finite $p$-group. The inertia group $I$
at the ramified prime has order $p$, since $E/L$ is unramified. Let $N$
be its normal closure in $G$. All inertia groups over that prime are
conjugate, so $E^N/K$ is everywhere unramified. A nontrivial finite
$p$-group has a quotient of order $p$; class field theory and the assumed
class-group vanishing therefore give $G/N=1$. Thus $G$ is normally
generated by a single element. Its Frattini quotient has dimension at
most one over $\mathbb F_p$, so $G$ is cyclic. Consequently $N=I=G$
and $E=L$. This proves $\operatorname{Cl}(L)\{p\}=0$.
There is no extra infinite-prime ramification for odd $p$.

Apply this argument successively to $F_{r+1}^+/F_r^+$: each extension
has degree $p$ and its only ramified finite prime is the unique prime
above $p$. All the real-layer $p$-class groups are zero. Extension and
norm for the quadratic extension $F_r/F_r^+$ identify
$\operatorname{Cl}(F_r^+)\{p\}$ with
$\operatorname{Cl}(F_r)\{p\}^+$, because their compositions are
multiplication by $2$ and $1+c$, respectively. Taking inverse limits
gives $X^+=0$.

Fix an even character $\theta\ne1,\omega^2$ and set
$\chi=\omega^{-1}\theta$. The hypotheses in
Wake--Wang-Erickson, \S2.1, are then satisfied with tame level $N=1$.
The reflected character
\[
 \theta'=\omega^2\theta^{-1}=\omega\chi^{-1}
\]
is also even and different from $1,\omega^2$; its associated odd
character is $\chi'=\chi^{-1}$. Since $X_{\theta'}=0$, their
Proposition~2.2.6 says that
\[
 \mathfrak X_{(\chi')^{-1}}^{\#}(1)
   =\mathfrak X_\chi^{\#}(1)
\]
is free of rank one over the relevant Iwasawa algebra. The operation
$\#$ inverts the Galois action, and $(1)$ applies the cyclotomic twist.
They transport the coefficient ring by automorphisms, so undoing them
preserves freeness and cyclicity. Hence $\mathfrak X_\chi$ is cyclic,
and its quotient $X_\chi$ is cyclic. In addition $X_\theta=0$.
Corollaries~5.1.1--5.1.2 now show that the cuspidal Eisenstein Hecke
algebra is Gorenstein and that Sharifi's map
\[
 \varpi:H^-_\theta/I_\theta H^-_\theta
       \longrightarrow X_\chi(1)
\]
is an integral isomorphism
\cite[\S2.1, Proposition~2.2.6, Corollaries~5.1.1--5.1.2;
preprint pp.~2--4 and~11]{WakeWangErickson2017}.

For completeness, the relevant Fukaya--Kato input has no hidden
simple-root hypothesis. Gorensteinness implies their condition
$C(\mathfrak h)$, namely rational cyclicity of the Eisenstein quotient.
Put $Q=(H^-_\theta/I_\theta H^-_\theta)/(p\text{-primary torsion})$.
Over $\mathbb Q_p$, the module $Q$ is free of rank one over
$\Lambda/(\xi)[1/p]$, so its endomorphism $1-\Upsilon\varpi$
is multiplication by a scalar. Their derivative identity and
Lemma~10.2.2 make that scalar nilpotent.
The module $Q$ is $p$-torsion free, so this endomorphism is nilpotent
on $Q$ itself. Its finite geometric series makes
$\Upsilon\varpi$ an integral automorphism. After composing $\Upsilon$ with this inverse, one obtains a splitting
of the target Iwasawa module as $Q\oplus K$. The characteristic-ideal
and Fitting-ideal argument in their \S7.1.3 makes $K$ finite.
The target has no $p$-torsion, so $K=0$. This gives their integral
Theorem~7.2.8. Their \S7.1.10 applies this
surjectivity at each finite cyclotomic level
\cite[\S\S7.1.3, 7.1.8--7.1.10, Theorem~7.2.8,
\S\S7.2.10--7.2.12 and~10.2.2--10.2.3;
preprint pp.~84--86, 89 and~107]{FukayaKato2012}.

In particular, at $r=1$ the $\theta$-component of
\[
 H^2(\operatorname{Spec}\mathbb Z[\mu_p,1/p],\mu_p^{\otimes2})
 \simeq (\operatorname{Cl}(F_1)/p)\otimes\mu_p
\]
is generated by cups of the $p$-units $1-\zeta_p^a$.
The class-group character in this component is
$\chi=\omega^{-1}\theta$; thus this is exactly the coefficient twist
required for the unit pairing. The excluded characters
$\theta=1,\omega^2$ correspond to $\chi=\omega^{-1},\omega$.
Their class-group components are zero by Stickelberger's theorem;
this exclusion is also explicitly checked in
\cite[proof of Lemma~5.2.7, p.~13]{WakeWangErickson2017}.
The remaining odd $\theta$-components of the displayed target vanish
because Vandiver gives the zero plus part of the class group.
We have therefore proved surjectivity of the full $p$-unit cup pairing.

The parity argument already proved above converts this into precisely
the hypothesis of Theorem~\ref{thm:pole-general}: the minus part of
the $p$-units modulo $p$ is the one-dimensional root-of-unity line,
its alternating self-cup is zero, and mixed-parity cups have zero
target. Thus plus Kummer classes suffice, and untwisting
$\mu_p=\mathbb F_p(1)$ makes them the odd-character part of
$H^1(P,\mathbb F_p)$. Theorem~\ref{thm:pole-general} gives the
kernel-field equality.

For $p=3$, the class group of $\mathbb Q(\mu_3)$ is zero, so the
same Kummer and Brauer calculation gives $H^2(P,\mathbb F_3)=0$.
The cup condition is vacuous and Theorem~\ref{thm:pole-general}
again applies. Finally, the equivalence of the Vandiver hypothesis
with the stated positive $K_{4k}$ vanishings is the known comparison
(K2)--(K3), proved in Section~\ref{sec:k-groups} using
\cite[VI.8.2, VI.8.3 and VI.10.9]{Weibel2013} and
\cite[Corollary~1.5, Proposition~3.7 and Remark~3.11]{Kurihara1992}.
\end{proof}

\begin{corollary}[K17-KERNEL-LT2POWER31]\label{cor:kernel-lt2power31}
For every prime $p<2^{31}$, one has $\Omega_p^*=\Omega_p$.
\end{corollary}

\begin{proof}
Hart--Harvey--Ong verify Vandiver for every odd prime below $2^{31}$
\cite[\S1, item~(2), p.~1]{HartHarveyOng2017}; apply
Theorem~\ref{thm:vandiver-kernel}. The prime $2$ is supplied by
Theorem~\ref{thm:pro2-kernel}.
\end{proof}

This numerical corollary uses the cited published computation. It
requires neither a new modular-symbol calculation nor any condition
on the multiplicities of the Iwasawa roots. For primes beyond this
range the same deduction applies whenever Vandiver is established.
The unconditional all-prime statement remains open in this argument.

% Checked algebraic deduction; whole-proof archive audit in cup_ext_archive_audit.md.
% This lemma alone does not assert a kernel theorem without its module hypotheses.
\subsection{Finite-level cups detect the actual companion poles}

The following comparison concerns the actual class-two extension of
the Galois group. It supplies the naturality needed to use a cup-product
theorem; an equality of characteristic ideals would not supply it.

\begin{lemma}[K17-CUP-EXT]\label{lem:cup-pole-bridge}
Let $p$ be odd, let $P=H\rtimes\Gamma$, where $H$ is free pro-$p$
and $\Gamma=\langle\gamma\rangle\simeq\mathbb Z_p$, and let an
involution $c$ preserve $H$ and fix $\gamma$. Put
\[
 \Lambda=\mathbb Z_p[[\Gamma]],\qquad R=\Lambda[1/p],\qquad
 H^{\rm ab}=U\oplus T,
\]
where $c=-1$ on $U$ and $c=1$ on $T$. Suppose that $U$ is finitely
generated over $\Lambda$, that $U[1/p]$ is free of finite rank over
$R$, and that $T$ is finite free over $\mathbb Z_p$. Suppose also
that each even character component $T_e[1/p]$ under the commuting
prime-to-$p$ action is cyclic over $R$. Write $q_e(z)$ for its monic
characteristic and minimal polynomial for $\gamma$.

Let the coefficient character satisfy $\chi(\gamma)=u\in1+p\mathbb Z_p$
and $\chi(c)=-1$, with $H$ acting trivially on its twists. Assume that,
for every $n\geq1$, the even part of
\[
 H^2(P_n,\mathbb Z/p^n(2)),\qquad
 P_n=H\rtimes\Gamma^{p^n},
\]
is generated by cups of $c$-even classes in
$H^1(P_n,\mathbb Z/p^n(1))$. A cofinal diagonal of the two level
indices suffices.

Then the actual odd--odd companion numerators satisfy the
pole-nonvanishing condition: for every root $\alpha$ of every
nonzero $q_e$, at least one coefficient has nonzero restriction to
\[
 D_\alpha=(1+T_1)(1+T_2)-\alpha=0.
                                                               \tag{CE.1}
\]
Here coefficients are taken in the bounded-denominator power-series
coordinates of an $R$-basis of $U[1/p]$. Repeated roots are allowed.
For the subsequent pole-separation application the arithmetic module
input supplies rank one in every odd character.
\end{lemma}

\begin{proof}
Since $2$ is a unit, the class-two quotient of $H$ has its integral
Baker--Campbell--Hausdorff Lie structure. Quotient its central
commutator module by all brackets having a $T$ input. The remaining
central module is
\[
 N=\widehat{\bigwedge}^{\,2}_{\mathbb Z_p}U.
\]
Complex conjugation is $1$ on $N$. Thus the minus part of the
resulting Lie module maps isomorphically to $U$, providing a
canonical $\Gamma$-equivariant odd section. Its plus part is the
continuous module extension
\[
 0\longrightarrow N\longrightarrow E\longrightarrow T
 \longrightarrow0.                                          \tag{CE.2}
\]
Choose a rational cyclic generator $v$ of $T_e[1/p]$ and a lift
$y\in E[1/p]$ in the same character component, using the
prime-to-$p$ character projector. The element
\[
 q_e(\gamma)y=C_e\in N[1/p]                                 \tag{CE.3}
\]
is precisely the odd--odd quadratic companion error. In the quotient
just taken, brackets involving even lifts vanish, so an ordered
companion product becomes its additive characteristic-polynomial
expression. A derived element has the same degree-two logarithmic
and Magnus coefficient.

Finite generation clears the denominators of the chosen isomorphism
$U[1/p]\simeq R^r$ and its inverse: for some $a$,
$p^a\Lambda^r\subset U\subset p^{-a}\Lambda^r$ in these
coordinates. Completed tensor and exterior squares inherit fixed
denominator bounds. Consequently $N[1/p]$ is the bounded-denominator
alternating two-variable coefficient module, rather than a module of
series with arbitrary coefficient denominators.

Pass to a finite splitting field $L/\mathbb Q_p$, with integers
$\mathcal O$, and suppose that every coefficient of $C_e$ vanishes
on $D_\alpha=0$. All roots satisfy $\alpha-1\in\mathfrak m_{\mathcal O}$,
because the action of the pro-$p$ group $\Gamma$ on the finite-rank
lattice $T$ is continuous. The change of variable
\[
 S=(1+T_1)(1+T_2)-\alpha,\qquad
 T_2=\frac{\alpha+S}{1+T_1}-1
\]
is an integral formal-coordinate change. Its specialization kernel
is $(S)$, and division preserves a common denominator bound.
Since $D_\alpha$ is symmetric, an alternating numerator has an
alternating quotient. Thus, in $N_L$,
\[
 C_e=(\gamma-\alpha)A.                                      \tag{CE.4}
\]
Wrong-character coefficients are identically zero, so it suffices
to assume vanishing for the matching odd-character pairs.

Set $r_e(z)=q_e(z)/(z-\alpha)$. Cyclicity implies that
$w=r_e(\gamma)v$ is nonzero and spans an $\alpha$-eigenline
$W\subset T_L$. For a repeated root it is the socle of the cyclic
generalized eigenspace. Equations (CE.3)--(CE.4) show that
\[
 \widetilde w=r_e(\gamma)y-A
\]
is a $\Gamma$-equivariant lift of $w$. Hence (CE.2) splits over
$W$ after inverting $p$.

We now compare this split with the actual cup map. The lattice
$T_{\mathcal O}\cap W$ is saturated, hence an $\mathcal O$-direct
summand of $T_{\mathcal O}$. Choose an $\mathcal O$-linear section
$s:T_{\mathcal O}\to E_{\mathcal O}[1/p]$ extending the equivariant
section on that lattice. A fixed power $p^b$ makes this section
integral. Write the class-two logarithm of $h\in H$ as
\[
 x(h)+s(t(h))+z(h),\qquad
 x(h)\in U_{\mathcal O},\quad t(h)\in T_{\mathcal O},\quad
 z(h)\in N_{\mathcal O}[1/p].
\]
The same $p^b$ clears every $z(h)$. The section defect is
\[
 \beta_\delta(t)=\delta s(t)-s(\delta t),\qquad\delta\in\Gamma,
                                                               \tag{CE.5}
\]
and it vanishes on $W$.

Fix $n$ and two even cohomology classes $f,g$. Average their cocycles
by $(1+c)/2$. The resulting invariant cocycles vanish on the chosen
$\Gamma_n=\Gamma^{p^n}$ section: $c$ fixes its elements and is
$-1$ on the coefficient module. On $H$ they are additive maps
from $U$ to $\mathcal O/p^n(1)$, vanish on $T$, and are
$\Gamma_n$-equivariant. Put
\[
 \psi=f\wedge g:N_{\mathcal O}\longrightarrow\mathcal O/p^n(2),
 \quad
 \psi([x_1,x_2])=f(x_1)g(x_2)-f(x_2)g(x_1).
\]
The continuous cochain
\[
 k(h)=-\frac{p^b}{2}f(h)g(h)-\psi(p^b z(h))                    \tag{CE.6}
\]
is integral. The class-two product has central coordinate
$z_1+z_2+[x_1,x_2]/2$. Its substitution in (CE.6) gives
\[
 dk(h_1,h_2)=p^b f(h_1)g(h_2).
\]
Indeed the two summands contribute respectively
$p^b(f_1g_2+f_2g_1)/2$ and $p^b(f_1g_2-f_2g_1)/2$.

Extend $k$ to $P_n$ by $k(h\delta)=k(h)$. The cup has zero mixed
terms on the chosen section, so its Hochschild--Serre edge is, up
to the convention for the edge sign, the equivariance defect of
$k$. Its quadratic summand is equivariant; the remaining term is
\[
 \pm\psi(p^b\beta_\delta(t)).                               \tag{CE.7}
\]
This vanishes on $W$. A convention that conjugates $t$ makes no
difference because $W$ is $\Gamma$-stable. Thus the same fixed
$p^b$ multiple of every cup has zero restriction to the dual of
$W$, at every finite level.

For completeness, the target and its transitions are explicit.
Both $H$ and $\Gamma_n$ have $p$-cohomological dimension one,
and the minus part $U$ contributes only odd parity. Therefore
\[
 H^2(P_n,\mathcal O/p^n(2))^+
 \simeq H^1(\Gamma_n,\operatorname{Hom}(T,\mathcal O/p^n)(2)).
                                                               \tag{CE.8}
\]
Let $V=\operatorname{Hom}_{\mathbb Z_p}(T,\mathcal O)(2)$, with
$\gamma$ operator $B$. Normalized cocycle evaluation identifies
the right side with
\[
 V\big/\bigl(p^nV+(B^{p^n}-1)V\bigr).                         \tag{CE.9}
\]
Corestriction for the procyclic index-$p$ subgroup sends its
generator value to the same value in the larger group's
coinvariants. Thus the transitions are quotient maps, without an
extra factor of $p$. Continuity implies $B^{p^n}-1\to0$; the
denominators in (CE.9) are cofinal with the $p$-adic neighbourhoods.
Their inverse limit is $V$.

Restriction to the primitive lattice in $W$ gives, after inverting
$p$, a nonzero quotient
\[
 V_L\longrightarrow W^*(2),\qquad
 \gamma|_{W^*(2)}=u^2/\alpha,
\]
of character $\chi^2e^{-1}$, or exponent $2-e$ in cyclotomic
character notation. The restriction maps on the integral Hom
modules are surjective, and the corresponding finite coinvariant
quotients form a separated inverse system. By (CE.7) and the
finite-level cup-surjectivity hypothesis, the image of $p^bV$ in
every one of these quotients is zero. Its image in $W^*(2)$ must
therefore be zero. This contradicts the nonzero torsionfree
quotient. No compatible choice of finite-level cup expressions is
needed, since $b$ was fixed before their selection. The contradiction
proves (CE.1) for every root, including repeated ones.
\end{proof}

% Complete whole-H rational-coordinate proof and arithmetic application.
% Independent review: greenberg_magnus_adaptation.md and sibling audit.
% Include cup_ext_pole_section.tex before this section.
\subsection{Greenberg's one-variable hypothesis suffices}
\label{sec:greenberg-kernel}

The integral freeness and cyclicity used in the preceding pole argument
can be replaced by their versions after inverting $p$. The change must
be made in the coefficient algebra of the entire free pro-$p$ group:
a subgroup with finite index on abelianizations need not be open.

\begin{theorem}[K17-GREENBERG-KERNEL]\label{thm:greenberg-kernel}
Let $p$ be an odd prime, put $F_r=\mathbb Q(\mu_{p^r})$ and
$F_\infty=\bigcup_r F_r$, and let
\[
 X=\varprojlim_r\operatorname{Cl}(F_r)\{p\}
\]
with transition maps given by norms. If $X^+$ is finite, then
\[
 \Omega_p^*=\Omega_p.
\]
The hypothesis is the one-variable CM form of Greenberg's conjecture.
No simplicity of Iwasawa roots is assumed, and no converse is asserted.
\end{theorem}

We first prove the rational-coordinate version of pole separation and
the precise cohomological comparison needed for its arithmetic use.

\begin{lemma}[K17-RATIONAL-MAGNUS]\label{lem:rational-magnus}
Let $p$ be odd and let $H$ be a second-countable free pro-$p$ group
with commuting continuous actions of
$\Gamma=\langle\gamma\rangle\simeq\mathbb Z_p$ and
$\Delta\simeq\mathbb F_p^\times$. Let $c\in\Delta$ have order two,
and put
\[
 \Lambda=\mathbb Z_p[[\Gamma]],\qquad R=\Lambda[1/p],\qquad
 M=H^{\rm ab}=U\oplus T,
\]
where $c=-1$ on $U$ and $c=1$ on $T$. Assume that $M$ is finitely
generated over $\Lambda$, that $U[1/p]$ is free of rank one over $R$
in every odd $\Delta$-character, and that $T$ is finite free over
$\mathbb Z_p$, with every nonzero $T_e[1/p]$ cyclic over $R$.
Assume also that abelian eigenclasses have exact $\Delta$-eigenlifts,
including inside arbitrarily small $\Delta$-stable neighbourhoods.

Let $\rho:H\to\mathcal U(\mathbb Q_p)$ be continuous and
$\Gamma\times\Delta$-equivariant, where $\mathcal U$ is the
weight-completed free prounipotent group on letters $s_m$ in all odd
weights $m\ge3$. Suppose its image is Zariski dense in every finite
weight quotient and, after a Levi normalization,
\[
 \gamma s_m=u^m s_m,\qquad \delta s_m=\omega(\delta)^m s_m,
\]
where $u$ topologically generates $1+p\mathbb Z_p$.

For each nonzero $T_e$, choose an integral multiple $v_e$ of a rational
cyclic generator, an exact eigenlift $y_e\in H$, and its actual
conjugates $y_{e,j}=\gamma^j y_e\gamma^{-j}$. Let $q_e(z)$ be the
characteristic polynomial of $\gamma$ on $T_e$. In a rational
$R$-basis of $U[1/p]$, let $C_{e;a,b}(T_1,T_2)$ be the odd--odd
quadratic coefficients of the actual companion error. Suppose that
for every root $\alpha$ of every $q_e$, there is a matching ordered
odd pair $(a,b)$ for which
\[
 C_{e;a,b}\left(T,\frac{\alpha}{1+T}-1\right)\ne0.
                                                               \tag{GM.1}
\]
Then $\rho$ is injective. All coefficient series here have a common
finite denominator in each fixed assigned degree; no bound uniform
over all degrees is required.
\end{lemma}

\begin{proof}
\emph{An actual compact basis.}
A compact torsionfree $\mathbb Z_p$-module is a product of copies of
$\mathbb Z_p$: its Pontryagin dual is a divisible $p$-primary group,
hence a direct sum of copies of $\mathbb Q_p/\mathbb Z_p$.
The abelianization of free $H$ is torsionfree and metrizable.
Using the integral character idempotents, choose pointed convergent
$\mathbb Z_p$-bases in the $U_a$, and finite bases in the $T_e$.
The open quotient map $H\to M$ permits lifts of the convergent
basis vectors tending to $1$. Apply the exact-eigenlift construction
inside a nested sequence of $\Delta$-stable open neighbourhoods;
it preserves both the abelian classes and their convergence.

Let $F$ be the free pro-$p$ group on this pointed sequence and the
finite even set. Its abelianization is the product module on those
basis vectors, so the resulting map $F\to H$ is an isomorphism on
abelianizations. Its compact image is Frattini-surjective and hence
equals $H$. Indeed a proper closed subgroup of any pro-$p$ group is
contained in a proper open subgroup and then in an index-$p$ subgroup,
so cannot be Frattini-surjective. Since $H$ is free, the surjection
has a continuous section. The section is inverse on abelianizations;
its closed image is Frattini-surjective in $F$ and is all of $F$.
Thus $F\simeq H$. This argument does not use an infinite-rank
Hopficity assertion.

Give the odd Magnus variables $X_i$ assigned degree $1$, and the
finite even variables $Y_l$ assigned degree $2$. This Magnus
expansion is injective. To see this directly, kill a tail of the
convergent free sequence. Every finite quotient kills such a tail,
so $H$ is the inverse limit of these finite-rank free groups.
An ordinary finite-rank Magnus expansion detects every nonidentity
element in some finite assigned degree. At each fixed degree the
coefficient module is a finite sum of completed tensor products of
$U$ and $T$, using these assigned degrees.

\emph{Odd rational coordinates.}
Choose a character-preserving isomorphism
\[
 U[1/p]\simeq\bigoplus_{a\ {
m odd}}Re_a.
\]
Finite generation over $\Lambda$ clears the denominators of this
isomorphism and its inverse. Thus there are continuous integral
maps between $U$ and $\Lambda^{(p-1)/2}$ whose two compositions are
multiplication by the same power of $p$. Their completed tensor
powers have compositions equal to the corresponding powers of that
scalar. The tensor modules are torsionfree: in the product bases
they are products of copies of $\mathbb Z_p$. Hence the maps are
injective and become inverse isomorphisms after inverting $p$.
Every fixed degree therefore has coordinates in
\[
 \mathbb Z_p[[\Gamma^s]][1/p]
   =\mathbb Z_p[[T_1,\ldots,T_s]][1/p],                       \tag{GM.2}
\]
for its singleton positions and their character types. These are
the Amice transforms of bounded measures on $\Gamma^s$.
The denominator bound may depend on the degree, which causes no
problem at a first nonzero degree.

The target linear coefficient of weight $m$ is a continuous additive
map from $M$ to $\mathbb Q_p$. Parity kills $T$, and character
equivariance kills every $U_a$ except $a\equiv m\pmod{p-1}$.
On the remaining rank-one $R$-module, $\Gamma$-equivariance makes
it a scalar multiple of evaluation at $T=u^m-1$: check this on
polynomials in $\gamma$ and then pass to bounded series by
continuity. The scalar is nonzero by Zariski density, since the
weight-$m$ abelian coordinate of the target cannot vanish on a
dense image. Rescale $s_m$ by this scalar. Such weightwise rescaling
is an automorphism of the completed free Lie algebra, because each
finite weight quotient has finitely many generators. The singleton
kernel is now exactly evaluation at $u^m-1$ in its matching type.

\emph{Even companion coordinates.}
Put $\lambda_e=\operatorname{rank}_{\mathbb Z_p}T_e$. The classes
\[
 v_e,\gamma v_e,\ldots,\gamma^{\lambda_e-1}v_e
\]
are a rational basis, though they need not be an integral basis.
The polynomial $q_e$ is monic integral, has unit constant term,
and reduces to $(z-1)^{\lambda_e}$ modulo $p$, since the image of
$\Gamma$ in $\operatorname{GL}(T_e/p)$ is a finite $p$-group.
It is also the minimal polynomial, by rational cyclicity.

The actual conjugates $y_{e,j}$ have assigned degree at least $2$
and target generator length at least $2$. Their assigned-degree-two
parts have the form
\[
 Z_{e,j}=\sum_l b_{j,l}Y_{e,l}+A_{e,j}(X),                    \tag{GM.3}
\]
where $b$ is a finite invertible rational matrix and $A_{e,j}$ is
a completed quadratic odd tensor. Before the odd coordinate change
these coefficients are integral; afterwards they retain bounded
denominators. There are no other terms of assigned degree two.
The quadratic tensor is alternating, by the commutator description
or primitiveness of the first Magnus coefficient.

Replace the finite variables $Y$ by $Z$. This is a homogeneous
triangular automorphism of the completed coefficient algebra:
it fixes $X$ and its inverse is $Y=b^{-1}(Z-A(X))$.
In each fixed degree only finitely many even variables occur,
so both directions preserve a common finite denominator bound.
In particular a first nonzero degree remains nonzero.
This is a change in the rational tensor algebra, not a claim that
the companion classes form an integral group basis.

By definition of $Z_{e,j}$, its leading target is the length-two
part of the actual element $y_{e,j}$. For an ordered target pair
$(m,n)$, put $z=u^{m+n}$. Its coefficient satisfies
\[
 B_{e,j;m,n}=z^jB_{e,0;m,n}.
\]
Write $q_e(z)=z^{\lambda_e}-\sum_jc_{e,j}z^j$ and form
\[
 E_e=(\gamma^{\lambda_e}y_e\gamma^{-\lambda_e})
       \left(\prod_jy_{e,j}^{c_{e,j}}\right)^{-1}.
\]
Its abelian class is zero, so $E_e$ is in the closed derived
subgroup. Its assigned-degree-two part is the odd alternating
tensor $C_e$. There are no target length-one terms in any even
factor, so comparison of length two gives exactly
\[
 q_e(z)B_{e,0;m,n}
   =C_{e;a,b}(u^m-1,u^n-1),\qquad
 a\equiv m,\ b\equiv n\pmod{p-1}.                           \tag{GM.4}
\]
This is the same companion class as the actual extension in
Lemma~\ref{lem:cup-pole-bridge}. Exact eigenlifts make its assigned
degree-two character $e$; hence wrong-character pairs have zero
coefficient. One can also obtain this zero from (GM.4) on the full
residue grid and the successive analytic identity theorem.

\emph{Partial fractions and continuity.}
Pass to a finite splitting field $L$, with integers $\mathcal O$.
All roots $\alpha$ satisfy $\alpha-1\in\mathfrak m_{\mathcal O}$.
The finite-dimensional rational-function space with basis
$z^j/q_e(z)$ has alternative basis
\[
 (z-\alpha)^{-k},\qquad
 1\le k\le\operatorname{mult}_e(\alpha).
\]
The corresponding constant invertible change in the finite $Z$
variables preserves assigned degree and bounded denominators in
every fixed degree. Its even letter $(e,\alpha,k)$ has kernel
\[
 \frac{C_{e;a,b}(T_1,T_2)}
      {((1+T_1)(1+T_2)-\alpha)^k}.                          \tag{GM.5}
\]
The numerator is common to all letters from the same companion
component, even when roots repeat.

There are finitely many positive integers $w$ such that $u^w$ is a
root of some $q_e$. Restrict all test weights to a sufficiently
large ordinary tail, excluding every vanishing pair denominator.
In each fixed odd residue class this still includes sequences
$m_0+(p-1)p^j$ with an interior $p$-adic accumulation point, for
a sufficiently large fixed representative $m_0$.

Substitution is legitimate at every finite arithmetic-weight cutoff.
In the actual compact basis, truncate its convergent sequence and
apply finite-rank Magnus substitution. These truncations converge
to the identity on $H$, since every finite group quotient kills a
tail. The image of $H$ in the finite target quotient is compact,
so its coefficient functions are bounded and continuous. Passing
to the limit therefore gives the completed tensor substitution.
The bounded rational coordinate changes preserve this statement in
each fixed degree. In (GM.2) it is equivalently integration against
bounded measures and the finite Fubini rule. We do not invoke a
free pro-$p$ universal property for a noncompact rational group.

\emph{Separation of a first nonzero degree.}
Suppose $g\ne1$ and $\rho(g)=1$, and let $d$ be its first nonzero
assigned degree after these invertible changes. A larger source
degree cannot contribute to target length $d$. Its coefficient is
a finite sum over ordered tilings of $d$ positions by singleton
odd types and pairs labelled $(e,\alpha,k)$:
\[
 \begin{gathered}
 \sum_{\mathcal P}A_{\mathcal P}(T_i:\ i\text{ singleton})
       \prod_{(j,e,\alpha,k)\in\mathcal P}
       \frac{C_{e;a_j,a_{j+1}}(T_j,T_{j+1})}
            {D_{j,\alpha}^{k}}=0,\\
 D_{j,\alpha}=(1+T_j)(1+T_{j+1})-\alpha.
 \end{gathered}                                            \tag{GM.6}
\]
Every $A_{\mathcal P}$ has a common bounded denominator. Clear
the finitely many denominators. Vanishing on the Cartesian grid
of large odd test weights and the successive one-variable identity
theorem make (GM.6) an identity in the fraction field of
$\mathcal O[[T_1,\ldots,T_d]][1/p]$.

Select a nonzero pattern with the largest number of pairs. Fix its
pair positions $J$ and its roots and even characters, and among
nonzero patterns with those data choose a componentwise maximal
pole-order vector $(k_j)$. At each selected pair choose the ordered
odd types in (GM.1), and at each singleton its source type. These
choices are independent because the pairs are disjoint.

Set $S_j=D_{j,\alpha_j}$ for $j\in J$. Retain as base variables
the singleton variables and the left variable in each selected pair.
The substitutions
\[
 T_{j+1}=\frac{\alpha_j+S_j}{1+T_j}-1
\]
are automorphisms of the complete integral power-series ring.
Let $K_0$ be the fraction field of the bounded-denominator base
ring. The identity embeds in $K_0[[S_j:j\in J]][S_j^{-1}:j\in J]$,
with finite lower bounds on exponents. All other denominators are
units in $K_0[[S_j]]$: a different root at the same position has
nonzero constant root difference; a different pair position joins
two distinct selected blocks, whose surviving variables are
independent, so its product is not identically a prescribed root.
Wrong even characters have already vanished.

Take the coefficient of $\prod_{j\in J}S_j^{-k_j}$.
A pattern missing a selected position or using a different root
there has no negative exponent in that variable. A pattern
containing all selected pairs and an extra pair has zero coefficient
by the first maximality choice. Among remaining patterns, a
contribution requires all pole orders to dominate $k_j$; the
componentwise maximal choice excludes strict domination and an
incomparable vector fails in at least one variable. Singleton types
must match. Thus exactly one coefficient remains:
\[
 A_{\mathcal P}(T_i:\ i\text{ singleton})
 \prod_{j\in J}
 C_{e_j;a_j,b_j}\left(T_j,\frac{\alpha_j}{1+T_j}-1\right).
\]
The first factor is nonzero and its coefficient ring injects into
$K_0$; every other factor is nonzero by (GM.1). Their product is
nonzero in a field, a contradiction. If $J$ is empty, the same
argument is the ordinary singleton moment identity. This proves
injectivity of $\rho$.
\end{proof}

\begin{lemma}[K17-ETALE-PROP]\label{lem:etale-prop}
Let $S$ be a connected normal affine scheme, let
$G=\pi_1^{\mathrm{\acute et}}(S)$, and let $P$ be its maximal
pro-$p$ quotient. For a finite $p$-primary local coefficient module
$A$ whose action factors through $P$, the natural maps give
\[
 H^1(P,A)\simeq H^1_{\mathrm{\acute et}}(S,A),\qquad
 H^2(P,A)\hookrightarrow H^2_{\mathrm{\acute et}}(S,A).
\]
They are compatible with cup products and with finite external
automorphisms. Consequently generation of a character summand of
the target $H^2$ by cups of $H^1$ classes gives the same generation
in group cohomology and makes the displayed injection an
isomorphism on that summand.
\end{lemma}

\begin{proof}
Put $N=\ker(G\to P)$. An element of $H^1(N,A)^P$ is a continuous
$P$-equivariant homomorphism $N\to A$, since $N$ acts trivially
on $A$. If its image were nonzero, its kernel would be $G$-normal
and the quotient of $G$ by that kernel would be an extension of
$P$ by a nonzero finite $p$-group. This is a pro-$p$ quotient:
every finite quotient has $p$-power order. It would be strictly
larger than the maximal pro-$p$ quotient, a contradiction.
Thus $H^1(N,A)^P=0$. The five-term Hochschild--Serre sequence gives
\[
 H^1(P,A)\simeq H^1(G,A),\qquad
 H^2(P,A)\hookrightarrow H^2(G,A).
\]
No finite-generation hypothesis on $N$ was used.

The full universal finite \'etale tower over $S$ has zero $H^1$
with finite locally constant coefficients: after trivializing the
coefficients, an $H^1$ torsor is finite \'etale and becomes trivial
at a later stage of the same tower. Finite \'etale objects and
torsors descend under these affine limits. The Cartan--Leray
spectral sequence for the tower therefore gives
\[
 H^1(G,A)\simeq H^1_{\mathrm{\acute et}}(S,A),\qquad
 H^2(G,A)\hookrightarrow H^2_{\mathrm{\acute et}}(S,A).
\]
This is only its low-degree consequence, and asserts no vanishing
of the higher rows. The sequence, the profinite-tower passage and
the simply connected $H^1$ vanishing are given explicitly in
\cite[Theorem~14.9, Remark~14.10, Example~14.11 and
Proposition~14.13, pp.~96--98]{MilneLEC2013}.
All comparison maps are natural and respect cups.
Lift each cup generator using the $H^1$ isomorphism and add its
group cups; this proves surjectivity onto the asserted summand.
Injectivity then identifies every group class in that summand with
the lifted sum.
\end{proof}

\begin{proof}[Proof of Theorem~\ref{thm:greenberg-kernel}]
The prime $3$ is already covered by the regular-prime result, so
assume $p\ge5$. Put
\[
 P=\operatorname{Gal}(\Omega_p/F_1),\quad
 H=\operatorname{Gal}(\Omega_p/F_\infty),\quad
 \Gamma=\operatorname{Gal}(F_\infty/F_1),\quad
 M=H^{\rm ab}=U\oplus T.
\]
Choose the finite $\Delta$ lift by the coprime complement theorem.
The fibre over a generator of $\Gamma$ is a $\Delta$-stable
$H$-torsor because the cyclotomic quotient is abelian.
The coprime identity $H^1(\Delta,H)=1$, obtained from the finite
complement theorem and compactness, gives a $\Delta$-fixed lift
$\gamma$. Its procyclic closure maps isomorphically to $\Gamma$
and supplies the commuting section used above.
The arithmetic structure used earlier, from Iwasawa and
Ferrero--Washington, gives free pro-$p$ $H$.
The finite-level Kummer sequence makes $H^1(P,\mathbb F_p)$
finite dimensional, so $P$ is finitely generated and $H$ is
second-countable. Adjust a finite generating set of $P$ by powers
of the chosen $\Gamma$ section to obtain elements of $H$.
Together with $\gamma$ they generate $P$; their closed normal
closure is $H$, since the remaining quotient is the given
procyclic quotient. Their abelian classes therefore generate $M$
over $\Lambda=\mathbb Z_p[[\Gamma]]$.
The actual-image results of Hain--Matsumoto and Brown give the
continuous dense rational representation and Levi normalization in
Lemma~\ref{lem:rational-magnus}; this is the same actual
representation as in the preceding pole theorem, not an assumed
faithful motivic representation. The exact eigenlift construction
is \cite[Lemma~2.1, pp.~3--4]{Sharifi2002}; performing it inside
stable neighbourhoods preserves convergence. We verify the changed
module and cup hypotheses in full.

\emph{Rational odd module.}
The ring $R=\Lambda[1/p]$ is a PID. Indeed $\Lambda$ is a
two-dimensional regular local UFD, and its only height-two prime
contains $p$. Thus $R$ is a one-dimensional noetherian UFD and
hence a PID. It remains the ring of bounded-denominator series,
not $\mathbb Q_p[[T]]$.

For an odd character $a\ne\omega,\omega^{-1}$, apply the full
duality calculation in Wake--Wang-Erickson's Proposition~2.2.6
with $\theta=\omega a^{-1}$. It gives
\[
 \operatorname{Ext}^1_\Lambda(M_a^{\#}(1),\Lambda)=X_\theta,
 \qquad
 \operatorname{Ext}^2_\Lambda(M_a^{\#}(1),\Lambda)=0.           \tag{GM.7}
\]
Here $M$ is their $p$-ramified module $\mathfrak X$, whereas $X$
on the right is unramified. Since $\theta$ is even, Greenberg
makes the right-hand side zero after inverting $p$. Ext commutes
with this localization for finitely generated modules. Over a PID,
the Ext group into the ring of a finitely generated torsion summand
is nonzero unless that summand is zero. Thus $M_a[1/p]$ is free.
Its rank is one by the local/global unit exact sequence in the
same proof: local units have rank one in this character, global
minus units contribute only roots of unity, and the unramified
quotient is torsion. The inversion $\#$ and Tate twist transport
the ring by automorphisms, so undoing them preserves freeness
\cite[Proposition~2.2.6 and proof, pp.~3--4]{WakeWangErickson2017}.

The two excluded characters can be calculated directly.
Coleman's exact sequence at tame level one is
\[
 0\longrightarrow\mathbb Z_p(1)\longrightarrow U_{\rm loc}
 \longrightarrow\mathbb Z_p[[\Delta\times\Gamma]]
 \longrightarrow\mathbb Z_p(1)\longrightarrow0
\]
\cite[Theorem~5.5.31, p.~148; Proposition~6.1.1,
pp.~151--152]{SharifiIwasawaNotes}.
Global minus units are exactly the root-of-unity term, and
Stickelberger gives $X_\omega=X_{\omega^{-1}}=0$; the latter
exceptional statement is explicitly recorded in
\cite[proof of Lemma~5.2.7, p.~13]{WakeWangErickson2017}.
The $\omega^{-1}$ component of $M$ is consequently free of rank
one. In the $\omega$ component, quotienting the local units by
their root-of-unity kernel identifies $M_\omega$ with
\[
 \ker(\Lambda\longrightarrow\mathbb Z_p(1))
     =(\gamma-u)\Lambda,
\]
again free of rank one. The two characters are distinct for
$p\ge5$. The valuation direction at the unique prime above $p$
is even and does not affect this calculation. We have therefore
proved the required rational freeness for every odd character.

\emph{Rational even module.}
Kummer duality identifies $T=M^+$ up to finite kernel and cokernel
with $(X^-)^\iota(1)$
\cite[Proposition~3.4.7 and Corollary~3.4.8,
p.~92]{SharifiIwasawaNotes}.
Ferrero--Washington makes it finitely generated over $\mathbb Z_p$,
and freeness of $H$ excludes $\mathbb Z_p$-torsion. Thus $T$ is
finite free. Greenberg implies that $X^-$ is pseudocyclic
character by character, by the exact local/global unit sequence and
Iwasawa adjunction in
\cite[Corollary~2.1.3 and Conjecture~2.1.5 with the following
explanation, p.~8]{WakeWangEricksonPseudo2018}.
Finite kernels and cokernels vanish after inverting $p$, so every
$T_e[1/p]$ is cyclic. This allows repeated roots; it does not assert
integral cyclicity without Vandiver.

\emph{Integral symbol surjectivity under Greenberg.}
Fix an even character $\theta\ne1,\omega^2$ and put
$\chi=\omega^{-1}\theta$. Wake--Wang-Erickson's
Corollary~8.1.3(2) assumes exactly
\[
 \begin{gathered}
 X_\chi(1)[1/p]\text{ cyclic over }(\Lambda/\xi_\chi)[1/p],\\
 X_\theta\otimes_\Lambda(\Lambda/\xi_\chi)[1/p]=0.
 \end{gathered}                                            \tag{GM.8}
\]
The localization in the second condition is essential.
Both conditions follow from the just stated consequences of
Greenberg. The conclusion is an integral isomorphism
\[
 \varpi:(\mathcal H^-_\theta/I_\theta\mathcal H^-_\theta)
           /(p\text{-primary torsion})
       \xrightarrow{\ \sim\ }X_\chi(1),                    \tag{GM.9}
\]
so the original $\varpi$ is surjective
\cite[Corollary~8.1.3 and complete proof,
pp.~48--49]{WakeWangEricksonPseudo2018}.
The source proves this with repeated roots: the derivative identity
makes the error nilpotent in each cyclic local factor, rather than
requiring the derivative to be invertible. Thus (GM.9) does not
insert a simple-root hypothesis.

Fukaya--Kato, \S7.1.10, deduce from this surjectivity that for every
$r\ge1$ the finite-level $\theta$ component is generated by the
symbols
\[
 \{1-\zeta_{p^r}^{a},1-\zeta_{p^r}^{b}\}.
\]
Their \S7.1.9 identifies its reduction modulo $p^r$ with the
corresponding class-group quotient, and their \S6.1.5 shows that
at tame level one the only excluded even characters are
$1,\omega^2$
\cite[\S6.1.5, p.~76; \S\S7.1.8--7.1.10,
p.~86]{FukayaKato2012}.
Those excluded characters correspond to the unramified components
$X_{\omega^{-1}}$ and $X_\omega$, which are zero. Norm projections
from $X^-$ onto each finite minus class group are surjective, so
the omitted finite targets are zero as well
\cite[Proposition~3.4.1 and proof,
pp.~90--91]{SharifiIwasawaNotes}.
Components with unit characteristic power series are also zero:
the source's $X^-$ is $p$-torsionfree, and a finite such module must
vanish. We apply the nonunit running assumption of the modular
source only to the nonzero components. Hence every even target is
covered.

\emph{All finite group levels and coefficient conventions.}
Let
\[
 P_n=H\rtimes\Gamma^{p^n},\qquad
 S_n=\operatorname{Spec}\mathcal O_{F_{n+1}}[1/p].
\]
The maximal pro-$p$ quotient of $\pi_1(S_n)$ is exactly $P_n$.
Indeed the normal closure over $F_1$ of any finite Galois
$p$-extension of $F_{n+1}$ unramified outside $p$ remains a
$p$-extension: its conjugate compositum over $F_{n+1}$ has
$p$-group Galois group, and $F_{n+1}/F_1$ has $p$-power degree.
Thus it lies in the fixed field $\Omega_p$ used throughout.

At $F_{n+1}$, the finite-level symbol theorem gives generation
with coefficients $p^{n+1}$. Reduction gives generation with
coefficients $p^n$, exactly as in Lemma~\ref{lem:cup-pole-bridge}.
The class-group quotient description in FK \S7.1.9 makes this
reduction explicit. The Brauer contribution is zero: these fields
have no real places, and only one finite prime is omitted, so the
global invariant condition forces its lone possible local
invariant to vanish.

For a $p$-unit $v$, the quotient $v/c(v)$ is a global unit of
absolute value one at every embedding and hence is a root of
unity. Consequently the negative $c$ part of the $p$-unit group
modulo $p^n$ is cyclic. Its alternating self-cup is zero because
$2$ is a unit; mixed-parity cups lie in the odd target. Thus the
even target is generated by cups of $c$-even unit classes. This
argument does not identify the whole of $H^1$ with units and does
not use Vandiver.

Apply Lemma~\ref{lem:etale-prop}. Each Kummer class has a unique
lift to $H^1(P_n,\mathbb Z/p^n(1))$, and group $H^2$ injects into
\'etale $H^2$. Lifting the even symbol generators makes this
injection surjective on the even part and gives precisely
\[
 H^2(P_n,\mathbb Z/p^n(2))^+
  =\bigl\langle f\smile g\mid
       f,g\in H^1(P_n,\mathbb Z/p^n(1))^+\bigr\rangle.       \tag{GM.10}
\]
This treats every $n\ge1$; no compatibility of expressions as sums
of cups is needed. Scalar extension to the integers of a finite
splitting field preserves the assertion, since that ring is finite
free over $\mathbb Z_p$.

Lemma~\ref{lem:cup-pole-bridge} now gives the nonzero companion
restriction (GM.1) for every root of every even component, including
repeated roots. Its dual root is $u^2/\alpha$ and its dual character
is $\omega^{2-e}$, so the inversion and both Tate twists match
the symbol target. Lemma~\ref{lem:rational-magnus} proves
injectivity on the entire $H$.

Finally an element in the kernel of the original pro-$p$ tripod
action is cyclotomically trivial, hence lies in $H$. Functorial
rational completion sends it to $1$ under $\rho$, and the proved
injectivity makes it trivial. This is exactly
$\Omega_p^*=\Omega_p$.
\end{proof}

The theorem is conditional only on the stated one-variable
Greenberg assertion. In particular it does not assert Greenberg
for unverified primes. Vandiver supplies the stronger condition
$X^+=0$ by the earlier persistence proof, so
Theorem~\ref{thm:vandiver-kernel} and its verified numerical range
remain unconditional consequences of their stated computations.

\subsection{A second explicit integral defect: weight twelve at 691}

\begin{corollary}[T-INTEGRAL691]
At $p=691$, every $K_{4k}(\mathbb Z)\{691\}$ vanishes, but the
actual integral tripod Lie algebra is not generated by its odd
Soul\'e classes. With the notation and normalization of
McCallum--Sharifi, one has
\[
 \mathfrak g_{691,12}/\mathfrak h_{691,12}
 \simeq (\mathfrak g_{691}^{\mathrm{ab}})_{12}
 \simeq\mathbb Z/691,
\]
and
\[
 [\sigma_3,\sigma_9]\equiv50[\sigma_5,\sigma_7]
       \pmod{691\mathfrak g_{691,12}}.
\]
Here $\mathfrak h$ denotes the integral Lie subalgebra generated by
the specified odd classes. The displayed coefficient $50$ belongs
to that normalization.
\end{corollary}

\begin{proof}
Vandiver at $691$ lies in the already cited verified range, so
VK-EQUIV gives the $K$-group assertion. Sharifi's cup-product
surjectivity theorem, \cite[Theorem~1.2, p.~2; Corollary~5.9,
p.~27]{Sharifi2007}, gives nontriviality in the irregular component
$(691,12)$. The same paper explicitly makes the application to
Ihara's weight-twelve relation on p.~28.
McCallum--Sharifi \cite[Theorem~9.11, p.~30]{McCallumSharifi2003}
then gives strict inclusion $\mathfrak h_{691,12}\subsetneq
\mathfrak g_{691,12}$ and the displayed congruence. Since $12$ is
the least irregular index, Theorem~9.7 identifies this quotient
with the degree-twelve abelianization. Proposition~9.10 and its
proof, pp.~29--30, identify its nonzero image with the cyclic
group of order $691$ in the stable derivation algebra. This proves
the claimed exact quotient. These are applications of the cited
results, not a new computation of the stable derivation algebra.
\end{proof}

Cup-product surjectivity alone does not assert that the defect for
every irregular prime first occurs at its least irregular index.
Theorem~9.7 additionally requires a nonzero coefficient with
$i+j=m$ and $3\le i<j$; a nonzero pairing whose least positive
representatives sum to $m+p-1$ does not meet that hypothesis.
The two explicit examples $37$ and $691$ have the required
published low-weight computations.

\subsection{A cubic-companion criterion beyond one-variable Greenberg}
\label{sec:odd-torsion-cubic}


The next criterion allows rational Iwasawa torsion in the odd characters.
Its proof isolates a precise class-three nonvanishing condition; no arithmetic
verification of that new condition is asserted here.


\begin{theorem}[K17-CUBIC-ODD]
\label{thm:odd-torsion-cubic}

Let \(p\) be odd, \(\Gamma=\langle\gamma\rangle\cong\mathbb Z_p\), \(\Delta\cong\mathbb F_p^\times\), \(\Lambda=\mathbb Z_p[[\Gamma]]\), and \(R=\Lambda[1/p]\). Let \(H\) be a second-countable free pro-\(p\) group with commuting continuous \(\Gamma\)- and \(\Delta\)-actions. Write \(M=H^{\rm ab}\), and assume it is finitely generated over \(\Lambda\). Let \(T\) be its \(\Lambda\)-torsion submodule and \(Q=M/T\). Assume:

\begin{enumerate}
\item \(Q[1/p]\) is free of rank one over \(R\) in each odd \(\Delta\)-character and zero in the even characters.
\item Every nonzero even character component \(T_e[1/p]\) is cyclic over \(R\).
\item Every nonzero odd character component \(T_o[1/p]\) is cyclic over \(R\). There may be several such characters; each has its own companion polynomial. Allowing degrees greater than one and repeated roots is part of the statement. When all are zero, the argument reduces to the previously audited pair-pole theorem.
\item Exact \(\Delta\)-eigenlifts of abelian eigenclasses can be chosen in \(H\), also inside arbitrarily small \(\Delta\)-stable neighborhoods, as in the earlier free-basis construction.
\item There is a continuous homomorphism \(\rho\) to the free prounipotent group, completed with respect to its weights, over \(\mathbb Q_p\) on one generator \(s_m\) for every odd \(m\ge3\). The map is \(\Gamma\times\Delta\)-equivariant and its image is Zariski dense in every finite weight quotient. In rational Levi coordinates, \(\gamma s_m=u^m s_m\), with \(u\) a topological generator of \(1+p\mathbb Z_p\), and \(\delta s_m=\omega(\delta)^m s_m\).
\end{enumerate}

Assume also the explicit Pair-NV and Triple-NV conditions on the actual companions defined below. Then \(\rho\) is injective.

\end{theorem}
This is a possible sufficient gate when the odd module is no longer rationally free. It does not prove that the gate holds for the arithmetic group. In particular it does not assert Greenberg, its negation, or nonvanishing of any unspecified Massey product.

All coefficient rings below consist of power series with a single bounded denominator at each fixed assigned degree. The ring \(R\) is not identified with \(\mathbb Q_p[[T]]\).

\begin{proof}
\par\medskip\noindent\textbf{Integral torsion and an actual free basis.}

The module \(M\) is \(\mathbb Z_p\)-torsion free because \(H\) is free pro-\(p\). The finitely generated torsion module \(T\) has a nonzero annihilator. Removing its power-of-\(p\) factor does not change annihilation, by torsion freeness. Weierstrass preparation then gives an annihilating distinguished polynomial. Thus \(T\) is finite free over \(\mathbb Z_p\). It is closed: if \(f\) is such an annihilator, then \(T=\ker(f:M\to M)\). It is saturated over \(\mathbb Z_p\), because \(p x\in T\) implies \(fpx=0\), hence \(fx=0\), so \(x\in T\). Therefore \(Q\) is also \(\mathbb Z_p\)-torsion free.

A compact torsion-free \(\mathbb Z_p\)-module is a product of copies of \(\mathbb Z_p\), by dualizing the classification of divisible \(p\)-primary groups. The exact sequence \(0\to T\to M\to Q\to0\) consequently splits continuously as \(\mathbb Z_p\)-modules: the dual of \(Q\) is divisible and hence injective. Apply this within each integral \(\Delta\)-idempotent to obtain a \(\Delta\)-equivariant splitting. It need not be \(\Gamma\)-equivariant, and none will be assumed.

Choose a convergent product basis in each odd component of \(Q\), finite bases of \(T^+\) and \(T^-\), and compatible exact eigenlifts. The lifts of the convergent basis tend to one, since the quotient map is open and the eigenlift hypothesis applies inside shrinking invariant neighborhoods. They form an actual compact free basis of all of \(H\). Indeed the resulting map from the free pro-\(p\) group on that pointed convergent set induces an isomorphism on abelianizations; its image is Frattini-surjective and therefore all of \(H\). A continuous section exists by freeness of \(H\); its image is again Frattini-surjective in the source and hence all of the source. The map is an isomorphism.

Assign degrees \(1,2,3\), respectively, to variables lifting \(Q,T^+,T^-\). This weighted Magnus expansion is injective. Finite quotients kill a tail of the convergent free basis, so ordinary finite-rank Magnus injectivity detects every nonidentity element at a finite assigned degree. At each such degree its coefficient module is a finite sum of completed tensor products of the indicated compact modules.

The \(\Gamma\)-action preserves this filtration. A lift of \(T^+\) has no degree-one term, and its image under \(\Gamma\) still has abelian class in \(T^+\). A lift of \(T^-\) has no degree-one term and no degree-two term. The latter assertion follows from parity: complex conjugation acts positively on every degree-two monomial (two odd degree-one variables or one even variable), while an exact odd eigenlift has negative leading character. Its \(\Gamma\)-image has the same property because the actions commute. The images of degree-one variables have degree at least one. This proves preservation on generators and thus on the completed algebra.

\par\medskip\noindent\textbf{Removing finite linear resonances.}

For every nonzero odd character component, let \(q_o(z)\) be the characteristic polynomial of \(\gamma\) on \(T_o[1/p]\). The linear target coordinate at \(s_m\), restricted to this component, satisfies

\[
q_o(u^m)\,\ell_m(T_o)=0.
\]

There are only finitely many positive integers \(m\) for which any of these finitely many \(q_o(u^m)\) vanish, since the powers \(u^m\) are distinct. Quotient the target free Lie algebra by the ideal generated by the corresponding finitely many \(s_m\), and let \(\rho'\) be the resulting representation. The quotient is still freely generated by all the other odd generators; it is equivariant and the image remains dense in each finite weight quotient.

Now every linear target coordinate kills \(T\) and factors through \(Q\). An odd torsion eigenlift has only odd target lengths, so after its length-one part has been removed its image starts in length at least three. Even lifts start in length at least two. Thus each target image starts in length at least its source assigned degree \(1,2,\) or \(3\), exactly as required for a first-nonzero-degree argument.

Choose a character-preserving \(R\)-basis of \(Q[1/p]\). The isomorphism and its inverse have bounded denominators, since both integral modules are finitely generated over \(\Lambda\). They induce injective bounded changes on every fixed completed tensor degree; the tensor modules are torsion free over \(\mathbb Z_p\). Their orbit coordinates are bounded measures, with transforms in

\[
\mathbb Z_p[[T_1,\ldots,T_s]][1/p].
\]

In a matching odd character, the linear target map is a nonzero scalar multiple of evaluation at \(T=u^m-1\). The scalar is nonzero by density. Rescale each retained \(s_m\), independently, so that the singleton kernel is exactly evaluation. This is an automorphism in the arithmetic-weight completion, because each finite weight quotient has only finitely many generators.

\par\medskip\noindent\textbf{Even companions and the pair condition.}

For each even \(e\), multiply a rational cyclic generator \(v_e\) by a power of \(p\) to make it integral. Choose an actual eigenlift \(y_e\), and its actual conjugates \(y_{e,j}=\gamma^j(y_e)\), \(0\le j<\deg q_e\), where \(q_e\) is the characteristic polynomial of \(\gamma\) on \(T_e[1/p]\). These conjugate abelian classes form a rational basis, though not necessarily an integral basis.

Their assigned-degree-two parts are an invertible finite linear change of the old even variables plus quadratic expressions in the degree-one variables. Replacing the old even variables by these leading parts is a homogeneous triangular automorphism of the rational completed tensor algebra. Its inverse has a bounded denominator in each degree.

Write \(q_e(z)=z^{r_e}-\sum_j c_{e,j}z^j\). The actual companion error

\[
E_e=\gamma^{r_e}(y_e)
       \left(\prod_j y_{e,j}^{\,c_{e,j}}\right)^{-1}
\]

lies in \(H'\). Its degree-two part is a quadratic tensor in the free odd variables \(C_e\), since odd torsion variables have degree three. The actual coefficients \(c_{e,j}\) are integral. All roots of \(q_e\) lie in \(1+\mathfrak m\) in a finite splitting field, because a continuous pro-\(p\) action on its finite integral lattice is unipotent modulo \(p\).

For retained target weights \(m,n\), with \(z=u^{m+n}\), comparison of target length two gives

\[
q_e(z)B_{e,0;m,n}
  =C_{e;a,b}(u^m-1,u^n-1),\qquad B_{e,j;m,n}=z^jB_{e,0;m,n}.
\]

Here \(a,b\) are the odd residue characters of \(m,n\), and \(C_{e;a,b}\) is zero unless \(a+b=e\). This follows either from the associated graded \(\Delta\)-action or the identity theorem on a residue grid.

\emph{Pair-NV.} For every root \(\beta\) of every \(q_e\), there is a matching ordered pair \((a,b)\) for which

\[
C_{e;a,b}(T,\beta/(1+T)-1)\ne0.
\]

After a constant partial-fraction change of the finite even variables, their kernels are

\[
\frac{C_{e;a,b}(T_1,T_2)}
 {((1+T_1)(1+T_2)-\beta)^k},
\quad 1\le k\le {\rm mult}_{q_e}(\beta).
\]

The same numerator occurs for all pole orders in a given cyclic component. No simple-root assumption is made.

\par\medskip\noindent\textbf{The odd companion and the precise cubic condition.}

For each nonzero odd character \(o\), choose \(v_o,y_o,y_{o,j}\) in the same way. Every \(y_{o,j}\) has assigned degree at least three. Its degree-three part has the form

\[
Z_{o,j}=\sum_l b_{jl}Z_l+
  A_j(X,X,X)+B_j(X,Y)+D_j(Y,X),
\]

with \(b\) a finite invertible rational matrix. After the even change of variables, this is still a homogeneous triangular automorphism: solve for the old \(Z_l\) by subtracting the lower-variable expressions and multiplying by \(b^{-1}\). The coefficient denominators remain bounded in each fixed degree. This preserves a nonzero first assigned degree of the whole group element.

Let \(E_o\) be the actual companion error formed from \(q_o\). It belongs to \(H'\). Its degree-three part has no linear odd-torsion term, and consists only of \(XXX\), \(XY\), and \(YX\) terms. Substitute the already determined even kernels into this part. For an ordered odd character triple \((a,b,c)\) this gives a meromorphic function

\[
\mathcal C_{o;a,b,c}(T_1,T_2,T_3).
\]

Put \(z_i=1+T_i\). The function belongs to the localization of
\(\mathcal O[[T_1,T_2,T_3]][1/p]\) obtained by inverting the finitely many proper adjacent-pair factors

\[
z_1z_2-\beta,\qquad z_2z_3-\beta
\]

coming from the even companions. It has no full triple-product denominator. Its character is \(a+b+c=o\). Comparison of target length three gives

\[
q_o(z_1z_2z_3)B_{o,0}=\mathcal C_o,\qquad
B_{o,j}=(z_1z_2z_3)^jB_{o,0}
\]

on all retained weight grids where the denominators are nonzero.

\emph{Triple-NV.} For each nonzero odd character \(o\) and each root \(\alpha\) of \(q_o\), at least one matching ordered triple has nonzero restriction

\[
\mathcal C_{o;a,b,c}
 \left(T_1,T_2,\frac{\alpha}{(1+T_1)(1+T_2)}-1\right)\ne0
\tag{CU.1}
\]

in the fraction field of the bounded-denominator two-variable ring. The proper pair factors remain nonzero functions on this triple divisor, so this restriction is well defined in that fraction field.

A constant partial-fraction change of the odd companion variables yields kernels

\[
\frac{\mathcal C_{o;a,b,c}(T_1,T_2,T_3)}
 {(z_1z_2z_3-\alpha)^k},
\quad 1\le k\le {\rm mult}_{q_o}(\alpha).
\]

This retains repeated roots and the entire odd cyclic module.

\par\medskip\noindent\textbf{Justification of substitution and analytic identities.}

All constructions so far are in the actual compact free basis followed by bounded rational changes, not in a subgroup of finite abelian index. In a finite arithmetic-weight target quotient, ordinary finite-rank Magnus substitution applies after truncating the convergent free basis. Those truncations converge to the identity on \(H\). The continuous image in the finite target quotient is compact, so its coefficient functions are bounded. Thus the tensor evaluations converge, and every fixed-degree rational change preserves this convergence. Equivalently these are finite iterated integrals of bounded functions against bounded measures.

Choose a large ordinary integer lower bound on every test weight so that no singleton resonance remains and no pair or triple sum is an integer \(w>0\) with \(u^w\) a root of any companion polynomial. Only finitely many such \(w\) exist. Within each odd residue class there remain Cartesian grids with interior \(p\)-adic accumulation points, for example sufficiently large representatives followed by \(m_0+(p-1)p^n\).

For a nonidentity \(g\), choose its first nonzero assigned degree \(d\) after the invertible changes. If \(\rho'(g)=1\), its target length-\(d\) coefficient is the image of this homogeneous component alone: higher source degrees have greater target length and lower source degrees are zero. Decompose the component into ordered tilings of \(d\) positions by singletons, pairs and triples. The coefficient of a fixed tiling is a bounded-denominator function \(A\) in the singleton variables. Its image is \(A\) times the pair and triple kernels at the corresponding blocks.

For any fixed odd residue assignment the sum of these finitely many meromorphic terms vanishes on the large Cartesian grid. Clearing every pair and triple denominator gives a bounded-denominator analytic function. Repeated one-variable identity theorems imply that it is identically zero. We have an identity in the fraction field of the full bounded power-series ring.

\par\medskip\noindent\textbf{Complete separation of pair and triple poles.}

Among nonzero tilings, first maximize the number of triple blocks, then the number of pair blocks. Choose one such tiling, fix its block positions and root labels, and choose a componentwise maximal vector of its pair and triple pole orders. At each selected block choose the matching odd residue types supplied by Pair-NV or Triple-NV; these choices are independent since the blocks are disjoint. At every singleton choose its own fixed character.

For each selected block \(B\), of length two or three, set

\[
S_B=\prod_{i\in B}z_i-\alpha_B.
\]

Retain as base variables the singleton variables and all but the last variable of each selected block. Solving for that last variable,

\[
z_{\max B}=\frac{\alpha_B+S_B}
                   {\prod_{i\in B,\ i\ne\max B}z_i},
\]

is an invertible change of formal coordinates over the splitting integers. The constants \(\alpha_B-1\) are in the maximal ideal. Let \(K_{\rm base}\) be the fraction field of the bounded-denominator base-variable ring.

Here is the needed divisor fact in full. On the stratum \(S_B=0\), a coordinate monomial is constant exactly when its exponent vector is an integer linear combination of the indicator vectors of the selected blocks. This follows by eliminating the last variable of each block: the exponent of every remaining variable in a block is its original exponent minus the exponent of the eliminated variable; singleton exponents remain unchanged. Therefore an interval indicator with entries zero or one is constant exactly when it is a union of whole selected blocks and contains no singleton. An interval of length two or three cannot be a union of two selected blocks, since each selected block has length at least two. It is therefore constant only when it is exactly one selected block. All other adjacent-pair or adjacent-triple denominators are nonzero in \(K_{\rm base}\) at \(S=0\).

Thus every denominator other than a selected outer block factor is a unit in \(K_{\rm base}[[S_B]]\). This includes every proper pair denominator internal to a selected triple. A different root at the same interval also gives a unit, its constant term being a nonzero root difference. The identity embeds in the multivariate Laurent-series ring with finite negative bounds.

Extract the coefficient \(\prod_B S_B^{-k_B}\) for the selected pole-order vector. To supply all selected triple negative exponents, a term must contain each selected triple with its selected root: no pair denominator or other triple interval can do so. Maximality of the triple count forces exactly those triple blocks and forbids extras. Their internal pair denominators lie inside these blocks and cannot supply the pole of any selected external pair. Therefore all selected pairs with their roots must also occur. Maximality of the pair count now forbids extra pairs. We have isolated the same tiling and root labels.

For these remaining terms, contribution requires the pole-order vector to dominate the selected vector componentwise. A strict domination is excluded by the Pareto-maximal choice, and any incomparable vector is insufficient in some component. Wrong character labels vanish on the chosen residue grid: the sum of the two or three target residue characters determines the even or odd companion character, respectively. Thus several odd character components cause no extra ambiguity, even if their polynomials share roots. The assumption of one cyclic module per character avoids the separate problem of several indistinguishable same-character blocks. Consequently the sole possible contribution is

\[
A_{\rm selected}\,
\prod_{\text{selected pairs}}C_e|_{S_B=0}\,
\prod_{\text{selected triples}}\mathcal C_o|_{S_B=0}.
\]

The first factor is nonzero, and all the other factors are nonzero by the two conditions. Their variables are the independent base variables of disjoint blocks, so their embeddings in \(K_{\rm base}\) are injective. A field has no zero divisors. This contradicts the zero coefficient. If no nonlinear blocks occur, the same argument is just the ordinary singleton identity theorem.

Therefore \(\rho'\) is injective. Since \(\rho'\) is a quotient of \(\rho\), the original \(\rho\) is injective as well.

\end{proof}
\par\medskip\noindent\textbf{What the cubic condition measures.}

At assigned degree three, pull back to one odd component to obtain the genuine rational \(\Gamma\)-module extension

\[
0\longrightarrow W_3
\longrightarrow E_3
\longrightarrow T_o
\longrightarrow0,
\qquad
W_3=L_3(Q)\oplus[Q,T^+]
\]

in rational graded coordinates. The displayed direct sum is a vector-space description; the \(\Gamma\)-action may mix its second summand into the first via the quadratic even companion. The odd companion represents the extension class modulo \(q_o(\gamma)W_3\). Evaluating the even letters gives the meromorphic numerator used in (CU.1).

Changing an actual lift of the same cyclic class changes its leading target kernel by a function \(A\) having only proper pair denominators. Hence

\[
\mathcal C'_o-\mathcal C_o=q_o(z_1z_2z_3)A.
\]

Its restriction to a root divisor is unchanged. Replacing the rational cyclic generator by \(f(\gamma)v_o\), with \(f\) invertible modulo \(q_o\), multiplies the restricted numerator by the nonzero scalar \(f(\alpha)\), up to such a vanishing term. The property is therefore independent of these choices. Rank-one free-coordinate changes multiply by invertible one-variable factors; changing even coordinates changes the expression but not the actual triple pole of the odd lift.

This is a specified class-three extension condition. Calling it an ordinary triple Massey product without further work would omit its defining system and indeterminacy. The quadratic cup theorem alone does not determine it. In a class-three truncated free Lie presentation, the action on an odd degree-three letter can be changed by a matching element of \(L_3(Q)\subset W_3\), leaving the abelian module and even the ordinary class-two quotient unchanged. The degree-three extension is therefore extra data. In particular a zero cubic companion is compatible with the same lower-degree data. This is a formal compatibility observation, not an arithmetic counterexample.

\par\medskip\noindent\textbf{Arithmetic interpretation and the gates still open.}

For the arithmetic application take \(p\ge5\), with the smaller primes already covered separately. The earlier arithmetic source checks establish finite \(\Lambda\)-generation, freeness of \(H\), and the actual rational target. For nonexceptional odd character \(a\), the complete proof of Wake--Wang-Erickson \cite[Proposition~2.2.6 and proof, pp.~3--4]{WakeWangErickson2017}, gives

\[
\operatorname{Ext}^1_\Lambda(M_a^\#(1),\Lambda)
   =X_{\omega a^{-1}},
\qquad
\operatorname{Ext}^2_\Lambda(M_a^\#(1),\Lambda)=0.
\]

The left module is the \(p\)-ramified module (the source's fraktur symbol); the right module is unramified. The rank of \(M_a\) is one independently of Greenberg, by the local-unit exact sequence in that proof. After inverting \(p\), the PID decomposition shows that its torsion elementary divisors, after the indicated involution and twist, are exactly the elementary divisors of the even unramified module. A root \(\alpha\) on the odd torsion corresponds to the root \(u/\alpha\) on that even module. Thus cyclicity character by character of \(X^+[1/p]\) gives the contemplated cyclic odd companions; a single nonzero cyclic even component is the first concrete case. The exceptional odd characters were treated separately by Coleman and Stickelberger in the established Greenberg section; this extension does not remove those exceptions silently.

The even module must still satisfy the rational cyclicity and Pair-NV assumptions. Without Greenberg these do not follow merely from its characteristic ideal. Moreover odd-\(H^1\) cup surjectivity, by itself, no longer forces the free-\(Q\) quadratic component to be nonzero: classes supported on odd torsion may contribute to those cups. A possible stronger route is to use norm-compatible cyclotomic-unit classes, whose universal cohomology restrictions land in a torsion-free Iwasawa coefficient module and hence annihilate \(\Lambda\)-torsion. The exact integral Shapiro and coefficient-specialization argument for that route has not been supplied here, so it is not used in the theorem.

Most decisively, no source presently checked supplies (CU.1) for the arithmetic odd companion. Even if Pair-NV and all cyclicity hypotheses are provided, the new cubic residue remains an explicit independent gate. The divisor argument shows that mixed even/odd nested denominators do not themselves block this first extension; the missing input is arithmetic noncancellation in the actual class-three extension.

The displayed module comparison uses the inspected author version \cite{WakeWangErickson2017}. The full Proposition~2.2.6 proof was read again for this extension. The algebraic conclusion retains both pole conditions as hypotheses.

\subsection{Norm-compatible cyclotomic units detect the quadratic companions}
\label{sec:universal-cyclotomic-pair}

Odd Iwasawa torsion creates a distinction between arbitrary odd
cohomology classes and the cyclotomic classes used by the modular-symbol
theorem.  The latter annihilate that torsion on restriction to the
cyclotomic kernel.  This supplies the quadratic pole condition under
precise hypotheses of Wake--Wang-Erickson, without assuming that the
whole plus Iwasawa module is finite.  The cubic pole condition remains
a separate arithmetic hypothesis.

\begin{lemma}[K17-NORM-TORSION]\label{lem:norm-torsion}
Let \(P=H\rtimes\Gamma\) be a pro-\(p\) group, where
\(\Gamma=\langle\gamma\rangle\cong\mathbb Z_p\).  Let the coefficient
character \(\kappa\) be trivial on \(H\), and write
\(\kappa(\gamma)=u\in1+p\mathbb Z_p\).
Put \(P_n=H\rtimes\Gamma^{p^n}\),
\(\Lambda=\mathbb Z_p[[\Gamma]]\), and \(M=H^{\rm ab}\).
Suppose that classes
\[
 c_{n,r}\in H^1(P_n,\mathbb Z/p^r(1)),
 \qquad n\geq0,\quad r\geq1,
\]
are compatible with corestriction in \(n\) and coefficient reduction
in \(r\).  Then every restriction
\[
 c_{n,r}|_H:M\longrightarrow\mathbb Z/p^r(1)
\]
annihilates the entire \(\Lambda\)-torsion submodule of \(M\).
\end{lemma}

\begin{proof}
Write \(a_{n,r}\) for this restriction.  Since \(H\) acts trivially
on the coefficients, it is an actual continuous homomorphism and is
unchanged by replacing a cocycle by a cohomologous one.  The cocycle
identity gives
\[
 a_{n,r}(\gamma^{p^n}m)=u^{p^n}a_{n,r}(m).
\]
Choose a compatible Tate basis to write scalar coordinates, and define
\[
 \Phi_{n,r}(m)=
 \sum_{j=0}^{p^n-1}u^{-j}a_{n,r}(\gamma^j m)[\gamma^j]
 \ \in(\mathbb Z/p^r)[\Gamma/\Gamma^{p^n}].                 \tag{NC.1}
\]
Give this coefficient module the \(P\)-action with \(H\) trivial and
\(\gamma\) acting by multiplication by \(u[\gamma]^{-1}\).
Shifting the sum, including its endpoint, proves
\[
 \Phi_{n,r}(\gamma m)=u[\gamma]^{-1}\Phi_{n,r}(m).
                                                               \tag{NC.2}
\]
The restricted corestriction formula, with representatives
\(\gamma^{-kp^n}\), is
\[
 a_{n,r}(m)=
 \sum_{k=0}^{p-1}u^{-kp^n}a_{n+1,r}(\gamma^{kp^n}m).
                                                               \tag{NC.3}
\]
It follows that ordinary group-ring projection from level \(n+1\)
to level \(n\), which sums coefficients over each fiber, sends
\(\Phi_{n+1,r}\) to \(\Phi_{n,r}\).  Coefficient reduction is also
compatible.  The inverse limit of these actual homomorphisms is
therefore a continuous equivariant map
\[
                  \Phi:M\longrightarrow\Lambda^\#(1).     \tag{NC.4}
\]
This is integral Shapiro in restriction coordinates.  The inverse
group action and corestriction convention agree with
\cite[\S2.5, pp.~16--17]{SharifiReciprocity2017};
the finite formulas above prove the part needed here directly.
In particular, no interchange of cohomology and inverse limit or
vanishing of a \(\varprojlim^1\) term is used.

If \(f(\gamma)t=0\), where \(f\in\Lambda\) is nonzero, then
equivariance implies
\[
                   f(u\gamma^{-1})\Phi(t)=0.
\]
The substitution \(\gamma\mapsto u\gamma^{-1}\) is an involutive
continuous automorphism of \(\Lambda\): in the parameter
\(Z=\gamma-1\), it is \(Z\mapsto u/(1+Z)-1\).
Since \(\mathbb Z_p[[Z]]\) is a domain, \(\Phi(t)=0\).
The coefficient of the identity in (NC.1) is \(a_{n,r}(t)\).
Hence every \(a_{n,r}(t)\) is zero, as claimed.
\end{proof}

\paragraph{All finite cyclotomic inputs.}
For \(p\) odd, put \(F^{(n)}=\mathbb Q(\mu_{p^{n+1}})\) and choose
compatible roots of unity.  For every \(a\in\mathbb Z_p^\times\),
\[
 v_n(a)=1-\zeta_{p^{n+1}}^{\,a\bmod p^{n+1}}
 \quad\hbox{satisfies}\quad
 N_{F^{(n+1)}/F^{(n)}}v_{n+1}(a)=v_n(a).                    \tag{NC.5}
\]
Kummer maps commute with norm and corestriction.  Their finite
classes therefore satisfy Lemma~\ref{lem:norm-torsion}.
They are classes of the maximal pro-\(p\) quotient: the Tate
action factors through it, and the degree-one comparison in
Lemma~\ref{lem:etale-prop} is an isomorphism, natural for the
tower transfers.  Equivalently, the Galois closure of each Kummer
extension is a \(p\)-extension, since the base contains \(\mu_p\).

Nonunit exponents cause no loss.  Every nonzero exponent modulo
\(p^{n+1}\) has the form \(p^s b\), with \(0\leq s\leq n\) and
\(b\) a unit, and
\[
 1-\zeta_{p^{n+1}}^{\,p^s b}
  =\prod_{j=0}^{p^s-1}
       \left(1-\zeta_{p^{n+1}}^{\,b+jp^{n+1-s}}\right).
                                                               \tag{NC.6}
\]
Every factor has a unit exponent and is a value of a family in
(NC.5).  Formula (NC.6) is the identity
\(\prod_{j=0}^{p^s-1}(1-X\eta^j)=1-X^{p^s}\) for a primitive
\(p^s\)-th root \(\eta\); it takes place in the same finite field
and involves no division by the degree.  Thus the restriction of
every finite cyclotomic Kummer input to \(H\) annihilates all of
\(T=M_{\Lambda\text{-tors}}\).

Complex conjugation \(c\) commutes with the norms, so averaging
these families by \((1+c)/2\) preserves this property.
For a \(p\)-unit \(v\), the quotient \(v/c(v)\) is a global unit
of absolute value one at every embedding: the valuations at the
unique prime above \(p\) cancel.  It is therefore a root of unity.
The negative parts of the Kummer classes consequently lie in one
cyclic subgroup.  Their mutual cups vanish, because self-cup is
zero when \(2\) is invertible.  Mixed-parity cups lie in the odd
target.  The even target spanned by cyclotomic symbols is thus
spanned by cups of \(c\)-even classes whose restrictions factor
through \(Q=M/T\).

\begin{proposition}[K17-UNIVERSAL-PAIR]\label{prop:universal-pair}
Let \(p\) be odd and \(P=H\rtimes\Gamma\) as above, with \(H\)
free pro-\(p\).  Suppose that
\(\Delta\cong\mathbb F_p^\times\) acts on \(H\), commutes with
\(\Gamma\), and extends \(\kappa\) to \(P\rtimes\Delta\).
Let \(c\) be its involution, and assume \(\kappa(c)=-1\).
Assume that \(M\) is finitely generated over \(\Lambda\), that
\(T=M_{\Lambda\text{-tors}}\) is finite free over \(\mathbb Z_p\),
and that \(Q=M/T\) is torsion free over \(\mathbb Z_p\), has only
negative \(c\)-parity, and is free of finite rank over
\(R=\Lambda[1/p]\) after inverting \(p\).
Assume that every even character component \(T_e^+[1/p]\) is
cyclic over \(R\).

Suppose that, on a cofinal diagonal of group and coefficient
levels,
\[
 H^2(P_n,\mathbb Z/p^n(2))^+
\]
is generated by cups of even classes whose restrictions to \(H\)
annihilate \(T\).  Then the actual free-\(Q\) quadratic companions
satisfy Pair-NV: for every root \(\alpha\) of the characteristic
and minimal polynomial \(q_e\) of \(\gamma\) on
\(T_e^+[1/p]\), at least one matching coefficient of \(C_e\) is
nonzero on
\[
                     (1+T_1)(1+T_2)=\alpha.                \tag{NC.7}
\]
Repeated roots are allowed.  The coefficients are in the
bounded-denominator coordinates of an \(R\)-basis of \(Q[1/p]\).
\end{proposition}

\begin{proof}
\par\medskip\noindent\textbf{The actual extension.}
The class-two quotient of \(H\) has its integral
Baker--Campbell--Hausdorff Lie structure, since \(p\) is odd.
Its central commutator module is
\(\widehat{\bigwedge}^2_{\mathbb Z_p}M\).
Quotient this module by all brackets with a \(T\) input.
The remaining central module is
\[
                      N=\widehat{\bigwedge}^2_{\mathbb Z_p}Q.
\]
For example, exactness follows from a continuous
\(\mathbb Z_p\)-splitting of \(M\to Q\), which exists because
compact torsion-free \(\mathbb Z_p\)-modules are products of
copies of \(\mathbb Z_p\).  No equivariant splitting is asserted.

The central module \(N\) has positive \(c\)-parity.  Hence the
negative part of the quotient Lie module maps isomorphically
onto \(M^-\), giving a canonical \(\Gamma\)-equivariant section.
The lift of \(T^-\) is central, because brackets with \(T\)
were killed.  Quotient this central equivariant lift as well.
The odd part is now the canonical section from \(Q\), and the
even part is the extension
\[
          0\longrightarrow N\longrightarrow E
                 \longrightarrow T^+\longrightarrow0.     \tag{NC.8}
\]
For a cyclic vector \(v_e\) and an even lift \(y_e\) in its
character component, set
\[
                    q_e(\gamma)y_e=C_e\in N[1/p].           \tag{NC.9}
\]
This is the quadratic free-\(Q\) companion used in
Theorem~\ref{thm:odd-torsion-cubic}.  Brackets with torsion
lifts vanish in this quotient, so its ordered companion product
is the additive characteristic-polynomial expression.
The logarithm and Magnus quadratic coefficient agree on a
derived element.

\par\medskip\noindent\textbf{A vanishing residue gives a split socle.}
Pass to a finite splitting field \(L/\mathbb Q_p\), with integers
\(\mathcal O\).  Finite generation clears the denominators of the
coordinates \(Q[1/p]\cong R^d\) and their inverse.  Completed
tensor and exterior squares retain a fixed denominator bound.
Thus \(N[1/p]\) is the alternating bounded-denominator
two-variable coefficient module.  Continuity of the pro-\(p\)
action on \(T\) gives \(\alpha-1\in\mathfrak m_{\mathcal O}\).

Suppose every coefficient of \(C_e\) vanishes on (NC.7).
The integral formal-coordinate change
\[
 S=(1+T_1)(1+T_2)-\alpha,\qquad
 T_2=\frac{\alpha+S}{1+T_1}-1
\]
has specialization kernel \((S)\).  Division preserves the
common denominator bound and, since \(S\) is symmetric, the
alternating property.  Hence
\[
                        C_e=(\gamma-\alpha)A
                        \quad\hbox{in }N_L.
\]
Let \(r_e(z)=q_e(z)/(z-\alpha)\).  Cyclicity gives a nonzero
vector \(w=r_e(\gamma)v_e\) spanning the \(\alpha\)-socle,
including when that root is repeated.  Then
\[
                       \widetilde w=r_e(\gamma)y_e-A
\]
is an equivariant lift of \(w\).  Put \(W=Lw\) and
\(W_{\mathcal O}=W\cap T_{\mathcal O}^+\).
This lattice is saturated, hence a direct summand over
\(\mathcal O\).  Choose an \(\mathcal O\)-linear section
\(s:T_{\mathcal O}^+\to E_L\) extending the equivariant section
on \(W_{\mathcal O}\).  One fixed power \(p^b\) clears its
denominators.  Write the class-two logarithm of \(h\in H\) as
\[
 x(h)+s(t(h))+z(h),\qquad
 x(h)\in Q_{\mathcal O},\quad
 t(h)\in T_{\mathcal O}^+,\quad z(h)\in N_L.
\]
The same \(p^b\) makes every \(z(h)\) integral.  The section
defect
\[
            \beta_\delta(t)=\delta s(t)-s(\delta t)
            \quad(\delta\in\Gamma)
\]
vanishes on \(W\).

\par\medskip\noindent\textbf{The actual cup and its dual quotient.}
Average two even cocycles \(f,g\) under \(c\).
They vanish on the chosen \(\Gamma^{p^n}\) section, since \(c\)
fixes its elements and acts by \(-1\) on the coefficients.
Their restrictions to \(H\) factor through \(Q\), by hypothesis.
Let
\[
 \psi=f\wedge g:N_{\mathcal O}\longrightarrow
                  \mathcal O/p^n(2).
\]
The integral cochain
\[
       k(h)=-\frac{p^b}{2}f(h)g(h)-\psi(p^b z(h))            \tag{NC.10}
\]
has differential \(dk(h_1,h_2)=p^b f(h_1)g(h_2)\).
Indeed the central coordinate of the product is
\(z_1+z_2+[x_1,x_2]/2\).  The two contributions to the
differential are
\[
 \frac{p^b}{2}(f_1g_2+f_2g_1),\qquad
 \frac{p^b}{2}(f_1g_2-f_2g_1).
\]
Extending \(k\) by \(k(h\delta)=k(h)\), its equivariance
defect computes the Hochschild--Serre edge of the cup.
Up to the convention for its sign, this is
\[
                         \psi(p^b\beta_\delta(t)),          \tag{NC.11}
\]
which vanishes on \(W\).

Both \(H\) and \(\Gamma^{p^n}\) have \(p\)-cohomological
dimension at most one.  All odd components of \(M\), including \(T^-\),
give negative parity in \(\operatorname{Hom}(M,\mathcal O/p^n)(2)\).
Therefore
\[
 H^2(P_n,\mathcal O/p^n(2))^+
 \cong H^1\bigl(\Gamma^{p^n},
         \operatorname{Hom}(T^+,\mathcal O/p^n)(2)\bigr).
                                                               \tag{NC.12}
\]
Set \(V=\operatorname{Hom}_{\mathbb Z_p}(T^+,\mathcal O)(2)\),
and let \(B\) denote its \(\gamma\) operator.  Normalized
generator evaluation identifies the last group with
\[
                   V/\bigl(p^nV+(B^{p^n}-1)V\bigr).         \tag{NC.13}
\]
Corestriction sends a generator value to the same value in
the larger group's coinvariants.  Its transition is consequently
the quotient map, with no extra factor \(p\).
As \(B^{p^n}-1\to0\), the denominator submodules are cofinal
with the \(p\)-adic neighborhoods, and the inverse limit is \(V\).

Restriction to \(W_{\mathcal O}\) gives a nonzero torsion-free
quotient
\[
 V\longrightarrow
       \operatorname{Hom}_{\mathcal O}(W_{\mathcal O},\mathcal O)(2).
\]
It is surjective because \(W_{\mathcal O}\) is saturated.
Its rational root is \(u^2/\alpha\), and its character is
\(\kappa^2 e^{-1}\).
By (NC.11) and finite-level cup generation, \(p^b\) times every
finite target has zero image in the corresponding quotient of
this dual line.  The separated inverse limit would make \(p^b\)
annihilate the displayed nonzero torsion-free quotient,
a contradiction.  The same \(b\) works at all levels; no
compatible expressions as sums of cups are required.
This proves (NC.7) for all roots.
\end{proof}

\begin{corollary}[K17-WWE-PAIR]\label{cor:wwe-pair}
Let \(p\geq5\), and use the arithmetic cyclotomic group and
notation of the preceding kernel theorems.
For every even character \(\theta\ne1,\omega^2\), put
\(\chi=\omega^{-1}\theta\).
Let \(\xi_\chi\) be a characteristic series of \(X_\chi(1)\)
in the Wake--Wang-Erickson ring and twist convention.
For every nonzero component suppose that
\[
 \begin{gathered}
 X_\chi(1)[1/p]\text{ is cyclic over }
                         (\Lambda/\xi_\chi)[1/p],\\
 X_\theta\otimes_\Lambda(\Lambda/\xi_\chi)[1/p]=0.
 \end{gathered}                                             \tag{NC.14}
\]
Then every rational even torsion character is cyclic, and the
actual free-\(Q\) quadratic companions satisfy Pair-NV.
This conclusion allows nonzero odd \(\Lambda\)-torsion in \(M\).
\end{corollary}

\begin{proof}
The conditions are exactly those of
\cite[Corollary~8.1.3(2), p.~49]{WakeWangEricksonPseudo2018}.
Its conclusion is an integral isomorphism
\[
 \varpi:
 (\mathcal H^-_\theta/I_\theta\mathcal H^-_\theta)/
           (p\text{-primary torsion})
       \ \xrightarrow{\ \sim\ }\ X_\chi(1).
\]
In particular, the original \(\varpi\) is surjective.
The full proof allows repeated roots; it does not require the
derivative of the characteristic series to be invertible.
At tame level one the primitive-character and other setup
conditions exclude exactly \(1,\omega^2\) among the even
characters \cite[\S1.1, pp.~2--3]{WakeWangEricksonPseudo2018}.
We apply its nonunit-series running hypothesis only to nonzero
components.  A unit characteristic series makes \(X_\chi\)
finite, and its \(p\)-torsion freeness makes it zero.

Fukaya--Kato turn this surjectivity into finite-level generation
by \(\{1-\zeta^a,1-\zeta^b\}\) for every cyclotomic level
\cite[\S\S7.1.8--7.1.10, p.~86]{FukayaKato2012}.
Their exponents are nonzero and at least one is prime to \(p\);
(NC.6) handles every input, including a nonunit exponent.
Their \S7.1.9 identifies the reduction modulo \(p^r\) with
the corresponding class-group quotient.

The excluded even targets \(1,\omega^2\) correspond to
\(X_{\omega^{-1}},X_\omega\), which vanish by the
Stickelberger exception calculation; the first is explicitly
recorded in \cite[proof of Lemma~5.2.7, p.~13]{WakeWangErickson2017}.
Norm projections from \(X^-\) onto every finite minus class
group are surjective
\cite[Proposition~3.4.1 and proof, pp.~90--91]{SharifiIwasawaNotes}.
Thus their finite targets also vanish.  At \(F^{(n)}\), the
Fukaya--Kato coefficient index is \(n+1\); reduction gives
coefficients \(p^n\), as required in
Proposition~\ref{prop:universal-pair}.  The Brauer term is zero,
since there is one omitted finite prime and no real place.

By Lemma~\ref{lem:etale-prop}, Kummer classes lift uniquely to
group \(H^1\), and group \(H^2\) injects into etale \(H^2\).
Lifting the generating cup expressions makes the latter
injection surjective on the even part.  This comparison is
natural for the tower transfers and assumes no arithmetic
\(K(\pi,1)\) theorem.  The paragraph following
Lemma~\ref{lem:norm-torsion} supplies even generating inputs
whose restrictions annihilate \(T\).

The unconditional module calculations in the preceding
arithmetic sections give \(Q[1/p]\) rank one in every odd
character, and zero in even characters.  The first condition
in (NC.14), and the Kummer pseudoisomorphism
\[
                         M^+\sim(X^-)^\iota(1),
\]
give cyclicity of every rational even component
\cite[Proposition~3.4.7 and Corollary~3.4.8, p.~92]{SharifiIwasawaNotes}.
An even companion root \(\alpha\) of character \(\omega^e\)
corresponds to the target root \(u^2/\alpha\) of character
\(\omega^{2-e}\), agreeing with both the involution and Tate
twists above.  Proposition~\ref{prop:universal-pair} now applies.
\end{proof}

\paragraph{What this removes, and what remains.}
Condition (NC.14) retains the exact rational support condition
of its source; it is not replaced by an untwisted shared-root
test or by finiteness of all \(X^+\).
Greenberg implies these conditions, but no new arithmetic
verification of them with \(X^+[1/p]\ne0\) is asserted.
Combining Corollary~\ref{cor:wwe-pair} with
Theorem~\ref{thm:odd-torsion-cubic} removes Pair-NV as a separate
hypothesis whenever (NC.14) holds.  Rational cyclicity of each
odd torsion character and the actual Triple-NV condition must
still be supplied.  In particular the norm-family argument
does not establish arithmetic cubic noncancellation or a new
unconditional kernel range.

\subsection{The integral comparison at \texorpdfstring{$2$}{2}}

Theorem~\ref{thm:pro2-kernel} and rational Deligne--Ihara do not prove surjectivity of
\[
 \operatorname{Lie}_{\mathbb Z_2}(s_3,s_5,\ldots)
      \longrightarrow\mathfrak g_2.
\]
Nonzero integral odd classes exist as explained above. Choosing odd classes whose rational images form a free generating system gives an injective integral free source. Equality with the actual lattice requires an additional argument. Theorem~\ref{thm:2-allweights} below supplies that argument in every weight: it constructs primitive adapted odd generators, proves integral freeness, and uses the exact outer-filtration comparison. It does not prescribe unit values for a separately normalized Soul\'e family.

There is a useful rational normal form, but it does not remove that question. Choose a Levi subgroup containing the image of $c$. Since $chc=h^{-1}$, the element $a=\log h$ has only odd total weights. Writing $\gamma=vu$, the element $\exp(a/2)v$ is $c$-fixed. A further even unipotent conjugation can remove this factor because $u^w-1\neq0$ in $\mathbb Q_2$ for every positive weight $w$, giving a rational form $\gamma=\exp(-a/2)u$. Division by $2$ and inversion of $u^w-1$ need not preserve an integral lattice. This explains the denominators in the finite-weight lattice calculation. The subsequent parity induction uses actual group commutators to remove them where needed and supplies the integral free basis.

\subsection{Mixed Tate motives and the current \texorpdfstring{$\mathfrak{grt}_1$}{grt1} statement}

Deligne--Goncharov's category over $\mathbb Z$ has rational coefficients. In the attached LaTeX, \S1.6, Propositions~1.7 and~1.9, and Proposition~2.3 establish the rational Ext and free-Lie statements; \S2.4 records the noncanonical nature of free lifts \cite{DeligneGoncharov2005}. The base $\operatorname{Spec}\mathbb Z$ specifies ramification conditions. It does not replace $\mathbb Q$-coefficients by $\mathbb Z$ or $\mathbb Z_p$. Rational $\operatorname{Ext}^2$ vanishing erases torsion and is not an integral vanishing result.

The rational Grothendieck--Teichm\"uller status must also be updated from current primary sources. Hirose's October2026 manuscript, arXiv:2610.23719v1, Theorem~4 and \S1.2, gives a Hoffman spanning result over $\mathbb Z_{(2)}$ for the confluence quotient and derives rational equality of the de Rham mixed-Tate motivic Galois group with GRT \cite{Hirose2026}. We checked the statement and its stated implication, not the whole proof. This is a separate current rational result. It does not identify the arithmetic $\mathbb Z_p$-image lattice, and it does not invalidate the counterexample at $37$. An integral associator lattice, a $\mathbb Z_{(2)}$-span of motivic periods, and the lattice of an arithmetic pro-$p$ Galois image are different objects unless a comparison theorem is supplied.

The supplied round-two war-game file mentions B3 and S9 but does not reproduce their original round-one rows. We have therefore used their subject as specified in the prompt and have not assigned an invented exact locator to missing text.

% Bibliography keys needed in this fragment and pro2_section.tex:
% Sharifi2002: Relationships between conjectures on the structure of pro-p Galois groups unramified outside p; https://www.math.ucla.edu/~sharifi/publ/galstr.pdf
% Brown2012: Mixed Tate motives over Z, Annals175(2012)949-976; https://annals.math.princeton.edu/wp-content/uploads/annals-v175-n2-p10-p.pdf
% Brown2019: From the Deligne-Ihara conjecture to multiple modular values; https://arxiv.org/pdf/1904.00179
% HainMatsumoto2001: Weighted completion of Galois groups and Galois actions on the fundamental group of P1-{0,1,infinity}; https://arxiv.org/pdf/math/0006158
% GrasJaulent2008: Note on 2-rational fields; https://arxiv.org/pdf/0805.1168
% MinacPasiniQuadrelliTan2020: Koszul algebras and quadratic duals in Galois cohomology; https://arxiv.org/pdf/1808.01695
% IshiiKernel2026: On the kernels of the pro-p outer Galois representations associated to once-punctured CM elliptic curves; arXiv2312.04196v3
% IshiiKummer2024: On Kummer characters arising from Galois actions on the pro-p fundamental groups of once-punctured CM elliptic curves; arXiv2312.04175v2
% McCallumSharifi2003: A cup product in the Galois cohomology of number fields; https://www.math.ucla.edu/~sharifi/publ/pairing.pdf
% BuhlerHarvey2011: Irregular primes to 163 million; arXiv0912.2121
% Iwasawa1981: Riemann-Hurwitz formula and p-adic Galois representations for number fields, Tohoku33(1981)263-288; https://www.jstage.jst.go.jp/article/tmj1949/33/2/33_2_263/_pdf/-char/en
% SharifiIwasawaNotes: Iwasawa Theory, Chapter3; https://www.math.ucla.edu/~sharifi/notes/iwasawa-ch03.html
% DeligneGoncharov2005: Groupes fondamentaux motiviques de Tate mixte; attached LaTeX Sec1ccc/Sec2ccc.
% Hirose2026: Confluence relations, Hoffman basis over Z_(2), and the Deligne-Drinfeld conjecture; arXiv2610.23719v1.

\subsection{Exact integral comparison and a finite calculation at 2}

\begin{proposition}[K17-OUT-SAT]\label{prop:outer-sat}
For a finitely generated free pro-$p$ group $\Pi$ of rank at least
two, every positive graded quotient of its lower-central
outer-automorphism filtration is a torsion-free $\mathbb Z_p$-module.
For an arithmetic outer action covered by Hain--Matsumoto
Proposition~8.2, the integral pullback filtration and its rational
unipotent pullback therefore agree exactly in positive degrees.
This includes $p=2$.
\end{proposition}

\begin{proof}
Write $\Gamma_n$ for the closed lower central series of $\Pi$,
$L_n=\Gamma_n/\Gamma_{n+1}$ and $V=L_1$. The graded Lie algebra
$L$ is the free $\mathbb Z_p$-Lie algebra on $V$, for every prime
$p$, by Koch--Kukkuk--Labute \cite[Theorem~2.1, pp.~48--49;
author PDF pp.~4--5]{KochKukkukLabute1998}.
Put
\[
 A_n=\ker\bigl(\operatorname{Aut}\Pi\longrightarrow
                    \operatorname{Aut}(\Pi/\Gamma_{n+1})\bigr).
\]
The map sending $\alpha$ to
$x\mapsto\alpha(x)x^{-1}\bmod\Gamma_{n+2}$ is
$\mathbb Z_p$-linear in $x\in V$, additive in $\alpha$, and has
kernel $A_{n+1}$. These assertions follow from the commutator
identities and the fact that $\alpha$ acts trivially modulo
$\Gamma_{n+1}$. Thus it embeds $A_n/A_{n+1}$ in
$\operatorname{Hom}_{\mathbb Z_p}(V,L_{n+1})$.

Every outer automorphism trivial modulo $\Gamma_{n+1}$ has an
$A_n$ representative: lift its quotient inner conjugator to
$\Pi$ and remove it. An inner conjugation by $g$ belongs to
$A_n$ precisely when $g\in\Gamma_n$. Indeed a first nonzero
degree of $g$ below $n$ would centralize all free Lie generators.
The center of a free Lie algebra of rank at least two over any
field is zero in positive degree. To see this, embed it in the
free associative algebra. A polynomial commuting with a free
letter $X$ contains only powers of $X$, by comparing word
coefficients and moving initial $X$'s to the right. Commutation
with another independent letter then makes it scalar, hence zero
in positive degree. The characteristic-zero conclusion also
holds over $\mathbb Z_p$ by tensoring.

It follows that the positive outer graded quotient embeds in
\[
 \operatorname{Hom}_{\mathbb Z_p}(V,L_{n+1})/\operatorname{ad}(L_n).
\]
The ambiguity is exactly the inner degree-$n$ map; if a Johnson
class is $\operatorname{ad}(z)$, an integral lift of $z$ removes
it and puts the representative in $A_{n+1}$. The displayed
quotient is torsion-free. In fact $\operatorname{ad}$ is injective
modulo $p$, since the free Lie algebra over $\mathbb F_p$ has zero
center. Thus $\operatorname{ad}(z)$ divisible by $p$ forces $z$
divisible by $p$, so its image is primitive in the finite free
target. This also proves torsion-freeness for every arithmetic
subgroup with the induced filtration.

Let $I^n\subseteq J^n$ be the integral and rational pullbacks.
They agree at $n=1$, since identity on the free integral
abelianization is equivalent to identity after tensoring with
$\mathbb Q_p$. If $I^n=J^n$, Hain--Matsumoto
\cite[Proposition~8.2, p.~25]{HainMatsumoto2001} gives a finite
kernel on $\operatorname{gr}_I^n\to\operatorname{gr}_J^n$.
The source is torsion-free, so the kernel is zero. Every element
of $J^{n+1}\subset I^n$ maps to that kernel and is in $I^{n+1}$.
Induction proves the asserted exact equality.
\end{proof}

\begin{theorem}[K17-2-FINITE18]\label{thm:2-finite18}
There are primitive integral odd classes $\sigma_m$ in the actual
tripod image at $2$ such that, in every arithmetic weight at most
$18$, their Hall brackets generate the entire integral graded
Lie algebra. This is an independently checked finite-range result;
Theorem~\ref{thm:2-allweights} below proves the all-weight assertion
by a separate induction.
\end{theorem}

\begin{proof}
We give the mathematical reduction and the exact finite
certificate. Use the checked presentation $C_2*\mathbb Z_2$,
choose $\chi(\gamma)=5$, and put $h=c\gamma c\gamma^{-1}$.
After placing $c$ in a rational Levi torus, $a=\log h$ has only
odd weights because $chc=h^{-1}$. Write $\gamma=v u$ with
$u=5$. The element $z=\exp(a/2)v$ is $c$-fixed. A $c$-fixed
conjugation removes $z$ successively in positive even weights:
the equation in weight $w$ divides by $5^w-1\ne0$ in
$\mathbb Q_2$. We obtain
\[
 h=\exp(a),\qquad \gamma=\exp(-a/2)u.
\]
Every odd homogeneous component of $a$ has nonzero
indecomposable coordinate, by density and the compact orbit
generation of the cyclotomic kernel. Brown's rational freeness
therefore permits choosing these components as the free
generators $s_m$, so that
\[
 a=\sum_{m\ge3,\ m\ {\mathrm{odd}}}s_m,\qquad
 D=\operatorname{Ad}(\exp(-a/2))\circ\tau,
 \quad\tau(s_m)=5^m s_m.
\]
All changes here transport the actual compact image and its
lattice; they do not declare the $s_m$ to be integral generators.
By Proposition~\ref{prop:outer-sat}, its transported weight
filtration is exactly the original integral filtration.

At cutoff $W$, work in the finite-dimensional free associative
algebra over $\mathbb Q$ truncated by arithmetic weight. The
finite exponential and logarithm give exact products and
commutators, retaining every denominator. The desired compact
group $H_W$ is the closed group generated by the continuous
$D$-orbit of $\exp(a)$.

Order a Hall--Lyndon basis by increasing weight, putting the
primitive letter first within each odd weight. Store group
generators by their logarithms. If a candidate has first
coordinate $b$ and the current pivot has first coordinate $c$,
eliminate it by the corresponding group power whenever
$v_2(b/c)\ge0$. If its valuation is smaller, replace the pivot
and reduce the old pivot in turn. Normalize only by unit
$\mathbb Z_2$-powers. Add conjugates under $D$ and pair
commutators until the output stabilizes. Every constructed
element belongs to $H_W$.

Here is why the finite checks certify the whole compact group.
For a triangular final list $g_1,\ldots,g_N$, put
\[
 C_i=g_i^{\mathbb Z_2}g_{i+1}^{\mathbb Z_2}\cdots
                         g_N^{\mathbb Z_2}.
\]
Every commutator has greater weight than both inputs. Its checked
zero reduction expresses it in the later ordered factors.
Descending induction makes each $C_i$ a subgroup: $g_i$
normalizes the already closed subgroup $C_{i+1}$, first for
nonnegative integral powers and then for all $\mathbb Z_2$
powers by continuity. In particular
$g_i^{2^r-1}\to g_i^{-1}$ supplies inverse conjugation.
Each $C_i$ is compact by its continuous parametrization.
Triangular first coordinates make that parametrization
injective and successive reduction an exact membership test.

The final certificate verifies that $h$, every $D(g_i)$,
every $D^{-1}(g_i)$ and every potentially nonzero pair
commutator reduce to the identity. Consequently $C_1$ contains
the whole compact orbit and $H_W\subseteq C_1$. The construction
gives the reverse inclusion. Hence $C_1=H_W$.

Vanishing of coordinates below weight $n$ is equivalent to
vanishing of all preceding product exponents. Its complete
weight-$n$ lattice is thus the $\mathbb Z_2$-span of the leading
weight-$n$ pivot vectors. Select the primitive odd pivots and
compute all their Hall brackets. They are contained in this
lattice. The quotient of the two determinants measures its
index; equality of their $2$-adic valuations proves equality
of the lattices.

The exact rational computation at $W=18$ gives the following
table. Columns $d_n$ and $d_n^{\rm odd}$ are determinant
valuations in the stated rational Hall coordinates.
\[
\begin{array}{c|r|r|r@{\qquad}c|r|r|r}
n&\mathrm{rank}&d_n&d_n^{\rm odd}&
n&\mathrm{rank}&d_n&d_n^{\rm odd}\\\hline
3&1&0&0&5&1&3&3\\
7&1&7&7&8&1&3&3\\
9&1&10&10&10&1&7&7\\
11&2&18&18&12&2&20&20\\
13&3&31&31&14&3&31&31\\
15&4&52&52&16&5&66&66\\
17&7&96&96&18&8&114&114
\end{array}
\]
All other positive weights at most $18$ have zero rational
space and hence zero integral graded group. Every listed index
is one. The complete rational vectors, arithmetic operations,
closure checks and output are in
\path{evidence/work/pro2_integral_malcev.py} and
\path{evidence/work/pro2_integral_malcev_extended.json}.
An independent rerun through $14$ and independent two-way
integral change-of-basis checks through $18$ are recorded in
the accompanying audit JSON files. The algorithm and the
collection and comparison proofs were read in full. This is
an exact reproducible computation, not a formal-verifier
certificate.
\end{proof}

The observed primitive odd valuation at weight $m=2r+3$ is
$3r+v_2(r!)$. Its abelian lower bound is exact:
\[
 v_2\!\left(\prod_{j=0}^{r-1}
 (5^{2r+3}-5^{2j+3})\right)=3r+v_2(r!),
\]
since $v_2(5^{2s}-1)=3+v_2(s)$. Successive division by the
monic linear factors describes the preceding-moment kernel
in the Iwasawa algebra. Passing to all weights requires integral
correction of lower decomposable terms. The scalar Newton determinant
does not alone perform that correction; the following compact-basis
reduction and parity induction supply it.

\subsection{An exact all-weight reduction at 2}

Continue with the checked presentation and normal form at $2$.
Identify $H_{\mathrm{ab}}$ with
$\Lambda=\mathbb Z_2[[T]]$, where $\gamma=1+T$ and $[h]=1$.
For $i\geq0$ put
\[
 m_i=2i+3,\qquad r_i=5^{m_i}-1,\qquad
 f_i(T)=\prod_{j<i}(T-r_j),\qquad \delta_i=f_i(r_i).
\]
Thus $v_2(\delta_i)=3i+v_2(i!)$.

\begin{proposition}[K17-2-NEWTON-BASIS]\label{prop:newton-basis}
Suppose that for every $i\geq0$ the actual primitive odd
image ideal in weight $m_i$ attains the Newton lower bound
$v_2(\delta_i)$. Then there are primitive classes in those odd
weights that freely generate the entire integral tripod graded
Lie algebra over $\mathbb Z_2$.
This proposition isolates a sufficient lifting hypothesis. The finite
calculation above does not establish it; Theorem~\ref{thm:2-allweights}
below establishes it in every weight.
\end{proposition}

\begin{proof}
The linear weight-$m$ coordinate of an element with abelian
class $F(T)$ is $F(5^m-1)$ in the normalized rational
coordinates. Choose $g_i\in H$ of actual weight at least
$m_i$ attaining the indicated minimum. Its earlier odd linear
coordinates vanish. Successive division by the monic factors
$T-r_j$ therefore gives
\[
 [g_i]=f_i(T)b_i(T).
\]
Its weight-$m_i$ coordinate is $\delta_i b_i(r_i)$, so
$b_i(r_i)$ is a unit. Since $r_i\in4\mathbb Z_2$, this is
equivalent to $b_i(0)$ being a unit, hence
$b_i\in\Lambda^\times$. Conversely, any lift of such an
abelian class having actual weight at least $m_i$ attains the
minimum. This proves the precise equivalence of the two
lifting formulations.

The classes $f_i b_i$ form a compact $\mathbb Z_2$-basis of
$\Lambda$. To check this, consider
\[
 B:\prod_{i\geq0}\mathbb Z_2\longrightarrow\Lambda,
       \qquad (a_i)\longmapsto\sum_{i\geq0}a_i f_i b_i.
\]
The series defines a continuous map. Indeed $r_j\in4\mathbb Z_2$
implies that the coefficient of $T^k$ in $f_i$ is divisible by
$2^{2(i-k)}$ for $i>k$; multiplication by an integral $b_i$
retains the required coefficientwise convergence. Equivalently,
the columns have increasing $(2,T)$-adic order. Modulo $2$,
the columns are $T^i\overline b_i(T)$, with unit leading
coefficient. Successive coefficient comparison makes this
triangular formal-series map an isomorphism. Lift modulo $2$,
subtract, divide by $2$, and repeat to prove surjectivity of
$B$. These corrections converge in the product topology.
The same argument applied to a kernel element shows that all
its coefficients are divisible by every power of $2$ and hence
zero. Thus $B$ is an isomorphism of compact modules.

The $g_i$ converge to $1$ in $H$. Their images have weights
tending to infinity, so they converge to $1$ in every fixed
rational weight truncation. The checked faithful rational
realization embeds compact $H$ into a Hausdorff inverse limit
and is a homeomorphism onto its image. The convergence therefore
holds in $H$ itself.

Take the free pro-$2$ group on the pointed compact sequence
$\{g_i:i\geq0\}\cup\{1\}$, where $g_i\to1$. It is the inverse
limit of the finite-rank free groups with successive tail
generators killed: every map of the convergent sequence into a
finite group kills its tail. This verifies the universal
property and also gives Magnus injectivity by finite-rank
projection. Its map to $H$
is an isomorphism on abelianizations by the preceding compact
basis calculation. Hence it is Frattini-surjective and thus
surjective. Projectivity of the free pro-$2$ group $H$ supplies
a section. The section has closed Frattini-surjective image in
the source, so its image is the whole source. Consequently the
$g_i$ form a compact free basis. This uses projectivity and
Frattini theory, not infinite-rank Hopficity.

Give the integral Magnus variable $X_i=g_i-1$ weight $m_i$.
A nonidentity group element has a first nonzero weighted Magnus
term. That term is primitive for the associated graded tensor
coproduct, because the group element is grouplike. The primitive
submodule of the integral tensor algebra is the free
$\mathbb Z_2$ Lie algebra: the rational primitive statement and
the unit-triangular Hall--Lyndon basis show that this Lie module
is saturated. Each arithmetic weight involves only finitely
many variables and words.

Let $\sigma_i$ be the actual weight-$m_i$ leading class of
$g_i$. Its $s_{m_i}$ indecomposable coefficient is
$\delta_i b_i(r_i)\ne0$. Every other term is a bracket in
strictly lower odd weights. Thus the homogeneous substitution
$X_i\mapsto\sigma_i$ is a triangular automorphism of the
rational free Lie algebra. It cannot kill a first nonzero
integral Magnus term. Therefore actual rational weight equals
abstract weighted Magnus degree, and every integral leading
class is a $\mathbb Z_2$ Lie polynomial in the $\sigma_i$.
Conversely every Hall bracket is realized by a group
commutator. This proves integral freeness for the rational
pullback filtration. Proposition~\ref{prop:outer-sat} identifies
it with the actual integral outer filtration.
\end{proof}

The all-weight problem is therefore reduced to the existence
of these lifts, or equivalently attainment of the Newton
minimum in every odd weight. This gate is closed by
Theorem~\ref{thm:2-allweights} below. A raw Newton word can have lower
even or decomposable terms, whose removal requires actual
integral commutators. The scalar abelian calculation does not
perform those corrections. Nor does a finite list of attained
odd minima certify all intervening lattices: later Newton
directions may still project to lower commutator weights.

Complex conjugation does not make every $\gamma$-conjugate of
$h$ exactly inverted, since $c\gamma c=h\gamma$. In particular,
the rational diagonal action $\tau$ differs from actual $D$ by
the half-power conjugation $\operatorname{Ad}(\exp(a/2))$.
Its integrality cannot be assumed. The finite certificate uses
actual $D$ and retains every denominator.

\subsection{All weights in the class-two quotient}

The next result is stronger in weight range than the finite
calculation but concerns the quotient by rational generator
length at least three. Its induced filtration need not be the
quotient of the full arithmetic filtration.

\begin{theorem}[K17-2-CLASS2]\label{thm:2-class2}
Let $H^{(2)}$ be the actual compact image of $H$ in the rational
class-two quotient of the free odd Lie algebra, with its induced
arithmetic-weight filtration. Its integral graded Lie algebra
is the free class-two $\mathbb Z_2$ Lie algebra on one primitive
generator in each odd weight $m_i=2i+3$. The primitive coefficient
in the canonical coordinate $s_{m_i}$ is, up to a unit,
\[
 \delta_i=\prod_{j<i}(5^{m_i}-5^{m_j}),\qquad
 v_2(\delta_i)=3i+v_2(i!).
\]
\end{theorem}

\begin{proof}
Fix $r\geq1$ and cutoff $W=2r+2$. Put $a_i=5^{2i+3}$ for
$0\leq i<r$, and
\[
 P_i(X)=\prod_{j<i}(X-a_j),\qquad
 P_r(X)=\sum_{t=0}^r c_tX^t.
\]
The linear image $A$ of $H$ is the evaluation image of
$\Lambda$ at these $r$ nodes, using $\gamma$ as variable.
Successive monic division gives kernel $P_r\Lambda$.
The classes of $1,\gamma,\ldots,\gamma^{r-1}$ and of
$P_0,\ldots,P_{r-1}$ are two integral bases. Both $\gamma$
and its inverse preserve $A$, since $P_r$ has unit constant
coefficient.

The image $C$ of the closed commutator subgroup is central
and is the image of $\bigwedge^2 A$. Its basis vectors are
\[
 P_i(a_k)P_j(a_l)-P_j(a_k)P_i(a_l),\qquad
       i<j,\quad i+j\leq r-2,
\]
in the target coordinates $[s_{2k+3},s_{2l+3}]$.
They are triangular by increasing $i+j$, with diagonal
$\delta_i\delta_j$: a nonzero evaluation requires
$k+l\geq i+j$, and equality forces the same pair.
Pairs with $i+j>r-2$ vanish at this cutoff. These are actual
commutators of lifts of the Newton basis, and $C$ is stable
under $\gamma^{\pm1}$ by its exterior-square description.

Set $h_t=\gamma^t h\gamma^{-t}$ and form the ordered word
\[
 R_r=\prod_{t=0}^r h_t^{c_t}.
\]
Its linear image is zero. Let $E_r$ be its central logarithm
at the cutoff. The kernel of the linear projection of the
compact image is generated by $C$ and the $\gamma$-orbit of
$E_r$: modulo the closed commutator subgroup the relevant
classes are $P_r\Lambda$, generated by the conjugates of
$R_r$, and compactness supplies their closed span. Thus
$E_r\in C$ would identify the whole kernel with $C$.
Here ``central'' refers to the length-two derived subspace;
other central length-one directions at the cutoff are not
part of this kernel.

For nonnegative integral $t$, the coefficient of
$[s_m,s_n]$ in $\log h_t$ is the polynomial
\[
 b_t(X,Y)=-\frac12\bigl(Y^t S_t(X)-X^tS_t(Y)\bigr),
 \qquad S_t(X)=\sum_{j=0}^{t-1}X^j,\quad S_0=0,
\]
evaluated at $X=5^m$, $Y=5^n$. This follows by summing the
geometric series for
$D=\operatorname{Ad}(\exp(-a/2))\tau$ in class two.
The exact class-two product formula therefore gives
\[
 Z_r(X,Y)=\sum_t c_tb_t(X,Y)
 +\frac12\sum_{x<y}c_xc_y(X^xY^y-X^yY^x)               \tag{C2.1}
\]
for the central coefficient of $R_r$. This alternating
polynomial belongs to $\tfrac12\mathbb Z[X,Y]$ and has degree
at most $r$ in each variable. Monic division in the tensor
Newton basis gives
\[
 Z_r=\sum_{i<j}z_{ij}
       \bigl(P_i(X)P_j(Y)-P_j(X)P_i(Y)\bigr),
       \qquad z_{ij}\in\tfrac12\mathbb Z.             \tag{C2.2}
\]

We prove $z_{ij}\in\mathbb Z$ for $i+j\leq r-2$ by parity.
All nodes are odd. In $\mathbb F_2[X,Y]$, put
$A=X-1$, $B=Y-1$, and $Q=A+B+AB=XY-1$.
Then $P_r(X)=A^r$, $P_r(Y)=B^r$, $P_r(XY)=Q^r$.
The linear part of twice (C2.1) is
$(Q^r+A^r)/B+(Q^r+B^r)/A$, with polynomial quotients.
The cross term is $P_r(X)P_r(Y)+P_r(XY)$ because
$c_t^2=c_t$ modulo $2$. Hence
\[
 2Z_r=\frac{Q^r+A^r}{B}+\frac{Q^r+B^r}{A}
                      +A^rB^r+Q^r
       \in(A,B)^{r-1}\quad\text{over }\mathbb F_2.    \tag{C2.3}
\]
Every quotient displayed here has minimum total degree
$r-1$. Since $P_i(X)P_j(Y)$ reduces to $A^iB^j$, the
coefficient $2z_{ij}$ is even for $i+j<r-1$. Its known
half-integrality then proves the required integrality.

All remaining tensor Newton terms vanish at the cutoff.
Equation (C2.2) therefore puts $E_r$ in the actual lattice
$C$, and its whole $\gamma$-orbit stays there. This proves
that the full linear kernel is exactly $C$.

Divide $R_r$ by an actual product of commutators with the
same central image through $2r+2$. The new word has unchanged
abelian class $P_r$ and no class-two component below weight
$2r+3$; its next linear coefficient is $\delta_r$. Thus every
odd Newton minimum is attained in $H^{(2)}$. At even weights,
the just identified lattice $C$ has the leading basis
$\delta_i\delta_j[s_{m_i},s_{m_j}]$ by triangularity. At odd
weights there is no length-two component. These observations
prove the theorem.
\end{proof}

The exact script \path{evidence/work/pro2_integral_class2_newton.py}
checks the group-product formula against its geometric-series
form and eliminates in the Newton-wedge basis. Its finite
checks helped identify (C2.3); the polynomial identity proves
the all-weight assertion. The result does not assert that
these corrected words have no components of generator length
at least three below their intended odd weight in the full
image. The higher-length lifting is supplied by the separate induction
in Theorem~\ref{thm:2-allweights}.

\subsection{Integral parity correction in all weights at 2}

The next argument closes the lifting hypothesis of
Proposition~\ref{prop:newton-basis}, using actual compact
subgroups rather than estimates for higher BCH denominators.

\begin{theorem}[K17-2-ALL]\label{thm:2-allweights}
In the checked pro-$2$ arithmetic presentation and actual
tripod realization, there are primitive integral classes
$\sigma_{2i+3}$, $i\geq0$, that freely generate the entire
integral graded Lie algebra over $\mathbb Z_2$.
Their indecomposable coefficients in the canonical rational
coordinates attain
\[
 v_2(\delta_i)=3i+v_2(i!),\qquad
 \delta_i=\prod_{j<i}(5^{2i+3}-5^{2j+3}).
\]
The statement concerns the constructed primitive generators;
identification with a separately prescribed normalization
requires its own comparison.
\end{theorem}

\begin{proof}
We retain the previously checked inputs
\[
 P=C_2*\mathbb Z_2,\qquad
 h=c\gamma c\gamma^{-1},\qquad u=\chi(\gamma)=5,
\]
and the faithful rational normal form
\[
 h=\exp(a),\quad a=\sum_{m\geq3\ {\mathrm{odd}}}s_m,
 \qquad\gamma=\exp(-a/2)u,
\]
\[
 D=\operatorname{Ad}(\exp(-a/2))\tau,
 \qquad\tau(s_m)=u^m s_m.
\]
Let $F^nH$ denote the inverse image of rational weight at
least $n$. Thus $H=F^3H$, the filtration is separated, and
$[F^aH,F^bH]\subseteq F^{a+b}H$. Write
$H_{\leq N}=H/F^{N+1}H$. Both $c$ and $D$ preserve this
filtration. Complex conjugation acts by $(-1)^n$ on rational
weight $n$, while $D$ acts there by $u^n$. The inner factor
of $D$ first raises weight by $3$.

Complex conjugation acts as minus identity on the entire
$H_{\rm ab}=\Lambda$, where $[h]=1$.
Indeed $c(h)=h^{-1}$ and $c\gamma c=h\gamma$; inner
conjugation by an element of $H$ is trivial in $H_{\rm ab}$.
It follows that each class $\gamma^x[h]$ is negated, and
these classes generate $\Lambda$.

Set $m_i=2i+3$, $a_i=u^{m_i}$ and
$P_i(\gamma)=\prod_{j<i}(\gamma-a_j)$.
We construct $g_i$ with
\[
 g_i\in F^{m_i}H,\qquad [g_i]=P_i(\gamma),            \tag{A3}
\]
so that the closed subgroup
$S_i=\langle g_0,\ldots,g_i\rangle$ surjects onto
$H_{\leq m_i}$. Simultaneously we retain exact identities
\[
 D(g_j)=g_{j+1}w_jg_j^{a_j},\qquad
 w_j\in S_j'\cap F^{m_j+1}H,\qquad j<i,             \tag{A4}
\]
where primes denote closed commutator subgroups.
For $i=0$ take $g_0=h$. Through weight $3$ its closed
procyclic subgroup is stable under $D^{\pm1}$, since
$D(h)=h^{u^3}$ there and $u^3$ is a $2$-adic unit.
It contains the whole orbit of $h$, which generates $H$,
so the initial surjectivity holds.

Assume the construction through $i$, and put $m=m_i$,
$g=g_i$, $S=S_i$. Write
\[
 \log g=A+B+\text{terms of weight at least }m+2,
       \qquad\operatorname{wt}(A)=m,\quad
              \operatorname{wt}(B)=m+1.
\]
Parity and $2m>m+1$ give
\[
 c(g)g\in F^{m+1}H\cap H',\qquad
          (\log(c(g)g))_{m+1}=2B.                   \tag{A5}
\]
Membership in $H'$ uses the exact minus action on
$H_{\rm ab}$, and every BCH commutator lies above the
weight displayed.

The following central extension supplies the required
integral lift without assuming generation in the next even
weight. In $H_{\leq m+1}$, the induction hypothesis gives
\[
 H_{\leq m+1}=\operatorname{im}(S)K,\qquad
 K=F^{m+1}H/F^{m+2}H.
\]
The subgroup $K$ is central, since its commutator with $H$
lies in $F^{m+4}H$. Therefore the closed derived subgroup
of this quotient equals that of $\operatorname{im}(S)$.
Continuous maps of compact groups carry closed commutator
subgroups onto those of their images. We can choose
$w_0\in S'$ whose image through weight $m+1$ is $c(g)g$.
In particular $w_0\in F^{m+1}H$ and
$(\log w_0)_{m+1}=2B$.

Now let $q=D(g)g^{-u^m}$. Its weight-$m$ term cancels.
The inner part of $D$ contributes first at weight $m+3$,
and the BCH cross terms are also above weight $m+1$.
Consequently
\[
 q\in F^{m+1}H,\qquad
           (\log q)_{m+1}=u^m(u-1)B.               \tag{A6}
\]
The scalar $t=u^m(u-1)/2=2\cdot5^m$ is an integer.
Define
\[
 w_i=w_0^t,\qquad g_{i+1}=q w_i^{-1}.
\]
Equations (A5)--(A6) cancel the entire next even term,
so $g_{i+1}\in F^{m+2}H$. Since $w_i\in S_i'$, its
abelian class is
\[
 [g_{i+1}]=(\gamma-u^m)[g_i]=P_{i+1}(\gamma).
\]
The order of the product gives exactly (A4) for $j=i$.
Its first weight is exactly $m+2$, since its linear
coordinate there is the nonzero product
$P_{i+1}(u^{m+2})=\delta_{i+1}$.

It remains to propagate surjectivity. Put $N=m+2$ and
let $G$ be the image of $S_{i+1}$ in $H_{\leq N}$.
The exact recursions put $D(g_j)$ in this subgroup for
$j\leq i$. For its last generator, scalar action on
its first weight gives
\[
 D(g_{i+1})=g_{i+1}^{u^N}\pmod{F^{N+1}H}.
\]
Thus $D$ restricts to an endomorphism of $G$.
On the displayed generators of $G/\Phi(G)$ it has matrix
\[
 g_j\longmapsto a_jg_j+g_{j+1}\quad(j\leq i),
 \qquad g_{i+1}\longmapsto a_{i+1}g_{i+1}.
\]
All commutator corrections disappear. The matrix is
triangular with odd diagonal, so its reduction modulo
$2$ is invertible. Its images therefore span the Frattini
quotient; independence of the displayed generators is
unnecessary. The closed subgroup $D(G)$ is consequently
all of $G$. Since $D$ is an ambient automorphism, $G$ is
also stable under $D^{-1}$. It contains $h$, hence its
entire integer orbit and, by compactness and continuity,
its $\mathbb Z_2$ orbit. Those elements generate $H$,
so $G=H_{\leq N}$. This completes the induction.

We have constructed (A3) for every $i$, with the exact
Newton abelian class. Proposition~\ref{prop:newton-basis}
therefore gives the compact free basis and integral Hall
generation in all weights. Proposition~\ref{prop:outer-sat}
identifies the filtration with the actual integral tripod
outer filtration and proves the theorem.
\end{proof}

The argument does not assume that $S_i$ is stable under
$c$, or that its integral graded lattice already equals
the unknown next even layer. The lift in (A5) comes from
the equality of derived images after adjoining a central
kernel. The exact minus action on the whole abelianization
is essential; an even Iwasawa component at an odd irregular
prime prevents this argument from applying unchanged.

% CHECKED DEDUCTION: full self-reading, archive independent audit and root review passed.
\subsection{Integral odd generation at every Vandiver prime}
\label{sec:integral-classification}

The kernel theorem gives a further integral consequence without locating
the first exceptional weight.  It distinguishes every irregular Vandiver
prime, rather than only the two explicit examples above.

\begin{theorem}[K17-INTEGRAL-CLASSIFICATION]
\label{thm:integral-classification}
Let $p$ be an odd prime satisfying Vandiver, and let
$\mathfrak g_p$ be the actual integral tripod Lie algebra with its
arithmetic weight grading.
If $p$ is regular, $\mathfrak g_p$ is the free $\mathbb Z_p$-Lie
algebra on one generator in every odd weight $m\ge3$.
If $p$ is irregular, there are infinitely many even weights $n$ for which
\[
 \bigl(\mathfrak g_p/[\mathfrak g_p,\mathfrak g_p]\bigr)_n
\]
is a nonzero finite $p$-group.  In particular, no collection of
homogeneous odd integral classes generates $\mathfrak g_p$ as a
$\mathbb Z_p$-Lie algebra.
\end{theorem}

\begin{proof}
The regular-prime assertion is Sharifi's Theorem~1.1 combined with
Brown's rational Deligne--Ihara theorem
\cite[Theorem~1.1, p.~2]{Sharifi2002}\cite[\S\S2.1--2.3]{Brown2019}.
For the irregular case, put
\[
 K=\mathbb Q(\mu_{p^\infty}),\qquad
 H=\operatorname{Gal}(\Omega_p/K),\qquad
 M=H/\overline{[H,H]}.
\]
Theorem~\ref{thm:vandiver-kernel} identifies $H$ with the actual
cyclotomic-kernel image.  Its weight filtration $F^nH$ is therefore
separated.  Proposition~\ref{prop:outer-sat} identifies it with the
actual integral outer filtration.  Write $c$ for complex conjugation.

We first justify that $M^+$ is a nonzero finite free
$\mathbb Z_p$-module.  The Iwasawa and Ferrero--Washington inputs used
in the proof of Theorem~\ref{thm:pole-general} make $H$ free pro-$p$;
thus $M$ has no $\mathbb Z_p$-torsion.  The Kummer comparison
\[
 M^+\simeq (X_\infty^-)^\iota(1)
 \quad\hbox{up to finite kernel and cokernel}
\]
and $\mu(X_\infty^-)=0$ make $M^+$ finitely generated over
$\mathbb Z_p$; hence it is finite free
\cite[Theorem~2]{Iwasawa1981}
\cite[Proposition~3.4.7 and Corollary~3.4.8]{SharifiIwasawaNotes}.

For nonvanishing, let $A_0$ be the $p$-part of the class group of
$\mathbb Q(\mu_p)$ and let $A_\infty$ be the direct limit of the
cyclotomic-layer $p$-class groups under extension of ideals.
Irregularity makes $A_0$ nonzero, and Vandiver makes $A_0^+=0$,
so $A_0^-\ne0$.  The minus-part maps into this direct limit are
injective, and Kummer duality gives an actual isomorphism
\[
 A_0^-\hookrightarrow A_\infty^-\ne0,
 \qquad M^+\cong(A_\infty^-)^\vee(1)
\]
\cite[Propositions~3.4.1 and~3.4.7]{SharifiIwasawaNotes}.
Here $\vee$ is Pontryagin duality; a nonzero discrete torsion
$p$-group has nonzero Pontryagin dual.  Thus $M^+\ne0$, and its
finite free rank is positive.

Let $\pi:H\to M$ be the quotient map and give $M$ the image
filtration
\[
 M^n=\pi(F^nH).
\]
Since $F^1H=H$, this starts with $M^1=M$.
These are closed $\mathbb Z_p$-submodules, because $F^nH$ is
compact.  This filtration is separated.  Indeed, if $x$ belonged
to every $M^n$, the nonempty compact sets
\[
 F^nH\cap\pi^{-1}(x)
\]
would be nested and would have a common point.  That point would
belong to $\bigcap_nF^nH=\{1\}$, forcing $x=0$.
This is the compactness argument needed here; separatedness is
not being assumed to pass automatically to an arbitrary quotient.

There is a surjective graded Lie map to an abelian graded module,
and hence surjections in every degree
\[
 \mathfrak g_p\longrightarrow\operatorname{gr}M,
 \qquad
 (\mathfrak g_p^{\mathrm{ab}})_n
       \twoheadrightarrow M^n/M^{n+1}.                       \tag{IC.1}
\]
Surjectivity follows from the definition of $M^n$; all brackets
map to zero since $\pi$ kills the closed group commutator subgroup.
Complex conjugation acts on $\mathfrak g_{p,n}$ by $(-1)^n$,
so it acts by the same scalar on the quotient in (IC.1).
Since the idempotents $(1\pm c)/2$ are integral, the induced
filtration $M^{n,+}=M^n\cap M^+$ has
\[
 M^{n,+}/M^{n+1,+}=0\quad\text{for odd }n.                   \tag{IC.2}
\]

Each $\mathfrak g_{p,n}$ is a finite free $\mathbb Z_p$-module:
it embeds in the finite free outer Johnson module described in
Proposition~\ref{prop:outer-sat}.  In a fixed weight the bracket
submodule is a finite sum of images of tensor products of finite
free modules, so it is finitely generated and closed.  Thus
$(\mathfrak g_p^{\mathrm{ab}})_n$ is finitely generated.
After tensoring with $\mathbb Q_p$, Brown's theorem makes the
Lie algebra free on odd generators and gives
\[
 (\mathfrak g_p^{\mathrm{ab}})_n\otimes\mathbb Q_p=0
 \quad\text{for even }n.
\]
Exactness of localization and the finite sum defining brackets
justify commuting this abelianization with tensoring.
Therefore every even $(\mathfrak g_p^{\mathrm{ab}})_n$ is a
finite $p$-group, and so is every even quotient
$M^{n,+}/M^{n+1,+}$.

There must be infinitely many nonzero such even quotients.
Otherwise (IC.2) would make the filtration on $M^+$ stationary
from some point onward.  Its intersection is zero, so the
stationary submodule would be zero.  Starting from $M^+$ and
using the finitely many preceding finite quotients would then
make $M^+$ finite, contradicting its positive free rank.
By (IC.1), each of those infinitely many even weights has
nonzero $(\mathfrak g_p^{\mathrm{ab}})_n$.

Finally, a graded Lie algebra generated by homogeneous odd
classes has abelianization generated by their images, all in odd
weights.  Its even abelianization would be zero.  The preceding
nonvanishing therefore excludes every such integral generating
collection, proving the theorem.
\end{proof}

Together with the separately proved all-weight theorem at $2$,
this classifies the existence of free homogeneous odd integral
generators among the primes covered by the Vandiver hypothesis:
they exist at $2$ and at the odd regular primes, and fail at every
irregular Vandiver prime.  The negative assertion does not locate
the first defect weight or its order.  The explicit weight-$32$
example at $37$ and weight-$12$ example at $691$ retain that
additional information.  Kernel-field equality and primitive
full-domain Soul\'e characters can hold at all these irregular
Vandiver primes despite the infinitely many even integral
indecomposable quotients.

% CHECKED DEDUCTION: complete self-read and archive independent audit passed.
% Include after greenberg_kernel_section.tex and integral_classification_section.tex.
\subsection{Integral odd generation under the one-variable Greenberg hypothesis}
\label{sec:greenberg-integral}

The integral classification extends from Vandiver to the weaker hypothesis
used in Theorem~\ref{thm:greenberg-kernel}.  The additional arithmetic
point is that irregularity already forces a nonzero minus class group;
Vandiver is unnecessary for that implication.

\begin{corollary}[K17-GREENBERG-INTEGRAL]
\label{cor:greenberg-integral}
Let $p$ be an odd prime, let
\[
 X=\varprojlim_r\operatorname{Cl}(\mathbb Q(\mu_{p^r}))\{p\}
\]
with norm transition maps, and assume that $X^+$ is finite.
Let $\mathfrak g_p$ be the actual integral tripod Lie algebra with
its arithmetic weight grading.
If $p$ is regular, then $\mathfrak g_p$ is the free
$\mathbb Z_p$-Lie algebra on one generator in each odd weight
$m\ge3$.  If $p$ is irregular, then
\[
 \bigl(\mathfrak g_p/[\mathfrak g_p,\mathfrak g_p]\bigr)_n
\]
is a nonzero finite $p$-group for infinitely many even weights $n$.
In particular, under the stated hypothesis, homogeneous odd integral
generators exist exactly at the regular odd primes, and at those primes
they can be chosen freely.
\end{corollary}

\begin{proof}
The regular-prime assertion is the one in
Theorem~\ref{thm:integral-classification}; its positive proof uses
Sharifi's Theorem~1.1 and Brown's rational Deligne--Ihara theorem
\cite[Theorem~1.1]{Sharifi2002}\cite[\S\S2.1--2.3]{Brown2019}.
For an irregular prime, put $F_\infty=\mathbb Q(\mu_{p^\infty})$,
\[
 H=\operatorname{Gal}(\Omega_p/F_\infty),\qquad
 M=H/\overline{[H,H]},\qquad
 A_0=\operatorname{Cl}(\mathbb Q(\mu_p))\{p\}.
\]
Theorem~\ref{thm:greenberg-kernel} identifies $H$ faithfully with
the actual cyclotomic-kernel image.  Its integral weight filtration
$F^nH$ is separated, and Proposition~\ref{prop:outer-sat}
identifies it with the actual outer filtration.

Here is the precise nonvanishing input.  Kummer's criterion says that
irregularity is equivalent to $p$ dividing the numerator of some
$B_k$ with $2\le k\le p-3$ even.  Ribet's Main Theorem~(1.1) gives
\[
 p\mid B_k
 \quad\Longleftrightarrow\quad
 (A_0/pA_0)^{(\chi^{1-k})}\ne0,
 \tag{GI.1}
\]
where $\chi$ is the mod-$p$ cyclotomic character
\cite[\S1, p.~151, Main Theorem~(1.1)]{Ribet1976}.
Since $1-k$ is odd, complex conjugation acts as $-1$ on this
eigenspace.  Thus $A_0^-\ne0$.  Conversely $A_0^-\ne0$ implies
that $p$ divides the full cyclotomic class number, hence that $p$
is irregular.  This is the exact irregularity/minus-class criterion
used here; it does not require $A_0^+=0$.

Let $A_\infty$ be the direct limit of the layer $p$-class groups
under extension of ideals.  The established Kummer comparisons and
minus-part injectivity give
\[
 A_0^-\hookrightarrow A_\infty^-,\qquad
 M^+\cong(A_\infty^-)^\vee(1),\qquad
 M^+\sim(X^-)^\iota(1),
 \tag{GI.2}
\]
where $\sim$ denotes a pseudo-isomorphism and $\vee$ is
Pontryagin duality
\cite[Propositions~3.4.1 and~3.4.7, Corollary~3.4.8]{SharifiIwasawaNotes}.
As in the proof of Theorem~\ref{thm:integral-classification},
Iwasawa's freeness input makes $H$ free pro-$p$, so $M$ is
$\mathbb Z_p$-torsion free; Ferrero--Washington's $\mu=0$
theorem and (GI.2) make $M^+$ finitely generated over $\mathbb Z_p$.
The injection and actual duality in (GI.2) make $M^+$ nonzero.
Consequently it is finite free of positive rank.  These steps use
no Vandiver assertion.

For completeness, the filtration argument is as follows.  Let
$M^n$ be the image of $F^nH$ in $M$.  These are closed submodules
and their intersection is zero: for a class $x$ in the intersection,
the nested nonempty compact fibres
$F^nH\cap\pi^{-1}(x)$ have a common point in
$\bigcap_nF^nH=\{1\}$, so $x=0$.
There are surjections
\[
 (\mathfrak g_p^{\mathrm{ab}})_n
 \twoheadrightarrow M^n/M^{n+1}.
 \tag{GI.3}
\]
Complex conjugation acts by $(-1)^n$ on these graded quotients.
Since $2$ is a unit in $\mathbb Z_p$, taking its plus eigenspace
is exact.  The filtration of $M^+$ therefore has zero quotients
in all odd weights.

Each integral graded Lie piece is finite free over $\mathbb Z_p$,
by its embedding in the outer Johnson module.  Its degreewise
abelianization is finitely generated, and Brown's rational theorem
makes that abelianization zero after inverting $p$ in every even
weight.  Thus each even $(\mathfrak g_p^{\mathrm{ab}})_n$, and
each even filtered quotient of $M^+$ in (GI.3), is a finite
$p$-group.  If only finitely many of the latter quotients were
nonzero, the separated filtration of $M^+$ would eventually
be zero and would express $M^+$ as an iterated extension of
finitely many finite groups.  This contradicts its positive
free rank.  Infinitely many even quotients, and hence infinitely
many even Lie indecomposable pieces in (GI.3), are nonzero.
An odd homogeneous generating collection would make all even
Lie indecomposables zero, completing the proof.
\end{proof}

The conclusion is conditional on the one-variable Greenberg hypothesis
when Vandiver is not known.  It gives neither the first exceptional
weight nor its order, and it makes no additional assertion about
primitivity of a prescribed normalization of the full-domain
Soul\'e characters.

\section{Cote 149, scans 79--85: Kummer theory and geometric tangential splittings}
\label{sec:archive-kummer}

\paragraph{Evidence and notation.}
All seven supplied TeX transcriptions were read in full.  The original
manuscript images were not supplied, so the editorial conclusions below are
mathematical readings, not verified readings of handwritten glyphs.  Write
\(D\) for the boundary of a smooth proper relative curve \(X/S\), put
\(U=X\setminus D\), and let \(g_i:S\longrightarrow D\) be a boundary section.
This avoids the manuscript's reuse of \(T\), which we reserve for a Tate module.
Write \(E_S=\pi_1(S)\), with geometric base points suppressed.

The passage studies a proposed cusp-normalizer extension
\begin{equation}\label{eq:arc-normalizer}
 1\longrightarrow T\longrightarrow N(L_i)
   \longrightarrow E_S\longrightarrow 1,
 \qquad T=\widehat{\mathbb Z}(1),
\end{equation}
and its obstruction class \(c_i\).  The geometric interpretation proposed in
scan 81 is the Kummer first Chern class of the conormal or normal line of
\(g_i(S)\).  The author explicitly leaves that sign convention unsettled.
Scan 85 then asks for an intrinsic profinite characterization of the
\emph{geometric} splittings inside the set of all splittings.  The obstruction
to the existence of a splitting and the recognition of that subset are
different questions.

In the calculations below, \(\mathbb Z_\ell(1)\) means the inverse system
\((\mu_{\ell^r})_r\), and \(H^j_{\mathrm{cont}}\) means continuous
cohomology of that inverse system.  We do not silently replace it by
ordinary cohomology of the inverse-limit sheaf on the small etale site.
For the full Tate module, the corresponding index set is all positive
integers ordered by divisibility.

\paragraph{F149-KUM: exact limiting Kummer sequences (proved).}
Let \(S\) be connected, normal and noetherian, and suppose that \(\ell\)
is invertible on \(S\).  Set
\[
 A=\Gamma(S,\mathcal O_S)^\times,\qquad
 P=\operatorname{Pic}(S),\qquad
 B=H^2_{\mathrm{et}}(S,\mathbb G_m).
\]
For an abelian group \(M\), put
\[
 M^\wedge_\ell=\varprojlim_r M/\ell^rM,
 \qquad
 T_\ell M=\varprojlim_r M[\ell^r],
\]
where the second transition map is multiplication by \(\ell\);
for \(A\), written multiplicatively, the first quotient is
\(A/A^{\ell^r}\).  There are natural short exact sequences
\begin{align}
 0&\longrightarrow A^\wedge_\ell
   \longrightarrow \varprojlim_r H^1_{\mathrm{et}}(S,\mu_{\ell^r})
   \longrightarrow T_\ell P\longrightarrow 0,\label{eq:arc-kum-h1}\\
 0&\longrightarrow P^\wedge_\ell
   \longrightarrow \varprojlim_r H^2_{\mathrm{et}}(S,\mu_{\ell^r})
   \longrightarrow T_\ell B\longrightarrow 0.\label{eq:arc-kum-h2}
\end{align}
Moreover,
\[
 H^1_{\mathrm{cont}}(S,\mathbb Z_\ell(1))
 \simeq \varprojlim_r H^1_{\mathrm{et}}(S,\mu_{\ell^r}).
\]
If \(P[\ell^\infty]\) has bounded exponent, then \(T_\ell P=0\), the
system \(H^1_{\mathrm{et}}(S,\mu_{\ell^r})\) is Mittag--Leffler, and
\begin{align*}
 H^1_{\mathrm{cont}}(S,\mathbb Z_\ell(1))&\simeq A^\wedge_\ell,\\
 H^2_{\mathrm{cont}}(S,\mathbb Z_\ell(1))
 &\simeq \varprojlim_r H^2_{\mathrm{et}}(S,\mu_{\ell^r}).
\end{align*}
In particular the extra hypothesis holds when \(P\) is finitely generated.

\paragraph{Proof.}
The etale Kummer sequence and Hilbert 90 give
\begin{align*}
 0&\longrightarrow A/A^{\ell^r}
   \longrightarrow H^1_{\mathrm{et}}(S,\mu_{\ell^r})
   \longrightarrow P[\ell^r]\longrightarrow 0,\\
 0&\longrightarrow P/\ell^rP
   \longrightarrow H^2_{\mathrm{et}}(S,\mu_{\ell^r})
   \longrightarrow B[\ell^r]\longrightarrow 0.
\end{align*}
The transition from level \(r+1\) to level \(r\) is induced by raising
roots of unity to the \(\ell\)-th power.  In the Kummer diagram the first
copy of \(\mathbb G_m\) maps by that power and the second by the identity.
Consequently the transitions on the left are quotient maps, and the
transitions on the right are multiplication by \(\ell\).  Both left
systems are surjective, so their first derived inverse limits vanish.
Taking inverse limits proves
\eqref{eq:arc-kum-h1}--\eqref{eq:arc-kum-h2}, including the terminal zeros.

If \(\ell^bP[\ell^\infty]=0\), the transition
\(P[\ell^m]\longrightarrow P[\ell^r]\) is zero whenever \(m\geq r+b\).
The image of \(H^1(S,\mu_{\ell^m})\) at level \(r\) is then contained in
the unit term \(A/A^{\ell^r}\).  It contains that entire term because
the unit quotient transitions are surjective.  These stabilized images
prove the Mittag--Leffler assertion and also give \(T_\ell P=0\).
For a countable inverse system, the continuous-cohomology comparison is
\begin{equation}\label{eq:arc-jannsen}
 0\longrightarrow \varprojlim\nolimits^1_r
 H^{j-1}_{\mathrm{et}}(S,\mu_{\ell^r})
 \longrightarrow H^j_{\mathrm{cont}}(S,\mathbb Z_\ell(1))
 \longrightarrow \varprojlim_r
 H^j_{\mathrm{et}}(S,\mu_{\ell^r})\longrightarrow 0.
\end{equation}
The groups \(H^0(S,\mu_{\ell^r})\) are finite: a section is determined
at the generic point of the connected normal scheme and satisfies a
polynomial of degree \(\ell^r\).  Their inverse system is therefore
Mittag--Leffler.  This removes the degree-one derived-limit term without
a Picard hypothesis.  The preceding calculation removes the
degree-two term under the bounded-torsion hypothesis.  This proves all
the assertions.

The finite-level inputs are \emph{SGA 4}~\cite{SGA4}, Expose IX,
Theorie de Kummer 3.2 and Theoreme 3.3; the primary French reedition was
checked at printed pp.~14--15.  Formula \eqref{eq:arc-jannsen} is Jannsen~\cite{Jannsen1988},
\emph{Continuous etale cohomology}, \emph{Math.\ Ann.} 280 (1988),
equation (3.1), p.~216; the inverse-system Kummer diagram is
equation (3.27), p.~222.  Both locators were checked in the primary text.

\paragraph{Full completion and its precise scope.}
If \(S\) is a scheme over \(\mathbb Q\), all positive integers are
invertible.  The same proof, on the cofinal sequence \(n=r!\), gives
\begin{align*}
 0&\longrightarrow A^\wedge
 \longrightarrow H^1_{\mathrm{cont}}(S,\widehat{\mathbb Z}(1))
 \longrightarrow \varprojlim_n P[n]\longrightarrow 0,\\
 0&\longrightarrow P^\wedge
 \longrightarrow \varprojlim_n H^2_{\mathrm{et}}(S,\mu_n)
 \longrightarrow \varprojlim_n B[n]\longrightarrow 0,
\end{align*}
where \(A^\wedge=\varprojlim_n A/A^n\) and
\(P^\wedge=\varprojlim_n P/nP\).  If the torsion of \(P\) has bounded
exponent, the same stabilization argument applies to the factorial
system and identifies the middle degree-two group with continuous
cohomology.  This is the pro-\(\mathbb N\) completion.  It need not be
a compact profinite group when the quotients \(A/A^n\) are infinite.

Thus the uncertain terminal arrow in scan 82, formula (21), can be
resolved mathematically as an arrow to zero with the usual transitions.
The \(K(\pi,1)\) hypothesis in scan 80 does not by itself justify dropping
every derived inverse limit.  In the finite-generation setting the
argument above supplies the missing justification.  These observations
are not a claim to have visually read the uncertain arrow.

The requirement that the relevant integers be invertible is substantive.
For example, over \(k=\mathbb F_p(t)\), the element \(t\) acquires no
\(p\)-th root in any separable algebraic extension: such a root would
generate the nontrivial purely inseparable extension defined by
\(X^p-t\).  Hence \(x\mapsto x^p\) on \(\mathbb G_m\) is not an epimorphism
on the etale site of \(k\).  Kummer theory for noninvertible exponents
uses the flat topology; its finite flat torsors cannot simply be
reinterpreted as etale fundamental-group torsors.  In particular, the
example over \(\operatorname{Spec}\mathbb Z[1/\ell]\) below is a
\emph{pro-\(\ell\)} example, not an unqualified full-\(\widehat{\mathbb Z}\)
etale calculation on that mixed-characteristic scheme.

\paragraph{F149-SPLIT: the two torsors (proved).}
Fix a split extension \(1\to T\to E\to G\to1\), with \(T\) abelian,
and choose a continuous splitting \(s_0\).  Every other splitting is
uniquely \(s(g)=a(g)s_0(g)\), where
\[
 a(gh)=a(g)\,{}^ga(h).
\]
This is exactly the continuous cocycle equation.  Conjugating a
splitting by an element of \(T\) changes \(a\) by a coboundary.
Consequently the \emph{conjugacy classes} of splittings form a torsor
under \(H^1_{\mathrm{cont}}(G,T)\).  Raw splittings are described by
\(Z^1_{\mathrm{cont}}\), not by \(H^1_{\mathrm{cont}}\).
For the prime-invertible systems used here, the degree-one group and
etale descriptions agree: finite \(\mu_n\)-torsors are finite etale,
and their degree-zero inverse system is Mittag--Leffler as above.
No degree-two \(K(\pi,1)\) comparison is needed for this assertion.

If a line bundle \(\mathcal L\) on \(S\) is trivial, its global bases form
a torsor under \(A=\Gamma(S,\mathcal O_S)^\times\).  Indeed, two bases
\(b,b'\) have a unique ratio \(u\in A\), and every \(u\in A\) occurs.
In a local tame parameter model, replacing \(t\) by \(ut\) changes the
root tower by roots of \(u\).  Its descent cocycle is precisely the
Kummer cocycle of \(u\).  Thus the map from bases to local splitting
classes is equivariant for
\[
 A\longrightarrow H^1_{\mathrm{cont}}(S,T).
\]
The same assertion for the proposed global normalizer follows when the
relative comparison specified below has been constructed.

Higher-order parameter changes introduce no additional finite-level
class: if \(n\) is invertible in a ring \(R\), the map
\(v\mapsto v^n\) is a bijection of \(1+tR[[t]]\).
To prove this, solve successively for the coefficients of
\(v=1+\sum_{j\geq1}v_jt^j\).  At stage \(j\) the new coefficient appears
as \(nv_j\), while every other term is determined by earlier
coefficients.  The invertibility of \(n\) gives existence and uniqueness
at every stage.  A parameter change \(t\mapsto ut+O(t^2)\) therefore has
the same Kummer ambiguity as its leading unit.

\paragraph{F149-FULLINJ: injectivity of the full unit completion (proved).}
If \(S\) is an integral scheme whose function field \(F\) is finitely
generated over \(\mathbb Q\), then
\[
 \Gamma(S,\mathcal O_S)^\times
 \longrightarrow
 \varprojlim_n\Gamma(S,\mathcal O_S)^\times/
       \bigl(\Gamma(S,\mathcal O_S)^\times\bigr)^n
\]
is injective.

\paragraph{Proof.}
It suffices to prove \(\bigcap_n(F^\times)^n=\{1\}\).  Let \(k\) be the
algebraic closure of \(\mathbb Q\) in \(F\); it is a number field.
Choose a normal projective \(k\)-model of \(F\).  If \(u=v_n^n\) for
every \(n\), the order of \(u\) at every prime divisor is divisible by
every positive integer, hence is zero.  A normal noetherian domain is
the intersection, in its fraction field, of its local rings at
height-one primes.  Applying this to affine opens, both \(u\) and
\(u^{-1}\) are global regular functions on the model.  Properness then
puts \(u\) in \(k^\times\); the chosen constant field is algebraically
closed in \(F\).  Each \(v_n\), being algebraic over \(k\), also lies in
\(k^\times\).  Its finite-prime valuations show that \(u\) and every
\(v_n\) are units of the ring of integers of \(k\).  Dirichlet's unit
theorem makes this unit group finitely generated.  The intersection
of all \(n\)-multiple subgroups of a finitely generated abelian group
is zero, by its decomposition into a free group and a finite group.
Thus \(u=1\), as required.

This proves the full-completion assertion used in scan 85 under the
stated arithmetic function-field hypothesis.  It does not prove the
analogous assertion for every single prime.  For every odd \(\ell\),
\(-1=(-1)^{\ell^r}\) for all \(r\), so \(-1\) maps to zero in the
\(\ell\)-adic completion.  The full and one-prime statements must not
be interchanged.

\paragraph{F149-UNIT: a nongeometric local splitting class for every prime
(proved).}
Let \(S_\ell=\operatorname{Spec}\mathbb Z[1/\ell]\).  Then
\[
 H^1_{\mathrm{cont}}(S_\ell,\mathbb Z_\ell(1))
 \simeq
 \begin{cases}
 \mathbb Z_\ell,&\ell\ne2,\\
 \mathbb Z/2\mathbb Z\oplus\mathbb Z_2,&\ell=2.
 \end{cases}
\]
In the local model of a standard tripod cusp with a fixed parameter,
there is a local splitting class whose valuation coordinate is
\[
 \alpha_\ell=\frac{1}{1-\ell^2}\in\mathbb Z_\ell\setminus\mathbb Z.
 \tag{F149-UNIT}
\]
It is not induced by any global unit or global tangent basis, despite
\(\operatorname{Pic}(S_\ell)=0\).

\paragraph{Proof.}
The localization \(\mathbb Z[1/\ell]\) is a principal ideal domain and
has unit group \(\{\pm1\}\times\ell^{\mathbb Z}\).  Thus its Picard group
is zero, and F149-KUM identifies the displayed cohomology with the
completion of this unit group.  For odd \(\ell\), raising to
\(\ell^r\) is an automorphism of \(\{\pm1\}\); this factor disappears
in every quotient.  For \(\ell=2\) it remains a copy of
\(\mathbb Z/2\mathbb Z\), with identity transitions.  In either case
the valuation quotients are \(\mathbb Z/\ell^r\mathbb Z\), with their
usual transitions.  This proves the group formula.

The denominator \(1-\ell^2\) is a unit in \(\mathbb Z_\ell\), while
its absolute value as a reduced integer denominator is
\(\ell^2-1>1\).  Hence \(\alpha_\ell\) is \(\ell\)-adically integral
but is not an integer.  Choose integers \(a_r\) with
\(a_r\equiv\alpha_\ell\pmod{\ell^r}\).  The Kummer classes of
\(\ell^{a_r}\) are compatible, since
\(a_{r+1}-a_r\) is divisible by \(\ell^r\), and they do not depend
on the choices.  If their inverse-limit class came from a global
unit \(\pm\ell^a\), its valuation coordinate would give
\(a=\alpha_\ell\) in \(\mathbb Z_\ell\), a contradiction.

The chosen standard parameter supplies the split local tame extension.
Its splitting classes are the torsor in F149-SPLIT; changing a global
tangent basis is exactly the Kummer action of a global unit.  The
constructed class lies outside that image.  The conclusion concerns
the local splitting torsor.  It neither assumes a general relative
normalizer theorem nor proves that a richer anabelian structure cannot
recognize the geometric subset.

\paragraph{F149-TORSOR: the completion mark in scan 85.}
The group acting on \emph{bases} is the uncompleted group
\(\mathbb G_m(S)\).  Scan 82, line 4, states this correctly.  When the
Kummer map is injective, the geometric image is a sub-torsor under
that uncompleted group, whereas all splitting classes are acted on by
\(H^1(S,T)\).  F149-UNIT exhibits explicitly why replacing the former
acting group by its completion changes the statement.  Therefore the
flagged uncertain \(\wedge\) in scan 85, line 8, should mathematically
be absent if the sentence refers to the geometric bases and their
image.  Here and below an uncertainty command in the supplied source
is a source annotation, not an assertion that the original glyph was
inspected.

\paragraph{F149-CHERN-LOCAL: the finite-level line-bundle obstruction
(proved locally).}
Let \(\mathcal L=J_i/J_i^2\) be the conormal line, and let \(n\) be
invertible on \(S\).  The local root-coordinate gluing obstruction is
the Kummer boundary
\[
 \delta_n(\mathcal L)\in H^2_{\mathrm{et}}(S,\mu_n),
\]
up to inversion if the normal rather than conormal convention is used.

\paragraph{Proof.}
Choose local frames and formal parameters \(t_a\), with leading
transition units \(u_{ab}\).  The \(u_{ab}\) represent the line bundle.
On an etale refinement of the overlaps, choose
\(v_{ab}\) with \(v_{ab}^n=u_{ab}\); such refinements exist because
\(n\) is invertible.  The failure of the root-coordinate identifications
to compose on triple overlaps is
\[
 v_{ab}v_{bc}v_{ca}\in\mu_n.
\]
This is the defining boundary cocycle for the Kummer connecting map.
If necessary one uses the corresponding refinement of the Cech
hypercover; no assertion that one fixed Zariski cover computes all of
\(H^2\) is required.  The higher-order formal unit factors have unique
\(n\)-th roots by the coefficient recursion in F149-SPLIT, and do not
change this obstruction.  Dualizing \(\mathcal L\) inverts the
transition units, hence negates its cohomology class.

The local tame-cover input is \emph{SGA 1}~\cite{SGA1}, Expose XIII,
Proposition 5.2 and Corollaire 5.3, reedition pp.~315--316
(original marginal pp.~428--430).  The relative statements are
Proposition 5.5 and Corollaire 5.6.  The discrete-valuation
Abhyankar lemma is Expose X, Lemme 3.6.  The local statements and
their Kummer description were checked in the primary reedition.

\paragraph{The relative comparison and the scope of recognition.}
The calculation just made does not identify the complete formal-tube
extension with \(N(L_i)\) in every generality suggested by scan 79.
That identification must specify the homotopy-exactness hypotheses,
the injection and normalizer of geometric inertia, and the descent
from the local tube.  A group-cohomology comparison with
\(H^2_{\mathrm{cont}}(S,T)\) must also be supplied where it is used.
The excerpt itself marks part of that setup as requiring verification.
Theorem~\ref{thm:archive-relative} below now supplies this comparison for a smooth characteristic-zero base under the exact geometric homotopy sequence assumed in scan~79. It retains that hypothesis and does not assert the mixed-characteristic or arbitrary-fibration variants.

Once that comparison identifies the obstruction of
\eqref{eq:arc-normalizer} with the full Kummer Chern class, and once
\(\operatorname{Pic}(S)\) is finitely generated, the line bundle is
trivial if and only if that obstruction vanishes.  Indeed,
\(\bigcap_n n\operatorname{Pic}(S)=0\) by the structure theorem for
finitely generated abelian groups, and F149-KUM injects its completion
into the relevant degree-two group.  For a single prime, prime-to-\(\ell\)
Picard torsion is lost, so this full-completion conclusion cannot
simply be repeated.

A basis of the conormal line is the precise first-order datum needed
for a tangent vector.  A globally chosen formal parameter or a single
global regular function defining the divisor is a stronger datum in
general; the word \emph{uniformizer} in the excerpt should not be used
to assert such a global lifting without its hypotheses.

Over a field of characteristic zero the local splitting torsor is
already established in Stix~\cite{StixCuspidal2012}, \emph{On cuspidal sections of algebraic
fundamental groups}, author version dated 24 February 2012,
Proposition 11, p.~7.  Proposition 16, p.~8, identifies inertia
normalizers for a hyperbolic curve.  Theorem 17, p.~8, and Corollary 19,
p.~9, address injection and abundance of cusp packets under their
stated section-injectivity hypotheses.  These checked results concern
the local packet and all its section classes; they do not select the
uncompleted-unit image inside it.  Corollary 32, p.~12, recognizes
cuspidality under the non-Tate-character hypothesis, which is again
different from the requested geometric tangent subset.

\paragraph{Mathematical corrections versus unresolved readings.}
There are four evidence-based editorial entries:
\begin{enumerate}
 \item Scan 85, line 8: the acting group for global tangent bases is
 \(\mathbb G_m(S)\), as proved above; the uncertain completion mark
 should be absent in that interpretation.
 \item Scan 83, line 5: the generic points of irreducible divisors of
 \(S\) are codimension-one points of \(S\), not of the relative curve
 \(X\).  The subsequent normalization of \(S\) in its function field
 confirms this referent.  This is a mathematical correction; without
 the image it may be an author slip rather than a transcription error.
 \item Scan 82, formula (21): its terminal map is surjective and may
 end in zero by F149-KUM.  The handwritten choice has not been checked.
 \item Scan 80, line 2: a literal inference from surjectivity of a
 normalizer to the existence of a splitting is invalid.  The exact
 sequence
 \[
 0\longrightarrow\mathbb Z_\ell
 \xrightarrow{\;\ell\;}\mathbb Z_\ell
 \longrightarrow\mathbb Z/\ell\mathbb Z\longrightarrow0
 \]
 is nonsplit: a section would send a nonzero element of order \(\ell\)
 to an element of order \(\ell\) in the torsion-free group
 \(\mathbb Z_\ell\).  The following pages themselves study a splitting
 obstruction.  This proves that the unconditional implication cannot
 be right; it does not determine the replacement French wording.
\end{enumerate}

No correction of the normal/conormal sign is justified, because the
author expressly flags that convention.  Nor is the marginal extension
direction in scan 79 a mathematically compelled transcription error.
Once two genuine exact extensions have the same kernel and quotient,
a morphism inducing the identity on both is automatically an
isomorphism: its kernel is contained in the common kernel, where it
is injective, and a lift of the quotient followed by correction in the
kernel proves surjectivity.  The marginal passage discusses a setup
before all those exactness hypotheses have been established.

The Hartogs step in scan 83 needs normality of \(S\), rather than
merely normality of its open subset \(S'\).  Since the surrounding
setup already makes \(S\) smooth, the occurrence of \(S'\) does not
refute the argument and is not promoted to a required emendation.
The standard extension criterion is \emph{EGA}~\cite{EGAIV2}
\(\mathrm{IV}_2\), Theoreme 5.10.5, pp.~115--116, combined with
the normality criterion 5.8.6, p.~108, and Proposition 5.9.8 and
Definition 5.9.9, pp.~111--112.  These precise primary locators have
now been checked.  Normality gives Serre's \(S_2\) condition, so
the depth of \(\mathcal O_S\) is at least two at every point of a
closed complement of codimension at least two.  Theorem 5.10.5
then says that \(\mathcal O_S\) is closed with respect to that
complement; Proposition 5.9.8 identifies this with
\(\mathcal O_S\simeq j_*\mathcal O_{S'}\).
The equality of sections used in scan 83 follows.

\paragraph{Relation to K17.}
The archive uses Tate twist one.  Its degree-one ambiguity is a
completion of units, and its degree-two obstruction is a first Chern
class of a line bundle.  In algebraic \(K\)-theory, these are connected
with units in \(K_1\) and with the determinant or first
gamma-filtration part of \(K_0\); \(\operatorname{Pic}(S)\) should
not be inserted into \(K_0(S)\) as an additive subgroup without
qualification.

The negative-twist module
\((\mathbb Q_\ell/\mathbb Z_\ell)(-e)\), its strict local conditions,
and a duality with \(H^2(S,\mathbb Z_\ell(e+1))\) belong to a
different calculation.  A comparison must specify that change of
twist, the cohomological degree, the local conditions and the duality
map.  The Kummer sequences alone supply none of those identifications.
In particular, a bound or even a vanishing assertion for higher
\(K_{4k}(\mathbb Z)\) does not by itself select global units inside
their completion.  F149-UNIT leaves that distinction explicit even
when the Picard group is zero.

The supplied SA06 E.4 intake consists of status summaries, several
expressly marked as not read in full by its author.  It has not been
used as a proof of an all-twist cohomology table.  The relevant
war-game entry for scan 85 is S24: intrinsic recognition of the
geometric tangential subset.  The local calculations here settle the
cohomological and editorial distinctions, and Theorem~\ref{thm:archive-relative} supplies the stated relative comparison in the smooth characteristic-zero setting. The subsequent F149-GENERIC and F149-MARKED-RECOGNITION prove recognition for the full marked generic diagram. F149-UNMARKED-CURVE shows why the specified smaller diagram does not suffice invariantly when its base group may move. These recognition statements are distinct from the normalizer comparison.

\subsection{The relative normalizer comparison under the manuscript's exactness hypothesis}

The relative tangential construction in Deligne's treatment supplies
the comparison left open in the local calculation. We state its
hypotheses explicitly, rather than asserting it for an arbitrary
fibration or in mixed characteristic.

\begin{theorem}[F149-RELATIVE]\label{thm:archive-relative}
Let $S$ be a connected smooth scheme of finite type over a field of
characteristic zero. Let $X\to S$ be a smooth proper relative curve,
let $D\subset X$ be finite etale over $S$, and suppose that the
geometric fibres of $U=X\setminus D$ are hyperbolic. Fix a section
$g:S\to D$ and assume the exact homotopy sequence required in scan~79:
\[
 1\longrightarrow\pi_1(U_{\bar s})
 \longrightarrow\pi_1(U)\longrightarrow\pi_1(S)\longrightarrow1.
                                                               \tag{R149.1}
\]
Let $\mathcal L=J/J^2$ be the conormal line of $g$, let
$N=\mathcal L^\vee$ be its normal line, and let $N^\times$
be the complement of the zero section in its total space.
Choose the compatible tangential base points, and write
$I\simeq\widehat{\mathbb Z}(1)$ for geometric cusp inertia.
Then tangential specialization identifies extensions
\[
 \pi_1(N^\times)\ \simeq\ N_{\pi_1(U)}(I).
\]
\[
 1\longrightarrow I\longrightarrow\pi_1(N^\times)
 \longrightarrow\pi_1(S)\longrightarrow1.                       \tag{R149.2}
\]
It also identifies the tame formal-tube extension with these groups.
At each invertible level $n$, the image in etale cohomology of the
extension pushed out to $\mu_n$ is the Kummer Chern class of the
conormal line, up to the normal/conormal and inertia-orientation sign
already specified in F149-CHERN-LOCAL.

If $S$ is an etale $K(\pi,1)$ and the torsion subgroup of
$\operatorname{Pic}(S)$ has bounded exponent, this is the corresponding equality in
$H^2_{\rm cont}(S,\widehat{\mathbb Z}(1))$ under the usual comparison.
If $\operatorname{Pic}(S)$ is finitely generated, the extension in
(R149.2) splits if and only if $\mathcal L$ is trivial. A global
conormal basis supplies the resulting geometric splitting, and changes
of basis act by the uncompleted unit Kummer map of F149-SPLIT.
\end{theorem}

\begin{proof}
Deligne constructs the relative tangential specialization functor
from tame covers of $U$ to covers of the punctured relative tangent
line in \S15.19, printed p.~239 of \cite{DeligneTangential1989}.
His \S15.26, p.~242, explicitly applies the formal-coordinate
description to that relative situation. The construction is
functorial over $S$ and restricts, at a geometric point of $S$, to
the usual inertia map of the punctured tangent line into the curve.
Here all finite etale covers of $U$ are tame along the boundary:
$X$ is regular, the boundary is a relative normal-crossings divisor,
and every residue characteristic is zero. The local description and
descent are the generalized Abhyankar lemma in \emph{SGA~1}, XIII,
5.2--5.6 \cite{SGA1}. Thus the specialization functor is an exact
functor of the relevant Galois categories. It induces
\[
 \phi:\pi_1(N^\times)\longrightarrow\pi_1(U)
\]
over $G=\pi_1(S)$, with restriction to geometric fibre inertia
equal to the cusp inclusion $I\hookrightarrow\pi_1(U_{\bar s})$.
The formal-tube category here is the category of covers tame along
the removed section, as in the cited relative construction; no
category of wild covers is inserted.

The line bundle compactifies as a $\mathbb P^1$-bundle with the
disjoint zero and infinity sections. Hence the right-exact homotopy
sequence of \emph{SGA~1}, XIII, Proposition~4.1 and the second
example in~4.4 (reedition pp.~306--307 and~309), gives
\[
 I\longrightarrow\pi_1(N^\times)\longrightarrow G
 \longrightarrow1.                                           \tag{R149.3}
\]
The set of allowed primes is all primes in characteristic zero,
so the source's relative prime-set quotient is the full group here.
The composite of the first arrow with $\phi$ is injective: cusp
inertia injects in the geometric hyperbolic fundamental group, and
that group injects by (R149.1). Thus the first arrow of (R149.3)
is injective as well. This proves the exact sequence in (R149.2).
In particular no additional unproved assertion about a boundary map
from $\pi_2(S)$ is used; the exact sequence (R149.1) is precisely
the retained hypothesis.

Geometric cusp inertia is self-normalizing in
$\pi_1(U_{\bar s})$, by the geometric form of Stix's
Proposition~16, p.~8 \cite{StixCuspidal2012}.
The global boundary section preserves its geometric conjugacy class.
For any element of $G$, lift it to $e\in\pi_1(U)$ using
(R149.1). There is then $a\in\pi_1(U_{\bar s})$ such that
$eIe^{-1}=aIa^{-1}$, and $a^{-1}e$ belongs to the normalizer.
Consequently that normalizer surjects onto $G$, and its intersection
with the geometric kernel is exactly $I$. It therefore has an exact
sequence with the same kernel and quotient as (R149.3).

The subgroup $I$ is normal in $\pi_1(N^\times)$, so $\phi$
takes values in this normalizer. The resulting map of extensions
induces the identity on $I$ and on $G$ and is an isomorphism.
Indeed its kernel lies in $I$, where it is trivial. To lift an
element of the target, first lift its class in $G$ and then correct
by an element of the common kernel $I$. This proves surjectivity;
the continuous bijection is a homeomorphism because the groups
are compact and Hausdorff. Deligne's relative formal-coordinate
equivalence identifies the tame tube group with the same source,
proving the asserted tube comparison.

For the obstruction, choose local conormal frames, represented by
parameters $t_a$, with leading transition units $u_{ab}$. At level
$n$ the fibre cover is obtained by adjoining an $n$th root of the
parameter. The choices of roots of $u_{ab}$ identify these covers
on overlaps; their failure to satisfy descent on triples is the
same $\mu_n$-valued boundary cocycle as in F149-CHERN-LOCAL.
More explicitly, the gerbe of $\mu_n$-torsors on $N^\times$
that restrict on every geometric fibre to the standard root cover
is the gerbe of lifts for the pushout extension. Its local objects
are exactly these root covers, and its transition identifications
have the displayed triple obstruction. This identifies its etale
class with the image of the group-extension class. Inverting the
inertia identification or dualizing the
line reverses its sign. Higher formal-coordinate terms have unique
$n$th roots and do not affect the calculation. This argument uses
etale refinements or hypercovers, as in the earlier proof, and does
not assume that a particular Zariski cover computes $H^2$.

If $S$ is an etale $K(\pi,1)$, finite locally constant cohomology
agrees with that of $G$. Apply this at every $\mu_n$ and pass to
the factorial cofinal system. The continuous-cohomology exact
sequence and the bounded Picard-torsion hypothesis verified in
F149-KUM give the claimed continuous degree-two identification.
This explains exactly the extra comparison used in scan~80.

Suppose now that $\operatorname{Pic}(S)$ is finitely generated.
If the extension splits, every finite pushout class is zero, so
every Kummer boundary $\delta_n(\mathcal L)$ is zero. The finite
Kummer sequence then gives $\mathcal L\in n\operatorname{Pic}(S)$
for every positive integer $n$. Their intersection is zero in a
finitely generated abelian group, hence $\mathcal L$ is trivial.
Conversely a conormal basis gives a section of $N^\times\to S$
(the dual normal vector taking value one on that basis). The
induced fundamental-group map splits (R149.3), hence also the
normalizer extension. The root-coordinate calculation of
F149-SPLIT identifies a change of basis with the Kummer action of
its actual global unit. This proves the last assertions.
\end{proof}

This settles the relative comparison in the smooth
characteristic-zero setting of the excerpt, with its explicit
homotopy-exactness hypothesis retained. It does not assert that an
arbitrary fibration has an injective geometric fundamental group,
or replace tame covers by wild covers in positive characteristic.
The intrinsic recognition of the uncompleted geometric tangent
subset in scan~85 remains a different question.

\subsection{Generic recognition and the necessary profinite input}
\label{sec:archive-recognition}

The last sentence of scan~85 asks for recognition of the distinguished
geometric subtorsor. There are two different input diagrams here.
The isolated normalizer extension does not determine that subtorsor.
The full relative curve group together with the generic Galois marking
used in scans~83--84 does determine it, by the following argument.
The conclusion is functorial determination and an exact membership
criterion; it is not a finite-quotient decision algorithm.

\begin{theorem}[F149-GENERIC]\label{thm:archive-generic}
Retain the hypotheses and the normalizer identification of
Theorem~\ref{thm:archive-relative}. Suppose that $S$ is integral and
normal, with function field $K$ finitely generated over $\mathbb Q$.
Let $V$ be the normal line of the marked boundary section, let
$I\simeq\widehat{\mathbb Z}(1)$ be its geometric inertia, and put
\[
 G=\pi_1(S),\qquad G_K=\operatorname{Gal}(\overline K/K),\qquad
 N=\pi_1(V^\times).
\]
For an $I$-conjugacy class $s$ of sections of $N\to G$, the following
conditions are equivalent:
\begin{enumerate}[label=\textup{(\roman*)}]
\item $s$ is induced by a nowhere-zero global section of $V$;
\item the pullback $s_K$ to $G_K$ is induced by a point of
      $V_K^\times(K)$.
\end{enumerate}
The realizing global section, when it exists, is unique.
No finite-generation hypothesis on $\operatorname{Pic}(S)$ is needed
for this equivalence.
\end{theorem}

\begin{proof}
The map $G_K\to G$ is surjective: a connected finite \'etale cover
of a normal integral scheme has connected generic fibre. The
fundamental exact sequence over the field $K$, the assumed exact
sequence over $S$, and their common geometric fibre identify
\[
 N_K:=\pi_1(V_K^\times)\simeq N\mathbin{\times}_G G_K.
                                                               \tag{RG.1}
\]
Indeed the natural map of these two extensions induces the identity
on the kernel $I$ and on $G_K$, and is therefore an isomorphism.
The same argument gives
$\pi_1(U_K)\simeq\pi_1(U)\times_GG_K$ for the full curve groups.

A global nowhere-zero section clearly gives its generic section,
so only the converse needs proof. Suppose $s_K$ is induced by a
nonzero rational vector $v\in V_K(K)$. On an affine open
$W\subset S$ trivializing $V$, choose a basis $b$ and write
$v=xb$ with $x\in K^\times$. The basis gives a reference section
of the normalizer extension over $W$. Its difference from the
restriction of $s$ is a continuous $I$-valued cocycle. Reducing
modulo $n$ gives a class in $H^1(W,\mu_n)$, whose generic restriction
is the Kummer class of $x$, up to the fixed normal/conormal sign.
That sign has no effect on the argument below.

For each codimension-one point $w\in W$, the local ring
$A=\mathcal O_{S,w}$ is a discrete valuation ring. Restricting the
class to $\operatorname{Spec}A$ and using $\operatorname{Pic}(A)=0$
in the finite Kummer sequence shows that the class of $x$ in
$K^\times/K^{\times n}$ is represented by an element of $A^\times$.
Thus $x=a_n y_n^n$ for some $a_n\in A^\times$, $y_n\in K^\times$,
and
\[
 \operatorname{ord}_w(x)=n\operatorname{ord}_w(y_n)
 \quad\hbox{for every positive integer }n.
\]
An integer divisible by every positive integer is zero. Hence
$\operatorname{ord}_w(x)=0$ at every codimension-one point of $W$.
Normality now implies both $x$ and $x^{-1}$ belong to
$\mathcal O(W)$: a noetherian normal domain is the intersection
of its height-one valuation rings inside its fraction field.
Equivalently, apply the normal Hartogs extension property in
\cite[Theorem~5.10.5 and proof, pp.~115--116]{EGAIV2} to the
locally free rank-one sheaf and its dual. The vector $v$ therefore
extends as a nowhere-zero section over every such $W$, and these
extensions agree on overlaps because they agree in $K$.
They give $v\in V^\times(S)$.

The two resulting section classes agree generically. Their
difference is in $H^1(G,I)$, and inflation
\[
 H^1(G,I)\longrightarrow H^1(G_K,I)
\]
is injective. This follows directly from cocycles: the kernel of
$G_K\to G$ acts trivially on $I$, and a cocycle inflated from $G$
which becomes a coboundary on $G_K$ is already that same
coboundary on $G$. Thus the section induced by $v$ is exactly $s$.

Finally, if two global nonzero vectors give the same section,
their ratio $a\in\mathcal O(S)^\times\subset K^\times$ has zero
full Kummer class in $H^1(G_K,I)$. The full-completion injectivity
F149-FULLINJ applies to this finitely generated characteristic-zero
field, so $a=1$. This proves uniqueness. No Picard-group assertion
was used in the converse.
\end{proof}

\begin{proposition}[F149-MARKED-RECOGNITION]
\label{prop:archive-marked-recognition}
Suppose that, in addition to the full curve extension and marked
geometric cusp inertia, the profinite input retains the diagram
\[
 \pi_1(U)\longrightarrow G\longleftarrow G_K
     \longrightarrow G_k,                                  \tag{RG.2}
\]
where $k$ is an identified number field, $K/k$ is a finitely
generated regular extension, and the last arrow is the actual
restriction map. This is the generic Galois marking used in
cote~149, scans~83--84, together with its arithmetic base.
Then the subset of geometric tangent section classes in scan~85
is functorially determined by this marked profinite diagram.
The assertion uses the full profinite groups in \textup{(RG.2)},
not just an isolated prime quotient or the extension $1\to I\to N\to G\to1$.
\end{proposition}

\begin{proof}
We specify both the reconstruction and why it is independent of
choices. The known anabelian inputs are Mochizuki's Theorems~A
and~B: a hyperbolic curve over a sub-$p$-adic field is fully faithful
for open fundamental-group homomorphisms over the base Galois
group, and the analogous full profinite statement holds for
function fields of arbitrary dimension over such a base.
Their precise original locators are
\cite[Theorem~A, p.~3; Theorem~B, p.~4;
Theorem~16.5 and the following profinite argument, p.~86;
Theorem~17.1 and proof, p.~87]{MochizukiLocal1999}.
Only the isomorphism consequences are needed here.
Number fields and their finitely generated extensions are sub-$p$-adic:
a finite presentation embeds into a finitely generated extension
of a finite extension of $\mathbb Q_p$, after choosing an
embedding of the number field into an algebraic closure of
$\mathbb Q_p$.

First reconstruct $K$ from $G_K\to G_k$ within the category of
regular finitely generated extensions of the identified $k$.
Theorem~B says exactly that two realizations and any isomorphism
of their marked Galois groups differ by a unique field isomorphism.
For transcendence degree zero this is just $K=k$.
Form, entirely from \textup{(RG.2)}, the fibre product
\[
 \Pi_K=\pi_1(U)\mathbin{\times}_G G_K.
\]
By (RG.1) and its curve analogue, this is the full arithmetic
fundamental group of $U_K$ over $G_K$. Theorem~A reconstructs
$U_K$ up to its unique corresponding isomorphisms. The marked
inertia class determines its marked cusp in the smooth proper
compactification. Its normal line $V_K$, and the tangential
specialization map from its nonzero rational vectors to section
classes of $N_K$, are then reconstructed along with the curve.
Tangential specialization is functorial by the same relative
construction of Deligne used in Theorem~\ref{thm:archive-relative}.

Here is the explicit resulting membership criterion. Choose any
nonzero vector $b\in V_K(K)$ in the reconstructed one-dimensional
space, and let $s_b$ be its tangential section. For a proposed
global splitting $s$, form
\[
 \delta_b(s_K)=[s_K]-[s_b]\in H^1(G_K,I).
\]
Then
\[
 s\text{ is geometric over }S
 \quad\Longleftrightarrow\quad
 \delta_b(s_K)\in
 \kappa_K(K^\times),                                      \tag{RG.3}
\]
where $\kappa_K$ is the full Kummer map of the reconstructed field,
using the inertia orientation of the reconstructed marked curve.
For the forward implication apply Theorem~\ref{thm:archive-generic}.
For the reverse implication, a Kummer representative $a\in K^\times$
gives the rational vector $ab$ (or its inverse convention), and
Theorem~\ref{thm:archive-generic} extends it over $S$.
Replacing $b$ by $ub$, $u\in K^\times$, translates
$\delta_b$ by $\kappa_K(u)$, so does not change the condition.

It remains to check that the criterion has not concealed a choice
of geometric realization. Given two realizations of the marked
field group, Theorem~B supplies the corresponding field isomorphism.
After transport by it, two realizations of $\Pi_K\to G_K$ are
related by the curve isomorphism supplied by Theorem~A. This
isomorphism extends to the smooth proper compactifications, sends
the marked cusp to the marked cusp, and its derivative bijects
their nonzero rational normal vectors. Naturality of tangential
specialization identifies their sets of geometric section classes.
The group isomorphisms are defined up to geometric conjugacy.
Once the marked inertia representatives are aligned, the remaining
geometric conjugacy lies in $I$, because cusp inertia is
self-normalizing in the geometric curve group
\cite[Proposition~16, p.~8]{StixCuspidal2012}.
Thus the induced set of $I$-conjugacy classes is unambiguous.
Consequently every isomorphism of the full marked diagram
preserves (RG.3). This proves functorial determination.

This argument assumes that the marked diagram is known to arise
from the geometric situation stated in the proposition. It neither
recognizes arbitrary profinite groups as such diagrams nor supplies
an effective decision from finitely many finite quotients. It also
does not recover the additional generic marking from $G$ alone.
\end{proof}

\begin{proposition}[F149-BARE-EXTENSION]
\label{prop:archive-bare-extension}
The isolated extension $1\to I\to N\to G\to1$, even with its
specified inertia module and cyclotomic action, does not in general
intrinsically determine the geometric splitting subset. More
precisely, that subset need not be preserved by automorphisms of
the extension acting identically on $I$ and $G$.
\end{proposition}

\begin{proof}
Take $S=\operatorname{Spec}\mathbb Q$ and the trivial normal line,
for example at the cusp $0$ of the tripod. A reference tangent
vector identifies the normalizer with
\[
 N=I\rtimes G_{\mathbb Q},\qquad I=\widehat{\mathbb Z}(1).
\]
For every continuous cocycle $z:G_{\mathbb Q}\to I$, the map
\[
 (t,g)\longmapsto(t+z(g),g)                                 \tag{RG.4}
\]
is a continuous group automorphism, with inverse obtained by
replacing $z$ by $-z$. The cocycle identity verifies multiplication,
and the map fixes both the kernel and the quotient. On section
classes it acts by translation by $[z]\in H^1(G_{\mathbb Q},I)$.
Such translations are transitive on all section classes.

The geometric subset, relative to the chosen tangent vector, is
$\kappa_{\mathbb Q}(\mathbb Q^\times)$, which is proper.
For an explicit extra class take
$\alpha\in\widehat{\mathbb Z}$ whose $2$-adic component is
$-1/3\in\mathbb Z_2$ and whose other prime components are zero.
The compatible classes of $2^{\alpha\bmod n}$ in
$\mathbb Q^\times/\mathbb Q^{\times n}$ define a full Kummer class.
If it came from $a\in\mathbb Q^\times$, the ordinary integer
$\operatorname{ord}_2(a)$ would equal $\alpha$ in
$\widehat{\mathbb Z}$, hence would equal $-1/3$ in $\mathbb Z_2$,
which no integer does. Thus this class is not geometric.
The continuous-cohomology/Kummer comparison F149-KUM realizes it
by a continuous cocycle $z$, and (RG.4) sends the geometric
reference section to this nongeometric class.

This proves the assertion about the isolated extension. It does
not give an automorphism of the full curve diagram (RG.2):
Proposition~\ref{prop:archive-marked-recognition} shows precisely
why retaining that diagram removes this ambiguity.
\end{proof}

The scope of scan~85 can therefore be stated precisely. Geometric
splittings are recognizable from the full relative curve group
with its generic arithmetic marking, by (RG.3); generic geometricity
already forces global geometricity by normality and all-exponent
Kummer ramification. They cannot be selected from the isolated
normalizer extension by an automorphism-invariant rule. Recovering
the generic marking from a single unmarked base fundamental group
is not part of this positive theorem. The following subsection
refutes functorial recovery from the specified smaller full diagram
when its base group is allowed to move.
These statements are independent of step~17 and of the higher
$K$-group vanishing problem.

\subsection{Why the generic marking cannot simply be omitted}
\label{sec:archive-unmarked-curve}

\begin{proposition}[F149-UNMARKED-CURVE]
\label{prop:archive-unmarked-curve}
Even the full relative curve diagram
\[
 I\ \subset\ N\ \subset\ \Pi\ \longrightarrow\ G
       \longrightarrow G_{\mathbb Q},                       \tag{UC.1}
\]
with its identified arithmetic quotient, marked cusp inertia and
cyclotomic action, need not determine the geometric tangent-section
subset invariantly under its profinite automorphisms. This remains
false if one retains a reference constant tangent section and a
rational base section of $G\to G_{\mathbb Q}$.

More precisely, take
\[
 S=\mathbb G_{m,\mathbb Q}=\operatorname{Spec}\mathbb Q[t,t^{-1}],
 \quad C=\mathbb P^1_{\mathbb Q}\setminus\{0,1,\infty\},
 \quad U=S\times_{\mathbb Q}C,
\]
with marked relative cusp $0$ and its constant normal basis $b$.
There is an automorphism of (UC.1), acting identically on $I$ and
$G_{\mathbb Q}$, which carries the section class associated with the
nowhere-zero normal vector $t b$ to a class not associated with
any global nowhere-zero normal vector.
The automorphism is allowed to act nontrivially on $G$; no assertion
is made with $G$ fixed pointwise, with its scheme realization fixed,
or with the generic Galois marking retained.
\end{proposition}

\begin{proof}
Write $\Gamma=G_{\mathbb Q}$. Use the rational point $t=1$ to
split the arithmetic exact sequence for $S$. The geometric
fundamental group of $S$ is
\[
 A=\pi_1(S_{\overline{\mathbb Q}})=\widehat{\mathbb Z}(1),
 \qquad G=A\rtimes\Gamma.                                  \tag{UC.2}
\]
The first statement follows from the punctured-curve presentation
in \cite[Expos\'e XIII, Corollaire~2.12, pp.~290--292]{SGA1};
the compatible covers $v^n=t$ identify the Galois action on $A$
with the cyclotomic action. They have the compatible points $v=1$
over $t=1$. These covers and their Kummer normalization are
\cite[Expos\'e XI, Proposition~6.1 and \S6.2--6.5,
pp.~233--235]{SGA1}.

The constant tangent basis $b$ gives a section of the cusp
normalizer of $\pi_1(C)\to\Gamma$. Choose the resulting tangential
basepoint and write
\[
 \pi_1(C)=\Delta\rtimes\Gamma,
 \quad \Delta=\pi_1(C_{\overline{\mathbb Q}}),
 \quad I\subset\Delta.
\]
The normalizer at that cusp is $I\rtimes\Gamma$.
Here the normal-line comparison and its section are exactly
F149-RELATIVE, applied over $\operatorname{Spec}\mathbb Q$;
geometric self-normalization is also
\cite[Proposition~16, p.~8]{StixCuspidal2012}.
The action of $\Gamma$ on $I=\widehat{\mathbb Z}(1)$ is again
cyclotomic.

For completeness, the full product group, rather than merely
its normalizer, is available. The natural map
\[
 \pi_1(U)\longrightarrow
        \pi_1(S)\mathbin{\times}_{\Gamma}\pi_1(C)             \tag{UC.3}
\]
is an isomorphism. On the geometric kernels it is the
K\"unneth isomorphism of
\cite[Expos\'e XIII, Proposition~4.6, pp.~310--311]{SGA1}.
Its hypotheses hold for the two smooth curves over
$\overline{\mathbb Q}$: in characteristic zero all primes are
allowed, and resolution in dimension at most one suffices.
The arithmetic fundamental exact sequence
\cite[Expos\'e IX, Th\'eor\`eme~6.1, p.~195]{SGA1} then makes
(UC.3) a homomorphism of short exact sequences which is an
isomorphism on the kernel and on the quotient $\Gamma$.
It is consequently an isomorphism of profinite groups.
With the chosen sections this gives
\[
 \Pi=(A\times\Delta)\rtimes\Gamma,
 \qquad N=(A\times I)\rtimes\Gamma.                         \tag{UC.4}
\]
Indeed, $A$ centralizes $\Delta$, and within $\Delta$ the
normalizer of $I$ is $I$; the chosen arithmetic section normalizes
$I$. Thus the second group in (UC.4) is the actual full
normalizer in $\Pi$. Projection to $G$ forgets the $I$ coordinate.
It also identifies $N$ with $I\rtimes G$, where $G$ acts on $I$
through its cyclotomic quotient $\Gamma$.

Choose the profinite unit
\[
 u\in\widehat{\mathbb Z}^{\times},\qquad
 u_2=3\in\mathbb Z_2^\times,\quad
 u_\ell=1\in\mathbb Z_\ell^\times\quad(\ell\ne2).           \tag{UC.5}
\]
Multiplication by $u$ on $A$ commutes with its cyclotomic
$\Gamma$-action. Therefore
\[
 \begin{split}
 \alpha_G(a,g)&=(ua,g),\\
 \alpha_\Pi(a,\delta,g)&=(ua,\delta,g)
 \end{split}                                               \tag{UC.6}
\]
are continuous group automorphisms, with inverse obtained by
replacing $u$ by $u^{-1}$. They commute with the projections in
(UC.1), fix its arithmetic quotient, and preserve $N$.
The restriction $\alpha_N$ fixes $I$ pointwise. They also fix
the chosen $\Gamma$ section, the constant tangent section, and
indeed all geometric cusp-inertia subgroups lying in $\Delta$.

Identify $A$ with $I$ through their common Kummer normalization.
Write an element of $N=I\rtimes G$ as $(c;(a,g))$.
Relative to the constant tangent section, the normal vector
$t b$ has cocycle
\[
 c_t(a,g)=a.                                               \tag{UC.7}
\]
This can be checked at every finite level using $v^n=t$:
the geometric inertia acts by its tautological $\mu_n$ quotient,
and the section $t=1$, lifted by $v=1$, has zero cocycle.
Equivalently, (UC.7) satisfies the cocycle identity directly,
since the two inertia modules have the same action of $g$.
Consequently the corresponding section is
$s_t(a,g)=(a;(a,g))$.

Transporting a section by an automorphism of the diagram means
$s\mapsto\alpha_N\circ s\circ\alpha_G^{-1}$. Thus the
transported $s_t$ is represented by
\[
 c'_t(a,g)=u^{-1}a.                                       \tag{UC.8}
\]
Every global nowhere-zero normal vector is of the form
$f b$, where
\[
 f\in\mathbb Q[t,t^{-1}]^\times=\mathbb Q^\times t^{\mathbb Z}.
                                                               \tag{UC.9}
\]
To see the last equality directly, the difference between the
highest and lowest exponent of a nonzero Laurent polynomial
is additive under multiplication over a field. A unit and its
inverse therefore both have exponent difference zero.
Write $f=a_0t^m$ with $a_0\in\mathbb Q^\times$, $m\in\mathbb Z$.
Kummer functoriality shows that its section cocycle restricts
to $A$ as multiplication by $m$; the constant factor has
zero restriction to $A$.

Restriction to $A$ is unchanged by $I$-conjugacy of sections:
$I$ is abelian and $A$ acts trivially on it, so every coboundary
restricts to zero. If (UC.8) were geometrically represented,
we would therefore have $m=u^{-1}$ in
$\widehat{\mathbb Z}$ for some ordinary integer $m$.
Its $2$-adic component would imply $3m=1$ in $\mathbb Z_2$,
hence $3m=1$ in $\mathbb Z$, which is impossible.
This proves that the transported section class is nongeometric.
A uniform reversal of normal/conormal inertia conventions changes
both slopes by the same sign and does not change the contradiction.
\end{proof}

\begin{corollary}[F149-NO-FUNCTORIAL-UNMARKING]
\label{cor:archive-no-unmarking}
On the class of the realizable diagrams in (UC.1), the full
generic arithmetic marking required in
F149-MARKED-RECOGNITION cannot be recovered functorially in a way
that both recovers its geometric realization and lifts every
isomorphism of the smaller diagram. In particular no
isomorphism-invariant recognition rule on the smaller data can
select all and only geometric tangent sections.
\end{corollary}

\begin{proof}
Such a recovery would lift the automorphism (UC.6) to an
isomorphism of full marked generic diagrams. By
F149-MARKED-RECOGNITION that isomorphism preserves the geometric
section subset. Proposition~\ref{prop:archive-unmarked-curve}
shows that its induced automorphism of the smaller diagram does
not preserve that subset, a contradiction.
The assertion concerns functorial recovery from the specified
profinite data. It does not exclude adding the scheme $S$, its
rational-point markings, or the generic Galois group as further
input, nor a noncanonical choice that is not respected by
all diagram isomorphisms.
\end{proof}

\section{Schemes, weights, and the limits of the comparison}
\label{sec:framework}

All arithmetic cohomology over the rational integers in this paper lives on
$U_\ell=\operatorname{Spec}\mathbb Z[1/\ell]$. The Tate lattice
$\mathbb Z_\ell(r)$ is the inverse system of finite etale sheaves
$\mu_{\ell^n}^{\otimes r}$, interpreted by continuous cohomology.
The discrete sheaf $\mathbb Q_\ell/\mathbb Z_\ell(-e)$ is its corresponding
divisible torsion coefficient object. These coefficient conventions,
the local conditions, and the compact/discrete duality are part of the
statements; replacing them by rational cohomology loses the finite groups
being bounded.

The weight convention is Deligne's \emph{Weil II}, Definitions 1.2.1--1.2.2,
p.~153, and 1.2.5(iv), p.~154 \cite{DeligneWeilII}: $\mathbb Q_\ell(1)$ has weight $-2$.
Thus the rational Tate object $\mathbb Q_\ell(-e)$ has weight $2e$,
and $\mathbb Q_\ell(e+1)$ has weight $-2(e+1)$. These statements refer
to rational sheaves and geometric Frobenius. They do not assign a
Frobenius purity theorem to a torsion module or prove its vanishing.
The finite $K$-groups at issue have zero rationalization even when their
integral vanishing is unknown. This elementary observation explains why
a weight or sign argument requires the additional integral congruence
and duality statements reconstructed above.

For the archive passage, \emph{SGA 4}, IX.3.2--3.3 supplies the
prime-invertible Kummer sequence and Hilbert 90. The first Chern class
is the Kummer boundary of a line bundle; \emph{SGA 5}, VII,
Proposition 3.4(i), identifies the degree-one term of the total Chern
class with that boundary \cite{SGA5}. The finite root-coordinate gluing proof in
the archive section fixes the precise form needed here. The tame
cover comparison is \emph{SGA 1}, XIII.5.2--5.6, with X.3.6 for the
discrete-valuation Abhyankar lemma. The extension of regular functions
across codimension at least two uses \emph{EGA} $\mathrm{IV}_2$,
5.8.6, 5.9.8--5.9.9 and 5.10.5. Recognition with the full generic
arithmetic marking additionally uses Mochizuki's anabelian theorems;
it is not a consequence of an FGA representability statement.

The $\lambda$-ring language belongs first to $K_0$: exterior powers
define $\lambda^i$, and the determinant or first gamma-filtration
piece records line bundles. An exact locator is \emph{SGA 6}, V,
Definition 2.1, Example 2.2(b), and Lemma 2.2.1 \cite{SGA6}. Higher $K_{4k}$-groups
require the higher etale comparison used in the $K$-group section. Neither
the $K_0$ operations nor a first Chern class is that comparison map.

Crystals and divided powers concern a different coefficient setting,
especially reduction in residue characteristic $\ell$. No new theorem
about a crystal, a divided-power envelope, or a crystalline comparison
is proved in this run. The etale Kummer calculation deliberately removes
$\ell$ from the base; for noninvertible exponents one must use flat
Kummer theory and then state a separate comparison with any crystalline
or prismatic object. The inseparable counterexample in the archive
section proves why the etale assertion itself cannot be transported
unchanged. Assigning one of the present finite Selmer estimates to a
crystalline object without such a map would add an unproved assertion.

Finally, the rational mixed-Tate category over $\operatorname{Spec}
\mathbb Z$ is the category used by Deligne--Goncharov and Brown. Its
base scheme controls ramification; its rational coefficients control
which extensions remain visible. The integral arithmetic lattice of a
pro-$\ell$ image is an additional object. The counterexample at $37$
shows that these distinctions change the answer to B3.

\subsection{Location of the new deductions in these frameworks}

The following grouping specifies the geometric object and the exact
framework used by each family of new findings. A source in the last
column supplies the stated language or comparison; it is not credited
with a later arithmetic or computational conclusion that it does not
contain.

\small
\begin{longtable}{p{.23\textwidth}p{.34\textwidth}p{.34\textwidth}}
\toprule
Finding family & Geometric object & Precise framework and scope\\
\midrule
\endhead
K17-S17, HOST, PT-Q; VK-EXP7 and POLY-CERT &
Strict-local-condition cohomology of
$\operatorname{Spec}\mathcal O_E[1/\ell]$ with negative Tate
coefficients; higher $K$-theory of $\operatorname{Spec}\mathbb Z$. &
SGA~4 IX.3.2--3.3 fixes finite Kummer coefficients; continuous limits
and arithmetic duality are the separately cited Jannsen/Milne inputs.
Weil~II 1.2.1--1.2.5 fixes rational weights, not torsion vanishing.\\
K17-HOST-GREENBERG-GATE &
Norm towers of finite arithmetic cohomology and finitely generated
Iwasawa modules. &
The same finite etale coefficient category; the countermodels test
growth inequalities only. No crystal or divided-power object is
identified with these modules.\\
K17-P2-MOMENT, VANDIVER-KERNEL, GREENBERG-KERNEL;
CUP-EXT and RATIONAL-MAGNUS &
The arithmetic action on the geometric pro-$\ell$ fundamental group
of $\mathbb P^1-\{0,1,\infty\}$; its weight filtration and
finite-level cup products. &
SGA~1 XIII.5.2--5.6 controls tame cusp inertia. SGA~4 IX.3.2--3.3
and the stated low-degree etale/group comparison place the cups.
The mixed-Tate rational target is Deligne--Goncharov \S1.6,
Propositions~1.7,1.9,2.3 and \S2.4, with Brown's cited theorem.\\
K17-OUT-SAT, 2-ALL and the two integral classifications &
An actual $\mathbb Z_\ell$ Lie lattice inside the rational
Tannakian image, with finite torsion in its abelianization. &
The above mixed-Tate framework supplies the rational comparison.
The integral lattice and its even torsion are additional arithmetic
data; SGA~6 V.2.1--2.2.1 on $K_0$ does not identify them.\\
K17-K12-THREE-RANKS, DOUBLE, CIRCUIT-SIMPLE,
ORIENT-SURVIVOR and SURVIVOR-CYCLE &
Quillen's rank filtration for $K(\operatorname{Spec}\mathbb Z)$,
Steinberg coefficients and integral Voronoi chains. &
These use the exact Quillen, Voronoi, building and Hurewicz sources
cited in their proofs. SGA~6 V describes the degree-zero
$\lambda$-ring; no claim identifies its operations with this
higher-homology computation.\\
F149-KUM, FULLINJ, CHERN and RELATIVE &
The line at a marked boundary section, its nonzero-vector scheme,
and the relative cusp-normalizer extension. &
Cote~149 scans~79--82 and~85; SGA~4 IX.3.2--3.3;
SGA~5 VII.3.4(i); SGA~1 XIII.4.1,4.4 and Deligne~15.19,15.26.
Normal and conormal orientations are kept distinct.\\
F149-GENERIC, MARKED-RECOGNITION and BARE-EXTENSION &
Geometric tangent splittings inside the full continuous
$H^1$ torsor, with the generic Galois marking retained. &
Cote~149 scans~83--85; EGA~$\mathrm{IV}_2$~5.10.5 for extension
from codimension one; finite Kummer at every exponent;
Mochizuki Theorems~A and~B for the extra anabelian reconstruction.
The isolated extension fails by the explicit cocycle automorphism.\\
\bottomrule
\end{longtable}
\normalsize



\clearpage
\part{Unresolved findings, failed implications, and continuation gates}
\section{Open assertions are not used as theorems}

\paragraph{OPEN-K-NUM: a specified new vanishing degree.}
The existential polynomial order bound and its explicit exponent do
not supply a small numerical constant. A new degree such as $K_{12}$
requires evaluated fixed-level cohomology bounds, or computation of the
two nonzero trace products and their gcd, or the three localized
Steinberg-homology vanishings in Proposition~\ref{prop:k12-three-ranks}.
The rank-$11$ term is eliminated independently. A dimension discrepancy
in the older enumeration leaves the announced $H_5(\mathrm{GL}_8;\mathrm{St}_8)$
endpoint without an independently checked source proof; it is retained
along with $H_4(\mathrm{GL}_9;\mathrm{St}_9)$ and
$H_3(\mathrm{GL}_{10};\mathrm{St}_{10})$ over $\mathbb Z[1/13!]$.
None of these three computations is claimed here.
Proposition~\ref{prop:k12-double-building} gives an alternative finite
criterion in terms of three explicitly identified double-building
homology groups for the last two ranks, with the rank-eight endpoint
still separate. These hypotheses remain unproved.
The small-index orientation elimination is proved, but
Proposition~\ref{prop:k12-orientation-survivor} gives an explicit
index-three chain generator with no orientation reverser. A full
vanishing proof must therefore address cycles and boundaries, not
just eliminate every cell by stabilizer parity.
Proposition~\ref{prop:k12-survivor-boundary} now supplies explicit
boundaries with coefficient one for the two displayed surviving
cycles. Proposition~\ref{prop:k12-trimodular-generated} further removes
every rank-ten degree-twelve cell whose minima generate the ambient
lattice and whose maximal minimum-basis index is at most three.
Theorem~\ref{thm:k12-ternary-cycles} classifies the proper index-three
glue cases as twenty-six explicit cycles, and
Proposition~\ref{prop:k12-index-three-boundaries} fills every one
with coefficient one. The rank-ten remaining chain gate is therefore
only cells of maximal index at least four.
Theorem~\ref{thm:k12-r8-nonsimplicial} removes every nonsimplicial
rank-eight endpoint cell, including ambient glue.
Corollary~\ref{cor:k12-r8-index-five} confines each remaining cell
to thirteen independent minimum rays and minimum-basis index at
most five, with an explicit necessary list of frame quotient types.
No incoming boundary matrix for those residual cells is computed.
Proposition~\ref{prop:k12-r8-glue-five} additionally rules out
proper minimum-generated lattice index five in this residual class.

\paragraph{OPEN-IRREG-KERNEL: irregular odd primes.}
The proved separation theorem handles cyclic even Iwasawa components of
arbitrary finite rank, including repeated and nonsplit roots.
Fukaya--Kato and Wake--Wang-Erickson supply its cup condition under
Vandiver, and Theorem~\ref{thm:greenberg-kernel} now extends the
conclusion to the weaker one-variable hypothesis that $X^+$ is finite.
It retains the entire free pro-$p$ group and permits bounded rational
Iwasawa coordinates and repeated roots. The remaining all-prime
question therefore reduces to Greenberg at the remaining primes, or
to a further argument with a positive-rank plus Iwasawa module.
Every prime below $2^{31}$ is already covered without that conjectural
hypothesis. The proof does not give a converse to Greenberg.
Theorem~\ref{thm:odd-torsion-cubic} permits cyclic torsion in each
odd character under explicit pair and triple companion nonvanishing
conditions. It resolves the separation of nested pair/triple poles,
but no new arithmetic prime is covered: arithmetic Triple-NV and
the additional module/cup hypotheses remain unproved.
Corollary~\ref{cor:wwe-pair} now supplies the pair condition under
the precise rational cyclicity and support-disjointness hypotheses
of Wake--Wang-Erickson, even with odd torsion present. Arithmetic
Triple-NV remains independent and unproved.

\paragraph{The former OPEN-2-LATTICE gate is closed.}
Theorem~\ref{thm:2-allweights} constructs every adapted Newton lift
by an integral parity correction. Proposition~\ref{prop:newton-basis}
then proves free generation in all weights, and
Proposition~\ref{prop:outer-sat} identifies the actual outer filtration.
The earlier finite and class-two computations remain separate checked
results; neither is extrapolated to obtain the theorem.

\paragraph{Scan-85 recognition: the generic marking is necessary in general.}
Theorem~\ref{thm:archive-generic} proves that a splitting over the
smooth normal arithmetic base is geometric exactly when its generic
restriction is geometric; no Picard finite-generation assumption is
needed for this equivalence.
Proposition~\ref{prop:archive-marked-recognition} determines that
subset from the full curve diagram with the marked generic absolute
Galois group over an identified number field. This is a functorial
recognition statement, not an effective finite-quotient algorithm.
Proposition~\ref{prop:archive-bare-extension} proves that the isolated
normalizer extension cannot determine the subset, even with its
cyclotomic inertia action. Proposition~\ref{prop:archive-unmarked-curve}
extends the negative result to the full relative curve diagram over
$\mathbb G_m/\mathbb Q$ when its base group may move. It preserves
the arithmetic quotient and cusp inertia but moves an integer Kummer
slope to a noninteger profinite slope.
Corollary~\ref{cor:archive-no-unmarking} therefore rules out functorial
recovery of the generic marking that lifts every smaller-diagram
isomorphism. It makes no claim when the base group is fixed pointwise.
Arbitrary
fibrations and wild mixed-characteristic extensions are outside the
proved relative statement.

\paragraph{Source and transcription scope.}
The broader host-field point theorem now has a closed source proof
composition, with its numerical layer bounds retained as premises;
the character theorem and rational-field specialization are also checked. Missing original
round-one B3/S9 text and missing handwritten archive images are not
reconstructed from conjecture. Mathematical editorial flags are
distinguished from verified readings in TRANSCRIPTION\_NOTES.tsv.

\section{Audit conventions and source reproducibility}
CHECK\_EVENTS.tsv records actual source readings, whole-proof checks,
independent re-proofs and exact computations. A cited theorem may be
used with its statement and hypotheses checked without a fresh audit of
its entire original proof; that distinction is retained. The source
contract accepts Alp\"oge's supplied theorems and does not reopen RH.
The pinned mirror, checksums and retrieved Lean excerpts are in the
evidence directory. Research notes also record primary URLs and exact
page or declaration locators. Earlier candidate notes are retained for
continuity but do not supersede the checked paper and latest ledger.

\begin{thebibliography}{99}

\bibitem{Kurihara1992}
M. Kurihara,
Some remarks on conjectures about cyclotomic fields and \(K\)-groups of \(\mathbb Z\),
\emph{Compositio Math.} 81 (1992), 223--236.
\url{https://www.numdam.org/item/CM_1992__81_2_223_0.pdf}.
Exact locators used: Lemma 1.2, p.225; Corollary 1.5, p.226;
Proposition 3.7, p.229; Remark 3.11, p.230;
(4), Proposition 5.1 and Remark 5.2, p.233.


\bibitem{Weibel2013}
C. A. Weibel,
\emph{The \(K\)-book: An Introduction to Algebraic \(K\)-theory},
Graduate Studies in Mathematics 145, AMS, 2013.
Author's Chapter VI:
\url{https://sites.math.rutgers.edu/~weibel/Kbook/Kbook.VI.pdf}.
Exact locators used: VI.8.2, VI.8.3, VI.9.8, VI.10.9.


\bibitem{RognesWeibel2000}
J. Rognes and C. Weibel,
Two-primary algebraic \(K\)-theory of rings of integers in number fields,
\emph{J. Amer. Math. Soc.} 13 (2000), 1--54.
\url{https://sites.math.rutgers.edu/~weibel/archive/papers-dir/RognesWeibel.pdf}.
Theorem 0.7, pp.4--5, proved as Theorem 7.11;
localization discussion p.1 and Theorem 0.6.


\bibitem{Rognes2000}
J. Rognes, \(K_4(\mathbb Z)\) is the trivial group,
\emph{Topology} 39 (2000), 267--281.
\url{https://doi.org/10.1016/S0040-9383(99)00007-5}.
Theorem 0.1, p.268.


\bibitem{DEKM2026}
M. Dutour Sikiri\'c, P. Elbaz-Vincent, A. Kupers and J. Martinet,
Voronoi complexes in higher dimensions, cohomology of
\(\operatorname{GL}_N(\mathbb Z)\) for \(N\ge8\) and the triviality of
\(K_8(\mathbb Z)\),
\emph{J. Inst. Math. Jussieu} 25 (2026), 565--590.
\url{https://doi.org/10.1017/S1474748025101394}.
Published online 24 November 2025.
Theorem 6.4; author preprint arXiv:1910.11598v1, p.19,
and Corollary 7.1, p.20:
\url{https://arxiv.org/abs/1910.11598}.


\bibitem{HartHarveyOng2017}
W. B. Hart, D. Harvey and W. Ong,
Irregular primes to two billion,
\emph{Math. Comp.} 86 (2017), 3031--3049.
\url{https://doi.org/10.1090/mcom/3211}.
Author preprint:
\url{https://arxiv.org/abs/1605.02398}.
Section 1, item (2), verifies Vandiver for every odd prime
\(\ell<2^{31}\).


\bibitem{MilneADT}
J. S. Milne, \emph{Arithmetic Duality Theorems},
second edition, BookSurge, 2006.
\url{https://www.jmilne.org/math/Books/ADT.pdf}.
Finite duality locators: I, Theorem 4.10; II, Theorem 2.9.


\bibitem{Jannsen1988}
U. Jannsen, Continuous \'etale cohomology,
\emph{Math. Ann.} 280 (1988), 207--245.
\url{https://doi.org/10.1007/BF01456052}.
Author's institutional copy:
\url{https://epub.uni-regensburg.de/26684/1/jannsen12.pdf}.
Exact inverse-limit sequence (3.1).


\bibitem{MilneANT}
J. S. Milne, \emph{Algebraic Number Theory},
course notes, version 3.08, 2020.
\url{https://www.jmilne.org/math/CourseNotes/ANT.pdf}.
Theorem 4.3, printed p.70.


\bibitem{EtingofGL2024}
P. Etingof, \emph{Lie Groups and Lie Algebras II},
MIT 18.755, Spring 2024.
Lecture 01, sections 27.2--27.3, printed pp.145--147:
\url{https://ocw.mit.edu/courses/18-755-lie-groups-and-lie-algebras-ii-spring-2024/mit18_755_s24_lec01.pdf}.
Lecture 02, section 28.1, printed pp.151--152:
\url{https://ocw.mit.edu/courses/18-755-lie-groups-and-lie-algebras-ii-spring-2024/mit18_755_s24_lec02.pdf}.
Highest-weight classification, determinant twists, and Weyl's dimension formula.


\bibitem{IshiiKummer2024}
S. Ishii,
On Kummer characters arising from Galois actions on the pro-\(p\)
fundamental groups of once-punctured CM elliptic curves,
\emph{Research in Number Theory} 11 (2025), article 38.
\url{https://doi.org/10.1007/s40993-025-00622-0}.
Author preprint arXiv:2312.04175v2:
\url{https://arxiv.org/pdf/2312.04175}.
Definition preceding Theorem 1.1 and Theorem 1.1, p.2.


\bibitem{IchimuraKaneko1989}
H. Ichimura and M. Kaneko,
On the universal power series for Jacobi sums and the Vandiver conjecture,
\emph{J. Number Theory} 31 (1989), 312--334.
Author's copy:
\url{https://www2.math.kyushu-u.ac.jp/~mkaneko/papers/04universal_power.pdf}.
Section 1.1, pp.313--314.


\bibitem{BuhlerHarvey2011}
J. P. Buhler and D. Harvey,
Irregular primes to 163 million,
\emph{Math. Comp.} 80 (2011), 2435--2444.
Author preprint:
\url{https://arxiv.org/abs/0912.2121}.
Section 4, pp.4--5 of version 2, for the Vandiver certificate criterion.


\bibitem{Kupers2019}
A. Kupers, A short proof that \(K_8(\mathbb Z)\cong0\),
author's note dated 21 October 2019.
\url{https://www.utsc.utoronto.ca/people/kupers/wp-content/uploads/sites/50/2021/01/k8zshorter.pdf}.
Theorem A and Remark 2.6, p.4.


\bibitem{BrueckEtAl2024}
B. Br\"uck, J. Miller, P. Patzt, R. J. Sroka and J. C. H. Wilson,
On the codimension-two cohomology of \(\operatorname{SL}_n(\mathbb Z)\),
\emph{Adv. Math.} 451 (2024), article 109795.
\url{https://doi.org/10.1016/j.aim.2024.109795}.
Author preprint arXiv:2204.11967v3:
\url{https://arxiv.org/abs/2204.11967}.
Theorem B: rational codimension-two vanishing for \(n\ge3\).
Proposition~8.5, p.~57; Lemmas~8.6--8.9 and proof of Theorem~B,
pp.~58--59, are the inspected inputs for the localized refinement.

% Source-check refinement for BrueckEtAl2024: Proposition8.5 p57, Lemmas8.6-8.9 pp58-59, proof of TheoremB p59; exact localized deduction in K17-ST2-LOCAL.

\bibitem{Sharifi2002}
R.~T.~Sharifi, \emph{Relationships between conjectures on the structure
of pro-$p$ Galois groups unramified outside $p$},
Proc. Sympos. Pure Math. \textbf{70} (2002), 275--284.
Author's 13-page version: \url{https://www.math.ucla.edu/~sharifi/publ/galstr.pdf}.


\bibitem{Brown2012}
F.~Brown, \emph{Mixed Tate motives over $\mathbb Z$},
Ann. of Math. \textbf{175} (2012), 949--976.
\url{https://annals.math.princeton.edu/wp-content/uploads/annals-v175-n2-p10-p.pdf}.


\bibitem{Brown2019}
F.~Brown, \emph{From the Deligne--Ihara conjecture to multiple modular values},
arXiv:1904.00179 (2019).
\url{https://arxiv.org/pdf/1904.00179}.


\bibitem{HainMatsumoto2001}
R.~Hain and M.~Matsumoto, \emph{Weighted completion of Galois groups
and Galois actions on the fundamental group of $\mathbb P^1-\{0,1,\infty\}$},
arXiv:math/0006158v2, 11 August 2001.
\url{https://arxiv.org/pdf/math/0006158}.


\bibitem{GrasJaulent2008}
G.~Gras and J.-F.~Jaulent, \emph{Note on $2$-rational fields},
J. Number Theory \textbf{129} (2009), 495--498;
arXiv:0805.1168.
\url{https://arxiv.org/pdf/0805.1168}.


\bibitem{MinacPasiniQuadrelliTan2020}
J.~Min\'a\v c, F.~Pasini, C.~Quadrelli, and Nguy{\~{\^e}}n Duy T\^an,
\emph{Koszul algebras and quadratic duals in Galois cohomology},
arXiv:1808.01695v3, 31 December 2020.
\url{https://arxiv.org/pdf/1808.01695}.


\bibitem{IshiiKernel2026}
S.~Ishii, \emph{On the kernels of the pro-$p$ outer Galois representations
associated to once-punctured CM elliptic curves},
arXiv:2312.04196v3, 6 January 2026.
\url{https://arxiv.org/pdf/2312.04196}.


\bibitem{McCallumSharifi2003}
W.~G.~McCallum and R.~T.~Sharifi,
\emph{A cup product in the Galois cohomology of number fields},
Duke Math. J. \textbf{120} (2003), 269--310.
Author's 34-page version dated 5 December 2003:
\url{https://www.math.ucla.edu/~sharifi/publ/pairing.pdf}.


\bibitem{Iwasawa1981}
K.~Iwasawa, \emph{Riemann--Hurwitz formula and $p$-adic Galois
representations for number fields},
Tohoku Math. J. \textbf{33} (1981), 263--288.
\url{https://www.jstage.jst.go.jp/article/tmj1949/33/2/33_2_263/_pdf/-char/en}.


\bibitem{SharifiIwasawaNotes}
R.~T.~Sharifi, \emph{Iwasawa Theory}, online notes, Chapter~3,
especially Corollary~3.3.14, Propositions~3.4.1 and~3.4.7,
and Corollary~3.4.8. Accessed 21 November 2026.
\url{https://www.math.ucla.edu/~sharifi/notes/iwasawa-ch03.html}.


\bibitem{DeligneGoncharov2005}
P.~Deligne and A.~B.~Goncharov,
\emph{Groupes fondamentaux motiviques de Tate mixte},
Ann. Sci. \'Ecole Norm. Sup. (4) \textbf{38} (2005), 1--56.
The supplied LaTeX source was used for \S1.6,
Propositions~1.7, 1.9 and~2.3, and \S2.4.


\bibitem{Hirose2026}
M.~Hirose, \emph{Confluence relations, Hoffman basis over $\mathbb Z_{(2)}$,
and the Deligne--Drinfeld conjecture},
arXiv:2610.23719v1, 16 October 2026.
\url{https://arxiv.org/pdf/2610.23719}.


\bibitem{Sharifi2007} R. Sharifi, \emph{Iwasawa theory and the Eisenstein ideal}, Duke Math. J. \textbf{137} (2007), 63--101. Locators refer to arXiv:math/0501236v4 (17 March 2007): Theorem 1.2 p.~2; Theorem 5.6 pp.~25--26; Theorem 5.7 and Corollaries 5.8--5.9 p.~27; weight-twelve application p.~28. \url{https://arxiv.org/abs/math/0501236}.

\bibitem{SGA1}
A.~Grothendieck et al., \emph{Rev\^etements \'etales et groupe fondamental
(SGA 1)}, Lecture Notes in Mathematics \textbf{224}, Springer, 1971;
annotated reedition, Documents Math\'ematiques \textbf{3},
Soci\'et\'e Math\'ematique de France, 2003.
Expos\'e IX, Th\'eor\`eme~6.1, p.~195;
Expos\'e XI, Proposition~6.1 and \S6.2--6.5, pp.~233--235;
Expos\'e XIII, Corollaire~2.12, pp.~290--292,
and Proposition~4.6, pp.~310--311.
Expos\'e XIII, Proposition~5.2 and Corollaire~5.3,
reedition pp.~315--316 (original marginal pp.~428--430);
relative versions, Proposition~5.5 and Corollaire~5.6;
Expos\'e X, Lemme~3.6.
Expos\'e XIII, Proposition~4.1, reedition pp.~306--307,
and Exemples~4.4, p.~309, for the relative homotopy sequence.
\url{https://arxiv.org/pdf/math/0206203}.


\bibitem{SGA4}
M.~Artin, A.~Grothendieck, and J.-L.~Verdier (eds.),
\emph{Th\'eorie des topos et cohomologie \'etale des sch\'emas
(SGA 4)}, Tome~3, Lecture Notes in Mathematics \textbf{305},
Springer, 1973. Expos\'e IX, Th\'eorie de Kummer~3.2 and
Th\'eor\`eme~3.3; French retypeset edition, printed pp.~14--15
(PDF pages~19--20).
\url{https://www.normalesup.org/~forgogozo/SGA4/tomes/tome3.pdf}.


\bibitem{StixCuspidal2012}
J.~Stix, \emph{On cuspidal sections of algebraic fundamental groups},
author version dated 24 February 2012.
Proposition~11, p.~7; Proposition~16 and Theorem~17, p.~8;
Corollary~19, p.~9; Corollary~32, p.~12.
\url{https://www.math.uni-frankfurt.de/~stix/research/preprints/STIXcuspsec20120802homepage.pdf}.


\bibitem{EGAIV2}
A.~Grothendieck, with the collaboration of J.~Dieudonn\'e,
\emph{\'El\'ements de g\'eom\'etrie alg\'ebrique IV:
\'Etude locale des sch\'emas et des morphismes de sch\'emas,
seconde partie}, Publ. Math. IH\'ES \textbf{24} (1965), 5--231.
Th\'eor\`eme~5.8.6, p.~108; Proposition~5.9.8 and
D\'efinition~5.9.9, pp.~111--112;
Th\'eor\`eme~5.10.5 and proof, pp.~115--116.
Primary scan:
\url{https://webusers.imj-prg.fr/~leila.schneps/grothendieckcircle/EGA/EGAIV.2.pdf}.
Numdam record: \url{https://archive.numdam.org/item/PMIHES_1965__24__5_0/}.

\bibitem{DeligneWeilII}
P.~Deligne, \emph{La conjecture de Weil. II},
Publ. Math. IHES \textbf{52} (1980), 137--252.
Definitions 1.2.1--1.2.2, p.~153, and 1.2.5(iv), p.~154.
\url{https://www.numdam.org/item/PMIHES_1980__52__137_0.pdf}.


\bibitem{SGA5}
A.~Grothendieck et al., \emph{Cohomologie $\ell$-adique et fonctions $L$},
SGA~5, Lecture Notes in Mathematics 589, Springer, 1977.
Expose VII, standing invertibility convention p.~282; Proposition 3.4(i), p.~291, proof p.~292; Corollary 3.7(i), p.~295. Original scanned pages visually checked.
\url{https://wstein.org/sga/sga5/SGA5.pdf}.


\bibitem{SGA6}
P.~Berthelot, A.~Grothendieck and L.~Illusie, eds.,
\emph{Theorie des intersections et theoreme de Riemann--Roch},
SGA~6, Lecture Notes in Mathematics 225, Springer, 1971.
Expose V, Definition 2.1, Example 2.2(b), Lemma 2.2.1.
Checked section transcription:
\url{https://grothendiecksga.com/read/sga6/fr/V-2.html}.

\bibitem{KochKukkukLabute1998}
H.~Koch, S.~Kukkuk and J.~Labute, \emph{Nilpotent local class field theory},
Acta Arithmetica \textbf{83} (1998), 45--64.
Theorem~2.1, pp.~48--49 (author PDF pp.~4--5).
\url{https://www.math.mcgill.ca/labute/papers/kkl.pdf}.

\bibitem{WakeWangErickson2017}
P.~Wake and C.~Wang-Erickson,
\emph{Ordinary pseudorepresentations and modular forms},
Proc. Amer. Math. Soc. Ser.~B \textbf{4} (2017), 53--71.
\url{https://doi.org/10.1090/bproc/29}.
Final author preprint arXiv:1510.01661v4, 30 January 2017:
\url{https://arxiv.org/pdf/1510.01661v4}.
Verified locators in that version: \S2.1, pp.~2--3;
Proposition~2.2.6, p.~4; Theorem~4.2.8, pp.~10--11;
Corollaries~5.1.1--5.1.2, p.~11; Lemma~5.2.7, p.~13.


\bibitem{FukayaKato2012}
T.~Fukaya and K.~Kato,
\emph{On conjectures of Sharifi}, author preprint,
15 August 2012, 121 pages.
\url{https://www.math.ucla.edu/~sharifi/sharificonj.pdf}.
All locators used here refer to this inspected author version:
\S\S6.1.4--6.1.8, pp.~76--77;
\S7.1.3, pp.~84--85;
Conjecture~7.1.8 and \S\S7.1.9--7.1.10, p.~86;
Theorem~7.2.8 and \S\S7.2.10--7.2.12, p.~89;
Lemma~10.2.2 and \S10.2.3, p.~107.

\bibitem{Arlettaz2000}
D.~Arlettaz,
\emph{Algebraic $K$-theory of rings from a topological viewpoint},
Publicacions Matem\`atiques \textbf{44} (2000), 3--84.
\url{https://mat.uab.es/pubmat/fitxers/download/FileType:pdf/FolderName:v44(1)/FileName:44100_01.pdf}.
Definition~6.5 and Theorem~6.6, p.~40;
the stable-to-unstable Hurewicz comparison, pp.~41--42.

\bibitem{MillerPatztPutman2025}
J.~Miller, P.~Patzt, and A.~Putman,
\emph{Homological vanishing for the Steinberg representation II:
reductive groups and integral conjectures},
arXiv:2509.01559v1 (1 September 2025), 74 pp.,
\url{https://arxiv.org/abs/2509.01559}.
Theorem~C and Theorem~D, p.~6; Corollary~E and Remark~1.18,
p.~7; Theorem~23.3 and Lemma~23.5, pp.~49--51;
Theorem~C$'$.1, p.~54; Lemma~25.3, Claim~2, pp.~56--57;
Theorem~D$'$, p.~64.

\bibitem{DeligneTangential1989}
P.~Deligne, \emph{Le groupe fondamental de la droite projective
moins trois points}, in \emph{Galois Groups over $\mathbb Q$},
Math. Sci. Res. Inst. Publ. \textbf{16}, Springer, 1989, 79--297.
\S15.19, p.~239 (relative tangential specialization);
Lemma~15.23, pp.~241--242; \S15.26, p.~242
(formal-coordinate interpretation and its relative tame version).
\url{https://publications.ias.edu/sites/default/files/61_LeGroupeFondamentalDroite.pdf}.

\bibitem{WakeWangEricksonPseudo2018}
P.~Wake and C.~Wang-Erickson,
\emph{Pseudo-modularity and Iwasawa theory},
Amer. J. Math. \textbf{140} (2018), no.~4, 977--1040.
\url{https://doi.org/10.1353/ajm.2018.0022}.
The inspected author version is arXiv:1505.05128v3,
7 September 2017, 54 pages:
\url{https://arxiv.org/pdf/1505.05128v3}.
Locators in that version: Corollary~2.1.3 and Conjecture~2.1.5, p.~8;
Corollaries~8.1.2--8.1.3, pp.~48--49;
Theorem~8.2.3, pp.~50--51.


\bibitem{MilneLEC2013}
J.~S.~Milne,
\emph{Lectures on \textup{\'Etale} Cohomology},
version~2.21, 22 March 2013, 202 pages.
\url{https://www.jmilne.org/math/CourseNotes/LEC.pdf}.
Theorem~14.9, Remark~14.10, Example~14.11 and Proposition~14.13,
printed pp.~96--98.

\bibitem{Ribet1976}
K.~A. Ribet,
\emph{A modular construction of unramified $p$-extensions of
$\mathbb Q(\mu_p)$},
Inventiones Mathematicae \textbf{34} (1976), 151--162.
Main Theorem~(1.1), p.~151.
\url{https://math.berkeley.edu/~ribet/Articles/invent_34.pdf}.
DOI: \href{https://doi.org/10.1007/BF01403065}{10.1007/BF01403065}.

\bibitem{MochizukiLocal1999}
S.~Mochizuki, \emph{The local pro-$p$ anabelian geometry of curves},
Invent. Math. \textbf{138} (1999), 319--423.
Author manuscript: Theorem~A, p.~3; Theorem~B, p.~4;
Theorem~16.5 and its following full profinite comparison, p.~86;
Theorem~17.1 and proof, p.~87.
\url{https://www.kurims.kyoto-u.ac.jp/~motizuki/The%20Local%20Pro-p%20Anabelian%20Geometry%20of%20Curves.pdf}.
DOI: \url{https://doi.org/10.1007/s002220050381}.

\bibitem{DutourHulekSchurmann2015}
M. Dutour Sikiri\'c, K. Hulek, and A. Sch\"urmann,
\emph{Smoothness and singularities of the perfect form and the second
Voronoi compactification of $\mathcal A_g$},
Algebraic Geometry \textbf{2} (2015), no.~5, 642--653.
doi:10.14231/AG-2015-027.
Theorem~1.1(i), printed p.~644; cone-dimension convention and proof,
pp.~645--646.


\bibitem{Heller1957}
I. Heller, \emph{On linear systems with integral valued solutions},
Pacific Journal of Mathematics \textbf{7} (1957), no.~3, 1351--1364.
Theorem~4.2 and proof, pp.~1358--1360.


\bibitem{LeePaatStallknechtXu2022}
J. Lee, J. Paat, I. Stallknecht, and L. Xu,
\emph{Polynomial upper bounds on the number of differing columns of
$\Delta$-modular integer programs},
Mathematics of Operations Research \textbf{48} (2023), no.~4, 2267--2286.
doi:10.1287/moor.2022.1339.
Primary preprint arXiv:2105.08160v5 (16 November 2022),
definitions pp.~1 and 3, Theorem~1 p.~4, proof \S4 pp.~17--19.

\bibitem{ElbazGanglSoule2013}
P. Elbaz-Vincent, H. Gangl and C. Soul\'e,
\emph{Perfect forms, $K$-theory and the cohomology of modular groups},
Advances in Mathematics \textbf{245} (2013), 587--624.
Theorem~4.3, p.~597; Theorem~8.1, p.~604.
\url{https://www.ihes.fr/~soule/documents/YAIMA4438.pdf}.
DOI: \url{https://doi.org/10.1016/j.aim.2013.06.014}.

\bibitem{AverkovSchymura2024}
G. Averkov and M. Schymura,
\emph{On the maximal number of columns of a $\Delta$-modular integer
matrix: bounds and computations},
Mathematical Programming \textbf{206} (2024), 61--89.
doi:10.1007/s10107-023-01964-2.
Section~5.2 and Tables~1 and~3, printed pp.~84--85;
data at \texttt{https://github.com/mschymura/delta-classification}.

\bibitem{K17PinnedLeanMirror}
Supplied public Lean source mirror,
\emph{riemann-hypothesis-lean},
pinned commit
\texttt{7cbf7a0c32fca1badf05931879a37c3f8a2b03b4}.
Source-content locators are given by complete relative file and
declaration names in the text; no fresh project build is claimed.
\url{https://github.com/CocoSgt/riemann-hypothesis-lean/tree/7cbf7a0c32fca1badf05931879a37c3f8a2b03b4}.

\bibitem{BergeMartinet2009}
A.-M. Berg\'e and J. Martinet,
\emph{On perfection relations in lattices},
Contemporary Mathematics \textbf{493} (2009), 29--49.
Locators used here refer to the 26-page author preprint:
Definition~2.4 p.~3; \S3 p.~6; Theorem~5.1 pp.~15--16;
Corollary~6.9 pp.~22--23.
\url{https://arxiv.org/abs/math/0611220}.


\bibitem{MartinetIndex2001}
J. Martinet, with an appendix by C. Batut,
\emph{Sur l'indice d'un sous-r\'eseau},
in \emph{R\'eseaux euclidiens, designs sph\'eriques et formes modulaires},
Monographies de L'Enseignement Math\'ematique \textbf{37}
(2001), 163--211.
Locators refer to the 51-page author PDF:
introductory theorem p.~3; Theorem~5.2 p.~16;
Theorem~10.1 pp.~39--40.
\url{https://jamartin.perso.math.cnrs.fr/Publications/watsonensmath.pdf}.


\bibitem{MartinetIndexErrata2014}
J. Martinet,
\emph{Erratum and complements to Monographie 37},
updated 3 April 2014.
Corrections and complements for the sixth paper, pp.~3--5;
appendix completing the rank-eight index bound, pp.~6--8.
\url{https://jamartin.perso.math.cnrs.fr/Publications/erratcompl37.pdf}.

\bibitem{SharifiReciprocity2017}
R.~T.~Sharifi,
\emph{Reciprocity maps with restricted ramification},
inspected author version dated 30 June 2017, 47 pages.
\url{https://www.math.ucla.edu/~sharifi/selmer.pdf}.
Section~2.5, printed pp.~16--17, gives the inverse group-action
convention and corestriction transitions for Iwasawa cohomology.
The finite restriction formulas used here are proved explicitly.

\end{thebibliography}
\end{document}
